% Options for packages loaded elsewhere
% Options for packages loaded elsewhere
\PassOptionsToPackage{unicode}{hyperref}
\PassOptionsToPackage{hyphens}{url}
%
\documentclass[
  letterpaper,
]{scrbook}
\usepackage{xcolor}
\usepackage[left=1in,textwidth=4.4in,marginparwidth=2.0in,marginparsep=0.3in]{geometry}
\usepackage{amsmath,amssymb}
\setcounter{secnumdepth}{2}
\usepackage{iftex}
\ifPDFTeX
  \usepackage[T1]{fontenc}
  \usepackage[utf8]{inputenc}
  \usepackage{textcomp} % provide euro and other symbols
\else % if luatex or xetex
  \usepackage{unicode-math} % this also loads fontspec
  \defaultfontfeatures{Scale=MatchLowercase}
  \defaultfontfeatures[\rmfamily]{Ligatures=TeX,Scale=1}
\fi
\usepackage{lmodern}
\ifPDFTeX\else
  % xetex/luatex font selection
\fi
% Use upquote if available, for straight quotes in verbatim environments
\IfFileExists{upquote.sty}{\usepackage{upquote}}{}
\IfFileExists{microtype.sty}{% use microtype if available
  \usepackage[]{microtype}
  \UseMicrotypeSet[protrusion]{basicmath} % disable protrusion for tt fonts
}{}
\makeatletter
\@ifundefined{KOMAClassName}{% if non-KOMA class
  \IfFileExists{parskip.sty}{%
    \usepackage{parskip}
  }{% else
    \setlength{\parindent}{0pt}
    \setlength{\parskip}{6pt plus 2pt minus 1pt}}
}{% if KOMA class
  \KOMAoptions{parskip=half}}
\makeatother
% Make \paragraph and \subparagraph free-standing
\makeatletter
\ifx\paragraph\undefined\else
  \let\oldparagraph\paragraph
  \renewcommand{\paragraph}{
    \@ifstar
      \xxxParagraphStar
      \xxxParagraphNoStar
  }
  \newcommand{\xxxParagraphStar}[1]{\oldparagraph*{#1}\mbox{}}
  \newcommand{\xxxParagraphNoStar}[1]{\oldparagraph{#1}\mbox{}}
\fi
\ifx\subparagraph\undefined\else
  \let\oldsubparagraph\subparagraph
  \renewcommand{\subparagraph}{
    \@ifstar
      \xxxSubParagraphStar
      \xxxSubParagraphNoStar
  }
  \newcommand{\xxxSubParagraphStar}[1]{\oldsubparagraph*{#1}\mbox{}}
  \newcommand{\xxxSubParagraphNoStar}[1]{\oldsubparagraph{#1}\mbox{}}
\fi
\makeatother

\usepackage{color}
\usepackage{fancyvrb}
\newcommand{\VerbBar}{|}
\newcommand{\VERB}{\Verb[commandchars=\\\{\}]}
\DefineVerbatimEnvironment{Highlighting}{Verbatim}{commandchars=\\\{\}}
% Add ',fontsize=\small' for more characters per line
\usepackage{framed}
\definecolor{shadecolor}{RGB}{241,243,245}
\newenvironment{Shaded}{\begin{snugshade}}{\end{snugshade}}
\newcommand{\AlertTok}[1]{\textcolor[rgb]{0.68,0.00,0.00}{#1}}
\newcommand{\AnnotationTok}[1]{\textcolor[rgb]{0.37,0.37,0.37}{#1}}
\newcommand{\AttributeTok}[1]{\textcolor[rgb]{0.40,0.45,0.13}{#1}}
\newcommand{\BaseNTok}[1]{\textcolor[rgb]{0.68,0.00,0.00}{#1}}
\newcommand{\BuiltInTok}[1]{\textcolor[rgb]{0.00,0.23,0.31}{#1}}
\newcommand{\CharTok}[1]{\textcolor[rgb]{0.13,0.47,0.30}{#1}}
\newcommand{\CommentTok}[1]{\textcolor[rgb]{0.37,0.37,0.37}{#1}}
\newcommand{\CommentVarTok}[1]{\textcolor[rgb]{0.37,0.37,0.37}{\textit{#1}}}
\newcommand{\ConstantTok}[1]{\textcolor[rgb]{0.56,0.35,0.01}{#1}}
\newcommand{\ControlFlowTok}[1]{\textcolor[rgb]{0.00,0.23,0.31}{\textbf{#1}}}
\newcommand{\DataTypeTok}[1]{\textcolor[rgb]{0.68,0.00,0.00}{#1}}
\newcommand{\DecValTok}[1]{\textcolor[rgb]{0.68,0.00,0.00}{#1}}
\newcommand{\DocumentationTok}[1]{\textcolor[rgb]{0.37,0.37,0.37}{\textit{#1}}}
\newcommand{\ErrorTok}[1]{\textcolor[rgb]{0.68,0.00,0.00}{#1}}
\newcommand{\ExtensionTok}[1]{\textcolor[rgb]{0.00,0.23,0.31}{#1}}
\newcommand{\FloatTok}[1]{\textcolor[rgb]{0.68,0.00,0.00}{#1}}
\newcommand{\FunctionTok}[1]{\textcolor[rgb]{0.28,0.35,0.67}{#1}}
\newcommand{\ImportTok}[1]{\textcolor[rgb]{0.00,0.46,0.62}{#1}}
\newcommand{\InformationTok}[1]{\textcolor[rgb]{0.37,0.37,0.37}{#1}}
\newcommand{\KeywordTok}[1]{\textcolor[rgb]{0.00,0.23,0.31}{\textbf{#1}}}
\newcommand{\NormalTok}[1]{\textcolor[rgb]{0.00,0.23,0.31}{#1}}
\newcommand{\OperatorTok}[1]{\textcolor[rgb]{0.37,0.37,0.37}{#1}}
\newcommand{\OtherTok}[1]{\textcolor[rgb]{0.00,0.23,0.31}{#1}}
\newcommand{\PreprocessorTok}[1]{\textcolor[rgb]{0.68,0.00,0.00}{#1}}
\newcommand{\RegionMarkerTok}[1]{\textcolor[rgb]{0.00,0.23,0.31}{#1}}
\newcommand{\SpecialCharTok}[1]{\textcolor[rgb]{0.37,0.37,0.37}{#1}}
\newcommand{\SpecialStringTok}[1]{\textcolor[rgb]{0.13,0.47,0.30}{#1}}
\newcommand{\StringTok}[1]{\textcolor[rgb]{0.13,0.47,0.30}{#1}}
\newcommand{\VariableTok}[1]{\textcolor[rgb]{0.07,0.07,0.07}{#1}}
\newcommand{\VerbatimStringTok}[1]{\textcolor[rgb]{0.13,0.47,0.30}{#1}}
\newcommand{\WarningTok}[1]{\textcolor[rgb]{0.37,0.37,0.37}{\textit{#1}}}

\usepackage{longtable,booktabs,array}
\newcounter{none} % for unnumbered tables
\usepackage{calc} % for calculating minipage widths
% Correct order of tables after \paragraph or \subparagraph
\usepackage{etoolbox}
\makeatletter
\patchcmd\longtable{\par}{\if@noskipsec\mbox{}\fi\par}{}{}
\makeatother
% Allow footnotes in longtable head/foot
\IfFileExists{footnotehyper.sty}{\usepackage{footnotehyper}}{\usepackage{footnote}}
\makesavenoteenv{longtable}
\usepackage{graphicx}
\makeatletter
\newsavebox\pandoc@box
\newcommand*\pandocbounded[1]{% scales image to fit in text height/width
  \sbox\pandoc@box{#1}%
  \Gscale@div\@tempa{\textheight}{\dimexpr\ht\pandoc@box+\dp\pandoc@box\relax}%
  \Gscale@div\@tempb{\linewidth}{\wd\pandoc@box}%
  \ifdim\@tempb\p@<\@tempa\p@\let\@tempa\@tempb\fi% select the smaller of both
  \ifdim\@tempa\p@<\p@\scalebox{\@tempa}{\usebox\pandoc@box}%
  \else\usebox{\pandoc@box}%
  \fi%
}
% Set default figure placement to htbp
\def\fps@figure{htbp}
\makeatother


% definitions for citeproc citations
\NewDocumentCommand\citeproctext{}{}
\NewDocumentCommand\citeproc{mm}{%
  \begingroup\def\citeproctext{#2}\cite{#1}\endgroup}
\makeatletter
 % allow citations to break across lines
 \let\@cite@ofmt\@firstofone
 % avoid brackets around text for \cite:
 \def\@biblabel#1{}
 \def\@cite#1#2{{#1\if@tempswa , #2\fi}}
\makeatother
\newlength{\cslhangindent}
\setlength{\cslhangindent}{1.5em}
\newlength{\csllabelwidth}
\setlength{\csllabelwidth}{3em}
\newenvironment{CSLReferences}[2] % #1 hanging-indent, #2 entry-spacing
 {\begin{list}{}{%
  \setlength{\itemindent}{0pt}
  \setlength{\leftmargin}{0pt}
  \setlength{\parsep}{0pt}
  % turn on hanging indent if param 1 is 1
  \ifodd #1
   \setlength{\leftmargin}{\cslhangindent}
   \setlength{\itemindent}{-1\cslhangindent}
  \fi
  % set entry spacing
  \setlength{\itemsep}{#2\baselineskip}}}
 {\end{list}}
\usepackage{calc}
\newcommand{\CSLBlock}[1]{\hfill\break\parbox[t]{\linewidth}{\strut\ignorespaces#1\strut}}
\newcommand{\CSLLeftMargin}[1]{\parbox[t]{\csllabelwidth}{\strut#1\strut}}
\newcommand{\CSLRightInline}[1]{\parbox[t]{\linewidth - \csllabelwidth}{\strut#1\strut}}
\newcommand{\CSLIndent}[1]{\hspace{\cslhangindent}#1}



\setlength{\emergencystretch}{3em} % prevent overfull lines

\providecommand{\tightlist}{%
  \setlength{\itemsep}{0pt}\setlength{\parskip}{0pt}}



 


% Math macros for the 6.1600 lecture notes.
% Trimmed from the original latexdefs.tex to the macros actually used in
% the text. Every definition here must be MathJax-compatible: this file is
% included in the PDF preamble AND (via macros/_macros.qmd, regenerated by
% scripts/gen-macros.sh) in every HTML page.
%
% Blackboard-bold sets and groups
\newcommand{\F}{\mathbb{F}}
\newcommand{\G}{\mathbb{G}}
\newcommand{\N}{\mathbb{N}}
\newcommand{\R}{\mathbb{R}}
\newcommand{\Z}{\mathbb{Z}}
\newcommand{\Zq}{\Z_q}
% Bitstrings
\newcommand{\zo}{\{0,1\}}
\newcommand{\bin}{\zo}
\newcommand{\zon}{\{0,1\}^n}
% Calligraphic letters
\newcommand{\A}{\mathcal{A}}
\newcommand{\B}{\mathcal{B}}
\newcommand{\calA}{\mathcal{A}}
\newcommand{\calB}{\mathcal{B}}
\newcommand{\calC}{\mathcal{C}}
\newcommand{\calD}{\mathcal{D}}
\newcommand{\calK}{\mathcal{K}}
\newcommand{\calM}{\mathcal{M}}
\newcommand{\calT}{\mathcal{T}}
\newcommand{\calX}{\mathcal{X}}
\newcommand{\calY}{\mathcal{Y}}
% Operators and symbols
\newcommand{\xor}{\oplus}
\newcommand{\cross}{\times}
\newcommand{\deq}{\mathrel{\mathop:}=}
\newcommand{\abs}[1]{\left| #1 \right|}
\newcommand{\set}[1]{\left\{ #1 \right\}}
\renewcommand{\vec}[1]{\mathbf{#1}}
\newcommand{\poly}{\operatorname{poly}}
\newcommand{\negl}{\operatorname{negl}}
\newcommand{\classP}{\mathsf{P}}
\newcommand{\classNP}{\mathsf{NP}}
% Cryptographic primitives
\newcommand{\INDCPA}{\text{IND-CPA}}
\newcommand{\Enc}{\mathsf{Enc}}
\newcommand{\Dec}{\mathsf{Dec}}
\newcommand{\Gen}{\mathsf{Gen}}
\newcommand{\Derive}{\mathsf{Derive}}
\newcommand{\Sign}{\mathsf{Sign}}
\newcommand{\Ver}{\mathsf{Ver}}
\newcommand{\MAC}{\mathsf{MAC}}
\newcommand{\MACSign}{\mathsf{MAC.Sign}}
\newcommand{\MACVerify}{\mathsf{MAC.Verify}}
\newcommand{\Funs}{\mathsf{Funs}}
\newcommand{\ToString}{\mathsf{ToString}}
% Keys and signatures
\newcommand{\sk}{\mathsf{sk}}
\newcommand{\pk}{\mathsf{pk}}
\newcommand{\vk}{\mathsf{vk}}
\newcommand{\sig}{\sigma}
% PDF-only macro definitions (not MathJax-compatible).
% The HTML build uses the rewrites in macros-html-overrides.tex instead.
%
% "Sample at random" arrow: a left arrow with a small R above it.
\newcommand{\rgetscmd}{\ooalign{$\leftarrow$\cr
    \hidewidth\raisebox{1.2\height}{\scalebox{0.5}{\ \textrm{R}}}\hidewidth\cr}}
\newcommand{\rgets}{\mathrel{\mathpalette\rgetscmd\relax}}
\newcommand{\getsr}{\rgets}

% Draft TODO notes. The Lua filter in scripts/todo-latex.lua wraps every
% .todo span in \todobox; xcolor is already loaded by the pandoc template.
% Notes reach the PDF only under the draft profile, so this never ships.
\definecolor{todocolor}{HTML}{B35A00}
\newcommand{\todobox}[1]{{\color{todocolor}#1}}
\makeatletter
\@ifpackageloaded{tcolorbox}{}{\usepackage[skins,breakable]{tcolorbox}}
\@ifpackageloaded{fontawesome5}{}{\usepackage{fontawesome5}}
\definecolor{quarto-callout-color}{HTML}{909090}
\definecolor{quarto-callout-note-color}{HTML}{0758E5}
\definecolor{quarto-callout-important-color}{HTML}{CC1914}
\definecolor{quarto-callout-warning-color}{HTML}{EB9113}
\definecolor{quarto-callout-tip-color}{HTML}{00A047}
\definecolor{quarto-callout-caution-color}{HTML}{FC5300}
\definecolor{quarto-callout-color-frame}{HTML}{acacac}
\definecolor{quarto-callout-note-color-frame}{HTML}{4582ec}
\definecolor{quarto-callout-important-color-frame}{HTML}{d9534f}
\definecolor{quarto-callout-warning-color-frame}{HTML}{f0ad4e}
\definecolor{quarto-callout-tip-color-frame}{HTML}{02b875}
\definecolor{quarto-callout-caution-color-frame}{HTML}{fd7e14}
\makeatother
\makeatletter
\@ifpackageloaded{bookmark}{}{\usepackage{bookmark}}
\makeatother
\makeatletter
\@ifpackageloaded{caption}{}{\usepackage{caption}}
\AtBeginDocument{%
\ifdefined\contentsname
  \renewcommand*\contentsname{Table of contents}
\else
  \newcommand\contentsname{Table of contents}
\fi
\ifdefined\listfigurename
  \renewcommand*\listfigurename{List of Figures}
\else
  \newcommand\listfigurename{List of Figures}
\fi
\ifdefined\listtablename
  \renewcommand*\listtablename{List of Tables}
\else
  \newcommand\listtablename{List of Tables}
\fi
\ifdefined\figurename
  \renewcommand*\figurename{Figure}
\else
  \newcommand\figurename{Figure}
\fi
\ifdefined\tablename
  \renewcommand*\tablename{Table}
\else
  \newcommand\tablename{Table}
\fi
}
\@ifpackageloaded{float}{}{\usepackage{float}}
\floatstyle{ruled}
\@ifundefined{c@chapter}{\newfloat{codelisting}{h}{lop}}{\newfloat{codelisting}{h}{lop}[chapter]}
\floatname{codelisting}{Listing}
\newcommand*\listoflistings{\listof{codelisting}{List of Listings}}
\usepackage{amsthm}
\theoremstyle{plain}
\newtheorem{lemma}{Lemma}[chapter]
\theoremstyle{plain}
\newtheorem{theorem}{Theorem}[chapter]
\theoremstyle{definition}
\newtheorem{definition}{Definition}[chapter]
\theoremstyle{remark}
\AtBeginDocument{\renewcommand*{\proofname}{Proof}}
\newtheorem*{remark}{Remark}
\newtheorem*{solution}{Solution}
\newtheorem{refremark}{Remark}[chapter]
\newtheorem{refsolution}{Solution}[chapter]
\makeatother
\makeatletter
\makeatother
\makeatletter
\@ifpackageloaded{caption}{}{\usepackage{caption}}
\@ifpackageloaded{subcaption}{}{\usepackage{subcaption}}
\makeatother
\makeatletter
\@ifpackageloaded{sidenotes}{}{\usepackage{sidenotes}}
\@ifpackageloaded{marginnote}{}{\usepackage{marginnote}}
\makeatother
\usepackage{bookmark}
\IfFileExists{xurl.sty}{\usepackage{xurl}}{} % add URL line breaks if available
\urlstyle{same}
\hypersetup{
  pdftitle={The Structure of Secure Systems},
  pdfauthor={Henry Corrigan-Gibbs, Srini Devadas, Yael Kalai, Nickolai Zeldovich},
  hidelinks,
  pdfcreator={LaTeX via pandoc}}


\title{The Structure of Secure Systems}
\author{Henry Corrigan-Gibbs, Srini Devadas, Yael Kalai, Nickolai
Zeldovich}
\date{}
\begin{document}
\frontmatter
\maketitle


\mainmatter
\bookmarksetup{startatroot}

\chapter*{Disclaimer}\label{disclaimer}
\addcontentsline{toc}{chapter}{Disclaimer}

\markboth{Disclaimer}{Disclaimer}

Charles Alfred Meurer, Still Life with Currency, 1895. Image from the
Art Institute of Chicago.

This set of notes is a work in progress. It may have errors and be
missing citations. It is certainly incomplete. Please let the staff know
of any errors that you find.

\section*{Contributors}\label{contributors}
\addcontentsline{toc}{section}{Contributors}

\markright{Contributors}

These notes are based on lectures by the 6.1600 course staff:

\begin{itemize}
\tightlist
\item
  \textbf{2024:} Srini Devadas and Yael Kalai
\item
  \textbf{2023:} Henry Corrigan-Gibbs and Nickolai Zeldovich
\item
  \textbf{2022:} Henry Corrigan-Gibbs, Yael Kalai, and Nickolai
  Zeldovich.
\item
  \textbf{2021:} Henry Corrigan-Gibbs, Srini Devadas, Yael Kalai, and
  Nickolai Zeldovich.
\end{itemize}

Ben Kettle, our course TA in 2022, was responsible for transcribing the
first set of these lecture notes in Fall 2022.

\bookmarksetup{startatroot}

\chapter{What is Security?}\label{what-is-security}

\section{Overview}\label{overview}

This course will give you an overview of the most important ``Big
Ideas'' on securing computer systems. Throughout the course, we will
touch on ideas from the fields of computer security, cryptography, and
(to some extent) computer systems.

\section{What is Security?}\label{what-is-security-1}

The goal of computer security is to ensure that a particular computer
system \emph{behaves correctly} even in the face of an \emph{adversary}
(or \emph{attacker}) whose goal is to foil the system.

\marginnote{\begin{footnotesize}

We will use the terms ``adversary'' and ``attacker'' interchangeably
throughout this course.

\end{footnotesize}}

To achieve this goal, we---as the designer of a computer system---will
need some kind of systematic plan. That is, we will have to carefully
define what it means for our system to \emph{behave correctly} and we
will have to specify the class of \emph{adversaries} against which we
want to defend.

For the purposes of this course, we will typically structure our plan in
terms of three components: a \emph{model} of the world in which our
system lives, a security \emph{goal}, an \emph{implementation}.

\begin{itemize}
\item
  \textbf{Model:} The system model specifies:

  \marginnote{\begin{footnotesize}

  Sometimes people call this the ``threat model.'' For the purposes of
  this chapter, think of the model as describing \emph{how} the attacker
  and system interact and the security goal as defining \emph{what} our
  implementation is trying to achieve.

  Many textbooks and courses use the term ``threat model'' to refer to
  some combination of the model and the security goal

  \end{footnotesize}}

  \begin{enumerate}
  \def\labelenumi{\arabic{enumi}.}
  \tightlist
  \item
    what the system is that we are trying to defend,
  \item
    what the attacker is, and
  \item
    how the system and attacker interact.
  \end{enumerate}

  For example, in network security we might think of the system and the
  attacker as being two computers that interact by sending messages over
  a network. In this model, the attacker can send maliciously crafted
  network packets to the victim system. (The attacker in this model
  \emph{cannot}, for example, swap out the hard drive of the victim
  system.)

  When we are studying hardware security, we might consider a different
  model: The system we are trying to defend is a CPU, and the adversary
  is an external device that can read all of the contents of the victim
  system's RAM. (The attacker in this model \emph{cannot}, for example,
  read the internal state of the victim's CPU.)

  When we are studying ATM security, the system we are trying to defend
  is an ATM machine. For an attacker, we might consider a person
  interacting with the ATM through it's normal interface: the attacker
  can insert a malformed ATM card into the machine and type arbitrary
  PINs into the device. (The attacker in this model cannot, for example,
  take a jackhammer to the ATM to extract the cash.)

  Modeling the adversary and the system in question is almost always the
  first step of thinking about system security. While we often only have
  an informal system model in mind, the more precise you---as a system
  designer---can be about your system model, the clearer your security
  properties can be.
\item
  \textbf{Goal:} The security goal defines what we want our system to
  achieve in our specified model.

  For example, in network security, we might want the property that
  ``only someone knowing Alice's secret password can execute shell
  commands on the machine.'' In hardware security, we might want the
  property that ``the attacker learns nothing about the data stored in
  RAM, apart from its size.'' In ATM security, we might want the
  property that ``an attacker can fraudulently authenticate as a victim
  with probability at most \(1/10000\).''

  As you will learn throughout the course, figuring out exactly what
  your security goal should be is often quite subtle and challenging.
\item
  \textbf{Implementation:} The implementation is how we achieve the
  goal. For example, mobile phones use PIN-based authentication to gate
  access to the private information on your phone.
\end{itemize}

Together, the model and the goal create our \emph{definition} of
security. As such, the model and goal cannot be ``wrong''---no more than
one person's definition of ``beautiful art'' can be wrong.

But it can be that a model and/or goal can be weak or bad. For example,
a poor threat model might not capture all of the attacks that a
real-world adversary can mount. For example, if we model an attacker and
defender as computers interactive at length over a network, no matter
how strong our security goal in that model is, we should not expect to
protect against an attacker that physically tampers with the defender's
hardware.

In addition, the security goal that we define might turn out to not be
exactly what we needed. For example, at first glance you might think of
defining the goal of an encryption system as: ``A secure encryption
system prevents an attacker from learning the encrypted message.'' An
encryption system that satisfies this security property might still leak
\emph{half of the bits of the encrypted message} to the attacker. This
example shows how goals that seem strong might turn out to be not nearly
strong enough.

Often, the process of designing the model and goal is iterative: when
the system designer discovers a surprising gap in the security goal (we
will see some shortly), she patches the goal and the implementation
accordingly.

\marginnote{\begin{footnotesize}

When you read about security failures in the news, it is worth trying to
understand whether the failure arose from a problem with the
\emph{model}, the \emph{goal}, or the \emph{implementation}.

\end{footnotesize}}

The implementation, on the other hand, can definitely be wrong---if the
implementation does not guarantee the desired security goal under the
designer's model, due to bugs or oversights or supply chain
vulnerabilities or anything else, the implementation has a mistake.

\textbf{The implementation is always public: Kerckhoff's principle.}
Throughout this course, we will always assume that the attacker knows
the implementation of whatever security mechanism we are using. (The
model and goal are always public as well.) The \emph{only} thing that we
keep secret from the adversary are the system's secret keys.

For example, if we are designing a bank vault, we will analyze its
security assuming that the adversary knows everything about the vault's
design. The only thing that we assume is secret is the passcode that
opens the vault.

This is known as \emph{Kerckhoff's principle}. The logic behind this way
of thinking is that: (1) it is often relatively easy for the attacker to
learn bits of information about an implementation and (2) it is much
easier to replace a set of cryptographic keys (if the attacker learns
them) than to redesign the entire system from scratch.

\subsection{Security is Hard}\label{security-is-hard}

Building secure systems is challenging. There are at least two broad
reasons for this.

\textbf{Secure systems must defend against worst-case behavior.} When we
are just thinking about functionality or correctness, we are often
satisfied with a system that performs well in the common case. If a
system behaves correctly on 99.9999\% of inputs, this is more than good
enough for many non-security-critical settings. In other words,
correctness is often about behavior in \emph{expected} or
\emph{average-case} situations.

In contrast, a secure system must defend against \emph{all possible
attacks} within the scope of the system model. If a computer system
behaves incorrectly on 0.0001\% of inputs, a clever attacker will figure
out a way to find these inputs and feed them to the system---breaking
the system's security goal. Thus, a secure system must behave correctly
in all cases---even on the 0.0001\% of inputs that a normal programmer
wouldn't worry about. In other words, security is concerned with
behavior in the \emph{worst-case}.

\marginnote{\begin{footnotesize}

Some developers do worry about correctness of their system for all
possible inputs and corner cases. In contrast, security-minded
programmers \emph{do} worry about all possible inputs and corner cases.

\end{footnotesize}}

For example, suppose that you are a car manufacturer and you want to
test a car stereo functions correctly. Testing that the stereo works
well \emph{on average} (i.e., in expected situations) is easy: turn the
stereo on and off 10,000 times, try playing some music through it, and
accept it as working if all of these checks pass. Testing that the
stereo works well \emph{in the worst case} is not as easy. For example,
an attacker might craft a malicious audio file that contains some
exploit code. If the stereo plays this file, the exploit code will
hijack the car and disables the car's breaks. How can we be sure that
our car's stereo is not vulnerable to such an attack? You will rarely
find these malicious packets by random testing---you will need more
powerful techniques, which we will discuss in this course.

\marginnote{\begin{footnotesize}

These attacks are actually possible in practice! See:
(\citeproc{ref-koscher2010experimental}{{Koscher et al.} 2010})

\end{footnotesize}}

To summarize: A secure system must defend against all possible attacks
within the threat model; being certain that a system satisfies this
strong security property is a challenge.

\textbf{An implementation can never defend against all possible
threats.} When we specify a system model, we delimit the set of
adversaries against which our implementation must defend. But real-world
adversaries can behave in ways that are outside of our model and thereby
violate our security goals.

For example, if we are running a course at a university, we might want
to make sure that only the course staff is able to access student
grades. This sounds like a simple goal but there are an essentially
unlimited number of possible ways to violate it! For example, an
attacker could:

\begin{itemize}
\tightlist
\item
  find a bug in the server software,
\item
  break into the professor's office,
\item
  compromise the professor's laptop via a software exploit,
\item
  steal the password to an administrator's account,
\item
  trick a teaching assistant into disclosing grades,
\item
  break the server's cryptography,
\item
  get a job at the university registrar's office getting permissions to
  read course grades.
\end{itemize}

And the list never ends. Because our threat model cannot capture all
possible threats, \textbf{security is never perfect}. There will
essentially always be \emph{some} attacker than can break your system.

This is why we need a threat model: the threat model defines what kinds
of attacks we worry about and which we decide are out of scope.

\textbf{Security is about economics.} Since a system cannot defend
against all possible threats, the designer needs to decide (via the
model and security goal) which threats are in scope and which threats
are out of scope.

Typically, the more threats are in scope, the more expensive the
implementation will be to build and maintain. For example, a bank vault
that is designed to protect against an attacker equipped with plastic
explosives will be more expensive than one that is designed to defend
against an attacker who only has a sledgehammer.

The \emph{cost} of defenses is what limits the defenses that we deploy
in practice. A designer typically balances the cost of new defenses
against the benefit that the defense brings in preventing various
attacks. For example, it usually does not make sense to spend \$1,000 on
a home safe to protect \$20 of jewelry. It might make sense to spend
\$1,000 on a home safe to protect \$50,000 of jewelry.

These cost-benefit trade-offs drive real-world decisions about which
defenses to deploy.

\section{Examples}\label{examples}

We give a handful of examples of security failures arising from poor
choices of security goals or threat model.

\todobox{\textbf{TODO (HCG):} Revise the following text.}

\subsection{Insufficient Attack
Models}\label{insufficient-attack-models}

\textbf{CAPTCHA.} CAPTCHAs were designed to be expensive to solve via
automation, but easy for a human to read. Indeed it might be expensive
to build an optical-character recognition system for CAPTCHAs in
general, but attackers who want to bypass CAPTCHAs do not do this.
Instead, they set up computer centers in countries where the cost of
labor is cheap. Attackers then pay people working in these centers to
solve CAPTCHAs. (\citeproc{ref-motoyama:recaptchas}{Motoyama et al.
2010}) The result is that it costs some fraction of a cent to solve a
CAPTCHA. The cost of solving a CAPTCHA is still non-zero, but the cost
is much lower than the system designers may have intended.

\textbf{DES.} There used to be an encryption standard called DES that
had \(2^{56}\) possible keys. At the time that it was designed, the U.S.
government standards agencies asserted that DES was secure against even
powerful attackers, but today a modern computer can try all \(2^{56}\)
keys at only modest cost. For what it's worth, academic researchers even
at the time of DES's design understood that 56-bit keys were not large
enough to prevent exhaustive cryptanalysis (\citeproc{ref-DH77}{Diffie
and Hellman 1977}). Because of the weakness of DES against modern
computers, everyone using the standard had to upgrade their block
ciphers.

\marginnote{\begin{footnotesize}

For example, \url{https://crack.sh/} uses an array of FPGAs to provide a
service that exhaustively checks all possible keys.

\end{footnotesize}}

\textbf{Model capturing a secure out-of-band channel: Two-factor
authentication (2FA) via SMS.} Many 2FA systems use a text message for
authentication, but an attacker then just needs to convince the clerk at
the AT\&T store to give them a new SIM card for your phone number.

\marginnote{\begin{footnotesize}

These attacks are often called ``SIM swapping'' or ``SIM hijacking''.

\end{footnotesize}}

\subsubsection{Assuming a correct compiler:
Xcode}\label{sec-intro-xcode}

\marginnote{\begin{footnotesize}

This prescient paper from 1984 anticipated the XcodeGhost attack:
(\citeproc{ref-thompson1984reflections}{Thompson 1984})

\end{footnotesize}}

iPhone apps are normally created and compiled on a developer's machine,
sent to Apple's App Store, and sent to iPhones from there. iPhone apps
are created using a tool called Xcode that is normally downloaded from
Apple servers. However, Xcode is a big piece of software and for
developers behind China's firewall, it was very slow to download.
Someone within China set up a much faster mirror of Xcode, and lots of
developers in China used the version of Xcode from this mirror. However,
this mirror was not serving exactly Apple's version of Xcode---instead,
it was serving a slightly modified version of Xcode that would inject
some malicious code into every app that was compiled with it. This took
a long time to detect.

\marginnote{\begin{footnotesize}

The Wikipedia article on XcodeGhost,
\url{https://en.wikipedia.org/wiki/XcodeGhost}, provides more details
about this attack.

\end{footnotesize}}

\subsection{Insufficient security
goal}\label{insufficient-security-goal}

\textbf{Business-class airfare.} An airline tried to add value to their
business-class tickets by allowing ticket-holders to change the ticket
(i.e., the departure date, origin, and destination) at any time with no
fee. One customer realized that they could board the flight \emph{then}
change their ticket. The customer could then take an unlimited number of
business-class flights for the price of one.

In this case, the airline's goal did not meet their real needs---perhaps
they needed to add an additional goal of the form ``every time someone
takes a flight, we get paid.''

\textbf{Sarah Palin's Email.} Sarah Palin had a Yahoo email account, and
Yahoo used security questions for password reset---their goal may have
been something like ``no one can reset a user's password unless they
know all of the answers to the user's security questions.'' (The
security questions are typically things like ``What is your mother's
maiden name?'') As it turned out, it was possible to find the answer to
all of Palin's account-recovery security questions on the Internet.
(\citeproc{ref-palin}{Zetter 2008}) Yahoo's implementation may have been
perfect, but their goal did not provide any meaningful security for
certain users.

\textbf{Instruction Set Architecture (ISA) Specification.} When defining
ISAs for processors, computer architects thought it would be acceptable
for each instruction to take a different number of cycles to execute.
This had big benefits for performance and for compatibility, but as
we'll talk about later in the semester, researchers have recently
exploited this timing variability to perform sophisticated attacks on
wide ranges of processors. (\citeproc{ref-hill:spectre-meltdown}{Hill et
al. 2019}) Even if a processor's implementation faithfully implements
the specification, the specification itself allows for certain types of
timing side-channel attacks.

\textbf{Complicated access-control policies.} A school in Fairfax,
Virginia used an online course-management software with a somewhat
complex access-control structure: each teacher is in charge of some
class, each class has many students, and each student has many files.
Teachers cannot access student's files, and there is also a
superintendent that has access to all the files. Teachers are able to
change their students' passwords, and are able to add students to their
class. It turned out that a teacher could add the superintendent as a
student, change the superintendent's password, and then access all files
via the superintendent's account. While each of the access-control
policies individually sounds reasonable, together they lead to a bad
outcome.

\textbf{Mat Honan's Gmail Account.} A journalist for wired named Mat
Honan had his Gmail account compromised via a clever attack.
(\citeproc{ref-honan}{Honan 2012}) Honan had a Gmail account. Gmail's
reset-password feature avoided using security questions, and instead
used a backup email account. The attacker triggered the reset-password
feature, which sent a reset link to Honan's Apple email account.

The attacker then attempted to gain access to Honan's Apple email
account. The attacker triggered the reset-password feature on the Apple
account. Apple's reset-password feature, in turn, required Honan's
address and the last four digits of the credit card number. The attacker
was able to find Honan's address publicly, but could not easily find his
credit-card number.

The attacker found this credit-card information through Honan's Amazon
account. Amazon, which knew his credit card number, required a full
credit-card in order to reset an account. However, Amazon allowed buying
something for a certain account without logging in, so long as you
provide a new credit card number. It also allowed \emph{saving} this new
credit card number to the user's Amazon account. So, the attacker made a
purchase on Amazon with the attacker's credit card. Then, the attacker
saved their own credit card number to Honan's Amazon account. Next, the
attacker triggered Amazon's reset-password feature, and---using the
credit-card number saved to the account---were able to reset Honan's
Amazon password and access his Amazon account. The attacker was then
able to see the last four digits of Honan's real credit card within his
Amazon account, use that to reset his Apple mail account, and then use
that to reset Honan's Gmail account.

Complex chains of systems like this can be very hard to reason about,
but these interactions ultimately are security-critical.

\subsection{Buggy implementations}\label{buggy-implementations}

Bugs, misconfigurations, and other mistakes are the most common cause of
security issues. A rule of thumb to keep in mind is that every 1000
lines of code will have around one bug. This is a very rough estimate,
but the basic idea is that more code will have more bugs. An effective
strategy to reduce security vulnerabilities is to reduce the amount of
code in your system.

\textbf{Missing Checks: iCloud.} Apple's iCloud performs many
functions---email, calendar, storage, and Find my iPhone. Each of these
had their own way of logging in, but across all of them a common goal
was to limit the attacker's ability to guess a user's password. To do
this, they added rate limiting to all the login interfaces, allowing
something like only 10 login attempts per hour---but they forgot the
Find my iPhone login interface. (\citeproc{ref-icloud}{Trew 2014})
Because this authentication code was duplicated all over the place,
there were many places to remember to add this rate limiting, but the
attacker only needed one weak login interface to brute-force a password.
In general, avoiding this repition will make it much easier to build a
secure system.

\textbf{Insecure Defaults.} When you set up a new service, they almost
always come with some defaults to make the setup simpler. Wi-Fi routers
come with default passwords, AWS S3 buckets come with default
permissions, and so on. These defaults can be convenient, but they are
very important to security because many people will forget or neglect to
change the default. Because of this, the default becomes the way that
the system operates. In order to build a secure system, it is important
that the default is secure.

\section{What are the general principles for secure system
design?}\label{what-are-the-general-principles-for-secure-system-design}

\subsection{Threat Models and Goals}\label{threat-models-and-goals}

\begin{itemize}
\tightlist
\item
  Create simple, general goals.
\item
  Avoid assumptions (such as ``no one else will be able to get a user's
  SIM card'') through better designs.
\item
  Learn and iterate.
\item
  Practice Defense in Depth: don't rely on one single defense for all
  your security---it is useful to use backup defenses to guard against
  bugs that will inevitably come up in one defense.
\end{itemize}

\subsection{Implementation}\label{implementation}

\begin{itemize}
\tightlist
\item
  A simpler system will lead to fewer problems.
\item
  Factor out the security-critical part into a small separate system or
  piece of code (for example, hardware security keys).
\item
  Reuse well-designed code, such as well-tested crypto libraries.
\item
  Understand and test the corner cases.
\end{itemize}

\part{Authentication}

\chapter{Authenticating People}\label{authenticating-people}

In this class, we will talk a lot about requests going to a computer
system. And a lot of security comes down to looking at that request and
deciding how to handle it. For this, it is crucial to know \emph{who}
issued the request. Then, we can decide whether the system should allow
the request.

Typically, a computer system performs two steps before processing a
request:

\begin{enumerate}
\def\labelenumi{\arabic{enumi}.}
\tightlist
\item
  \textbf{Authenticate:} Identify the person or machine (the
  ``\emph{principal}'') making the request.
\item
  \textbf{Authorize:} Decide if the principal is authorized to make the
  request.
\item
  \textbf{Audit:} Log some information about what requests your system
  authorized, so that you can identify malicious requests after the fact
  and/or clean up your system after attacks on the authentication
  system. (For example, a user who accidentally reveals their password
  to an attacker.)
\end{enumerate}

This chapter focuses on authentication. We will first discuss the attack
model and security goals. Then we will describe common implementations
that aim to achieve these goals.

\section{Authentication: Security
goals}\label{authentication-security-goals}

In the simplest model of authentication, we have a client and
server---two machines communicating over a network.

\marginnote{\begin{footnotesize}

Another common scenario has a human being authenticating to a computer:
think about typing a PIN into a phone or a password into a computer
terminal. We will discuss this setting more when we come to passwords.

\end{footnotesize}}

All authentication schemes try to prevent an attacker from impersonating
an honest user. To precisely define the security goal of an
authentication system though, we have to specify the attacker's power:
against which types of attack are we trying to defend?

\begin{figure}

\centering{

\includegraphics[width=0.9\linewidth,height=\textheight,keepaspectratio]{index_files/mediabag/figs/auth-attacks.pdf}

}

\caption{\label{fig-auth-attacks}There are (at least) three interesting
security goals for client-server authentication systems: (a) direct
attack, (b) eavesdropping attack, and (c) active attack.}

\end{figure}%

Figure~\ref{fig-auth-attacks} sketches three types of attacks on
authentication systems. In more detail these are, in increasing order of
strength are:

\begin{itemize}
\item
  \emph{Direct attack.} The attacker never sees the target user
  authenticate and then tries to impersonate the honest user. PIN-based
  authentication systems, e.g., on your phone or on an ATM machine,
  often aim to defend only against direct attacks. The screen-lock
  password that protects your laptop is also an authentication system
  that just attempts to protect against direct attack.

  That is, these systems do not protect against an attacker that can
  look over your shoulder while you are typing your password. These
  systems only aim to provide security when the attacker \emph{never}
  sees you (the honest user) authenticate.
\item
  \emph{Eavesdropping attack.} The attacker observes an honest user
  authenticating many times---i.e., the attacker sees all of the traffic
  between the client and server---and then tries to impersonate that
  user. One-time passwords, such as the six-digit authentication codes
  that the Google Authenticator app uses, aim to protect against
  eavesdropping attacks: an attacker who sees one of your one-time
  passwords will not be able to use it to authenticate as you; it is a
  \emph{one-time} password.
\item
  \emph{Active attack.} The attacker that compromises the server,
  interacts with the honest user, and after the server is restored to a
  good state, tries to authenticate as the honest user (\emph{active
  attack}).

  U2F security keys, and other schemes based on digital signatures
  (Chapter~\ref{sec-sig}), aim to protect against active attacks.
\end{itemize}

Systems that defend against active attacks provide the strongest form of
security, in that they also protect against eavesdropping attacks and
direct attacks. Systems that defend against eavesdropping also defend
against direct attacks. Systems that defend against direct attacks are
the weakest---they do not necessarily provide any protection against the
other types of attack.

\section{Protecting against direct attacks: Bearer tokens, PINs, and
passwords}\label{protecting-against-direct-attacks-bearer-tokens-pins-and-passwords}

\begin{figure}

\centering{

\includegraphics[width=0.7\linewidth,height=\textheight,keepaspectratio]{index_files/mediabag/figs/auth-direct.pdf}

}

\caption{\label{fig-auth-direct}A bearer-token-based authentication
scheme. To authenticate, the client sends its secret key \(k\) to the
server. {WARNING: In practice, the server should store the secret \(k\)
under a slow hash function (see Section~\ref{sec-auth-pwhash}).}}

\end{figure}%

The simplest form of authentication uses secret passwords or PINs. We
sometimes call these secret values \emph{bearer tokens}; whoever bears
(holds) a user's token can authenticate as that user. Authentication
with such schemes works as follows:

\begin{enumerate}
\def\labelenumi{\arabic{enumi}.}
\tightlist
\item
  The server holds a password (or PIN or random token string) for the
  client.
\item
  To authenticate, the client sends their password to the server.
\item
  The server checks the password against its stored password and accepts
  if they match.
\end{enumerate}

The benefit of password-based authentication schemes is that they are
simple and easy to implement. In addition, a human can play the role of
the client in a password-based authentication system (e.g., as you do
when you type a PIN into your phone). Fancier authentication systems
require the client to compute non-trivial cryptographic functions---not
functions that normal humans can compute in their brains.

Bearer-token-based schemes do provide security against direct attacks:
if an attacker has never seen the user authenticate, the attacker's best
strategy is to just guess the user's password. Thus, the security of
these schemes against direct attacks depends entirely on the adversary's
uncertainty about the password.

In some bearer-token-based systems, the server can assign a random
password to each user. For example, when you create an account for
certain web APIs, the API provider will give you a random secret key---a
bearer token. You will have to include this secret key with each API
request. Modern APIs use the stronger authentication mechanisms we
describe later in this chapter.

In the vast majority of password- and PIN-based login systems, the user
may choose their own password. This creates all sorts of
headaches\ldots{}

\subsection{What makes a good
password?}\label{what-makes-a-good-password}

\begin{figure}[H]

{\centering \includegraphics[width=1\linewidth,height=\textheight,keepaspectratio]{index_files/mediabag/figs/password-dist.pdf}

}

\caption{The most common passwords, according to NordPass
(\url{https://nordpass.com/most-common-passwords-list/}), sorted by
their frequency descending. A small number of common passwords
dominate.}

\end{figure}%

\marginnote{\begin{footnotesize}

\begin{longtable}[]{@{}cc@{}}
\caption{The most popular passwords in 2021, according to NordPass,
\url{https://nordpass.com/most-common-passwords-list/}.}\tabularnewline
\toprule\noalign{}
Rank & Password \\
\midrule\noalign{}
\endfirsthead
\toprule\noalign{}
Rank & Password \\
\midrule\noalign{}
\endhead
\bottomrule\noalign{}
\endlastfoot
1 & \texttt{123456} \\
2 & \texttt{123456789} \\
3 & 12345 \\
4 & qwerty \\
5 & password \\
6 & 12345678 \\
7 & 111111 \\
8 & 123123 \\
9 & 1234567890 \\
10 & 1234567 \\
11 & qwerty123 \\
12 & 000000 \\
13 & 1q2w3e \\
14 & aa12345678 \\
15 & abc123 \\
16 & password1 \\
17 & 1234 \\
18 & qwertyuiop \\
19 & 123321 \\
20 & password123 \\
\end{longtable}

\end{footnotesize}}

The security of a password-based authentication system rests entirely on
the attacker's inability to guess the password in a small number of
guesses.

\emph{Entropy} is a way to quantify an adversary's uncertainty about a
value sampled using a random process, or from a particular probability
distribution. If a distribution has \(b\) bits of entropy, then it will
take at least roughly \(2^b\) guesses for an attacker to correctly guess
a value sampled from this distribution.

The uniform distribution over 128-bit strings has 128 bits of entropy.
The distribution from which humans typically choose their passwords has
much less entropy---empirically, more like 20 bits.

Ideally, we would want all passwords to be equally as likely, from the
adversary's perspective. (These would be ``high-entropy'' passwords.) If
a system generated a truly random password for each user, each password
would be indeed equally likely. But then it becomes very difficult for a
user to remember their password, much less many random passwords for
many different services.

Typically, people have to remember their passwords, so we let them pick
their own passwords. When they do, it turns out that many people are
likely to choose the same password.

A typical password might be sampled from a distribution with roughly 20
bits of entropy---if an adversary is able to make \(2^{20}\) guesses at
the password, they can expect to guess the password correctly. One can
computer easily make \(2^{20}\) authentication requests to another in a
few minutes.

\subsection{Dealing with weak
passwords}\label{dealing-with-weak-passwords}

\emph{Every} system that uses password- or PIN-based authentication must
contend with the fact that most passwords are not that difficult for an
attacker to guess.

In this section, we describe some mitigation strategies: all are flawed,
but each is better than nothing. The goal of each is to make the
attacker's job slightly more difficult; by stacking a few of these
defenses on top of each other, we can substantially strengthen the
end-to-end system.

\textbf{Aggressively limit the number of guesses.} Therefore, when using
passwords as an authentication mechanism, an authentication system must
\emph{always} somehow limit the number of password guesses.

For example, some phones allow 10 guesses at the screen-lock PIN before
the device resets itself. Limiting the number of guesses effectively
prevents a \emph{single} account from being compromised---provided that
the password is not too too weak. One downside is guess limits create
the possibility for denial-of-service attacks: an attacker can
potentially make 10 guesses at your password and lock you out of your
phone or online-banking account.

In addition, in many physical computer systems have multiple authorized
users, each with their own password. If the guess limit is enforced only
on a per-user basis, then an attacker can often compromise \emph{some}
account on the machine if it is allowed 10 guesses at \emph{every}
user's password. Preventing these types of attacks requires some
additional measures: websites that use password authentication
rate-limit guesses by IP, or use CAPTCHAs, etc.

\textbf{Try to coerce users into picking stronger passwords.} Modern
websites will often provide the user with a ``password-strength
checker'' that tries to give the user some sense of how strong or weak
their password is. These strength meters are completely heuristic and
can be wildly wrong: they might say that \texttt{6175551212} is a great
password; if the attacker knows that 617-555-1212 is my phone number, it
is probably not a great password. These strength meters sometimes check
a user's password against public lists of popular passwords. Ensuring
that your password isn't in the million most popular ones gives you at
least some protection against untargeted attacks.

Two common strategies for encouraging users to choose strong passwords
that don't work terribly well are:

\begin{itemize}
\tightlist
\item
  \emph{Require longer passwords.} If someone tries to use
  \texttt{abc123} as a password but it's not long enough, they might use
  \texttt{abc123456}---but this doesn't really add much uncertainty.
  There are standard ways to lengthen passwords, and a clever attacker
  will try these first.
\item
  \emph{Prohibit using common English words in passwords.} It's not
  clear that this is a good idea. Five randomly chosen words from the
  dictionary will form a strong password, and prohibiting English words
  in passwords may make passwords much more difficult to remember.
\end{itemize}

\subsection{Avoiding weak passwords with a password
manager}\label{avoiding-weak-passwords-with-a-password-manager}

When using passwords to authenticate to a website, a user can install a
password manager on their computer that will generate random passwords
for them. Since the user doesn't need to remember these passwords, they
can be sampled truly at random from a high-entropy distribution. Once
the user authenticates to their computer (using a password, typically),
they can then access their randomly generated passwords and use them to
log in to their websites.

Internally, the password-manager software maintains a table of servers
and the corresponding passwords:

{\def\LTcaptype{none} % do not increment counter
\begin{longtable}[]{@{}ccc@{}}
\toprule\noalign{}
server & user & pw \\
\midrule\noalign{}
\endhead
\bottomrule\noalign{}
\endlastfoot
\texttt{amazon.com} & \texttt{alice} & \texttt{3xyt42...} \\
\texttt{mit.edu} & \texttt{alice4} & \texttt{a21\$z...} \\
\end{longtable}
}

Even when using a password manager, password-based authentication
schemes provide \emph{no security} against eavesdropping or active
attacks. If an attacker can observe you sending your password to the
server (e.g., with a phishing attack) it can still authenticate as you.

\subsection{Password hashing: Trying to get some protection against
server compromise}\label{sec-auth-pwhash}

Password-based authentication schemes provide no security against active
attacks, in which the attacker compromises the server. And yet, since
attackers manage to breach web servers quite often, we would really like
to provide some defense against server compromise.

\marginnote{\begin{footnotesize}

\emph{Forcing password changes.} A system may force their users to
change passwords on a regular schedule (e.g., every six months). If an
attacker has compromised the password database on the server, it only
has a limited amount of time to access the system before the passwords
change and it will get locked out. It is not clear that the cost of
requiring frequent password changes is worth the benefit.

\end{footnotesize}}

Since, as we have seen, passwords are easy to guess, avoiding
password-based authentication entirely is the safest option where
possible. When a system must use passwords for authentication, the
safest way to store them (e.g., on a server) is using a \emph{salted
cryptographic password-hashing function}. The goal is to make it as
difficult as possible for an attacker to recover the plaintext
passwords, given the hashed values stored on the server.

To describe how this works: when a user creates an account with password
pw, the server chooses a random 128-bit string, called a \emph{salt},
and the server stores the salt and the hash value
\(h=H(\text{salt}\|\text{pw})\), where \(H\) is a special
password-hashing function.

The server then stores a table that looks like this:

\marginnote{\begin{footnotesize}

A \emph{rainbow table} is a common data structure that an attacker can
use to invert unsalted password hashes. A rainbow table is essentially a
compressed table \((\mathtt{passwd}, H(\mathtt{passwd}))\) pairs, where
\(H\) is a common hash function, and \(\mathtt{passwd}\) ranges over a
large set of common passwords.

An attacker can download rainbow tables for common cryptographic hash
functions, such as MD5 or SHA1, from the web. If the password hash is
salted with a 128-bit salt, it will be infeasible to produce a table
that covers any reasonably large fraction of the
\((\mathtt{salt}\|\mathtt{passwd})\) pairs.

\end{footnotesize}}

{\def\LTcaptype{none} % do not increment counter
\begin{longtable}[]{@{}ccc@{}}
\toprule\noalign{}
user & salt & \(H(\text{salt}\|\text{pw})\) \\
\midrule\noalign{}
\endhead
\bottomrule\noalign{}
\endlastfoot
alice & \(r_a\) & \(h_a\) \\
bob & \(r_b\) & \(h_b\) \\
\end{longtable}
}

Later on, when the user sends a password \(\text{pw}'\) to the server to
authenticate, the server can use the salt and hash function to compute a
value \(h' = H(\text{salt}\|\text{pw}')\). If this hash value \(h'\)
matches the server's stored value \(h\) for this user, the server
accepts the password.

To explain the rationale for this design:

\begin{itemize}
\tightlist
\item
  The password-hashing function \(H\) is designed to be relatively
  expensive to compute---possibly using a large amount of RAM and taking
  a second or more of computation. This makes it more difficult for an
  attacker to brute-force invert the hash value, since each guess at the
  password requires a second of computation (instead of the microseconds
  required to compute a standard hash function, such as SHA256).
\item
  The use of a per-user random salt ensures that guesses at one user's
  password are useless in inverting another user's password hash.
  Salting also defeats \emph{precomputation attacks}, in which an
  attacker precomputes the hashes of many common passwords to speed up
  this hash-inversion step later on.
\end{itemize}

\subsection{Biometrics}\label{biometrics}

Biometrics are physical features like your fingerprints, your face, etc.
These are essentially a type of bearer token: whoever is able to produce
a face that looks like yours is able to authenticate as you.

Biometrics are very convenient to use for authentication, since you will
not forget them and cannot easily lose them. Biometrics most useful when
authenticating in person to a device, such as for phone unlock, or to
grant a person access to a secure vault. In these settings, the device
performing the authentication has a ``trusted input path'' that can
provide some assurance that a real human who owns that biometric is on
the other end. Biometrics are not so useful for authenticating over a
network because the network typically does not provide a trusted input
path (i.e., does not provide any assurance that the biometric readings
are coming from a real human), and the biometric data itself is not
particularly secret. In particular, if we used biometrics for network
authentication, an adversary who knows what your fingerprints looks like
could log in to your account. (Since biometrics are essentially
impossible to change, this is a major drawback.)

\section{Protecting against eavesdropping attacks: Challenge-response
protocols}\label{protecting-against-eavesdropping-attacks-challenge-response-protocols}

We have so far been talking about a human manually authenticating to a
device (ATM, phone, laptop, etc.) by physically entering a PIN or
password into the device. But we often log in to some server on the
network---Facebook, Gmail, MIT, and so on. In this scenario, we can get
much more creative with the authentication mechanism we use and the
security properties we can demand.

We now assume that our computer has some key \(k\) (e.g., a random
128-bit string), and the server also holds the same key \(k\). In this
setting, we can hope to provide security against eavesdropping attacks:
even if an attacker can observe the traffic between the client and
server, the attacker learns no information that can help it authenticate
as the client later on.

\begin{figure}

\centering{

\includegraphics[width=0.8\linewidth,height=\textheight,keepaspectratio]{index_files/mediabag/figs/auth-eavesdropping.pdf}

}

\caption{\label{fig-auth-eaves}A challenge-response-based authentication
system. The server and client share a secret key \(k\). To authenticate,
the client sends a function \(F\) of its secret key and the server's
challenge \(c\).}

\end{figure}%

Figure~\ref{fig-auth-eaves} describes an über-simplified
challenge-response authentication scheme that provides security against
eavesdropping attacks. The protocol takes place between a client and
server holding a shared secret key \(k\).

\begin{enumerate}
\def\labelenumi{\arabic{enumi}.}
\tightlist
\item
  The server chooses a long random string \(c\), which we call a
  \emph{challenge} and sends it to the authenticating client.
\item
  The client computes an authentication ``tag'' \(t \gets F(k, c)\),
  where \(F(\cdot, c)\) is hard to compute without knowing \(k\). (The
  function \(F\) here is a Message Authentication Code, which we will
  talk more specifically about in Chapter~\ref{sec-mac}.) The sends the
  MAC tag \(t\) to the server.
\item
  The server receives a tag \(t'\) from the client and ensures that
  \(t' = F(k,c)\). If so, the server considers the authentication
  successful.
\end{enumerate}

The security of this scheme derives from the fact that an attacker
cannot produce tags \(F(k,c)\) on new challenge values \(c\). (An
attacker can always try to replay an old tag it has seen, but since the
challenge changes with every authentication request, the server will
always reject the old tag.)

\subsection{Time-based One-Time Passwords
(TOTP)}\label{time-based-one-time-passwords-totp}

Time-based one-time passwords are a type of challenge-response
authentication protocol. The only difference from
Figure~\ref{fig-auth-eaves} is that in a TOTP scheme, the client and
server derive the challenge from the current time. The user has a
device, such as a phone, that shares a secret key \(k\) (e.g., a random
128-bit string) with the server. Both parties agree on a protocol by
which to generate this code---something like
\(F(k, \mathtt{gettimeofday()\ /\ 30})\). The phone can generate the
code, display it to the user, and the server can then verify the code by
recomputing it.

\subsection{Authenticating requests}\label{authenticating-requests}

Often, a client will want to send an \emph{authenticated message} to a
server. That is, the client often wants to simultaneously authenticate
to the server and send a request \(\mathsf{req}\), such as
\(\mathsf{req} = \mathtt{rm\ file.txt}\). To accomplish this, the client
can compute the challenge value as the hash of the server-provided
challenge \(c\) and the client's request. So the tag looks like:
\(t_{\mathsf{req}} \gets F(k, c \| \mathsf{req})\). Then the client
sends the pair \((t_{\mathsf{req}}, \mathsf{req})\) to the server. In
this way, the server can simultaneously authenticate the client and be
sure that the request \(\mathsf{req}\) came from the client.

An \textbf{unsafe} way for the client to simultaneously authenticate to
the server and send a request would be for the client to compute the MAC
tag \(t \gets \MAC(k,r)\) and then send \((t, \mathsf{req})\) to the
server. A network attacker could modify the client's request to
\((t, \mathsf{req'})\) en route to the server without the server being
able to detect this attack.

\subsection{Phishing attacks (attacker-in-the-middle
attacks)}\label{phishing-attacks-attacker-in-the-middle-attacks}

A \emph{phishing} attack is one in which an attacker tricks a user into
giving away their Gmail password, for example, by creating a website
that looks, for example, like the \texttt{gmail.com} login page. TOTP
passwords have a similar vulnerability: an attacker can simply ask the
user to give her the one-time code by pretend to be tech support, or the
user's employer, or a customer-service representative. In this setting,
TOTP codes are slightly better than standard passwords since the
attacker must use a stolen TOTP code within \$\approx\$30 seconds of
stealing it, which requires a much more sophisticated attack.

Phishing attacks take advantage of the fact that in password- and
TOTP-based authentication schemes, there is no binding between the
authentication process and the server's subsequent communication with
the client. The attacker here doesn't really break the authentication
scheme; the problem is that the authentication scheme didn't
authenticate enough. U2F, which we now discuss, handles that issue.

\section{Protecting against active attacks: Signatures and
U2F}\label{protecting-against-active-attacks-signatures-and-u2f}

\begin{figure}

\centering{

\includegraphics[width=0.7\linewidth,height=\textheight,keepaspectratio]{index_files/mediabag/figs/auth-active.pdf}

}

\caption{\label{fig-auth-active}An authentication scheme based on
digital signatures.}

\end{figure}%

To provide security against active attacks, we can use an authentication
scheme based on digital signatures, which we will discuss in
Chapter~\ref{sec-sig}. With these schemes, the client has a secret key
\(k\) and the server stores some hard-to-invert function of the key
\(F(k)\). (We call \(F(k)\) the ``public key.'') In particular, the
server \emph{does not store any secrets}---even if the attacker can
compromise the server and/or interact with the honest client, it cannot
learn the client's secret key \(k\) nor learn any information that it
can use to later authenticate as the client. To authenticate, the server
sends the client a challenge and the client produces a digital signature
on the challenge \(c\)---essentially a proof that the client knows the
secret key \(k\) and that it intended to sign the challenge \(c\).

The U2F USB security tokens that you may have seen use this form of
authentication. As an added bonus, they prevent phishing attacks by
binding the authentication process to the name of the server that the
client is trying to authenticate to. In particular, the U2F software on
the client passes the name of the server (e.g., \texttt{amazon.com}), in
addition to a server-provided random challenge \(c\), to the U2F token.
The token then produces a signature on the string
\(c\|\mathtt{amazon.com}\). If the attacker sets up \texttt{amason.com}
and gets the user to visit it, the U2F device will only generate a code
that is good for \texttt{amason.com} and not the real
\texttt{amazon.com}.

\section{Two-Factor Authentication}\label{two-factor-authentication}

Many systems use multiple forms of authentication to try to boost
security. In particular, as we have already seen, passwords are a weak
authentication mechanism: humans are bad at choosing strong passwords
and attackers have become good at stealing password databases and
recovering many users' passwords at once.

A common technique to harden password-based authentication systems is to
combine passwords with a second method of authentication---one with a
different failure mode. Common authentication schemes are:

\begin{itemize}
\tightlist
\item
  Something you know: password, PIN, etc
\item
  Something you have: USB key, phone, etc
\item
  Something you are: biometrics (fingerprint, face ID)\ldots{}
\end{itemize}

\chapter{Collision Resistance and File
Authentication}\label{collision-resistance-and-file-authentication}

In the last chapter, we focused on authenticating people---ensuring that
a person (or a request on behalf of that person) is likely who they
claim to be. In this chapter, we will focus on authenticating files,
code, and other data. When we say that we want to authenticate a file,
we mean that we want to verify that the file's contents are exactly as
they were when we or someone we trust last viewed them. The key new tool
we use to do so is a \emph{collision-resistant hash function}.

\section{Intuition: Collision
resistance}\label{intuition-collision-resistance}

For our purposes, a hash function \(H\) maps a bitstring of any length
onto a fixed-size space of outputs, so the type signature is
\(H: \bin^* \rightarrow \bin^{\lambda}\).

In order for a hash function to be \emph{collision-resistant}, we want
it to be the case that for any input, the generated output should be
``unique.'' Of course, it cannot really be unique---we are mapping
infinitely many inputs onto finitely many outputs---but we want it to be
\emph{computationally infeasible} to find a pair of distinct inputs that
have the same hash values (a ``collision'').

\textbf{Security goal.} A hash function \(H\) is \emph{collision
resistant} if it is ``computationally infeasible'' to find two distinct
strings \(x\) and \(x'\) such that \(H(x) = H(x')\).

Given a long message \(m\), it's hash \(H(m)\) under a
collision-resistant hash function is like a short ``fingerprint'' of the
message---the hash essentially uniquely identifies the message \(m\).
For that reason, collision-resistant hash functions let you authenticate
a long message \(m\) by authenticating the short fixed-length string
\(H(m)\). We often call the hash value \(H(m)\) a \emph{digest}.

\subsection{Applications}\label{applications}

\textbf{Secure File Mirroring.} Often a user wants to download large
files (e.g., software updates) from a far-away server. To speed up this
process, a company or Internet-service provider may set up local
\emph{mirrors} of the remote files. Users can then download the files
from the nearby mirror instead of the far-away server. However, without
additional security measures, the mirror may server users a different
file than the one the mirror fetched from the origin server. If the
mirror is malicious, it can, for example, trick the user into installing
a backdoored software update. (We saw an attack based on mirrors in
Section~\ref{sec-intro-xcode}.)

To protect against a malicious mirror, we can add some authentication on
the file that the mirror hosts. Say that the origin server publishes a
large software update \(f\). The origin server will send the file \(f\)
to its mirrors and the origin server itself will serve the hash digest
\(d\gets H(f)\) to anyone who asks for it. A user who wants to fetch the
update can download \(d\) from the origin server directly---this will be
fast since the digest is tiny. Then, the client can fetch the update
itself from a (potentially untrustworthy) mirror. When the client
receives a file \(\hat f\) from the mirror, it can check that
\(d = H(\hat f)\) to ensure that \(\hat f\) is the true software update.
If \(H\) is collision resistant, then if the hash value \(H(\hat f)\)
matches the origin server's digest \(d\), the files are almost certainly
identical.

\textbf{Subresource Integrity.} If a program fetches a file from some
content delivery network, it can store the hash of that file locally and
use it to verify that the contents of the file did not change since the
application was developed.

\textbf{Outsourced File Storage.} If you want to store your files on a
cloud provider, you want to be sure that the cloud provider does not
maliciously modify the files without you noticing. To make sure of this,
you can store \(H(\text{files})\) locally, which takes very little
storage space. Then, when you redownload your files locally, you can
recompute the hash to verify that they were not tampered with.

\section{Defining collision resistance (slightly more
formally)}\label{defining-collision-resistance-slightly-more-formally}

An adversary's goal in breaking a collision resistant hash function is
to find a collision---a pair of values \(m_0, m_1 \in \zo^*\) such that
\(m_0 \neq m_1\) and \(H(m_0) = H(m_1)\).

\begin{definition}[Collision
Resistance]\protect\hypertarget{def-crhf}{}\label{def-crhf}

A function \(H: \bin^* \rightarrow \bin^{\lambda}\) is
collision-resistant if for all ``efficient'' adversaries \(\calA\), we
have that:
\[ \Pr[H(m_0) = H(m_1), m_0 \neq m_1: (m_0, m_1) \leftarrow \calA()] \leq \text{``negligible''} \]

\end{definition}

In words, this means that the probability of finding a collision is so
small that no efficient adversary could hope to do it.

There are two ways of thinking about the terms ``efficient'' and
``negligible'' that we use in this definition---one mindset we use in
practice and the other mindset we use in theory.

\begin{itemize}
\tightlist
\item
  In theory\(\ldots\)

  \begin{itemize}
  \tightlist
  \item
    All of our cryptographic constructions are parameterized by an
    integer \(\lambda \in \{1, 2, 3, \dots\}\) that we call the
    \emph{security parameter}. So instead of a single
    collision-resistant hash function \(H\), we have a family of
    functions \(\{ H_1, H_2, H_3, \dots \}\), where the function
    \(H_\lambda\) has \(\lambda\)-bit output.
  \item
    An ``efficient'' algorithm is a randomized algorithm that runs in
    time polynomial in \(\lambda\).
  \item
    A ``negligible'' function is one that grows slower than the inverse
    of every polynomial---a function that is \(O(\frac{1}{\lambda^c})\)
    for all constants \(c \in \mathbb{N}\).
  \end{itemize}
\item
  In practice\(\ldots\)

  \begin{itemize}
  \tightlist
  \item
    We use a fixed hash function \(H\) with a fixed-length output, which
    might be as 256 or 512 bits.
  \item
    An ``efficient'' adversary is one that runs in time
    \(\leq 2^{128}\).
  \item
    A ``negligible'' probability is some very small constant, like one
    smaller than \(2^{-128}\).
  \end{itemize}
\end{itemize}

\subsection{Understanding which attacks are
feasible}\label{understanding-which-attacks-are-feasible}

Typically, we think of an attack that runs in more than \(2^{128}\) time
as infeasible and an event that happens with probability less than
\(2^{-128}\) is one that will never happen. These seemingly magic
constants come from empirical considerations:

{\def\LTcaptype{none} % do not increment counter
\begin{longtable}[]{@{}
  >{\raggedleft\arraybackslash}p{(\linewidth - 2\tabcolsep) * \real{0.1622}}
  >{\raggedright\arraybackslash}p{(\linewidth - 2\tabcolsep) * \real{0.8378}}@{}}
\toprule\noalign{}
\endhead
\bottomrule\noalign{}
\endlastfoot
\(2^{30}\) & operations/second on a laptop \\
\(2^{58}\) & ops/sec on Fugaku supercomputer \\
\(2^{68}\) & hashes/second on the Bitcoin network (as of Fall 2022) \\
\(2^{92}\) & hashes/yr on the Bitcoin network \\
\(2^{114}\) & hashes required to use enough energy to boil the ocean \\
\(2^{140}\) & hashes required to use one year of the sun's energy \\
\end{longtable}
}

See Lenstra, Kleinjung, and Thomé
(\citeproc{ref-lenstra:universal}{Lenstra et al. 2013}) for an
entertaining discussion of these constants.

{\def\LTcaptype{none} % do not increment counter
\begin{longtable}[]{@{}
  >{\raggedleft\arraybackslash}p{(\linewidth - 2\tabcolsep) * \real{0.1622}}
  >{\raggedright\arraybackslash}p{(\linewidth - 2\tabcolsep) * \real{0.8378}}@{}}
\toprule\noalign{}
\endhead
\bottomrule\noalign{}
\endlastfoot
\(2^{-1}\) & fair coin lands heads \\
\(2^{-13}\) & probability that a randomly sampled MIT grad is a Nobel
Prize winner \\
\(2^{-19}\) & probability of being struck by lightning next year \\
\(2^{-28}\) & probability of winning the Mega Millions jackpot \\
\(2^{-128}\) & will essentially never happen \\
\end{longtable}
}

\marginnote{\begin{footnotesize}

For most cryptosystems, there is a tradeoff between the attacker's
running time and success probability. For example, an attacker running
in time \(T\) can find a collision in a hash function with \(n\)-bit
output with probability \(T^2/2^n\). So, as the attack runs for more
time, it has a better chance of finding a collision.

\end{footnotesize}}

The takeaway is that if an attacker finds a collision with probability
\(2^{-128}\), we can be extremely sure that a collision will never
occur.

\section{Constructing a collision-resistant hash
function}\label{constructing-a-collision-resistant-hash-function}

The current standard for fast collision-resistant hashing is SHA256
(a.k.a. SHA2), which was designed by the NSA in 2001. The SHA2 hash
functions are designed using the following common two-step approach:

\marginnote{\begin{footnotesize}

We can also build collision-resistant hash functions that are secure
under ``nice'' cryptographic assumptions, such as the assumption that
factoring large numbers is hard. Unfortunately, hash functions based on
these nice assumptions tend to be very slow and, as a result, are almost
never used in practice.

\end{footnotesize}}

\begin{enumerate}
\def\labelenumi{\arabic{enumi}.}
\item
  Build a small collision-resistant hash function on a fixed-size domain
  \(H_{small}: \bin^{2\lambda} \rightarrow \bin^\lambda\). This step is,
  to some degree, ``more art than science.'' The standard practice is to
  design a hash function that defeats all known collision-finding
  attacks. If no known attack works well, we declare the candidate
  function to be collision resistant.
\item
  Use \(H_{\text{small}}\) to construct
  \(H: \bin^* \rightarrow \bin^\lambda\). This can be done very cleanly
  using the ``Merkle-Damgård'' approach described below. This step
  requires no additional assumptions: we can prove unconditionally that
  if \(H_{\text{small}}\) is collision resistant, then \(H\) is as well.
\end{enumerate}

\marginnote{\begin{footnotesize}

Another way to build collision-resistant hash functions is to use the
so-called ``sponge'' construction. It is similar to the approach
described here in that we start with a small primitive, which we assume
secure in some sense, and then we use the small primitive to build a
hash function on a large domain.

\end{footnotesize}}

\subsection{Merkle-Damgård}\label{merkle-damguxe5rd}

The Merkle-Damgård construction gives a way to construct a
collision-resistant hash function for all bitstrings (i.e., \(\zo^*\))
from a collision resistant hash function that maps \(2\lambda\)-bit
strings down to \(\lambda\)-bit strings, sketched out in
Figure~\ref{fig-merkle-damgard}.

\begin{figure}

\centering{

\pandocbounded{\includegraphics[keepaspectratio]{index_files/mediabag/figs/merkle-damgard.pdf}}

}

\caption{\label{fig-merkle-damgard}Sketch of the Merkle-Damgård
construction for a collision-resistant hash function.}

\end{figure}%

The Merkle-Damgård construction first splits the message into
\(\lambda\)-sized blocks \([m_1, \ldots, m_n]\) and successively hashes
them together. In the following pseudocode, the function
\(\mathsf{ToBlock}\) converts an integer, representing the length of the
input message in blocks, into a \(\lambda\)-bit string. Then the
Merkle-Damgård construction is:

\begin{figure}

\centering{

\(H(m_1, \dots, m_n)\): // Merkle-Damgård construction

\begin{itemize}
\tightlist
\item
  Let \(b \gets 0^\lambda\).
\item
  For \(i = 1, \dots, n\):

  \begin{itemize}
  \tightlist
  \item
    Let \(b \gets H_\text{small}(b, m_i)\).
  \end{itemize}
\item
  Let \(b \gets H_\text{small}(b, \mathsf{ToBlock}(n))\).
\item
  Output \(b\).
\end{itemize}

}

\caption{\label{fig-md}The Merkle-Damgård construction of a large-domain
collision-resistant hash function \(H\) from a small-domain
collision-resistant hash function \(H_\text{small}\).}

\end{figure}%

(Here, we are assuming that the message is at most
\(\lambda 2^{\lambda}\) bits long.)

\marginnote{\begin{footnotesize}

In practice, standard hash functions have limits on the length of the
messages that they can hash. For example, SHA256 can hash messages of
length up to \(2^{64}-1\) bits.

\end{footnotesize}}

We won't prove it here, but we can use the fact that \(H_{small}\) is
collision-resistant to prove that \(H\) must also be
collision-resistant. The basic idea of the proof is to show that given a
collision in \(H\), we can easily compute a collision in
\(H_{\text{small}}\).

\textbf{Note.} In the Merkle-Damgård construction of
Figure~\ref{fig-md}, we initialize the variable \(b\) to the all-zeros
string. The construction is collision-resistant if we omit the all-zeros
string and start by setting \(b \gets m_1\) and then continue by hashing
\(m_2, m_3, \dots\). The construction is \emph{not} collision resistant
if we omit the length block \(\mathsf{ToBlock}(n)\) that we hash in at
the end.

\subsection{The Birthday Paradox}\label{the-birthday-paradox}

\marginnote{\begin{footnotesize}

If you sample \(2^{\lambda/2}\) random \(10\lambda\)-bit strings and
hash them with a hash function that has \(\lambda\)-bit outputs, you
will find a collision among these inputs with constant probability.

\end{footnotesize}}

An important thing to understand when dealing with hash functions is the
``Birthday Paradox,'' which states that given a hash function with
\(\lambda\)-bit output, you can always find a collision in time
\(O(\sqrt{2^\lambda}) = O(2^{\lambda/2})\). So, if you want to force an
attacker to use at least \(2^{128}\) to find a collision, you must use a
hash function with at least 256 bits of output.

\subsection{Domain Separation}\label{domain-separation}

In many applications, we have a one-input CRHF (such as SHA256)
\(H: \bin^* \rightarrow \bin^{256}\) and we need to construct a
two-input CRHF \(H_2(x, y)\).

\textbf{Bad idea.} An obvious solution to construct the two-input hash
function \(H_2\) is to concatenate the two values, so that
\(H_2(x, y) = H(x || y)\). However, this construction allows two
different pairs of messages to hash to the same value:
\[H_2(\text{"key", "value"}) = H_2(\text{"ke", "yvalue"}).\] Both Amazon
and Flickr had a bug arising from this---they concatenated all
parameters before hashing, and had parameters such that two different
intents had the same concatenation (\citeproc{ref-flickr}{Duong and
Rizzo 2009}).

\subsection{Length-Extension}\label{length-extension}

Recall the concept of Message Authentication Codes (MAC) from the last
lecture---a code that can be sent along with a message to verify that
the message was not changed. (We will see the formal definition in
Chapter~\ref{sec-mac}.)

\textbf{Bad idea.} Poorly designed software uses
\(\MAC(k, m) = H(k \| m)\) as a very simple construction of a MAC.
However, this construction has an easy attack---given \(\MAC(k, m)\), it
is easy to compute \(\MAC(k, m\|m')\) without knowing the key \(k\) if
\(H\) is a hash function built with the Merkle-Damgård construction. To
do this, the attacker hashes the output of \(\MAC(k, m)\) with two more
blocks---a new message \(m''\) and another length block. Now, we have
computed \(\MAC(k, m\|m')\) where \(m'\) is the original length block
plus some custom message without knowing the key \(k\).

This problem here is that we were using a hash function that was
\emph{only} guaranteed to be collision resistant, but we assumed that it
had other properties (such as that it is guaranteed to be difficult to
compute the hash of an extension of the original message).
Figure~\ref{fig-merkle-damgard-extension} sketches out the
length-extension attack on the Merkle-Damgård construction from
Figure~\ref{fig-merkle-damgard}.

\begin{figure}

\centering{

\pandocbounded{\includegraphics[keepaspectratio]{index_files/mediabag/figs/merkle-damgard-extension.pdf}}

}

\caption{\label{fig-merkle-damgard-extension}Sketch of the extension
attack on the Merkle-Damgård construction, starting with a hash of
\(m_1 || m_2 || \ldots || m_n\) to compute the hash of
\(m_1 || m_2 || \ldots || m_n || \textsf{ToBlock}(n) || x\).}

\end{figure}%

\section{Applications: Merkle Trees}\label{applications-merkle-trees}

In many settings, an origin server has \(N\) files (e.g., Android app
binaries) and wants to serve these files from potentially untrustworthy
mirror servers (e.g., Akamai servers) distributed around the globe.

To do this, the origin server can put the \(N\) files at the leaves of a
binary tree. Then the server hashes together pairs of files, then hashes
each pair of hashes and so on until it eventually ends up with a single
root hash \(h\). The client fetches the root hash \(h\) from the origin
server directly.

Later on, the client can download any one of the \(N\) files from the
untrustworthy mirror server. The mirror can produce the file, along with
\(O(\log N)\) hashes---the sibling nodes of each node on every path from
the file's leaf to the root. The client can use the root hash \(h\) it
got from the origin server, along with the additional hashes from the
mirror server, to be convinced that the mirrored file it downloaded was
authentic.

\todobox{\textbf{TODO:} Add diagram from lecture.}

\chapter{Message Authentication Codes}\label{sec-mac}

So far, we have talked about authenticating \emph{people} and
authenticating \emph{files}. In this section, we will discuss
authenticating \emph{communication}. If we have two parties that are
communicating over the network, we want some way to guarantee to each
party that the message they received really came from the other party
and was not tampered with along the way.

At a first glance, this seems impossible. If there is some eavesdropper
Eve in between the two parties, they can just replace the message with
one of their own choosing and the other party will have no idea. To make
this possible, we need to relax the scenario a bit and add an
assumption---that the two parties share some secret key \(k\).

With this shared key \(k\) between the two parties, our goal will be to
add some ``tag'' onto the message that validates its authenticity.
Necessarily, this tag will be a function of this shared key \(k\). If
this were not the case, the eavesdropper would be able to compute a
valid tag herself---the secret \(k\) is the only information in this
scenario that Eve does not know.

\section{Defining message authentication
codes}\label{defining-message-authentication-codes}

\textbf{Syntax.} A message authentication code (MAC) over key space
\(\mathcal{K}\), message space \(\mathcal{M}\) , and tag space \(\calT\)
is an efficient algorithm \(\MAC \colon \calK \times \calM \to \calT\).
In order for a MAC to be useful, it must be \emph{secure}, in the
following sense. We first give the definition and then explain why it is
a useful one:

\begin{definition}[MAC Security: Existentially unforgeability against
adaptive chosen message
attacks]\protect\hypertarget{def-mac-sec}{}\label{def-mac-sec}

A MAC \(\MAC\) over key space \(\calK\) and message space
\(\mathcal{M}\) is secure (existentially unforgeable against adaptive
chosen message attacks) if any poly-time adversary \(\mathcal{A}\) wins
the following game with at most negligible probability:

\marginnote{\begin{footnotesize}

In practice, ``a poly-time adversary'' means ``any real-life
adversary''. But we need to place some mathematical bound on real-life
to make the proofs work out.

\end{footnotesize}}

\begin{itemize}
\tightlist
\item
  The challenger samples a MAC key \(k \getsr \calK\).
\item
  For \(i = 1, 2, \dots\) (polynomially many times)

  \begin{itemize}
  \tightlist
  \item
    The adversary sends any message \(m_i \in \calM\) to the challenger
  \item
    The challenger responds with \(\mathsf{MAC}(k, m_i)\).
  \end{itemize}
\item
  The adversary sends the challenger a message-tag pair \((m^*, t^*)\).
\item
  The adversary wins the game if \(\mathsf{MAC}(k, m^*) = t^*\) and
  \(m^* \notin \set{m_1, m_2, \ldots, m_n}\).
\end{itemize}

\marginnote{\begin{footnotesize}

A subtlety of this definition is that, even if the MAC scheme is secure
under this definition, it is possible for an adversary, given a valid
message-tag pair \((m,t)\) to produce a second valid message-tag pair
\((m,t')\) on the same message without knowing the secret key.

\end{footnotesize}}

\end{definition}

\subsection{Intuition for the security
definition}\label{intuition-for-the-security-definition}

To formulate our security notion, we need to define the adversary's goal
and the adversary's power.

The adversary's goal in this definition is to compute a valid MAC of
\emph{any} message \(m\in \mathcal{M}\) of its choice. It's not entirely
obvious why we care about the adversary producing a valid MAC on
\emph{any} message: ``If the adversary MACs a message that is jibberish,
they are unlikely to be able to do any harm with it,'' you might think.
But there will certainly be applications that authenticate messages that
violate whatever definition of ``non-jibberish'' we define. So allowing
the adversary to forge a MAC tag on any message makes the definition as
broadly applicable as possible.

As far as the adversary's power goes: we, as usual in cryptography,
restrict the adversary to be efficient (i.e., to run in polynomial
time). But in the MAC security game we also allow the adversary to
obtain MAC tags on messages of its choice. This captures the reality
that in many systems, an adversary can trick an honest system into
MACing adversarial messages. For example, if an email-backup system MACs
every email that a user receives, an adversary may be able to obtain MAC
tags on messages of its choice by sending emails to the backup system.

\marginnote{\begin{footnotesize}

This has some interesting implications---importantly, the adversary can
store these messages along with their MAC and replay them later.

\end{footnotesize}}

\subsection{MACs require
pseudorandomness}\label{macs-require-pseudorandomness}

The fact that it is even possible to construct a MAC seems a bit
surprising---in effect, for a MAC to satisfy the definition, the tag has
to effectively be random. But the only ``randomness'' that we have is
the key \(k\)---to generate tags for arbitrarily many messages, we need
much more randomness than one key's worth. This seems impossible. How
can we generate a large number of random-looking tags from only a single
short random key?

We get ourselves out of this conundrum by observing that the adversary
must be an \emph{efficient} algorithm. So while we cannot generate a
large number of truly random bits from a short key, we can---under
appropriate and reasonable cryptographic assumptions---generate a large
number of bits that \emph{look} truly random from the perspective of any
efficient algorithm. We call these bits \emph{pseudorandom}.

This surprising and powerful idea leads us to our next cryptographic
primitive\(\ldots\)

\section{Pseudorandom Functions}\label{pseudorandom-functions}

A pseudorandom function is defined over a keyspace \(\calK\), and input
space \(\calX\) and output space \(\calY\). To be useful a pseudorandom
function must satisfy the following security definition:

\begin{definition}[Pseudorandom Function,
PRF]\protect\hypertarget{def-prf}{}\label{def-prf}

A function \(F: \mathcal{K} \cross \calX \rightarrow \calY\) is a
pseudorandom function if all efficient algorithms \(\mathcal{A}\) win
the following game with probability
\(\tfrac{1}{2} + \text{``negligible''}\):

\begin{itemize}
\item
  The challenger samples a random bit \(b \gets \zo\) and a key
  \(k \getsr \calK\).
\item
  If \(b = 0\), the challenger sets \(f(\cdot) \deq F(k, \cdot)\).
\item
  If \(b = 1\), the challenger sets
  \(f(\cdot) \getsr \Funs[\calX, \calY]\).

  \marginnote{\begin{footnotesize}

  Here, \(\Funs\) is the set of all functions from \(\calX\) to
  \(\calY\).

  \end{footnotesize}}
\item
  Then for \(i = 1, 2, \dots\) (polynomially many times):

  \begin{itemize}
  \tightlist
  \item
    The adversary sends the challenger a values \(x_i \in \calX\).
  \item
    The challenger responds with \(y_i \gets f(x_i) \in \calY\).
  \end{itemize}
\item
  The adversary outputs a guess \(\hat{b}\) at the bit \(b\).
\item
  The adversary wins if \(b = \hat{b}\).
\end{itemize}

First, the challenger will sample a random \(b \leftarrow \bin\) and a
key \(k \leftarrow \mathcal{K}\).

\end{definition}

The adversary can trivially win this game with probability
\(\tfrac{1}{2}\) by just guessing the bit \(b\) at random. This
definition asserts that no efficient adversary can do much better than
that.

If we have such a pseudorandom function \(F\), we could easily construct
a MAC---we can just use the message as the input to the pseudorandom
function along with the key: \(\MAC(k, m) \deq F(k, m)\).

\subsection{Constructing pseudorandom functions from
one-wayness}\label{constructing-pseudorandom-functions-from-one-wayness}

It is not at all obvious that pseudorandom functions should exist at
all! They seem like a very magical primitive indeed.

One surprising fact is that if there exists \emph{any} function that is
``hard to invert,'' in a sense we will define, then pseudorandom
functions exist. For example, if you believe that factoring large
numbers is difficult (as many people do), then pseudorandom functions
exist.

In particular the following definition captures the notion of a function
that is hard to invert:

\begin{definition}[One-Way
Function]\protect\hypertarget{def-owf}{}\label{def-owf}

A function \(f \colon \calX \to \calY\) is a \emph{one-way function} if
for all efficient adversaries \(\calA\),
\[ \Pr[f(\calA(f(x))) = f(x) \colon x \getsr \calX ] \leq \text{``negligible''}.\]

\end{definition}

Having defined one-way functions, we now have the following surprising
and non-obvious result:

\marginnote{\begin{footnotesize}

Notice that if \(\classP = \classNP\), one-way functions do not exist,
and therefore psuedorandom functions do not exist.

\end{footnotesize}}

\begin{theorem}[]\protect\hypertarget{thm-prf-iff-owf}{}\label{thm-prf-iff-owf}

Psueodorandom functions exist if and only if one-way functions exist.
(\citeproc{ref-HILL99}{Hastad et al. 1999})

\end{theorem}

In practice, we assume that:

\begin{itemize}
\tightlist
\item
  the function \(f(x) \deq \text{SHA256}(x)\) is a one-way function
  where the domain is the set of 256-bit strings,
\item
  the function \(f(x) \deq \text{AES}(x, 0^{128})\) is a one-way
  function, where the domain is the set of 128-bit strings, and
\item
  the function \(f(x) \deq 2^x \bmod p\) is a one-way function on domain
  \(\{1, \dots, p\}\), for a sufficiently large prime \(p\).
\end{itemize}

\subsection{Pseudorandom functions in
practice}\label{pseudorandom-functions-in-practice}

In practice, we use the Advanced Encryption Standard (AES) as a
pseudorandom function. The AES function on key length
\(\kappa \in \{128, 192, 256\}\) has the type signature
\(\mathsf{AES}: \bin^\kappa \times \bin^{128} \rightarrow \bin^{128}\).
That is, it takes a 128-bit input and generates a 128-bit output.

\marginnote{\begin{footnotesize}

We don't have any mathematical proof that AES is a pseudorandom
function. However, it has undergone a tremendous amount of cryptanalysis
and the best attacks on AES are only marginally better than the obvious
brute-force attacks.

\end{footnotesize}}

\section{From pseudorandom functions to
MACs}\label{from-pseudorandom-functions-to-macs}

\textbf{MACs for short messages.} Using AES as a pseudorandom function
on a 128-bit domain, we can build a MAC for 128-bit messages as
described above : \(\MAC(k, m) \deq \mathsf{AES}_k(m)\). However, since
AES takes only 128 bits as input, using AES directly, we can only
authenticate 128-bit messages.

\textbf{Insecure ways to construct a MAC for long messages.} A bad way
to construct a MAC for long messages from a pseudorandom function \(F\)
for 128-bit messages is just to chop our message \(m\) up into 128-bit
blocks \(m = (m_1, m_2, \dots)\) and MAC each block separately. Our tag,
then, would look something like \(\left(F(k,m_1), F(k,m_1)\right)\).
However, there is a problem! Given the tag \(t = (t_1, t_2)\) for a
message \(m=(m_1, m_2)\), we can easily generate a valid tag
\(t' = (t_2, t_1)\) for a different message \(m'=(m_2, m_1)\).

\textbf{MACs for long messages: The easy way.} If we have a pseudorandom
function \(F\) with an input space of 256-bits, we can construct a MAC
on arbitrary-length messages using the ``hash-and-sign'' paradigm. In
particular, we use a collision-resistant hash function
\(H\colon \zo^* \to \zo^{256}\) (Definition~\ref{def-crhf}) and we
define the MAC on message space \(\zo^*\) as:
\[ \MAC(k,m) \deq F(k, H(m)).\]

\marginnote{\begin{footnotesize}

Notice that we cannot use AES as the pseudorandom function \(F\) in this
construction, since AES only takes a 128-bit input. In this case, we
would need a collision-resistant hash function
\(H \colon \zo^* \to \zo^{128}\), but it is always possible to find
collisions in hash functions with 128-bit output in time \(2^{64}\). So
such a MAC can never be secure against attackers running in time
\(2^{64}\).

\end{footnotesize}}

In practice, we typically do not construct MACs in this way because
collision-resistant hash functions are typically more expensive to
compute (per bit of input) than pseudorandom functions, such as AES.

\subsection{MACs for long messages: Cipher-Block Chaining
MAC}\label{macs-for-long-messages-cipher-block-chaining-mac}

\marginnote{\begin{footnotesize}

Applying the PRF to the last block using an independent random key is
important. If we do not use a new key, an adversary can mount a
length-extension attack. That is, if the adversary asks for
\(t = \mathsf{MAC}(k, m_1)\) and \(t' = \mathsf{MAC}(k, t)\), \(t'\) is
also a valid key for the original message with two zero blocks attached
\(\mathsf{MAC}(k, m_1 \| 0 \| 0)\). The chain of AES applications
becomes equivalent, since zero blocks are equivalent to skipping the XOR
and adding AES applications.

\end{footnotesize}}

A common and secure way to construct a MAC for long messages from a MAC
for short messages is to \emph{chain} the output of each of these calls
to the pseudorandom function. Given our chopped message
\((m_1, m_2, \ldots, m_n)\), we will generate \(t_1 = F(k,m_1)\) as
before. When generating \(t_2\), we will first XOR \(t_1\) into the
input: \(t_2 = F(k, m_2 \oplus t_1)\). This continues until the end of
the message, at which point have a tag \(t_n\). Finally, we apply the
PRF \emph{with a different key} \(k'\) to the value \(t_n\) and output
this tag \(t \gets F(k', t_n)\). This construction is called CBC-MAC or
CMAC.

\begin{figure}

\centering{

\pandocbounded{\includegraphics[keepaspectratio]{index_files/mediabag/figs/cbc-mac.pdf}}

}

\caption{\label{fig-cbc-mac}The CBC-MAC construction.}

\end{figure}%

CBC-MAC is going out of favor for two reasons:

\begin{enumerate}
\def\labelenumi{\arabic{enumi}.}
\tightlist
\item
  It is impossible to parallelize the MAC computation: the chaining
  procedure is inherently sequential so you cannot speed it up, even if
  you have a computer with many CPU cores.
\item
  Computing the MAC requires one PRF invocation \emph{per block} of the
  message. There are even faster MACs that require only one PRF
  invocation \emph{per message} total, plus a number of fast
  ``non-cryptographic'' operations per message block. These MACs can be
  faster than CBC-MAC on some processors. The GMAC construction we will
  see next is one example.
\end{enumerate}

\subsection{A parallelizable MAC: Carter-Wegman MAC}\label{sec-mac-cw}

We now describe a different way to authenticate long messages. This MAC
scheme is parallelizable and also requires only one single PRF
invocation per message authenticated (independent of the message
length). The construction is named the Carter-Wegman MAC, after its
inventors.(\citeproc{ref-CW81}{Wegman and Carter 1981}) Modern
encryption schemes, including AES-GCM (Section~\ref{sec-enc-gcm}) use a
Carter-Wegman-style MAC as a key ingredient.

For this construction, we will use the notation \(\Z_p\) to indicate the
set of integers modulo \(p\) with addition and multiplication modulo
\(p\). So \(x + y \in \Z_p\) means that we add \(x\) and \(y\) as
integers and reduce the result modulo \(p\). Typically, we will think of
\(p\) as a prime---of 64 bits, for example. The construction uses a
fixed a prime number \(p\) as a parameter, where \(p \approx 2^n\) for
security parameter \(n\).

\marginnote{\begin{footnotesize}

So in practice, we take \(p \approx 2^{128}\) for 128-bit security.

\end{footnotesize}}

\textbf{Universal hash function.}

Before we look at the construction of the Carter-Wegman MAC, we first
define an important building block: the notion of a \emph{universal hash
function}, or UHF for short. A universal hash function is keyed, and
provides collision-resistance when the adversary does not know the key.
Specifically, we say that \(H\) is a universal hash function if, when an
adversary chooses two messages \(m\) and \(m'\) where \(m\neq m'\),

\[
Pr[H(k, m) = H(k, m')] \leq \text{negl.}
\]

Intuitively, a universal hash function is a weaker primitive than a
collision-resistant hash function: the adversary does not know the
precise hash function that will be applied to their messages, because
the adversary does not know what key will be used.

One simple and practical construction of a universal hash function is
based on polynomials. Given a long message \(m\), break it up into
fixed-size chunks \(m_0, m_1, \ldots, m_{l-1}\). Then, the hash of that
message is defined as:

\marginnote{\begin{footnotesize}

In practice, these fixed-size chunks are going to be 128 or 256 bits
long.

\end{footnotesize}}

\[
H(k, m_0 || m_1 || \ldots || m_{l-1}) = (m_0 + m_1 k + m_2 k^2 +
  \ldots + m_{l-1} k^{l-1}) \mod p
\]

We can give some intuition for why \(H\) is a universal hash function
(i.e., collision-resistant for a randomly chosen key). In order for a
pair of messages \(m\) and \(m'\) to collide, it must mean that
\(H(k, m) = H(k, m')\), which in turn means that
\(H(k, m) - H(k, m') = 0\). Expanding the definition of \(H\) as a
polynomial, this means that

\[
(m_0 - m'_0) + (m_1 - m'_1) k + (m_2 - m'_2) k^2 + \ldots + (m_{l-1} - m'_{l-1}) k^{l-1} = 0
\]

which is another way of saying that \(k\) is a root of that
degree-\(l-1\) polynomial. We know that there can be at most \(l-1\)
roots of a degree-\(l-1\) polynomial, but there are \(p \approx 2^n\)
possible choices for \(k\), so the probability that our randomly-chosen
\(k\) happens to be one of those roots is \(\frac{l-1}{p}\), which is
negligible.

\textbf{MAC construction.}

The MAC uses a pseudorandom function
\(F \colon \calK \times \Z_p \to \Z_p\). The keyspace for the MAC is
\(\calK\), so the MAC key consists of a key for the pseudorandom
function. The message space for the MAC is \(\calM = \Z_p^{\leq L}\),
the set of vectors of integers of \(\Z_p\) elements of length at most
\(L\) where \(L \ll p\). Here, assume that the message vector has length
at least \(1\).

\marginnote{\begin{footnotesize}

Here, the input space of the pseudorandom function \(F\) is the set of
integers in \(\{0, \dots, p-1\}\). Given a pseudorandom function on
bitstrings, it is indeed possible to construct one that operates on
numbers in \(\Z_p\) like this by interpreting each number as a
bitstring.

\end{footnotesize}}

One other difference is that this MAC construction is \emph{randomized}.
So there are now two algorithms:

\begin{itemize}
\tightlist
\item
  \(\MACSign(k, m) \to t\), which takes as input a key \(k\) and message
  \(m\) and outputs a MAC tag \(t\), and
\item
  \(\MACVerify(k, m, t) \to \zo\), which takes as input a key \(k\),
  message \(m\), tag \(t\), and outputs an accept/reject bit.
\end{itemize}

The security definition here is essentially the same as for
deterministic MACs, except that we use different algorithms to generate
and verify the MAC tags.

The Carter-Wegman MAC construction is then:

\(\MACSign(k, m \in \Z^{\leq L}_p) \to t\).

\begin{itemize}
\item
  Compute \(v \gets F(k, 0) \in \Z_p\).
\item
  Parse the message into chunks as
  \((m_1, \dots, m_{\ell}) \gets m \in \Z_p^\ell\).
\item
  Compute
  \(M(v) \gets m_1 v + m_2 v^2 + m_3 v^3 \cdots + m_{\ell} v^{\ell} \in \Z_p\).

  \marginnote{\begin{footnotesize}

  Essentially we are viewing the blocks of the message \(m\) as
  coefficients of a degree-\(t\) polynomial \(M(\cdot)\). We then
  evaluate this polynomial at the secret point \(v\) determined by the
  MAC key.

  \end{footnotesize}}

  \todobox{\textbf{TODO (HCG):} Check the definition of the message
  polynomial \(M\). Should there be an additional \(v^{t+1}\) monomial?}
\item
  Sample a nonce \(r \getsr \Z_p\).
\item
  Output \(t \gets \big(r, F(k, r) + M(v)\big) \in \Z^2_p\) as the MAC
  tag.
\end{itemize}

\(\MACVerify(k, m, t) \to \zo\).

\begin{itemize}
\tightlist
\item
  Compute \(v \gets F(k, 0) \in \Z_p\).
\item
  Parse the message into chunks as
  \((m_1, \dots, m_{\ell}) \gets m \in \Z_p^\ell\).
\item
  Compute
  \(M(v) \gets m_1 v + m_2 v^2 + m_3 v^3 + \cdots + m_{\ell} v^{\ell}\in \Z_p\).
\item
  Parse the tag \((r, z) \gets t \in \Z^2_p\).
\item
  Output ``1'' if and only if \(z - F(k, r) = M(v)\).
\end{itemize}

\textbf{Security intuition.} The security argument here goes as follows:

\marginnote{\begin{footnotesize}

For a detailed treatment of Carter-Wegman security see Boneh and Shoup's
textbook, \emph{A Graduate Course in Applied Cryptography}, Section 7.4.

\end{footnotesize}}

\begin{itemize}
\item
  First, we appeal to the PRF security property to argue that we can
  replace the values \(F(k_F, r)\) used to generate the tags with truly
  random values.
\item
  Next, we show that as long as the \(\MACSign\) algorithm never samples
  the same nonce \(r\) twice, the masking values \(F(k,r)\) are
  independent random values that complete hide the values \(M(v)\). So,
  the adversary learns no information on the secret point \(v\) by
  making MAC queries.
\item
  Now, say that the adversary finds a forged message-tag pair
  \((m^*, t^*)\). There are two cases:

  \begin{itemize}
  \tightlist
  \item
    Either the forgery uses a fresh random nonce \(r^*\) that did not
    appear as the response to any of the adversary's MAC queries. In
    this case, the forgery is only valid with probability \(1/p\).
  \item
    Alternatively, the forger could use a random nonce \(r^*\) that is
    equal to the nonce \(r\) returned from one of the adversary's MAC
    queries. In this case, we have the following relations, where
    message \(m\) polynomial \(M\) was the message the adversary queried
    of the challenger: \[\begin{aligned}
      F(k, r) &= M(v) - z\\
      F(k, r) &= M^*(v) - z^*\\
      0 &= \big(M(v) - M^*(v)\big) + (z - z^*).
    \end{aligned}\] Since \(m \neq m^*\), we know \(z \neq z^*\). So
    \((M(\cdot) - M^*(\cdot)) + (z^* - z)\) is a non-zero polynomial of
    degree at most \(t\). Since such a polynomial can have at most
    \(\ell \leq L\) zeros in \(\Z_p\), and since the adversaries view is
    independent of the evaluation point \(v \in \Z_p\), the probability
    that the adversary's forgery is valid is at most \(\ell/p\).
  \end{itemize}

  In either case, the adversary's probability of forging is
  \(O(L)/p = \poly(\lambda) \cdot \negl(\lambda) = \negl(\lambda)\) on
  security parameter \(\lambda\).
\end{itemize}

\chapter{Digital Signatures}\label{sec-sig}

In the last section, our strategy for authentication depended on two
parties sharing a secret key. In that discussion, we completely left out
of the picture how these parties should exchange this secret key. Our
implication was that they went to some private room and exchanged the
key in secret, but in many cases this is not practical: if they could
whisper a key, why not just whisper the message?

Luckily, there is a way to get around this requirement for a shared
secret using \emph{public-key cryptography}.(\citeproc{ref-DH76}{Diffie
and Hellman 1976})

\marginnote{\begin{footnotesize}

The original Diffie-Hellman paper from 1976, which introduced public-key
cryptography, is a fascinating read.

\end{footnotesize}}

\section{Definitions}\label{definitions}

The basic idea of public-key cryptography, applied to authentication, is
that each party will generate two linked keys---a secret signing key and
a public verification key. The verification key will be good enough to
verify that a signature is valid, but not to generate new signatures.

\begin{definition}[Signature
Scheme]\protect\hypertarget{def-signature-scheme}{}\label{def-signature-scheme}

A signature scheme is associated with a message space \(\calM\) and
three efficient algorithms \((\Gen, \Sign, \Ver)\).

\marginnote{\begin{footnotesize}

In theoretical papers, people will write \(\Gen(1^\lambda)\) to indicate
that the key-generation algorithm takes as input a length-\(\lambda\)
string of ones. This is just a hack to make the input given to \(\Gen\)
\(\lambda\) bits long so that the \(\Gen\) algorithm can run in time
polynomial in its input length: \(\poly(\lambda)\). If we express
\(\lambda\) in binary, then \(\Gen(\lambda)\) gets a
\(\log_2 \lambda\)-bit input and can only run in time
\(\poly(\log \lambda)\). This distinction is really unimportant, but if
you see the \(1^\lambda\) notation, you will now know what it means.

\end{footnotesize}}

\begin{itemize}
\tightlist
\item
  \(\Gen(\lambda) \to (\sk, \vk)\). The key-generation algorithm as
  input a security parameter \(\lambda \in \N\) and outputs a secret
  signing key \(\sk\) and public verification key \(\vk\). The algorithm
  \(\Gen\) runs in time \(\poly(\lambda)\).
\item
  \(\Sign(\sk, m) \to \sigma\). The signing algorithm takes as input a
  secret key \(\sk\) and a message \(m \in \calM\), and outputs a
  signature \(\sigma\).
\item
  \(\Ver(\vk, m, \sigma) \to \zo\). The signature-verification algorithm
  takes as input a public verification key \(\vk\), a message
  \(m \in \calM\), and a signature \(\sig\), and outputs \(\bin\),
  indicating acceptance or rejection.
\end{itemize}

\end{definition}

For a signature scheme to be useful, a correct verifier must always
accept messages from an honest signer. Formally, we have:

\begin{definition}[Digital signatures:
Correctness]\protect\hypertarget{def-sig-correctness}{}\label{def-sig-correctness}

A digital-signature scheme \((\Gen, \Sign, \Ver)\) is \emph{correct} if,
for all messages \(m \in \calM\):
\[ \Pr\big[\Ver(\vk, m, \Sign(\sk, m)) = 1 \colon (\sk, \vk) \gets \Gen(\lambda) \big] = 1. \]

\end{definition}

The standard security notion for digital signatures is very similar to
that for MACs (Definition~\ref{def-mac-sec}). The only difference here
is that a digital-signature scheme splits the single secret MAC key into
two keys: a secret signing key and a public verification key. Otherwise
the definition is essentially identical.

\begin{definition}[Digital signatures: Security -- existential
unforgeability under chosen message
attack]\protect\hypertarget{def-sig-sec}{}\label{def-sig-sec}

A digital-signature scheme \((\Gen, \Sign, \Ver)\) is \emph{secure} if
all efficient adversaries win the following security game with only
negligible probability:

\begin{itemize}
\tightlist
\item
  The challenger runs \((\sk, \vk) \gets \Gen(\lambda)\) and sends
  \(\vk\) to the adversary.
\item
  For \(i = 1, 2, \dots\) (polynomially many times)

  \begin{itemize}
  \tightlist
  \item
    The adversary sends a message \(m_i \in \calM\) to the challenger.
  \item
    The challenger replies with \(\sigma_i \gets \Sign(\sk, m_i)\).
  \end{itemize}
\item
  The adversary outputs a message-signature pair \((m^*, \sigma^*)\).
\item
  The adversary wins if \(\Ver(\vk, m^*, \sigma^*) = 1\) and
  \(m^* \not \in \{m_1, m_2, \dots\}\).
\end{itemize}

\end{definition}

Notice that this security definition does not guarantee that an attacker
cannot forge a new signature on a message that it has already seen a
signature of. Namely, given a valid message-signature pair
\((m, \sigma)\) an adversary may be able to produce additional valid
message-signature pairs on the same message:
\((m, \sigma'), (m, \sigma''), \dots\).

In some applications, we want to prohibit an attacker from finding
\emph{any} new message-signature pair. We call this security notion
``\emph{strong} existential unforgeability under chosen message
attack.'' The definition is the same as in 1 except that we require the
adversary to find a valid-message signature pair \((m^*, \sigma^*)\)
such that
\((m^*, \sigma^*) \not \in \{ (m_1, \sigma_1), (m_2, \sigma_2), \dots \}\).
Standard digital-signature schemes, such as the elliptic-curve digital
signature algorithm (EC-DSA) or the RSA algorithm with full-domain
hashing (RSA-FDH), are believed to have this strong security property.

\section{Constructing a Signature
Scheme}\label{constructing-a-signature-scheme}

In the following sections, we will show how to construct a
digital-signature scheme from any one-way function
(Definition~\ref{def-owf}).

We will generate a signature scheme that is secure, but that has
relatively large signatures and public keys: to achieve security against
attackers running in time \(2^\lambda\), this signature scheme has
signatures of \(O(\lambda^2)\) bits. Widely used modern digital
signature schemes (e.g., EC-DSA) have signatures of \(O(\lambda)\) bits.

\marginnote{\begin{footnotesize}

One benefit of the signature scheme that we present here is
that---unlike EC-DSA, RSA, DSA, and other widely used signature
schemes---this one is plausibly secure even against \emph{quantum}
adversaries. There is ongoing work to standardize signature schemes
secure against quantum adversaries; see
\url{https://csrc.nist.gov/projects/pqc-dig-sig}

\end{footnotesize}}

We will construct this scheme in three stages:

\begin{enumerate}
\def\labelenumi{\arabic{enumi}.}
\item
  Construct a signature scheme for signing a \emph{single bit}.
\item
  Construct a \emph{one-time secure} signature scheme for signing a
  \emph{fixed length} messages. With this scheme, an attacker who sees
  two signatures under the same signing key can forge signatures. In
  addition, the secret signing key for this scheme will be larger than
  the size of the message being signed.
\item
  Construct a \emph{one-time secure} scheme for \emph{arbitrary-length
  messages}. Here, we construct a one-time signature scheme whose secret
  signing key is independent of the length of the signed message.
\item
  Construct a \emph{many-time secure} scheme (i.e., a fully secure one
  under

  \begin{enumerate}
  \def\labelenumii{\arabic{enumii})}
  \tightlist
  \item
    for \emph{arbitrary-length messages}. This last scheme is a fully
    secure and fully functional digital-signature scheme.
  \end{enumerate}
\end{enumerate}

\section{Constructing a Signature Scheme for Signing a Single
Bit}\label{constructing-a-signature-scheme-for-signing-a-single-bit}

This signature scheme is not useful on its own, and is given only as a
step towards the final construction. It uses as a building block a
one-way function \(f:\calX\rightarrow\calY\). Recall that \(f\) a
one-way function if it is easy to compute but hard to invert; namely
there is an efficient algorithm that given \(x\in\calX\) outputs
\(f(x)\), and at the same time any efficient algorithm given \(y=f(x)\)
for a random \(x\leftarrow \cal X\), finds an inverse \(x'\in\calX\)
such that \(f(x')=f(x)\) with only negligible probability.

\marginnote{\begin{footnotesize}

In this construction, we leave the security parameter \(\lambda\)
implicit. To be fully formal, \(\Gen\) would take \(\lambda\) an input.
The one-way function \(f\) and its domain \(\calX\) would both depend on
\(\lambda\). So we would write \(f_\lambda\) and \(\calX_\lambda\).

\end{footnotesize}}

\begin{itemize}
\tightlist
\item
  \(\Gen() \to (\sk, \vk)\). Choose two random elements \(x_0, x_1\)
  from \(\calX\) and let \((y_0,y_1)=(f(x_0),f(x_1))\). Output
  \(\sk=(x_0,x_1)\) and \(\vk=(y_0,y_1)\).
\item
  \(\Sign(\sk, b) \to \sigma\). Parse \(\sk=(x_0,x_1)\) and output
  \(\sigma=x_b\),
\item
  \(\Ver(\vk, b, \sig) \to \zo\). Parse \(\vk=(y_0,y_1)\) and output
  \(1\) if and only if \(f(\sig)=y_b\). (Otherwise, the signing routine
  rejects.)
\end{itemize}

\section{One-time-secure Signatures (Lamport
Signatures)}\label{sec-lamport}

In this section we give a very simple and elegant construction of a
one-time-secure digital signature scheme, due to
Lamport.(\citeproc{ref-L79}{Lamport 1979}) The construction is a
straightforward generalization of the signature scheme constructed
above: Each message is signed bit-by-bit, where each bit is signed using
a fresh and independently generated secret key.

Before giving the construction, we define one-time security for
digital-signature schemes. This signature scheme is not generally useful
on its own, but is useful as a building block.

\begin{definition}[Digital signatures: One-Time
Security]\protect\hypertarget{def-sig-once}{}\label{def-sig-once}

A digital-signature scheme \((\Gen, \Sign, \Ver)\) over message space
\(\calM\) is \emph{one-time secure} if all efficient adversaries win the
following game with negligible probability:

\begin{itemize}
\tightlist
\item
  The challenger generates \((\sk, \vk) \leftarrow \Gen(\lambda)\) and
  sends \(\vk\) to the adversary.
\item
  The adversary sends the challenger \emph{single} message
  \(m \in \calM\).
\item
  The challenger responds with \(\sig = \Sign(\sk, m)\).
\item
  The adversary outputs \((m^*, \sig^*)\).
\item
  The adversary wins the game if \(\Ver(\vk, m^*, \sig^*) = 1\) and
  \(m^* \neq m\).
\end{itemize}

\end{definition}

\textbf{Lamport signatures.} We now construct a one-time secure
signature scheme for messages in \(\bin^n\), for some fixed message
length \(n \in \N\). To do this, we will define the following
algorithms, which make use of a one-way function
\(f \colon \calX \to \calY\):

\begin{itemize}
\tightlist
\item
  \(\Gen() \to (\sk, \vk)\). Choose \(2n\) random elements from
  \(\calX\), the domain of the one-way function \(f\). Arrange these
  values in to a \(2 \times n\) matrix, which forms the secret signing
  key \(\sk\). The public verification key just consists of the \(2n\)
  images of these values under the one-way function \(f\):
  \[ \sk \gets \begin{pmatrix}
  x_{10} & \ldots & x_{n0} \\
  x_{11} & \ldots & x_{n1} \\
  \end{pmatrix},\quad \vk \gets \begin{pmatrix}
  f(x_{10}) & \ldots & f(x_{n0}) \\
  f(x_{11}) & \ldots & f(x_{n1}) \\
  \end{pmatrix}.\]
\item
  \(\Sign(\sk, m) \to \sigma\) outputs \((x_{1m_1}, \ldots x_{nm_n})\),
  where \(m_1 \dots m_n\) are the individual bits of the length-\(n\)
  message \(m \in \zo^n\).
\item
  \(\Ver(\vk, m, \sig) \to \zo\) parses the the message into bits
  \(m = m_1\dots m_n \in \zon\) and the signature \(\sig\) into its
  individual symbols \(\sig = (x_1^*, \ldots x_n^*) \in \calX^n\). The
  signing routine accepts if, for all \(i \in \{1, \dots, n\}\):
  \begin{equation}\protect\phantomsection\label{eq-lamport}{
  f(x^*_i) = \vk_{i,m_i}.
  }\end{equation} In other words, the routine accepts if applying the
  one-way function \(f\) to each symbol of the signature matches the
  corresponding value in the verification key. (Otherwise, the signing
  routine rejects.)
\end{itemize}

This signature scheme has relatively large keys: the verification key,
in particular consists of \(2n\) values, where each is of length
\(\Omega(\lambda)\) bits. So the total length is roughly \(2n\lambda\)
bits---much longer than the \(n\)-bit message being signed.

In addition, notice that an adversary who sees signatures on even two
messages can forge signatures on messages of its choice. In particular:

\begin{itemize}
\tightlist
\item
  The adversary first asks for a signature on the message \(m_0 = 0^n\).
  It receives \(\sig_0 = (x_{10}, \ldots, x_{n0})\).
\item
  The adversary then then asks for a signature on the message
  \(m_1 = 1^n\). It receives \(\sig_1 = (x_{11}, \ldots, x_{n1})\).
\item
  At this point, the adversary has the entire secret signing key!
\end{itemize}

However, we will show that this scheme is indeed one-time secure.

\textbf{Claim.} The Lamport signature scheme is one-time secure under
the assumption that \(f\) is a one-way function.

\marginnote{\begin{footnotesize}

Remember that if \(\classP = \classNP\), one-way functions, and also
digital signature schemes, do not exist. So any proof of security of a
digital-signature scheme will require some sort of cryptographic
assumption.

\end{footnotesize}}

In cryptography, we generally prove these security claims by
\emph{reduction}: we will show that if there exists an efficient
adversary \(\calA\) that breaks the security of our scheme, then we can
construct an efficient adversary \(\calB\) that breaks one of our
assumptions. If we do this, we have reached a contradiction to one of
our assumptions, so the first adversary cannot exist.

\begin{proof}
Suppose there exists an adversary \(\A\) that wins the one-time-security
game of Definition~\ref{def-sig-once} with non-negligible probability
\(\epsilon\). That is, the adversary can produces \((m^*, \sig^*)\) such
that \(\Ver(\vk, m^*, \sig^*) = 1\) and \(m \neq m^*\) given only
\(\sig = \Sign(\sk, m)\). We can then construct an adversary \(\B\) that
can use \(\A\) to invert the one-way function.

\marginnote{\begin{footnotesize}

Lamport's construction shows that if one-way functions exist, then so do
digital signatures. Can you show that if digital signatures exist, then
so do one-way functions?

\end{footnotesize}}

In particular, our adversary \(\calB\) will use algorithm \(\calA\) as a
subroutine to invert the one-way function. We will show that if
\(\calA\) wins in the one-time signature security game often, then
algorithm \(\calB\) will invert the one-way function often, which is a
contradiction.

Assume our one-way function is of the form \(f \colon \calX \to \calY\)
and that the Lamport signature scheme is on \(n\)-bit messages. The
one-way-function adversary \(\calB\) operates as follows:

\begin{itemize}
\tightlist
\item
  The adversary \(\calB\) is given a point \(y \in \calY\) and its task
  is to produce a preimage of \(y\) under \(f\).
\item
  The adversary \(\calB\) generates a signing keypair as follows:

  \begin{itemize}
  \tightlist
  \item
    It runs the key-generation algorithm for the Lamport signature
    scheme \((\sk, \vk) \gets \Gen()\).
  \item
    The adversary chooses a random value \(i^* \getsr \{1, \dots, n\}\)
    and a random bit \(\beta^* \getsr \zo\).
  \item
    The adversary sets \(\vk_{i^*,\beta^*} \gets y\). That is, it
    inserts the one-way-function point it must invert into a random
    location in the verification key.
  \end{itemize}
\item
  The adversary then sends the verification key \(\vk\) to the
  Lamport-signature adversary \(\calA\).
\item
  The adversary \(\calA\) asks for the signature on a message
  \(m = m_1 m_2 \dots m_n \in \zon\).
\item
  If \(m_{i^*} = \beta^*\), then algorithm \(\calB\) cannot produce a
  valid signature on the message \(m\) and it outputs FAIL.
\item
  Otherwise, the algorithm \(\calB\) returns the signature
  \(\sigma = (\sk_{1,m_1}, \dots, \sk_{n,m_n}) \in \calX^n\) to
  algorithm \(\calA\).
\item
  Algorithm \(\calA\) then produces a forged message-signature pair
  \((m^*, \sigma^*)\), where \(m \neq m^*\).
\item
  Algorithm \(\calB\) parses \(m^* = m^*_1 \dots m^*_n \in \zon\) and
  \(\sigma^* = \sigma^*_1 \dots \sigma^*_n \in \calX^n\). Then:

  \begin{itemize}
  \tightlist
  \item
    If \(m_{i^*} = m_i\), algorithm \(\calB\) outputs FAIL.
  \item
    Otherwise, algorithm \(\calB\) outputs
    \(x \gets \sigma^*_{i^*} \in \calX\).
  \end{itemize}
\end{itemize}

First, notice that whenever \((m^*, \sigma^*)\) is a valid
message-signature pair and whenever algorithm \(\calB\) does not output
FAIL, algorithm \(\calB\) outputs a preimage \(x \in \calX\) of point
\(y \in \calY\) under the one-way function \(f\). That is because, by
the verification relation (Equation~\ref{eq-lamport}) for Lamport
signatures,
\[f(x) = f(\sigma^*_{i^*}) = \vk_{i^*,m^*_{i^*}} = \vk_{i^*, 1 - m_i} = \vk_{i^*, \beta^*} = y.\]

Now, we must show that algorithm \(\calB\) does not output FAIL too
often. Since algorithm \(\calB\) chooses the values \(i^*\) and
\(\beta^*\) at random, and since the adversary \(\calA\) behavior is
\emph{independent} of these values, we can say:

\begin{itemize}
\tightlist
\item
  the probability of the first failure event is \(1/2\), since there are
  two possible choices of \(m_{i^*}\) and only one of these is bad, and
\item
  the probability of the second failure event is at most \(1/n\), since
  \(m\) and \(m^*\) must differ in at least one of \(n\) bits, and there
  is a \(1/n\) probability that this differing bit is at index \(i^*\).
\end{itemize}

The events that \(\calA\) breaks the signature scheme and that either of
these failures occur are all \emph{independent}. Then if \(\calA\)
breaks the one-time security of the Lamport signature scheme with
probability \(\epsilon\), our one-way-function adversary \(\calB\)
inverts the one-way function with probability
\[ \epsilon_\text{one-way} = \epsilon \cdot \frac{1}{2} \cdot \frac{1}{n}.\]

The probability of either bad is at most \(1/2 + 1/n\), by the union
bound. Therefore if algorithm \(\calA\) breaks one-time security of
Lamport's scheme with probability \(\epsilon\), If \(\epsilon\) is
non-negligible, then \(\epsilon_\text{one-way} = \epsilon/2n\) is also
non-negligible, and we have a contradiction.
\end{proof}

\section{A one-time signature scheme for arbitrary-length
messages}\label{a-one-time-signature-scheme-for-arbitrary-length-messages}

In the Lamport signature scheme (Section~\ref{sec-lamport}), the length
of the keys scales with the size of the message being signed. To adapt
our scheme from above into a scheme that works on arbitrary-length
messages without the key growing arbitrarily large, we will use a
strategy called \emph{hash-and-sign}. In essence, the signing algorithm
will pass the message through a hash function to generate a fixed-size
digest before applying a signature scheme that works only on
fixed-length messages.

\marginnote{\begin{footnotesize}

Essentially all signature schemes used in practice use this
hash-and-sign construction.

\end{footnotesize}}

This paradigm is called ``hash and sign,'' and is very common. In
practice, hashing is computationally cheap operation while signing turns
out to be computationally relatively expensive. So it is common to hash
a message before signing it in order to reduce the size of the message
that must be signed.

The following claim gives the general construction:

\textbf{Claim (Hash-and-sign paradigm).} Given a collision-resistant
hash function \(h: \bin^* \rightarrow \bin^{n}\) and a signature scheme
\((\Gen, \Sign, \Ver)\) for message space \(\calM = \bin^n\) (such as
the one in Section~\ref{sec-lamport}), there exists a signature scheme
\((\Gen', \Sign', \Ver')\) for \(\calM' = \bin^*\) as follows:

\begin{itemize}
\tightlist
\item
  \(\Gen' \deq \Gen\). The key-generation algorithm is unchanged.
\item
  \(\Sign'(\sk, m) \deq \Sign(sk, h(m))\). We hash the message using the
  hash function \(h\) before passing the hashed message to the original
  signing function.
\item
  \(\Ver'(\vk, m, \sig) \deq \Ver(\vk, h(m), \sig)\). We use the
  original \(\Ver\) to check that the tag matches \emph{hash} of the
  original message.
\end{itemize}

\textbf{Security Intuition.} Suppose that there exists an efficiency
adversary that breaks \((\Gen', \Sign', \Ver')\). In particular, given
\(((m_1, \sig_1, \ldots, (m_t, \sig_t))\), the adversary is able to
construct a valid message-signature pair \((m^*, \sig^*)\) such that
\(m^* \notin \set{m_1, \ldots, m_t}\). There are then two cases:

\begin{enumerate}
\def\labelenumi{\arabic{enumi}.}
\tightlist
\item
  \(h(m^*) \in \set{h(m_1), \ldots, h(m_t)}\). If this is the case,
  there is some \(i \in [t]\) such that \(h(m^*) = h(m_i)\). However,
  \(h\) is collision-resistant as in the definition, so this is a
  contradiction!
\item
  \(h(m^*) \notin \set{h(m_1), \ldots, h(m_t)}\). Since the message that
  we pass to the underlying signature scheme is \(h(m)\), this means
  that the adversary has found a valid signature for \(h(m)\) under the
  original scheme \((\Gen, \Sign, \Ver)\) after seeing only signatures
  of \(h(m_1)\neq h(m), \ldots, h(m_t)\neq h(m_t)\). This breaks the
  security of the underlying signature scheme, which is a contradiction!
\end{enumerate}

\textbf{Application to Lamport.} We can apply the hash-and-sign paradigm
to the Lamport signature scheme from Section~\ref{sec-lamport}. We can
fix our input to the Lamport scheme at, for example, 256 bits, and then
run messages through a hash function that outputs 256 bits before
passing them to the Lamport scheme. This gives a \emph{one-time-secure}
signature scheme for messages of arbitrary length.

\textbf{Security implications of hash and sign.} In practice,
hash-and-sign can actually \emph{increase} the security of our signature
scheme, in a certain sense. As shown in case 2 above, it is absolutely
crucial that the hash function used is collision-resistant: if not, an
adversary can find messages that cause collisions, and then a signature
for one message will also be a valid signature for the other. However,
in practice we often think of hash functions like SHA256 as behaving
like \emph{random oracles}. That is, for a hash function
\(h \colon \zo^* \to \zo^\lambda\) and a string \(x \in \zo^*\) we think
of the value \(h(x)\) as being an independently sampled and uniformly
random value from the co-domain of the hash function, \(\zo^\lambda\).
(Of course, a real-world hash function is \emph{never actually} a random
oracle. A random oracle from \(h \colon \zo^* \to \zo^\lambda\) would
take infinitely many bits to describe, while real-world hash functions
have finite size (and polynomial-size descriptions).)

\marginnote{\begin{footnotesize}

Recall that we faced a similar problem in our MAC construction. Why not
use hash and MAC there? The reason is that we used AES as our PRF, which
takes input of \(\bin^{128}\). As explained by the birthday paradox, it
is possible to find a collision in an output space of size \(2^{128}\)
in only time \(2^{64}\)! This does not provide sufficient security for
practical use, as it would be quite practical to find collisions. If we
had a version of AES that outputted 256 bits, we could indeed apply hash
and sign.

Another reason to not use hash and MAC is that MACs can be faster to
compute than collision-resistant hash functions.

\end{footnotesize}}

Recall that the standard security definition for digital signatures (1)
allows the attacker to request signatures on messages of its choice. If
we pass a message through a hash function before signing it using an
underlying signature scheme scheme, however, we effectively randomize
the message---the adversary can no longer control the input to the
underlying signature scheme. This allows us to define another meaningful
definition of security:

\begin{definition}[Digital signatures: security against random message
attacks]\protect\hypertarget{def-sig-random-message}{}\label{def-sig-random-message}

Any efficient adversary given the public verification key and a list of
random message-signature pairs
\(((m_1, \sig_1), \ldots, (m_t, \sig_t))\) cannot generate a forged
message-signature pair \((m^*, \sig^*)\) such that
\(\Ver(\vk, m^*, \sig^*) = 1\) and
\(m^* \notin \set{m_1, \ldots, m_t}\).

\end{definition}

Note that this definition is \emph{not} good enough on its own---an
adversary often does have the ability to generate signatures for
messages of his choice. However, paired with a hash function modeled as
a random oracle, this definition becomes very useful---if the inputs are
passed through the hash function before they are passed to the signature
scheme, they become effectively random. We can even relax the definition
further without losing practicality: by the same logic, with hash
function in front of the signature scheme, what the adversary needs to
sign is really not a message of their choise, but is the \emph{hash} of
a message of their choice---effectively a random value.

\begin{definition}[Digital signatures: random security against random
message
attacks]\protect\hypertarget{def-signatures-random-security}{}\label{def-signatures-random-security}

Any efficient adversary given \(\vk\) and a list of random
message-signature pairs \(((m_1, \sig_1), \ldots, (m_t, \sig_t))\) and
random \(m^* \notin \set{m_1, \ldots, m_t}\) cannot generate \(\sig^*\)
such that \(\Ver(\vk, m^*, \sig^*) = 1\).

\end{definition}

It is possible to formally argue that given a signature scheme
satisfying Definition~\ref{def-signatures-random-security} and a random
oracle, we can construct a scheme satisfying existential unforgeability
under chosen message attacks.

\section{From one-time security to many-time
security}\label{sec-sig-manytime}

After applying the hash-and-sign strategy above to our Lamport scheme,
we have a signature scheme that is one-time secure for messages of
arbitrary length. In order for the scheme to be useful and satisfy our
security definition, we need to be able to sign polynomially many
messages with a single key pair. To do this, we will use a construction
very similar to the Merkle tree construction we have seen before.

Informally, we will build up a binary tree of Lamport keys of depth 256.
The signing key in each of the \(2^{256}\) leaf nodes will be used to
actually sign messages; we will use a random leaf node to sign a
message, so the fact that there are \(2^{256}\) of them means that the
probability of accidentally choosing the same leaf twice is negligible
(\(2^{-128}\)). The signing key in an intermediate nodes (and in the
root) will be used to sign the (public) verification keys of the two
corresponding child nodes in the tree. Figure~\ref{fig-lamport-tree}
shows a sketch of this tree.

\begin{figure}

\centering{

\pandocbounded{\includegraphics[keepaspectratio]{index_files/mediabag/figs/lamport-tree.pdf}}

}

\caption{\label{fig-lamport-tree}Sketch of the tree of Lamport signature
keys used for the many-time secure signature construction.}

\end{figure}%

Importantly, every signing key in the tree will be used to sign only one
message ever: the key at non-leaf nodes will only ever sign the pair of
its children, and the key at leaf nodes will only ever be used to sign a
single message.

The signature, in this scheme, will consist of the signatures and
verification keys along the path from the root to the chosen leaf node.
The signature will also include information about which child node was
chosen.

This tree, of course, is impractically large, but we can solve that
problem by lazily constructing it using a pseudorandom function. That
is, instead of actually building up all of the leaves of the tree, we
will build up the tree (i.e., generate the signing keys, verification
keys, and signatures) on-demand, and furthermore, we will build it in a
deterministic way using the pseudorandom function, so that we don't have
to remember what parts we might have already computed in the past.

To make the construction more precise, we will assume that we are given:

\begin{itemize}
\tightlist
\item
  a pseudorandom function \(f\) with keyspace \(\calK\), and
\item
  a one-time secure signature scheme \((\Gen_0, \Sign_0, \Ver_0)\).
\end{itemize}

We will need the ability to run the (non-deterministic) \(\Gen_0\)
algorithm on specific randomness, so as to make it deterministic. For a
PRF key \(k \in \calK\) and string \(s\), we will write \(\Gen_0^{F(k,
s)}()\) to indicate the process of running the key-generation algorithm
\(\Gen_0\) using randomness derived from the output of the PRF
\(F(k, s)\). We will assume the use of SHA256 (with a 256-bit digest
length) as a collision-resistant hash function. Using these building
blocks, we will construct a many-time secure signature scheme
\((\Gen, \Sign,
\Ver)\) for arbitrary-length messages (i.e., on message space
\(\zo^\ast\)), where all of the algorithms are efficient
(\(\poly(\cdot)\) running time).

Our construction is as follows:

\begin{tcolorbox}[enhanced jigsaw, arc=.35mm, bottomrule=.15mm, breakable, colback=white, colframe=quarto-callout-note-color-frame, left=2mm, leftrule=.75mm, opacityback=0, rightrule=.15mm, toprule=.15mm]

\begin{itemize}
\tightlist
\item
  \(\Gen() \to (\sk, \vk)\):

  \begin{itemize}
  \tightlist
  \item
    Sample a fresh PRF key \(k \getsr \calK\).
  \item
    Set
    \((\sk_\epsilon, \vk_\epsilon) \gets \Gen_0^{F(k, \text{“”})}()\).
  \item
    Output \((\sk, \vk) \gets (k, \vk_\epsilon)\).
  \end{itemize}
\item
  \(\Sign_t(k, m) \to \sigma\):

  \begin{itemize}
  \tightlist
  \item
    Choose a random 256-bit value
    \(r = (r_1 \dots r_{256}) \in \zo^{256})\).
  \item
    Compute \((\sk_r, \vk_r) \gets \Gen_0^{F(k, \text{r})}()\)
  \item
    Compute \(\sigma_r \gets \Sign_0(\sk_r, \text{SHA256}(m))\).
  \item
    Compute \(\sigma_\varepsilon\), \(\vk_0\), \(\vk_1\),
    \(\sigma_{r_1}\), \(\vk_{r_1 0}\), \(\vk_{r_1 1}\),
    \(\sigma_{r_1 r_2}\), \(\ldots\),
    \(\sigma_{r_1 r_2 \ldots r_{255}}\),
    \(\vk_{r_1 r_2 \ldots r_{255} 0}\),
    \(\vk_{r_1 r_2 \ldots r_{255} 1}\) as shown in
    Figure~\ref{fig-lamport-tree}.
  \item
    Output
    \(\sigma \gets (r, \sigma_\varepsilon, \vk_0, \vk_1, \sigma_{r_1}, \vk_{r_1 0}, \vk_{r_1 1}, \ldots, \sigma_r)\).
  \end{itemize}
\item
  \(\Ver_t(\vk_\epsilon, m, \sigma) \to \zo\):

  \begin{itemize}
  \tightlist
  \item
    Parse
    \((r, \sigma_\varepsilon, \vk_0, \vk_1, \sigma_{r_1}, \vk_{r_1 0}, \vk_{r_1 1}, \ldots, \sigma_r) \gets \sigma\).
  \item
    Output ``1'' if and only if

    \begin{itemize}
    \tightlist
    \item
      \(\Ver_0(\vk_r, \sigma_r, \text{SHA256}(m)) = 1\) and
    \item
      \(\Ver_0(\vk_x, \sigma_x, \vk_{x0} || \vk_{x1}) = 1\) for every
      prefix \(x\) of \(r\) (from \(\varepsilon\) to
      \(r_1 r_2 \ldots r_{255}\)).
    \end{itemize}
  \end{itemize}
\end{itemize}

\end{tcolorbox}

\section{Choosing Signature Schemes}\label{choosing-signature-schemes}

The signature scheme we presented in Section~\ref{sec-sig-manytime} is
not particularly efficient in terms of signature size.

\begin{longtable}[]{@{}
  >{\raggedleft\arraybackslash}p{(\linewidth - 8\tabcolsep) * \real{0.2000}}
  >{\centering\arraybackslash}p{(\linewidth - 8\tabcolsep) * \real{0.1600}}
  >{\centering\arraybackslash}p{(\linewidth - 8\tabcolsep) * \real{0.1733}}
  >{\centering\arraybackslash}p{(\linewidth - 8\tabcolsep) * \real{0.2133}}
  >{\centering\arraybackslash}p{(\linewidth - 8\tabcolsep) * \real{0.2533}}@{}}
\caption{Statistics about various signature schemes used in
practice}\label{tbl-sig-schemes}\tabularnewline
\toprule\noalign{}
\begin{minipage}[b]{\linewidth}\raggedleft
\textbf{Algorithm}
\end{minipage} & \begin{minipage}[b]{\linewidth}\centering
\(\vk\) size
\end{minipage} & \begin{minipage}[b]{\linewidth}\centering
\(\sig\) size
\end{minipage} & \begin{minipage}[b]{\linewidth}\centering
signatures/sec
\end{minipage} & \begin{minipage}[b]{\linewidth}\centering
verifications/sec
\end{minipage} \\
\midrule\noalign{}
\endfirsthead
\toprule\noalign{}
\begin{minipage}[b]{\linewidth}\raggedleft
\textbf{Algorithm}
\end{minipage} & \begin{minipage}[b]{\linewidth}\centering
\(\vk\) size
\end{minipage} & \begin{minipage}[b]{\linewidth}\centering
\(\sig\) size
\end{minipage} & \begin{minipage}[b]{\linewidth}\centering
signatures/sec
\end{minipage} & \begin{minipage}[b]{\linewidth}\centering
verifications/sec
\end{minipage} \\
\midrule\noalign{}
\endhead
\bottomrule\noalign{}
\endlastfoot
\textbf{SPHINCS+-128} & 32 B & 8000 B & 5 & 750 \\
\textbf{RSA 2048} & 256 B & 256 B & 2,000 & 50,000 \\
\textbf{ECDSA256} & 32 B & 64 B & 42,000 & 14,000 \\
\end{longtable}

Many deployed systems today use the ECDSA256 signature scheme. Legacy
application still use RSA signatures, though because of their relatively
large public-key and signature sizes, few new applications use these
schemes.

\marginnote{\begin{footnotesize}

Hashing is much, much faster than signing---the commonly used SHA256
hashing algorithm can compute around 10,000,000 hashes per second. This
is one reason the hash-and-sign paradigm is so useful.

\end{footnotesize}}

\chapter{RSA Signatures}\label{sec-rsa}

In this chapter, we will discuss the RSA digital-signature scheme. The
RSA paper (\citeproc{ref-RSA}{Rivest et al. 1978}) was tremendously
influential because it gave the first constructions of digital
signatures and public-key encryption. (We will talk about public-key
encryption in detail later on.)

The RSA cryptosystem is going out of style for a few reasons: generating
RSA keys is relatively expensive and the keys are relatively large (4096
bits for RSA versus 256 bits for more modern elliptic-curve-based
cryptosystems). In addition, a large-scale quantum computer could---in
theory, at least---break RSA-style cryptosystems.

The RSA cryptosystem is worth studying for a few reasons:

\begin{itemize}
\item
  RSA's security is related to the problem of factoring large integers,
  which is (arguably) the most natural ``hard'' computational problem
  out there.
\item
  RSA gives the only known instantiation of a \emph{trapdoor one-way
  permutation}, which we will define shortly.
\item
  RSA has a number of esoteric properties that are useful for advanced
  cryptographic constructions. For example, it gives a ``group of
  unknown order.'' See
  \href{https://toc.cryptobook.us/book.pdf\#page=436}{Boneh-Shoup,
  Chapter 10.9} for details.
\item
  RSA signatures are used on the vast majority of public-key
  certificates today.\footnote{As of today, around 94\% of certificates
    in the Certificate Transparency logs use RSA signatures:
    \url{https://ct.cloudflare.com/}.}

  The most commonly used type of RSA signatures (``PKCS \#1 v1.5'') is
  more complicated---and no more secure---than the construction we
  describe here, but that construction is still used for historical
  reasons.
\end{itemize}

\section{Trapdoor one-way
permutations}\label{trapdoor-one-way-permutations}

\begin{figure}

\begin{minipage}[t]{0.45\linewidth}

\raisebox{-\height}{

\pandocbounded{\includegraphics[keepaspectratio]{index_files/mediabag/figs/rsa-owf.pdf}}

}

\end{minipage}%
%
\begin{minipage}[t]{0.55\linewidth}

\raisebox{-\height}{

\pandocbounded{\includegraphics[keepaspectratio]{index_files/mediabag/figs/rsa-trapdoor.pdf}}

}

\end{minipage}%

\caption{\label{fig-rsa-owf-trapdoor}A one-way function (at left) is
hard to invert on random inputs. A trapdoor one-way function/permutation
(at right) is a function keyed by a public key. The function is easy to
invert given the secret key and is hard to invert otherwise.}

\end{figure}%

RSA implements a \emph{trapdoor one-way function}. Informally, a
trapdoor one-way function is a function that is easy to compute in the
forward direction but that is hard to invert \emph{except} to someone
knowing a secret key. So it is like a one-way function with a
``trapdoor'' that allows efficient inversion.

RSA actually implements a trapdoor one-way \emph{permutation}---that is,
it maps an input space onto itself with no collisions.

\subsection{Definition}\label{definition}

Formally, a trapdoor one-way permutation over input space \(\calX\) is a
triple of efficient algorithms:

\marginnote{\begin{footnotesize}

If we wanted to be completely formal, the input space would be
parameterized by the security parameter \(\lambda\). So we would have a
family of input spaces \(\{\calX_\lambda\}_{\lambda \in \N}\)---one for
each choice of \(\lambda\). This way the input space can grow with
\(\lambda\).

\end{footnotesize}}

\begin{itemize}
\item
  \(\Gen(1^\lambda) \to (\sk, \pk)\). The key-generation algorithm takes
  as input the security parameter \(\lambda \in \N\), expressed as a
  unary string, and outputs a secret key and a public key.
\item
  \(F(\pk, x) \to y\). The evaluation algorithm \(F\) takes as input the
  public key \(\pk\) and an input \(x \in \calX\), and outputs a value
  \(y \in \calX\).

  \marginnote{\begin{footnotesize}

  In the RSA construction, the input space \(\calX\) depends on the
  public key, but we elide that technical detail here.

  \end{footnotesize}}
\item
  \(I(\sk, y) \to x'\). The inversion algorithm \(I\) takes as input the
  secret key \(\sk\) and a point \(y \in \calX\), and outputs its
  inverse \(x \in \calX\).
\end{itemize}

\textbf{Correctness.} Informally, we want that for keypairs
\((\sk, \pk)\) output by \(\Gen\), we have that \(F(\pk, \cdot)\) and
\(I(\sk, \cdot)\) are inverses of each other. More formally, for all
\(\lambda \in \N\), \((\sk, \pk) \gets \Gen(1^\lambda)\), and
\(x \in \calX\), we require: \[ I(\sk, F(\pk, x))= x.\]

\textbf{Security.} Security requires that \(F(\pk, \cdot)\) is hard to
invert (in the sense of a one-way function) on a randomly sampled input
in the input space \(\calX\), even when the adversary is given the
public key \(\pk\). That is, for all efficient adversaries \(\calA\),
there exists a negligible function \(\negl(\cdot)\) such that
\[ \Pr\left[
\calA(\pk, F(\pk, x)) = x
\colon \begin{aligned}
  (\sk, \pk) &\gets \Gen(1^\lambda)\\
  x &\getsr \calX
\end{aligned}\right] \leq \negl(\lambda).\]

When we use the RSA cryptosystem, we make the assumption that the RSA
function is hard to invert given only the public key:

\begin{definition}[RSA
Assumption]\protect\hypertarget{def-rsa-assumption}{}\label{def-rsa-assumption}

The RSA function \((\Gen, F, I)\) is a trapdoor one-way permutation.

\end{definition}

\textbf{IMPORTANT:} Just as a one-way function is only hard to invert on
a \emph{randomly sampled input}, a trapdoor one-way function is only
hard to invert on a randomly sampled input. Many of the cryptographic
failures of RSA come from assuming that the RSA one-way function is hard
to invert on non-random inputs.

\subsection{Digital signatures from trapdoor one-way
permutations}\label{digital-signatures-from-trapdoor-one-way-permutations}

This construction is called ``full-domain hash.''
(\citeproc{ref-BR93}{Bellare and Rogaway 1993}) The construction makes
use of a hash function \(H\) and resulting signature scheme is secure,
provided that we model the hash function \(H\) as a ``random oracle.''

In other words, to argue security, it is not sufficient to show that
\(H\) is, for example, collision resistant. Instead, we can only prove
security provided that we pretend that \(H\) is a truly random
function---i.e., in the random-oracle model. When we instantiate the
hash function \(H\) with some concrete cryptographic hash function, such
as SHA256, we hope that the resulting signature scheme is still secure.
In practice, this approach works quite well.

One way to think about it is that if a signature scheme is secure in the
random-oracle model, then the concrete signature scheme is in some sense
secure against attacks that do not exploit the peculiarities of the hash
function.

In the construction, we use:

\begin{itemize}
\tightlist
\item
  a trapdoor one-way permutation \((\Gen, F, I)\), and
\item
  a hash function \(H \colon \zo^* \to \calX\), which we model as a
  random oracle in the security analysis.
\end{itemize}

\textbf{Construction.} We construct a digital-signature scheme
\((\Gen, \Sign, \Ver)\) as follows:

\begin{itemize}
\tightlist
\item
  \(\Gen\) -- Just run the key-generation algorithm for the trapdoor
  one-way permutation.
\item
  \(\Sign(\sk, m) \to \sigma\). Hash the message down to an element
  \(h\) of the input space \(\calX\) of the trapdoor one-way permutation
  using the hash function \(H\). Then invert the trapdoor one-way
  permutation at that point:

  \begin{itemize}
  \tightlist
  \item
    Compute \(h \gets H(m)\).
  \item
    Output \(\sigma \gets I(\sk, h)\).
  \end{itemize}
\item
  \(\Ver(\pk, m, \sigma) \to \zo.\)

  \begin{itemize}
  \tightlist
  \item
    Compute \(h' \gets H(m)\).
  \item
    Accept if \(F(\pk, \sigma) = h'\).
  \end{itemize}
\end{itemize}

Notice that the use of a hash function here is \textbf{critical} to
security, since (in the random oracle) it means that forging a signature
is as hard as inverting \(F\) on a random point in its co-domain.
Without the hash function, forging a signature is only as hard as
inverting \(F\) on an attacker-chosen point in its co-domain, which
could be easy.

\textbf{Correctness.} For all \(\lambda \in \N\),
\((\sk, \pk) \gets \Gen(1^\lambda)\), and \(m \in \zo^*\), we have:
\[\begin{aligned}
\Ver(\pk, m, \Sign(\sk, m)) &= 1\{ F(\pk, I(\sk, H(m))) = H(m) \}\\
&= 1\{ I(\sk, F(\pk, I(\sk, H(m)))) = I(\sk, H(m)) \}
\end{aligned}\] and by correctness of the trapdoor one-way permutation:
\[\begin{aligned}
&= 1\{ I(\sk, H(m)) = I(\sk, H(m)) \} = 1.
\end{aligned}\]

\textbf{Security.} The intuition here is that if the adversary cannot
invert \(F\), it cannot find the preimage of \(H(m)\) under \(F\) for
any message on which it has not seen a signature. See
\href{https://toc.cryptobook.us/book.pdf\#page=550}{Boneh-Shoup Chapter
13.3} for the full security analysis.

\section{The RSA construction: Forward
direction}\label{the-rsa-construction-forward-direction}

The algorithms for key-generation and for evaluating the RSA permutation
in the forward direction are not too complicated.

In what follows, we present RSA with public exponent \(e=5\). The same
construction works with many other choices of \(e\), just by replacing
all of the ``5''s below with some other small prime: 3, 7, 13, etc. A
popular choice of the public exponent \(e\) in practice is
\(e=2^{16}+1\). The complexity of computing the RSA function in the
forward direction scales with the size of \(e\), so we prefer small
choices of \(e\).

\begin{itemize}
\item
  \(\Gen(1^\lambda) \to (\sk, \pk)\).

  \marginnote{\begin{footnotesize}

  In practice, we usually take the bitlength of primes to be
  \(\lambda=1024\) or \(\lambda = 2048\).

  \end{footnotesize}}

  \begin{itemize}
  \item
    Sample two random \(\lambda\)-bit primes \(p\) and \(q\) such that
    \(p \equiv q \equiv 4 \pmod 5\).

    \marginnote{\begin{footnotesize}

    Standard RSA implementations require the weaker condition that the
    public exponent \(e\) shares no prime factors with \(p-1\) and
    \(q-1\). Using the stronger condition here simplifies the inversion
    algorithm.

    \end{footnotesize}}
  \item
    Set \(N \gets p \cdot q\).
  \item
    Output \(\sk \gets (p, q)\), and \(\pk = N\).
  \end{itemize}
\item
  \(F(\pk = N, x) \to y\).

  \begin{itemize}
  \tightlist
  \item
    The input space for the RSA function is \(\calX = \Z^*_N\)---the set
    of elements in \(\{0, 1, 2, \dots, N-1\}\) relatively prime to
    \(N\).
  \item
    Output \(y \gets x^5 \bmod N\).
  \end{itemize}
\end{itemize}

\begin{refremark}
The key-generation algorithm relies on us being able to sample large
random primes. One perhaps surprising fact is that there are many many
large primes. In particular, if you pick a random \(\lambda\)-bit
number, the probability that it is prime is roughly \(1/\lambda\).

\marginnote{\begin{footnotesize}

For more on this, look up the Prime Number Theorem.

\end{footnotesize}}

We can sample a random \(\lambda\)-bit prime by just picking random
integers in the range \([2^\lambda, 2^{\lambda + 1})\) until we find a
prime. We can test for primality in \(\approx \lambda^4\) time using the
Miller-Rabin primality test. We also need that there are infinitely many
primes congruent to \(4 \bmod 5\), but fortunately there are. Generating
RSA keys is expensive---it can take a few seconds even on a modern
machine.

\label{rem-sampling-primes}

\end{refremark}

Notice that computing the RSA function in the forward direction is
relatively fast: it just requires three multiplications modulo a
2048-bit number \(N\). That is, to compute \(x^5 \bmod N\), we compute:
\[ (x^2)^2 \cdot x = x^5\mod N.\]

Before describing the RSA inversion algorithm, we discuss why the RSA
trapdoor one-way permutation should be hard to invert without the secret
key.

\subsection{Why should the RSA function be hard to
invert?}\label{why-should-the-rsa-function-be-hard-to-invert}

To invert the RSA function, the attacker's is effectively given a value
\(y \getsr \Z_N\) and must find a value \(x\) such that
\(x^5 = y \bmod N\). Or, put another way, the attacker's task is
essentially the following:

\begin{itemize}
\tightlist
\item
  \textbf{Given:} A polynomial \(p(X) \deq X^5 - y \in \Z_N[X]\), for
  \(y \getsr \Z_N\).
\item
  \textbf{Find:} A value \(x \in \Z_N\) such that \(p(x) = 0 \in \Z_N\).
\end{itemize}

So the attacker must find the root of a polynomial modulo a composite
integer \(N\).

The premise of RSA-style cryptosystems is that we only know of
essentially two ways to find roots of polynomials modulo \(N\):

\begin{itemize}
\item
  \textbf{Factor \(N\) into primes} and find a root modulo each of the
  primes. (We will say more on this in a moment.) Since the best
  algorithms for factoring run in time roughly
  \(2^{\sqrt[3]{\log N}} = 2^{\sqrt[3]{\lambda}}\), this approach is
  infeasible at present without knowing the factorization of \(N\).

  \marginnote{\begin{footnotesize}

  In Appendix~\ref{sec-fact} we present a factoring algorithm that runs
  in sub-exponential time \(2^{\sqrt{\log N \log \log N}}\).

  \end{footnotesize}}
\item
  \textbf{Find a root over the integers} and reduce it modulo \(N\).

  \marginnote{\begin{footnotesize}

  Actually, it suffices to find a root over the rational numbers, but
  the distinction isn't important here.

  \end{footnotesize}}

  For example, it is easy to find a root of polynomials such as:
  \[\begin{aligned}
    X + 4 &= 3 &&\mod N,&\\
    X + 2Y &= 5 &&\mod N,&\\
    X^2 &= 9 &&\mod N\text{, and}&\\
    X^2-3x+2 &= (X-2)(X-1) = 3 &&\mod N.
  \end{aligned}\]

  \marginnote{\begin{footnotesize}

  There are many clever attacks for solving polynomial equations modulo
  composites that work in certain special cases, but for most purposes
  these are the two known attacks.

  \end{footnotesize}}

  When \(y \getsr \Z_N\), the probability that \(y\) is a perfect 5-th
  power, and thus that there is an integral root to \(X^5 - y\), is
  \(\sqrt[5]{N}/N \approx 2^{-4\lambda/5}\), which is negligible in the
  security parameter \(\lambda\). So solving this equation over the
  integers is a dead end.

  \textbf{Is inverting the RSA function as hard as factoring the
  modulus?} No one knows---the question has been open since the
  invention of RSA. We do know that finding roots of certain polynomial
  equations, such as \(p(X) \deq X^2 - y \bmod N\) for random
  \(y \getsr \Z_N\) \emph{is} as hard as factoring the modulus \(N\).
  But for RSA-type polynomials, the answer is unclear.
\end{itemize}

\section{The RSA construction: Inverse
direction}\label{the-rsa-construction-inverse-direction}

To understand how the inversion algorithm works, we will need some
number-theoretic tools.

\subsection{Tools from number theory}\label{tools-from-number-theory}

For a natural number \(N\), let \(\phi(N)\) denote the number of
integers in \(\Z_N = \{1, 2, 3, \dots, N\}\) that are relatively prime
to \(N\). When \(p\) is prime \(\phi(p) = p-1\). The function
\(\phi(\cdot)\) is called \emph{Euler's totient function}.

\marginnote{\begin{footnotesize}

Two natural numbers are \emph{relatively prime} if they share no prime
factors.

\end{footnotesize}}

When \(N=pq\) is the product of two distinct primes,
\(\phi(N) = (p-1)(q-1)\). That is so because all numbers in \(\Z_N\) are
relatively prime to \(N\) except \(N\) and the multiples of \(p\) and
\(q\): \[ p, 2p, 3p, \dots, (q-1)p,\ \ \ q, 2q, 3q, \dots, (p-1)q.\] So
there are \(N - (q-1) - (p-1) - 1 = (p-1)(q-1)\) numbers in \(\Z_N\)
relatively prime to \(N\).

\begin{theorem}[Euler's
Theorem]\protect\hypertarget{thm-euler}{}\label{thm-euler}

Let \(N\) be a natural number. Then for all \(a \in \Z^*_N\),
\[ a^{\phi(N)} = 1 \mod N.\]

\end{theorem}

\begin{proof}
Consider the sets \(\Z^*_N\) and \(\{ax \bmod N \mid x \in \Z^*_N\}\).
These sets are equal, so the product of the elements in the two sets is
equal. Let \(X \gets \prod_{x \in \Z^*_N} x \bmod N\). Then
\[ X = a^{\phi(N)}X \bmod N \quad \Rightarrow\quad 1 = a^{\phi(N)} \bmod N.\]
\end{proof}

\begin{lemma}[]\protect\hypertarget{lem-inv}{}\label{lem-inv}

Let \(p\) and \(q\) be distinct primes congruent to \(4\) modulo \(5\).
Define the integer \(d = \frac{\phi(N) - 4}{5} + 1\). Then
\(5d \equiv 1 \bmod \phi(N)\).

\end{lemma}

\begin{proof}
Observe that
\[ p \equiv 4 \bmod 5 \quad \Rightarrow \quad \phi(N) - 4 \equiv 0 \bmod 5,\]
so \(\frac{\phi(N) - 4}{5}\) is an integer and thus \(d\) is well
defined. Then \(5 d = \phi(N) - 4 + 5 = 1 \bmod \phi(N)\).
\end{proof}

\subsection{Inverting the RSA
function}\label{inverting-the-rsa-function}

With the number theory out of the way, we can now describe how to invert
the RSA function. All we have to do is to show how to compute a fifth
root of \(y \bmod N\).

\begin{itemize}
\tightlist
\item
  \(I(\sk, y) \to x\).

  \begin{itemize}
  \tightlist
  \item
    The secret key \(\sk\) consists of the prime factors \(p,q\) of
    \(N\). Recall that \(\phi(N) = (p-1)(q-1)\).
  \item
    Compute the integer \(d \gets \frac{\phi(N) - 4}{5} + 1\), as in

    \begin{enumerate}
    \def\labelenumi{\arabic{enumi}.}
    \setcounter{enumi}{1}
    \tightlist
    \item
    \end{enumerate}

    \marginnote{\begin{footnotesize}

    We sometimes call \(d\) the \emph{private exponent} in RSA.

    \end{footnotesize}}
  \item
    Return \(y^d \bmod N\).
  \end{itemize}
\end{itemize}

It is not obvious why the inversion algorithm is correct. Say that
\(y = x^5 \bmod N\). Then: \[\begin{aligned}
  y^d &= (x^5)^d &&\bmod N\\
      &= x^{5d} &&\bmod N\\
      &= x^{k \cdot \phi(N) + 1} &&\bmod N,
\end{aligned}\] for some \(k \in \Z\), by 2, and \[\begin{aligned}
      &= x \cdot (x^{\phi(N)})^k &&\bmod N\\
      &= x &&\bmod N,
\end{aligned}\] by Theorem~\ref{thm-euler}. We could write
\(5d = k \phi(N) + 1\) because from 2, we know that
\(5d \equiv 1 \bmod \phi(N)\).

\textbf{Using other public exponents.} For our RSA-inversion algorithm
to work, we need only to compute the multiplicative inverse \(e\) modulo
\(\phi(N)\). That is, we need to compute an integers \(d\) such that
\(ed \equiv 1 \bmod \phi(N)\). Such an inverse always exists when \(e\)
and \(\phi(N)=(p-1)(q-1)\) are relatively prime. RSA implementations
typically use the extended Euclidean algorithm to compute the
multiplicative inverse of \(e\) modulo \(\phi(N)\). That algorithm is
more general, but the one we used in 2 is simpler and is self-contained.

\textbf{Inverting RSA is easy on a negligible fraction of points.}
Recall the RSA is If the preimage under the RSA function of a point
\(y\) is very very small, then If \(x < N^{1/5}\), then computing \(x\)
given \(y = x^5 \bmod N\) is \emph{easy}.

\textbf{Is inverting RSA as hard as factoring the modulus \(N\)?} The
inversion algorithm we showed here requires knowing the prime factors of
the modulus \(N\). Inverting RSA is thus \emph{no harder than} factoring
\(N\).

Is inverting RSA \emph{as hard as} factoring \(N\)? In particular, if we
have an efficient algorithm \(\calA\) that inverts RSA, can we use
\(\calA\) to factor the modulus \(N\)? No one knows!

Most cryptographers, I would guess, believe that inverting the RSA
function is as hard as factoring. But for all we know, it could be that
computing fifth roots modulo \(N\) is \emph{easier} than factoring the
modulus.

\chapter{Public-key Infrastructure}\label{public-key-infrastructure}

In the last chapter, we discussed digital signatures, which allow us to
authenticate messages without a shared secret. For example, if I have
the public signature-verification key of the university dean, I can
verify that signed emails from the dean really came from her and not
from someone pretending to be her. But to verify the signature on the
dean's message, I need to know her signature-verification key \(\vk\).
How can I (the recipient) obtain this verification key without a secure
channel to the dean (the sender)?

Unfortunately, there are no perfect solutions to this problem. In this
section, we will discuss some of the approaches that we use in practice.

\section{Public-key infrastructure}\label{public-key-infrastructure-1}

The goal of a public-key infrastructure is to facilitate the mapping of
\textbf{human-intelligible names} to \textbf{signature-verification
keys}. Examples of human-intelligible names that we map to keys are:
email addresses, domain names, legal entities, phone numbers, and
usernames (e.g., within a company).

We can think of the public-key infrastructure as implementing the
following (grossly simplified) API:
\[ \mathsf{IsKeyFor}(\vk, \mathtt{name}) \to \zo.\] That is, given a
verification key \(\vk\) and a name \(\mathtt{name}\), the public-key
infrastructure gives a way to check whether this mapping is valid.

We now discuss some ways to implement this API.

\section{Option 1: Use verification keys as
names}\label{option-1-use-verification-keys-as-names}

One option is to just refer to everyone by the bytes of their
signature-verification key. This way, there is no need to do a messy
name-to-verification-key translation at all.

This is not practical for humans generally: it would be difficult to
remember your friends' names if you had to call them by random 32-byte
strings! However, some digital services such as Bitcoin indeed use
verification keys as identities: when you want to transfer Bitcoin to
someone, you send the coins to an account identified by their public
key. The public key \emph{is} name of that account.

Using keys as names has a two major problems:

\begin{itemize}
\tightlist
\item
  Verification keys are hard to remember. Things like email addresses,
  domain names, kerberos usernames, phone numbers, and so on, are much
  easier for humans to remember.
\item
  There is no way to update the name-to-key mapping. If you lose the
  secret key associated with your name/account, there is no way to
  ``update'' the key to a new value. In practice, people lose their
  secret keys all the time, so supporting key updates is critical in
  most systems.
\end{itemize}

\section{Trust on first use (TOFU)}\label{trust-on-first-use-tofu}

Another strategy is to avoid having any global mapping from names to
verification keys. Instead, a client can just accept the first
verification key that it sees associated with a given name. The secure
shell system (\texttt{SSH}) uses TOFU for key management by default.

In particular, the key-validation logic looks like this:

\begin{tcolorbox}[enhanced jigsaw, arc=.35mm, bottomrule=.15mm, breakable, colback=white, colframe=quarto-callout-note-color-frame, left=2mm, leftrule=.75mm, opacityback=0, rightrule=.15mm, toprule=.15mm]

\(\mathsf{keymap} \gets \{ \}\).

\(\mathsf{IsKeyFor}(\vk, \mathtt{name}):\)

\begin{itemize}
\tightlist
\item
  If \(\mathsf{keymap}[\mathtt{name}]\) is undefined:

  \begin{itemize}
  \tightlist
  \item
    Set \(\mathsf{keymap}[\mathtt{name}] \gets \vk\).
  \item
    Return true..
  \end{itemize}
\item
  Else: Return \(\mathsf{keymap}[\mathtt{name}] == \vk\).
\end{itemize}

\end{tcolorbox}

That is, the client will accept the \emph{first} verification key it
sees associated with a particular name. Later on, the client will only
accept the same verification key for that name.

TOFU is very simple to implement and provides a meaningful security
guarantees. There are two drawbacks:

\begin{itemize}
\tightlist
\item
  If the first key that client receives for a particular name is
  incorrect/attacker-generated, the attacker can forge signatures under
  that name.
\item
  It is not clear with TOFU how to handle key updates. In most systems
  that use TOFU, whenever the sender's key changes, the system notifies
  the user and allows them to accept or reject the new key. The burden
  is then on the user to figure out whether the sender really did change
  their signing keypair, or whether there is an attack in progress.
\end{itemize}

\section{Certificates}\label{certificates}

Another option is to rely on a few parties to manage the mapping of
names to public keys. These entities are called \emph{Certificate
Authorities} (CAs). Your operating system and web browser typically come
bundled with a set of roughly 100 public signature-verification keys,
owned by each of 100 CAs.

Whenever the owner of website \texttt{example.com}, for example,
generates a new public key \(\vk_\mathtt{example.com}\), the website
owner can ask a certificate authority to certify that
\(\vk_\mathtt{example.com}\) really belongs to \texttt{example.com}. The
certificate authority does this by signing the pair
\((\mathtt{example.com}, \vk_\mathtt{example.com})\) using its own
signing keypair \(\vk_\text{CA}\) to generate a signature
\(\sigma_\text{CA}\). This signed attestation
\((\mathtt{example.com}, \vk_\mathtt{example.com}, \sigma_\text{CA})\)
is called a \emph{certificate}.

\marginnote{\begin{footnotesize}

In practice the structure of certificates is much more complicated than
we are showing here, and include all sorts of additional metadata. But
the basic idea is the same as we describe here.

\end{footnotesize}}

When a client connects to \texttt{example.com}, the server at
\texttt{example.com} will supply the client with the certificate
\((\mathtt{example.com}, \vk_\mathtt{example.com},\sigma_\text{CA})\).
As long as this certificate was signed by a CA that the client trusts
(i.e., a CA for which the client has a verification key), the client can
validate the certificate and conclude that the verification key
\(\vk_\mathtt{example.com}\) really belongs to \texttt{example.com}.

In pseudocode the logic for verifying certificates looks like this:

\begin{tcolorbox}[enhanced jigsaw, arc=.35mm, bottomrule=.15mm, breakable, colback=white, colframe=quarto-callout-note-color-frame, left=2mm, leftrule=.75mm, opacityback=0, rightrule=.15mm, toprule=.15mm]

\(\mathsf{caKeys} \gets \{\vk_\text{Verisign}, \vk_\text{Let's Encrypt}, \dots \}\).

\(\mathsf{IsKeyFor}((\vk, \sigma), \mathtt{name}):\)

\begin{itemize}
\tightlist
\item
  For each \(\vk_\text{CA}\) in \(\mathsf{caKeys}\):

  \begin{itemize}
  \tightlist
  \item
    If \(\Ver(\vk_\text{CA}, (\mathtt{name}, \vk), \sigma) = 1\): Return
    true.
  \end{itemize}
\item
  Return false.
\end{itemize}

\end{tcolorbox}

A very nice feature of certificate-based public-key infrastructure is
that the client does not need to communicate with the CA to validate a
name-to-key mapping. The client only needs to perform one
signature-verification check.

\marginnote{\begin{footnotesize}

In order to use TLS on a website you own, you need to convince one of
the certificate authorities to give you a certificate---i.e., to sign
your \((\mathtt{name}, \vk)\) pair. To do so, the CA will have some
protocol to follow---typically, you will send your
\((\mathtt{name}, \vk)\) pair to the CA, who will then ask you to verify
that you own the name somehow. This is called \emph{domain validation}.
In the case of web certificates, the CA may verify ownership by
requiring you to upload a file file to your server, to add a new DNS
record with a random value, or something similar. Once the CA is
convinced that you own the domain, the CA will reply with a certificate:
a signature over the tuple \((\mathtt{name}, \vk)\). This
\((\mathtt{name}, \vk, \sig_{\text{CA}})\).

\end{footnotesize}}

Certificates in practice works quite well:

\begin{itemize}
\tightlist
\item
  The client only needs to store \(\approx 100\) CA verification keys,
  and yet the client can validate the name-to-key mappings for millions
  of websites.
\item
  A client can choose which CAs to trust (though in practice, clients
  typically delegate this decision to software vendors).
\item
  The client never needs to interact with the CA.
\end{itemize}

However, certificates still have some drawbacks:

\begin{itemize}
\tightlist
\item
  If an attacker compromises \emph{any} CA, they can generate
  certificates for any domain.
\item
  Certificate authorities often perform quite minimal validation of
  domain ownership.
\item
  If a server's private key gets stolen, there is no great plan for
  \emph{revoking} or updating a name-to-key mapping.
\end{itemize}

\subsection{Domain Validation}\label{domain-validation}

Domain Validated (DV) certificates are quite common. CA validation uses
a simple technical check for domain ownership. For example, the CA asks
the requester amazon.com to put a nonce in a file on an amazon.com
server. Then, the CA retrieves that file, and checks the nonce. The
(reasonable) assumption is that only the server owner could have created
such a file. Alternately, the CA could ask the requester to create
certain Domain Name Services (DNS) records under amazon.com, e.g.,
acme-challenge.amazon.com.

This check is often automated, and is fast, convenient, cheap. An
example DV CA is \textbf{Let's Encrypt}. It is easy to use and free; the
requester doesn't have to be a real company.

DV certificates only have low-level DNS ownership guarantees: ``CA
validated that owner of certificate owns amazon.com''. However, the
guarantee closely matches what CAs can be expected to validate in real
life.

A standard protocol for validating domain ownership is ACME\footnote{\url{https://tools.ietf.org/html/rfc8555\#section-8.3}}.

\subsection{Revocation}\label{revocation}

In many cases, a CA will want to delete or change a name-to-key mapping.
This process is called \emph{certificate revocation}. There are several
possible reasons for this:

\begin{itemize}
\tightlist
\item
  The owner of a verification key may have their corresponding secret
  key be lost or stolen.
\item
  A company may want to rotate keys, for example to update to a new
  cryptographic algorithm.
\item
  A website may go out of business and another entity buys their domain
  name.
\item
  Software bugs may lead a user to generate an insecure keypair that
  they later want to revoke. (\citeproc{ref-yilek2009private}{Yilek et
  al. 2009}; \citeproc{ref-nemec2017return}{Nemec et al. 2017})
\end{itemize}

In a scheme that uses certificates, this seems like a hard problem:
since there is no interaction with the CA to verify a certificate, there
is no way for a CA to ``take back'' a certificate. There are again no
excellent solutions to this, but there are a few strategies used in
practice.

\textbf{Expiration Dates.} One pragmatic way to handle revocation is to
add an expiration date to each generated certificate---if this
expiration date has passed, the client will reject the certificate. This
way, a server will need to re-authenticate to the CA that they own the
name that they claim to own periodically. So, for example, an attacker
that steals a website's secret key will only be able to use it until the
certificate expires. In practice, certificates used on the Internet
typically expiration dates between 90 days and 1-2 years.

\textbf{Software Updates.} Another solution is for the browser (or
client, more generally) to maintain a list of revoked certificates. On
every connection, the browser will check whether the provided
certificate is in this local revocation list and refuse it if so. Since
browsers today check for updates very frequently, this strategy can
respond to a stolen secret key quickly. However, there is a large
storage cost since now every browser needs to store this (potentially
large) list of revoked certificates.

\textbf{CA Revocation List.} To avoid depending on browser manufacturers
to update this revocation list, another strategy is to ask the CA for it
directly. One way to do this is similar to the above: periodically query
the CA to download its updated revocation list and check that each
certificate is not in this list. This method has fallen out of favor, in
part because clients (e.g., behind corporate firewalls) cannot connect
directly to the CAs to download these revocation lists.

\part{Transport Security}

\chapter{Introduction to Encryption}\label{introduction-to-encryption}

So far, we have explored methods for authenticating data---verifying
that it has not been modified in transit. However, the
integrity-protection mechanisms we have discussed provide no
\emph{confidentiality}: a network eavesdropper can still view everything
that we send over the network.

In this chapter, we will discuss \emph{encryption}, which allow two
parties to exchange messages over an insecure network while hiding the
contents of their communications.

We will cover encryption in a sequence of steps:

\begin{itemize}
\tightlist
\item
  First, we will construct an encryption scheme with a \emph{weak form
  of security} for \emph{fixed-length messages} for settings in which
  the sender and recipient have a \emph{shared secret key}.
\item
  Next, we will show how to extend this scheme to support
  \emph{variable-length messages}.
\item
  Then, we will show how to improve the scheme to have a \emph{strong
  form of security}.
\item
  Finally, we will show how to implement encryption in settings in which
  the sender and recipient \emph{have no shared secrets}.
\end{itemize}

At the conclusion of this part, we will discuss how deployed systems use
encryption, and we will think about some problems that encryption does
not solve.

\section{Background}\label{background}

\textbf{The need for encryption.} The internet is a massive network of
wifi access points, routers, switches, undersea cables, DNS servers, and
much more. There are many, many devices for a potential adversary to
compromise and many vantage points from which an attacker can observe
network traffic. Every single hop your packets take is a potential point
of compromise.

To make matters worse, most standard network protocols provide \emph{no
authentication or encryption}: in Ethernet, IP, DNS, email, HTTP, and
others, an adversary is free to modify and read the traffic we send and
receive. Protecting confidentiality typically requires augmenting these
standard network protocols with some form of authentication and
encryption.

\textbf{Systems using encryption.} Encryption shows up in a large number
of deployed systems. A few examples are:

\begin{itemize}
\tightlist
\item
  \textbf{Messaging apps}, such as WhatsApp, Signal, and iMessage,
  encrypt traffic between app users so that the server cannot easily
  read it.
\item
  \textbf{Network protocols}, such as SSH and HTTPS use encryption to
  protect traffic between a service's clients and servers.
\item
  \textbf{File-storage systems} use encryption to protect data at rest.
  So if a thief steals your laptop, they will not easily be able to read
  the encrypted files on your hard disk.
\end{itemize}

\section{Encryption Scheme Syntax}\label{encryption-scheme-syntax}

Encryption schemes are defined with respect to a key space \(\calK\), a
message space \(\calM\), and a ciphertext space \(\calC\). For now,
think of \(\calK = \calM = \zo^n\) and \(\calC = \zo^{2n}\), for a
security parameter \(n\).

\marginnote{\begin{footnotesize}

Remember that in practice, we will often take the size of the keyspace
\(\abs{\calK}\) to be at least \(2^{128}\) to prevent brute-force
key-guessing attacks.

\end{footnotesize}}

An encryption scheme then consists of two algorithms:

\begin{itemize}
\item
  \(\Enc: \calK \times \calM \rightarrow \calC\)
\item
  \(\Dec: \calK \times \calC \rightarrow \calM\)

  \marginnote{\begin{footnotesize}

  As we will see later on, the decryption algorithm \(\Dec\) can
  sometimes output ``FAIL'' or \(\bot\).

  \end{footnotesize}}
\end{itemize}

The definition of correctness for an encryption scheme just states that
if you encrypt a message \(m\) with a key \(k\), and then you decrypt
the resulting ciphertext with the same key \(k\), you end up with the
same message \(m\) that you started with:

\begin{definition}[Encryption Scheme,
Correctness]\protect\hypertarget{def-enc-correctness}{}\label{def-enc-correctness}

An encryption scheme is correct if, for all keys \(k \in \calK\) and all
messages \(m \in \calM\), \(\Dec(k, \Enc(k, m)) = m\).

\end{definition}

Defining security for encryption scheme is tricky business. Many of the
most obvious security definitions are insufficient:

\begin{itemize}
\tightlist
\item
  \textbf{Bad definition:} \emph{``An encryption scheme is secure if it
  is infeasible for an attacker to recover the plaintext message given
  only a ciphertext.} This definition admits encryption schemes in which
  the ciphertext leaks half of the plaintext bits.
\item
  \textbf{Bad definition:} \emph{``An encryption scheme is secure if it
  is infeasible for an attacker to recover any bit of the plaintext
  message given only a ciphertext.} This definition admits encryption
  schemes in which the ciphertext leaks the parity of the plaintext
  bits.
\end{itemize}

The starting-point (weak) security definition we will use is called
\emph{indistinguishability under adaptive chosen plaintext attack}
(\(\INDCPA\)). Intuitively, a scheme is CPA-secure if an attacker cannot
tell which of two chosen messages are encrypted, even after seeing
encryptions of many attacker-chosen messages.

\begin{definition}[Encryption Scheme, CPA Security
(weak)]\protect\hypertarget{def-cpa}{}\label{def-cpa}

Formally, an encryption scheme \((\Enc, \Dec)\) over message space
\(\calM\) and key space \(\calK\) is CPA-secure if all efficient
adversaries win the following game with probability at most
\(\tfrac{1}{2}\) + ``negligible:''

\begin{itemize}
\item
  The challenger samples \(b \rgets \bin\) and \(k \rgets \calK\).
\item
  Polynomially many times: \emph{// Chosen-plaintext queries}

  \begin{itemize}
  \tightlist
  \item
    The adversary sends the challenger a message \(m_i \in \calM\)
  \item
    The challenger replies with \(c_i \gets \Enc(k, m_i)\).
  \end{itemize}
\item
  The adversary then sends two messages \(m^*_0, m^*_1 \in \calM\) to
  the challenger. (We require \(\abs{m^*_0} = \abs{m^*_1}\).)

  \marginnote{\begin{footnotesize}

  Standard encryption systems do not hide the length of the message
  being encrypted. So, if the message space \(\calM\) contains messages
  of different lengths, our security definition requires the adversary
  to distinguish the encryption of two messages of the \emph{same}
  length.

  \end{footnotesize}}
\item
  The challenger replies with \(c^* \gets \Enc(k, m^*_b)\).
\item
  The adversary outputs a value \(b' \in \bin\). The adversary wins if
  \(b=b'\).
\end{itemize}

\end{definition}

One potentially surprising consequence of the CPA-security definition
for encryption schemes is:

\emph{For an encryption scheme to be secure in any meaningful sense, the
encryption algorithm \textbf{must} be randomized}.

\marginnote{\begin{footnotesize}

In contrast, secure MACs can be---and typically are---deterministic!

\end{footnotesize}}

If the encryption algorithm is deterministic (i.e., not randomized), an
attacker can win the CPA security game in Definition~\ref{def-cpa} with
probability 1. To do so:

\begin{itemize}
\tightlist
\item
  The attacker first asks for encryption of a message \(m_0\) and
  receives a ciphertext \(c_0\).
\item
  Then, the attacker attacker choose a message \(m_1 \neq m_0\) and
  sends \((m^*_0, m^*_1) = (m_0, m_1)\) to the challenger and receives
  the challenge ciphertext \(c^*\).
\item
  If \(c^* = c_0\), the attacker outputs 0. Otherwise, the attacker
  outputs 1.
\end{itemize}

Deterministic encryption schemes are not only broken in theory, they are
also broken in practice. For example, if you encrypt the pixels of an
image using a deterministic encryption scheme, the encrypted image
essentially reveals the plaintext image.

\textbf{Why CPA security is a ``weak'' form of security.} There are two
reasons why CPA security is a weak or insufficient definition of
security for an encryption scheme:

\begin{itemize}
\tightlist
\item
  First, the CPA security definition guarantees nothing about message
  \emph{integrity}: an attacker can modify the ciphertext and
  potentially change the meaning of the encrypted message in a
  meaningful way (even if the attacker does not know the encrypted
  message!).
\item
  Second, the CPA security definition guarantees nothing in the event
  that the adversary can obtain \emph{decryptions} of ciphertexts of its
  choosing. In practice, attackers can often obtain decryptions of
  chosen ciphertexts. for example, in a system where a client sends
  encrypted messages to a server and the server does something in
  response, an attacker can send encrypted queries to the server and
  observe its behavior to learn some function of the decrypted contents
  of the message. Later on, we will expand our definition to include
  these \emph{chosen ciphertext attacks}.
\end{itemize}

\section{One-time Pad}\label{one-time-pad}

The one-time pad is perhaps the simplest encryption scheme. Its
keyspace, message space, and ciphertext space are all the set of
\(n\)-bit strings. The algorithms are then:

\begin{itemize}
\tightlist
\item
  \(\Enc(k, m) \to c\). Compute \(c \gets (k \oplus m)\).
\item
  \(\Dec(k, c) \to m\). Compute \(m' \gets (k \oplus c)\).
\end{itemize}

The encryption scheme is correct since
\(\Dec(k, \Enc(k, m)) = k \oplus (k \oplus m) = m\). A less obvious fact
is that it is also \emph{one-time secure}: if an attacker sees only one
message with encrypted a one-time-pad key \(k\), it learns nothing about
the underlying plaintext.

\marginnote{\begin{footnotesize}

The one-time pad is actually one-time secure in a very strong sense: it
protects confidentiality even against a computationally unbounded
attacker---one that can perform arbitrary amounts of computation.

\end{footnotesize}}

\textbf{The ``two-time pad'' attack.} The one-time pad is insecure if
the same key \(k\) is every used to encrypt two messages. In particular,
if two ciphertexts are ever computed using the same key, we have:
\[\begin{aligned}
  c_1 &= k \xor m_1 \\
  c_2 &= k \xor m_2 \\
  c_1 \oplus c_2 &= (k \xor m_1) \xor (k \xor m_2) = m_1 \xor m_2.
\end{aligned}\] The attacker then learns the XOR of the two encrypted
messages, which is often enough to leak all sorts of sensitive
information about the plaintext. For example, if the attacker knows some
bits of \(m_1\), it can learn some bits of \(m_2\).

\textbf{Why the one-time pad is useful.} The one-time pad encryption
scheme feels in some sense useless: if a secure channel exists through
which Alice and Bob can exchange an \(n\)-bit secret key, they may as
well use that channel to exchange the message itself! There is some
merit still in the one-time pad: it is possible to exchange the key
ahead of time, and then send encrypted messages later on. Diplomats
indeed used the one-time pad in this way throughout the 20th century.
They would exchange huge books of keying material (random strings) and
then use these keys to communicate securely over long distances.

In practice, it is much more convenient to be able to exchange a
\emph{short} key and then use it to encrypt many \emph{long} messages.

\section{A Weak Encryption Scheme}\label{sec-lec8-weakenc}

What we effectively need is a way to generate many bits of
random-looking keys (i.e., keys for the one-time pad encryption scheme)
from a single short random string.

To generate this, we can use a pseudorandom function! In particular, we
would like a pseudorandom function of the form
\(F \colon \bin^n \times \bin^n \rightarrow \bin^n\).

\marginnote{\begin{footnotesize}

If you forget what a pseudorandom function is, refer back to
Definition~\ref{def-prf}.

\end{footnotesize}}

\marginnote{\begin{footnotesize}

In practice, we will use AES or another block cipher as a pseudorandom
function, so we will have \(n = 128\) or so.

\end{footnotesize}}

Using such a pseudorandom function, we can construct a new encryption
scheme that replaces the truly random keys with pseudorandom keys
generated from the pseudorandom function. The keyspace and message space
for this encryption space are \(\zo^n\) and the ciphertext space is
\(\zo^{2n}\). The algorithms are:

\begin{itemize}
\tightlist
\item
  \(\Enc(k, m) \to c\):

  \begin{itemize}
  \tightlist
  \item
    Sample a random value \(r \rgets \bin^n\). We call this the
    ``nonce.''
  \item
    Compute a one-time key \(k \gets F(k, r)\) using the pseudorandom
    function.
  \item
    Use this key to compute the ciphertext using the one-time pad:
    output \(c \gets (r, k \xor m) \in \zo^{2n}\).
  \end{itemize}
\item
  \(\Dec(k, c)\):

  \begin{itemize}
  \tightlist
  \item
    Parse the ciphertext \(c\) into \((r, c')\).
  \item
    Compute \(m \gets c' \xor F(k, r)\).
  \end{itemize}
\end{itemize}

\marginnote{\begin{footnotesize}

While this encryption scheme is CPA-secure, it provides no message
integrity. For any string \(\Delta \in \zo^n\), an attacker can modify a
ciphetext \((r, c')\) to \((r, c' \oplus \Delta)\). This ciphertexts now
decrypts to the original message XORd with the attacker-chosen string
\(\Delta\).

\end{footnotesize}}

Correctness holds by construction. The security argument goes in three
steps:

\begin{itemize}
\item
  First, we argue that if \(n\) is large enough, the probability that
  the encryption algorithm ever chooses the same \(r\) value twice is
  negligible.

  \marginnote{\begin{footnotesize}

  By the Birthday Paradox, after encrypting \(T\) messages, the
  probability of a repeated \(r\) value is \(O(T^2/2^n)\), which is
  negligible in the key length \(n\).

  \end{footnotesize}}
\item
  Second, we argue that as long as the encryption algorithm never
  chooses the same \(r\) value twice, we can replace the values
  \(F(k, r)\) with truly random strings. By the security of the
  pseudorandom function, the adversary will not notice this change.
\item
  At this point, we can appeal to the security of the one-time pad
  scheme to argue that the adversary has no chance of winning the
  security game (Definition~\ref{def-cpa}).
\end{itemize}

For security to hold, it is crucial that the probability that the
encryption algorithm uses the same encryption nonce \(r\) twice be
negligibly small. If the encryption algorithm ever selects the same
random nonce \(r\) twice, the pad \(F(k, r)\) will be identical in two
ciphertexts. An attacker can then apply the two-time-pad attack to
recover the XOR of the two messages.

By the Birthday Paradox, if we sample the nonce from a space of size
\(2^{128}\), we can expect a collision in 128-bit random values once
around \(2^{64}\) have been generated. Therefore, any single
encryption-key must be used for \(\ll 2^{64}\) messages. Cryptographic
standards typically limit the number of bytes that users can encrypt
with the same key to prevent these sorts of problems.

\section{Encrypting longer messages: Counter mode}\label{sec-enc-ctr}

The CPA-secure encryption scheme of Section~\ref{sec-lec8-weakenc} only
allowed encrypting messages of a fixed length. We now show how to use
\emph{counter-mode encryption} to extend this scheme to support messages
of arbitrary length. That is, we will construct an encryption scheme
\(\Enc \colon \calK \times \zo^* \to \zo^*\) for messages of any length.

\marginnote{\begin{footnotesize}

In practice, encryption schemes place some (large) bound on the length
of encrypted messages. For example, the AES-GCM cipher has a maximum
message limit of just under 64 GiB.

\end{footnotesize}}

Counter-mode encryption works much as the CPA-secure encryption scheme
we have already seen in Section~\ref{sec-lec8-weakenc}, except that we
split the message into blocks and encrypt each block separately.

We will use the function
\(\ToString \colon \{0, \dots, 2^n-1\} \to \zon\), which converts an
integer in \(\{0, \dots, 2^n -1\}\) to an \(n\)-bit string in the
natural way.

The encryption scheme uses a pseudorandom function
\(F \colon \calK \times \zon \to \zon\). The secret encryption key is a
key \(k \in \calK\) for the pseudorandom function. To encrypt a message,
we choose a fresh random value \(r \getsr \{0, \dots, 2^n-1\}\) and XOR
the \(i\)th block of the message with the value
\(F(k, \ToString(r+i \bmod 2^n))\).

\begin{itemize}
\tightlist
\item
  \(\Enc(k, m):\)

  \begin{itemize}
  \item
    Split the message \(m\) into blocks of \(n\) bits:
    \((m_1, m_2, m_3, \dots, m_\ell)\). The last message block
    \(m_\ell\) may be shorter than \(n\) bits.
  \item
    Sample a random nonce \(r \getsr \{0, \dots, 2^n-1\}\).

    \marginnote{\begin{footnotesize}

    The nonce is sometimes called an ``initialization vector'' or
    ``IV.''

    \end{footnotesize}}
  \item
    For \(i = 1, \dots, \ell\): Compute
    \(c_i \gets m_i \oplus F(k, \mathsf{ToString}(r+i \bmod 2^n))\). (If
    the last message block \(m_\ell\) is less than \(n\) bits long,
    truncate the last ciphertext block \(c_\ell\) to the length of
    \(m_\ell\).)
  \item
    Output the ciphertext \(c = \big(r, c_1, \dots, c_\ell)\).
  \end{itemize}
\item
  \(\Dec(k, c):\)

  \begin{itemize}
  \tightlist
  \item
    Parse the ciphertext \(c\) as \((r, c_1, \dots, c_\ell)\), where all
    values but the last are exactly \(n\) bits long.
  \item
    For \(i = 1, \dots, \ell\): Compute
    \(m_i \gets c_i \oplus F(k, \ToString(r+i \bmod 2^n))\). (Truncate
    \(m_\ell\) to the length of \(c_\ell\).)
  \item
    Output the message \(m = (m_1 \| \dots \| m_\ell)\).
  \end{itemize}
\end{itemize}

As long as the space of random values \(r\) is large enough to ensure
that the encryption routine never evaluates the pseudorandom function
\(F(k, \cdot)\) on the same input twice, this scheme will be CPA secure.

You may notice that this encryption scheme reveals the length of the
encrypted message to the attacker! This indeed is a potential risk, but
in some sense it is required: if we were to hide the length of the
message, we would need to set some maximum message length and pad it up
to this length. If we did this, encrypting a single word would
necessarily result in a ciphertext equal in length to the ciphertext of
encrypting a movie! This would greatly decrease the practicality of our
encryption scheme. For applications where hiding the length is
especially important, the messages can be padded to ensure they are all
the same length before they are passed to the encryption scheme.

\chapter{Authenticated Encryption}\label{authenticated-encryption}

We have just constructed an encryption scheme with weak (CPA) security:
one that provided security given that the adversary could see
\emph{encryptions} of messages of her choice. The ``gold standard''
security notion for encryption schemes allows the attacker to receive
both encryptions of messages of its choice \emph{and} decryptions of
ciphertexts of its choice. Our security notion then says that even an
attacker with this power should not be able to distinguish which of two
chosen plaintext message a given ciphertexts encrypts. This new and
strong notion of security for encryption schemes is called
\emph{security against adaoptive chosen-ciphertext attacks}.

\textbf{Motivation: Chosen-ciphertext security.} It is not clear why
chosen-ciphertext security is the right security notion to consider. In
real-life applications, why would we ever allow an attacker to obtain
\emph{decryptions} of ciphertexts of its choosing? It turns out that, in
many settings, attackers can indeed trick honest parties into decrypting
adversarially ciphertexts and revealing their contents.

As a simple example, imagine a server that receives encrypted requests
from the network, decrypts them, and either:

\begin{itemize}
\tightlist
\item
  returns an error if the decrypted request is malformed, or
\item
  silently processes the request otherwise.
\end{itemize}

In this case, an attacker can send ciphertexts to the server and learn
information about their decryptions by noticing whether or not the
server returned an error message.

\marginnote{\begin{footnotesize}

If the server returns an error when the decrypted message is illformed,
we often call it a \emph{padding oracle}. Many many real-world protocols
using non-chosen-ciphertext-secure encryption schemes fall victim to
this sort of attack.

\end{footnotesize}}

CPA-secure encryption schemes provide \emph{no} security guarantees in
this setting. Even if the server leaks a single bit to the attacker
about the decrypted value (such as whether the decrypted ciphertext is a
well-formed request or not), the attacker could potentially learn the
entire secret key!

As a concrete application-level example, consider what could happen if
SSH (the secure shell protocol) used the counter-mode encryption scheme,
as described in the previous lecture, without any authentication. A user
might send the command
\texttt{echo\ secret\ \textgreater{}\ secret-file} to write their secret
string, \texttt{secret}, to their own private file \texttt{secret-file}.
The encryption of this command, sent over the network, would be the XOR
of the command and the PRF-generated pseudorandom bytes. Suppose that
the adversary knows the user is running this command, but doesn't know
the contents of the 6-byte secret string. Consider what happens if the
adversary XORs the encrypted message with the hexadecimal bytes
\texttt{00\ 00\ 00\ 00\ 00\ 00\ 00\ 00\ 00\ 00\ 00\ 00\ 00\ 00\ 5c\ 11\ 0e\ 02\ 4a\ 04\ 58\ 04\ 05\ 05\ 06}.
The leading zeroes mean that the first part of the command,
\texttt{echo\ secret\ \textgreater{}}, will remain unchanged. When the
server decrypts the latter part of the command (originally
\texttt{secret-file}), however, it will obtain the XOR of
\texttt{secret-file} and the bytes that the adversary XOR'ed in, which
happens to turn into the string \texttt{/tmp/public}. As a result, if
the server now runs this command, the user's secret data will be written
to a file \texttt{/tmp/public}, which might be available to an adversary
that also has an account on that same server and can look at files in
\texttt{/tmp}.

Chosen-ciphertext security guarantees that such attacks are ineffective.

\section{Defining Authenticated
Encryption}\label{defining-authenticated-encryption}

Authenticated-encryption schemes simultaneously provide message
integrity (as a MAC does) and confidentiality (as CPA-secure encryption
does). Additionally, authenticated-encryption schemes remain secure even
when an attacker can see encryptions and decrypts of ciphertexts of her
choice. Perhaps unsurprisingly, the standard way to construct an
authenticated-encryption scheme is to combine a CPA-secure encryption
scheme with a MAC in a careful way.

We now formally define our strong security notion:
\emph{indistinguishability under chosen ciphertext attacks}, also known
as IND-CCA2 security or CCA security.

\begin{definition}[CCA security for encryption
(strong)]\protect\hypertarget{def-cca}{}\label{def-cca}

An encryption scheme is \emph{secure against adaptive chosen-ciphertext
attacks} if every efficient adversary wins the following game with
probability at most \(\tfrac{1}{2} + \text{``negligible''}\):

\begin{itemize}
\tightlist
\item
  The challenger samples \(b \rgets \bin\) and \(k \rgets \calK\).
\item
  The adversary can make either of the following queries to the
  challenger repeatedly:

  \begin{itemize}
  \tightlist
  \item
    \emph{Chosen-plaintext queries}

    \begin{itemize}
    \tightlist
    \item
      The adversary sends the challenger a message \(m_i \in \calM\)
    \item
      The challenger replies with \(c_{m_i} \gets \Enc(k, m_i)\).
    \end{itemize}
  \item
    \emph{Chosen-ciphertext queries}

    \begin{itemize}
    \tightlist
    \item
      The adversary sends the challenger a \emph{ciphertext}
      \(c_j \notin \set{c_{m_0}, \ldots, c_{m_i}}\)
    \item
      The challenger replies with \(m_{c_j} \gets \Dec(k, c_j)\).
    \end{itemize}
  \end{itemize}
\item
  The adversary then sends two messages \((m^*_0, m^*_1) \in \calM^2\)
  to the challenger, where \(\abs{m^*_0} = \abs{m^*_1}\).
\item
  The challenger replies with \(c^* \gets \Enc(k, m^*_b)\).
\item
  The adversary can make more chosen-plaintext queries and more
  chosen-ciphertext queries. (The adversary may not make a
  chosen-ciphertext query on the challenge ciphertext \(c^*\).)
\item
  The adversary outputs a value \(b' \in \bin\).
\item
  The adversary wins if \(b=b'\).
\end{itemize}

\end{definition}

\subsection{Encrypt then MAC}\label{encrypt-then-mac}

We typically achieve CCA security using the ``encrypt-then-MAC''
construction:

\begin{itemize}
\item
  First, encrypt the message using a CPA-secure encryption scheme on key
  \(k_\mathsf{Enc}\).
\item
  Next, MAC the \emph{ciphertext} using a secure MAC scheme and an
  independent key \(k_\mathsf{MAC}\).

  \marginnote{\begin{footnotesize}

  As we discuss below, it is possible to derive both keys
  \(k_\mathsf{Enc}\) and \(k_\mathsf{MAC}\) from a single key \(k\)
  using a pseudorandom function.

  \end{footnotesize}}
\item
  Output the ciphertext and the MAC tag.
\end{itemize}

The decryption routine first checks the MAC tag, then decrypts the
ciphertext.

Using independent keys \((k_\mathsf{Enc}, k_\mathsf{MAC})\) is important
in encrypt-then-MAC, as in many other cryptographic constructions. For
example, a CPA-secure encryption scheme using \(n\)-bit keys can reveal
the low order \(n/2\) bits of its secret key in the ciphertext. And a
secure MAC scheme using \(n\)-bit keys can reveal the high-order \(n/2\)
bits of its secret key in each MAC tag. Used independently, the
encryption scheme and the MAC scheme are both secure. Used together with
the same key \(k\), the attacker learns all \(n\) bits of the key \(n\)
and can break both primitives!

So, in general, you should always use independent keys for different
primitives. To reduce the amount of keying material parties need to
store, it is actually sufficient to store a single secret key \(k\) for
a pseudorandom function \(F\) and derive all subsequent keys from the
pseudorandom strings \(F(k, 0), F(k, 1), \dots\).

\begin{theorem}[Informal]\protect\hypertarget{thm-encrypt-then-mac}{}\label{thm-encrypt-then-mac}

The Encrypt-then-MAC construction yields a CCA-secure encryption scheme,
provided that: the underlying encryption scheme is CPA-secure and the
underlying MAC scheme is secure (existentially unforgeable against
adaptive chosen message attacks).

\end{theorem}

\textbf{Warning! Only use encrypt-then-MAC.} There are a number of bad
ways to combine encryption and MACs to attempt to build
authenticated-encryption schemes. MAC-then-encrypt is one way.
Encrypt-and-MAC (i.e., MAC the message instead of the ciphertext) is
another. Neither of these constructions is necessarily CCA-secure when
used with a CPA-secure encryption scheme and a secure MAC scheme. So the
\textbf{\emph{only}} flavor of authenticated encryption you should use
is encrypt-then-MAC.

\marginnote{\begin{footnotesize}

Reminder: You should never need to implement authenticated-encryption
schemes yourself. Instead use an off-the-shelf implementation that does
the hard work for you. AES-GCM is one popular and widely implemented
authenticated encryption scheme

\end{footnotesize}}

\section{AES-GCM (Galois Counter Mode)}\label{sec-enc-gcm}

One of the most widely used authenticated-encryption constructions is
AES-GCM. It follows the encrypt-then-MAC paradigm. It uses AES as a
pseudorandom function for counter-mode encryption from the previous
lecture. It uses a Carter-Wegman-style MAC -- see previous lecture -- as
the MAC scheme.

There are a few optimizations that AES-GCM uses beyond what we have
described:

\begin{itemize}
\item
  AES-GCM derives both the encryption and MAC keys from a single short
  key using a pseudorandom function.
\item
  AES-GCM implements a fast form of the Carter-Wegman MAC that does not
  need arithmetic modulo a big prime \(p\), as the scheme described in
  Section~\ref{sec-mac-cw} does. Instead, of defining the MAC using
  \(\Z_p\) (integers mod a 128-bit prime \(p\)), GCM works with 128-bit
  strings. The GCM mode of operation replaces addition modulo \(p\) with
  XOR of 128-bit strings and it replaces multiplication modulo \(p\)
  with a somewhat complicated operation on 128-bit strings.

  \marginnote{\begin{footnotesize}

  If you are interested, to implement the multiplication operation:
  think of both 128-bit strings as polynomials with 128 coefficients in
  \(\Z_2 = \{0,1\}\). Multiply the polynomials, reducing the
  coefficients modulo \(2\). Then reduce the resulting polynomial modulo
  some fixed polynomial of degree-\(128\). Then interpret the result as
  a \(128\)-bit string.

  \end{footnotesize}}

  (Formally, the scheme works over the field \(\F_{2^{128}}\) of order
  \(2^{128}\).) This gives a big performance boost with no loss in
  security.
\end{itemize}

\chapter{Key Exchange and Public-key
Encryption}\label{key-exchange-and-public-key-encryption}

So far, we have talked about encryption systems that require the sender
and recipient to \emph{share a secret key}. In this chapter, we discuss
how the sender and recipient can agree on a shared secret even if they
only ever communicate over an open (insecure) network.

\section{Key exchange}\label{key-exchange}

A key-agreement scheme over keyspace \(\calK\) is defined by efficient
functions \((\Gen, \Derive)\):

\begin{itemize}
\tightlist
\item
  \(\Gen() \to (\sk, \pk)\). The \(\Gen\) algorithm generates a secret
  key and a public key for one party.
\item
  \(\Derive(\sk_A, \pk_B) \to k\). The \(\Derive\) algorithm takes as
  input one party's secret key \(\sk_A\) and the other party's public
  key \(\pk_B\) and outputs a shared key \(k \in \calK\).
\end{itemize}

\marginnote{\begin{footnotesize}

Formally, \(\Gen\) also takes as input the security parameter.

\end{footnotesize}}

Given this syntax, let us see how two parties can use a key-agreement
scheme \((\Gen, \Derive)\) to agree on a shared secret:

\begin{enumerate}
\def\labelenumi{\arabic{enumi}.}
\tightlist
\item
  Alice runs \((\sk_A, \pk_A) \gets \Gen()\) and sends \(\pk_A\) to Bob.
\item
  Bob runs \((\sk_B, \pk_B) \gets \Gen()\) and sends \(\pk_B\) to Alice.
\item
  Alice and Bob both then compute a key \(k \in \calK\) as:

  \begin{itemize}
  \tightlist
  \item
    Alice computes \(k \gets \Derive(\sk_A, \pk_B)\).
  \item
    Bob computes \(k \gets \Derive(\sk_B, \pk_A)\).
  \end{itemize}
\end{enumerate}

\textbf{Correctness.} Correctness for a key-agreement scheme just says
that the two parties should always agree on the same shared secret:

\begin{definition}[Key Agreement
Correctness]\protect\hypertarget{def-key-agreement-correctness}{}\label{def-key-agreement-correctness}

We say that a key-agreement scheme is \emph{correct} if for all
\((\sk_A, \pk_A) \gets \Gen()\) and \((\sk_B, \pk_B) \gets \Gen()\), it
holds that: \[ \Derive(\sk_A, \pk_B) = \Derive(\sk_B, \pk_A). \]

\end{definition}

\textbf{Security.} The standard notion of security for a key-agreement
scheme only considers \emph{passive attacks}: we consider adversaries
that can view the network traffic but cannot modify it. In practice, we
can combine key-agreement schemes with authentication schemes (e.g.,
digital signatures) to prevent active attacks by in-network adversaries.

Our security definition for key agreement says that even if an adversary
sees both parties' public keys, it should not be able to distinguish the
shared secret from random. That is, the following probability
distributions \(D_\text{real}\) and \(D_\text{random}\) should be
computationally indistinguishable:

\marginnote{\begin{footnotesize}

When we say that two distributions are \emph{computationally
indistinguishable}, we mean that if we give the adversary a sample from
one or the other, it can guess which sample it got with probability at
most \(1/2 + \text{``negligible''}\).

\end{footnotesize}}

\[\begin{aligned}
  D_\text{real} &\deq \left\{
(\pk_A, \pk_B, k) \colon  \begin{aligned}
  (\sk_A, \pk_A) &\getsr \Gen()\\
  (\sk_B, \pk_B) &\getsr \Gen()\\
  k &\getsr \Derive(\sk_A, \pk_B)
\end{aligned}
\right\} \\
  D_\text{random} &\deq \left\{ (\pk_A, \pk_B, k) \colon  \begin{aligned}
  (\_, \pk_A) &\getsr \Gen()\\
  (\_, \pk_B) &\getsr \Gen()\\
  k &\getsr \calK
\end{aligned}
  \right\}
\end{aligned}\]

\section{Diffie-Hellman key exchange}\label{sec-dh}

We now give a simple and beautiful key-exchange protocol, due to Diffie
and Hellman.

\marginnote{\begin{footnotesize}

So far, we have been able to construct our cryptographic entities from
constructs like a PRF, which meant that we could use ``unstructured''
algorithms like AES to compute them. We so far only know how to
construct key-exchange schemes from more structured problems (e.g.,
based on number theory).

\end{footnotesize}}

The version we will see uses large primes \(p\) and \(q\) a public
parameters, along with a number \(g \in \Z^*_p\), called the
``generator.'' (\textbf{Warning:} The parameters \(p\), \(q\), and \(g\)
must have some particular relation. So do not attempt to pick these
parameters yourself.) Typically, we will have
\(p \approx q \approx 2^{2048}\).

\marginnote{\begin{footnotesize}

For a prime \(p\), the notation \(\Z^*_p\) just denotes the non-zero
integers modulo \(p\). So when we write \(ab \in \Z^*_p\), we mean
\(a \cdot b \bmod p\).

\end{footnotesize}}

The keyspace of the Diffie-Hellman scheme is \(\calK = \Z^*_p\) and the
algorithms are as follows:

\begin{itemize}
\tightlist
\item
  \(\Gen() \to (\sk, \pk)\).

  \begin{itemize}
  \tightlist
  \item
    Sample \(\sk \getsr \{1, \dots, q\}\).
  \item
    Set \(\pk \gets g^{\sk} \in \Z^*_p\).
  \item
    Output \((\sk, \pk)\).
  \end{itemize}
\item
  \(\Derive(\sk_A, \pk_B) \to k\).

  \begin{itemize}
  \tightlist
  \item
    We have \(\sk_A \in \{1, \dots, q\}\) and \(\pk_B \in \Z^*_p\).
  \item
    Output \(k \gets (\pk_B)^{\sk_A} \in \Z^*_p\).
  \end{itemize}
\end{itemize}

Before we argue correctness and security, let us consider the
computational efficiency of the scheme:

\textbf{Efficiency.} In order for the algorithm to be useful, Alice and
Bob must be able to compute \(g^x \in \Z^*_p\) efficiently, for
\(x \in \{1, \dots, q \approx 2^{2048}\}\). However, trying to compute
\(g^x\) where \(x\) is 2048 bits long naively would certainly not be
efficient: it would require \(x \approx 2^{2048}\) multiplications!
However, we can compute this exponentiation much more efficiently using
the following strategy:

\begin{itemize}
\tightlist
\item
  \textbf{Compute powers of \(g\).} Write
  \(\ell \gets \lceil \log_2 p \rceil\). Then compute
  \((g, g^2, g^4, g^{8}, g^{16}, \dots, g^{2^\ell})\), where all of
  these are in \(\Z^*_p\). It is possible to compute \(g^{2^{i}}\) with
  a single multiplication modulo \(p\) as
  \(g^{2^i} = (g^{2^{(i-1)}})^2 \in \Z^*_p\). So this step takes only
  \(\ell = 2048\) multiplications.
\item
  \textbf{Compute exponentiation.} Write the bits of the exponent as
  \(x = x_0 \cdots x_{\ell-1}\). Then compute:
  \[ g^x = g^{\sum_{i = 0}^{\ell-1} x_i 2^i} = \prod_{i=0}^{\ell-1} x_i(g^{2^i}) \in \Z^*_p\]
  This step again takes only \(\ell\) multiplications.
\end{itemize}

\marginnote{\begin{footnotesize}

In many applications, the generator \(g\) is fixed in advance. In this
case, the implementation can precompute and store these powers of \(g\).

\end{footnotesize}}

\textbf{Correctness.} Correctness holds since
\(g^{xy}=g^{yx} \in \Z^*_p\) for all \(x,y \in \Z\): \[\begin{aligned}
  \Derive(\sk_A, \pk_B) = (\pk_B)^{\sk_A} = (\pk_A)^{\sk_B} = \Derive(\sk_B, \pk_A).
\end{aligned}\]

\textbf{Security.} To argue security, we must rely on a new
computational assumption: essentially we just assume that the
key-agreement scheme is secure.

\begin{definition}[Decision Diffie-Hellman (DDH)
assumption]\protect\hypertarget{def-ddh}{}\label{def-ddh}

The \emph{decision Diffie-Hellman assumption} states that, for a
suitable choice of \(p\), \(q\), and \(g\), the following distributions
are computationally indistinguishable:

\marginnote{\begin{footnotesize}

To make the statement fully formal, we need to let \(p\), \(q\), and
\(g\) grow with the security parameter.

\end{footnotesize}}

\[\begin{aligned}
    D_\text{real} &\deq \{ (g, g^x, g^y, g^{xy}) \in (\Z^*_p)^4\colon x, y \,\,\,\,\,\getsr \{1, \dots, q\}\}\\
    D_\text{ideal} &\deq \{ (g, g^x, g^y, g^z) \,\,\,\in (\Z^*_p)^4 \colon x, y, z \getsr \{1, \dots, q\}\}
\end{aligned}\]

\end{definition}

In practice, we typically first run Diffie-Hellman key agreement, have
the two parties run the shared secret that they get through a
cryptographic hash function, and then use the hashed value as an
encryption key. If we model the hash function as a random oracle,
security can rely on a slightly weaker assumption---the
``computational'' Diffie-Hellman (CDH) assumption. The CDH assumption
asserts, informally, that given \((g, g^x, g^y)\), it is infeasible to
compute \(g^{xy}\). More formally, we have:

\begin{definition}[Computational Diffie-Hellman (CDH)
assumption]\protect\hypertarget{def-cdh}{}\label{def-cdh}

The \emph{computational Diffie-Hellman assumption} states that, for a
suitable choice of \(p\), \(q\), and \(g\), and for all adversaries
\(\calA\):
\[ \Pr[\calA(g, g^x, g^y) = g^{xy} \colon x, y, \getsr\{1, \dots, q\} ] \leq \text{``negligible.''}\]

\end{definition}

\section{The discrete-log problem}\label{the-discrete-log-problem}

The DDH assumption (Definition~\ref{def-ddh}) is no harder than the
following problem, which asserts that computing \(x\) given
\((g, g^x) \in (\Z_p^*)^2\) is computationally infeasible:

\begin{definition}[Discrete-log
assumption]\protect\hypertarget{def-dlog}{}\label{def-dlog}

The \emph{discrete-log assumption} states that, for \(p\), \(q\), and
\(g\), and for all efficient adversaries \(\calA\):
\[ \Pr[\calA(g, g^x) = x \colon x \getsr \{1, \dots, q\} ] \leq \text{``negligible''}.\]

\end{definition}

Given an algorithm for the discrete-log problem, we can use it to solve
the DDH problem. Given a DDH challenge \((g, g^x, g^y, g^{z})\), compute
the discrete log of \(g^x\) and test whether \(g^z = {(g^y)}^x\). For
certain choices of \(p\), \(q\), and \(g\), the best known algorithm for
the DDH problem is to first solve the discrete-log problem in
\(\Z^*_p\).

How hard it to solve the discrete-log problem then?

\begin{enumerate}
\def\labelenumi{\arabic{enumi}.}
\tightlist
\item
  The most basic attack is to enumerate all \(p\) possible values of
  \(x\) and check whether the corresponding \(g^x\) matches. This will
  take time \(p \approx 2^{2048}\).
\item
  There is a slightly more clever algorithm, called ``Baby Step Giant
  Step,'' that is able to compute \(x\) in time \(\sqrt{p}\).
\item
  In \(\Z^*_p\), a much better attack is the Number Field Sieve. This
  algorithm is able to compute \(x\) in (roughly) time
  \(\exp((\log p)^{1/3} (\log \log p)^{2/3})\)---sub-exponential time!
  The existence of this attack is the reason we require \(p\) to be
  \(2048\) bits long to get \(128\)-bit security.
\end{enumerate}

\marginnote{\begin{footnotesize}

An attack that runs in time \(\sqrt p\) runs in time
\(2^{\frac 12 \log p}\). In contrast, the Number Field Sieve runs in
time roughly \(2^{\sqrt[3]{\log p}}\). This is much much much faster
than the \(\sqrt p\)-time algorithms, since the exponent grows much more
slowly.

\end{footnotesize}}

\section{Generalizations of
Diffie-Hellman}\label{generalizations-of-diffie-hellman}

We have described the Diffie-Hellman protocol in terms of
\(\Z^*_p\)---multiplication of integers modulo \(p\). In particular, the
protocol uses a \emph{set} of elements (here, \(\Z^*_p\)) and a
\emph{binary operation on elements} (here, multiplication modulo \(p\)).
There is a natural generalization of the Diffie-Hellman protocol that
works with other sets of elements and binary operations that operate on
them.

\marginnote{\begin{footnotesize}

More precisely, we can define Diffie-Hellman key exchange with respect
to any \emph{finite cyclic group}---a concept from abstract algebra. A
cyclic group just consists of a \emph{set} of elements and a
\emph{binary operation on elements}. The set and operation need to
satisfy certain mathematical properties---associativity, etc.

Given a group \(\G\), we can then define a discrete-log assumption on
\(\G\) and whenever discrete-log is hard on \(\G\), we can use \(\G\) to
construct cryptosystems.

\end{footnotesize}}

The most widely used version of the Diffie-Hellman protocol today uses
\emph{elliptic-curve groups}. The set of elements in an elliptic-curve
group is a set of points \((x,y) \in \Z^2\) in two-dimensional space,
where \(0 \leq x, y \leq p\) for some prime \(p \approx 2^{256}\). The
binary operation on elements is some geometric operation on two points
that yields a third point. Even though the underlying objects are now
points instead of integers modulo \(p\), the Diffie-Hellman protocol
looks exactly the same in this setting.

The appeal of elliptic-curve cryptosystems is that the best known
discrete-log algorithm on certain elliptic-curve groups is Baby Step
Giant Step, which runs in \(\sqrt p\) time. So, we can use
elliptic-curve groups of size \(2^{256}\) and the Diffie-Hellman public
keys take only \(\approx 256\) bits to represent. In contrast, when
working in \(\Z^*_p\), we need to work modulo a prime
\(p \approx 2^{2048}\) to defeat the Number Field Sieve attack, so
Diffie-Hellman public keys in this setting take \(\approx 2048\) bits to
represent.

\section{Defining Public-Key
Encryption}\label{defining-public-key-encryption}

The definition for a public-key encryption scheme will be similar to the
definitions we saw for symmetric-key encryption:

\begin{definition}[Public-Key Encryption
Scheme]\protect\hypertarget{def-pke}{}\label{def-pke}

A public-key encryption scheme over message space \(\calM\) consists of
three efficient algorithms \((\Gen, \Enc, \Dec)\):

\begin{itemize}
\tightlist
\item
  \(\Gen() \to (\sk, \pk)\): Generates a keypair with secret key \(\sk\)
  and public key \(\pk\).
\item
  \(\Enc(\pk, m) \to c\): Uses public key \(\pk\) to encrypts a message
  \(m \in \calM\) to a ciphertext \(c\).
\item
  \(\Dec(\sk, c) \to m\): Uses secret key \(\sk\) to decrypt ciphertext
  \(c\) into message \(m \in \calM\).
\end{itemize}

\end{definition}

\begin{definition}[Public-Key Encryption -
Correctness]\protect\hypertarget{def-pke-correctness}{}\label{def-pke-correctness}

For all keypairs \((\sk, \pk) \gets \Gen()\) and for all messages
\(m \in \calM\), it holds that \[\Dec(\sk, \Enc(\pk, m)) = m.\]

\end{definition}

\begin{definition}[Public-Key Encryption - Security against
chosen-plaintext attacks
(Weak)]\protect\hypertarget{def-pke-cpa}{}\label{def-pke-cpa}

~

\marginnote{\begin{footnotesize}

The literature sometimes calls security against chosen-plaintext attacks
(``CPA security'') \emph{semantic security}.

\end{footnotesize}}

A public-key encryption scheme \((\Gen, \Enc, \Dec)\) is \emph{secure
against chosen-plaintext attacks} if all efficient adversaries \(\calA\)
win the following game with probability \(\leq \tfrac{1}{2} + \negl\):

\begin{itemize}
\tightlist
\item
  The challenger generates \((\sk, \pk) \gets \Gen()\) and
  \(b \getsr \bin\) and sends the public key \(\pk\) to \(\A\).
\item
  The adversary \(\A\) sends \(m_0, m_1 \in \calM\) to the challenger,
  where \(\abs{m_0} = \abs{m_1}\).
\item
  The challenger responds with \(\Enc(\pk, m_b)\)
\item
  The adversary outputs \(b'\) and wins if \(b' = b\).
\end{itemize}

\end{definition}

Note that since this is a public-key encryption scheme, the adversary
can use the public key to generate encryptions of any message of their
choice. As in the secret-key setting, here it is crucial that the
encryption algorithm be randomized---otherwise two encryptions of the
same message are identical and the attacker can break CPA security.

For symmetric-key encryption, we also defined a stronger notion of
security that we called security against chosen-ciphertext attacks (CCA
security). We can extend this definition of CCA security to the
public-key setting.

This security definition asserts that the attacker should not be able to
distinguish the encryption \(c^*\) of two messages of its choice, even
if the attacker can obtain decryptions of any ciphertext except \(c^*\).
This is a very strong notion of security (since the security definition
gives the attacker a lot of power) and it is the notion of security that
we typically demand in practice.

\begin{definition}[Public-Key Encryption - Security against
chosen-ciphertext attacks
(Strong)]\protect\hypertarget{def-pke-cca}{}\label{def-pke-cca}

A public-key encryption scheme \((\Gen, \Enc, \Dec)\) is \emph{secure
against chosen-ciphertext attacks} if all efficient adversaries
\(\calA\) win the following game with probability
\(\leq \tfrac{1}{2} + \negl\):

\begin{itemize}
\tightlist
\item
  The challenger generates \((\sk, \pk) \gets \Gen()\) and
  \(b \getsr \bin\) and sends the public key \(\pk\) to the adversary
  \(\A\).
\item
  The adversary may make polynomially many \emph{decryption queries}:

  \begin{itemize}
  \tightlist
  \item
    The adversary \(\A\) sends ciphertext \(c\) to the challenger.
  \item
    The challenger responds with \(m \gets \Dec(\sk, c)\)
  \end{itemize}
\item
  At some point, the adversary sends a pair of messages
  \(m_0, m_1 \in \calM\) to the challenger, where
  \(\abs{m_0}=\abs{m_1}\).
\item
  The challenger returns \(c^* \gets \Enc(\pk, m_b)\).
\item
  The adversary can continue to make decryption queries, provided that
  it never asks the challenger to decrypt the ciphertext \(c^*\).
\item
  The adversary outputs \(b'\) and wins if \(b' = b\).
\end{itemize}

\end{definition}

\section{ElGamal Encryption Scheme}\label{elgamal-encryption-scheme}

In public-key encryption, we want a sender to be able to encrypt a
message to a recipient. The goal of public-key encryption is to perform
encryption without a shared secret key.

ElGamal's encryption scheme essentially uses Diffie-Hellman key exchange
to allow the sender and recipient to agree on a shared secret key \(k\),
and then has the sender encrypt her message with a symmetric-key
cryptosystem under key \(k\).

\marginnote{\begin{footnotesize}

ElGamal's scheme was not the first public-key encryption scheme. The
first scheme, RSA, was much more complicated despite Diffie-Hellman
already having been published.

\end{footnotesize}}

\begin{definition}[Hashed ElGamal
Encryption]\protect\hypertarget{def-elgamal}{}\label{def-elgamal}

Let \(p\), \(q\), and \(g\) be integers of the sort we use for
Diffie-Hellman key exchange (Section~\ref{sec-dh}). In particular,
\(p \approx q \approx 2^{2048}\) and \(g \in \Z^*_p\).

\marginnote{\begin{footnotesize}

In practice, we typically use elliptic-curve groups to instantiate
ElGamal encryption, instead of \(\Z^*_p\). But the general principle is
exactly the same.

\end{footnotesize}}

Let \(H: \Z^*_p \rightarrow \bin^*\) be a hash function (modelled as a
random oracle). Let \((\Enc', \Dec')\) be a symmetric-key
authenticated-encryption scheme. Then define:

\begin{itemize}
\tightlist
\item
  \(\Gen() \to (\sk, \pk)\):

  \begin{itemize}
  \tightlist
  \item
    Choose \(a \gets \{1, \dots, q\}\).
  \item
    Compute \(A \gets g^a \in \Z^*_p\).
  \item
    Output \((\sk, \pk) \gets (a, A)\).
  \end{itemize}
\item
  \(\Enc(\pk, m) \to c\):

  \begin{itemize}
  \tightlist
  \item
    Choose \(r \gets \{1, \dots, q\}\).
  \item
    Compute \(R \gets g^r \in \Z^*_p\).
  \item
    Compute \(k \gets H(\pk^r \in \Z^*_p)\).
  \item
    Output \(c \gets (R, \Enc'(k, m)\).
  \end{itemize}
\item
  \(\Dec(\sk, c) \to m\):

  \begin{itemize}
  \tightlist
  \item
    Parse \((R, c') \gets c\).
  \item
    Compute \(k \gets H(R^a \in \Z^*_p)\).
  \item
    Output \(m \gets \Dec'(k, c')\).
  \end{itemize}
\end{itemize}

\end{definition}

\subsection{Performance}\label{performance}

The performance of this scheme is limited by the exponentiations---the
symmetric encryption scheme is quite fast (gigabytes per second), but a
single exponentiation can take a millisecond on modern processors. For
encryption, this scheme requires two exponentiations (\(g^r\) and
\(\pk^r\)). Decryption requires one exponentiation \((R^\sk)\). To speed
up the encryption routine, we can precompute powers of \(g\):
\(g^2, g^4, g^8, g^{16}, \dots\), which saves a factor of two in
exponentiations.

\marginnote{\begin{footnotesize}

Hashed ElGamal encryption is one of the most common public-key
encryption schemes used today.

\end{footnotesize}}

\subsection{Security}\label{security}

If we model the hash function \(H\) as a random oracle, we can prove
security of ElGamal encryption from (a) the computational Diffie-Hellman
assumption (Definition~\ref{def-cdh}) and (b) the CCA security of the
underlying authenticated-encryption scheme \((\Enc', \Dec')\).

\chapter{Encryption in Practice}\label{encryption-in-practice}

So far, we have established several constructions that allow us to hide
the contents of transmissions: we created chosen-plaintext- and
chosen-ciphertext-secure encryption schemes that worked with and without
a shared key. We will now discuss a few practical applications of
transport encryption, and why it is often difficult to get right.

\section{File Encryption}\label{file-encryption}

Perhaps the most straightforward use of encryption is file encryption.

\textbf{Example: WhatsApp Encrypted Backup.} WhatsApp allows the app's
users to back up their messages and contacts to the cloud. This way, a
user can recover her messages if she loses or breaks her phone. To hide
the user's data from WhatsApp's cloud servers, WhatsApp uses encrypted
backup. To achieve this, the user's device generates a 128-bit AES key
\(k\) at the time of backup and encrypt the message data (photos,
messages, etc.) using AES-GCM\((k, \cdot)\) before sending the
\emph{ciphertext} to the server. In order to allow you to restore your
backup on a new phone, the app allows you to export 64 decimal digits
that encode the AES key used. When restoring your backup, you will enter
these digits and your phone will fetch the ciphertext from the server
and use these digits to reconstruct the key and decrypt your messages.

\marginnote{\begin{footnotesize}

See
\href{https://scontent.fphl1-1.fna.fbcdn.net/v/t39.8562-6/241394876_546674233234181_8907137889500301879_n.pdf?_nc_cat=108&ccb=1-7&_nc_sid=ad8a9d&_nc_ohc=ANW9FNYPigIAX93N8Oj&_nc_ht=scontent.fphl1-1.fna&oh=00_AT9IVosU4uWqddd5woEO4eUWw3mbd76IgvwjcqDApb3p9A&oe=63565D66}{\textbf{this
document}} for details on how WhatsApp encrypts backups. (The document
also describes a more complicated backup scheme that uses passwords for
encryption.)

\end{footnotesize}}

This is a fairly simple application of file security. However, file
encryption can be much more tricky: many applications require or desire
features beyond simple encryption and decryption.

\subsection{Case Study: PDF v1.5
Encryption}\label{case-study-pdf-v1.5-encryption}

One instance of this desire for extra features that ended up going wrong
was a previous version of the PDF standard, PDF
v1.5.(\citeproc{ref-muller2019practical}{Müller et al. 2019}) This
standard provided several features:

\begin{enumerate}
\def\labelenumi{\arabic{enumi}.}
\tightlist
\item
  It is possible to password-encrypt some or all of the document. (The
  encryption scheme does not matter, but think of it as a secure
  authenticated-encryption scheme.) For example, a PDF could have an
  unencrypted title page and have the rest of the pages be encrypted.
\item
  A PDF document can contain a form that the PDF reader submits to a
  server via HTTP.
\item
  A form in a PDF document can reference other parts of a document.
\item
  The PDF reader may submit a form when an event happens: when the PDF
  is opened, closed, decrypted, etc.
\end{enumerate}

Each of these features individually seems innocuous. However, when
combined, there is a clever attack that allows an attacker to learn the
contents of the encrypted portion of a PDF.

The attack works against a PDF document with an \emph{unecrypted
(public) title page} and an \emph{encrypted (secret) body}. We assume
that the attacker can modify the PDF file on its way to the victim.

To mount the attacker, the attacker intercepts the PDF on the way to the
victim and replaces the title page with an ``evil'' title page. The evil
title page:

\begin{itemize}
\tightlist
\item
  contains an invisible PDF form,
\item
  reads the contents of the decrypted pages into a PDF form element, on
  the event that decryption of the PDF body succeeds, and then
\item
  submits the form via HTTP to \texttt{evil.com}
\end{itemize}

So when the victim recipient enters her password into the PDF reader to
decrypt the document, the evil title page is exfiltrates the decrypted
content to \texttt{evil.com} via a PDF form.

\textbf{What went wrong?} The core issue here was that the unencrypted
contents of the document were not authenticated. That is, an attacker
could modify the unencrypted pages of the document without detection. A
better design would have been to either use a MAC over all pages of the
document or to use a primitive called ``Authenticated Encryption with
Associated Data'' to authenticate the encrypted data together with the
unencrypted data.

One downside of MACing the entire document is that, before rendering
even a single page of the PDF, the reader would have to compute a MAC
over the entire PDF. If the PDF consists of many thousands of pages,
this could be expensive. Using a more sophisticated cryptographic
construction: per-page MACs, plus MACs on the document metadata, etc.,
we might be able to construct a scheme that protects document integrity
while allowing partial document rendering. But the design of such a
scheme would have to carefully defend document integrity against the
sort of attacks we describe here.

\section{Stream Encryption: Transport Layer Security
(TLS)}\label{stream-encryption-transport-layer-security-tls}

The standard Internet data transfer protocols---TCP and UDP---provide no
integrity or confidentiality whatsoever. To protect the data that we
send over the network, we use the \emph{transport-layer security} (TLS)
protocol. The TLS protocol runs on top of TCP and aims to provide an
encrypted ``tunnel'' between a client and server.

\marginnote{\begin{footnotesize}

HTTPS is simply HTTP run over TLS.

\end{footnotesize}}

While designing a stream-encryption protocol may seem straightforward at
first glance, the task is much more subtle than it would seem. However,
as is often the case in security, features and practical requirements
make the situation much more complex.

\subsection{Downgrade Attack}\label{downgrade-attack}

We now give one example of a \emph{downgrade attack}---the sort of
subtle security issue that can arise in protocol design.

The current version of TLS is TLS 1.3. An older version of the TLS
standard, called SSLv3, is vulnerable to devastating attacks that allow
an attacker to recover the plaintext traffic. However, for backwards
compatibility, many TLS clients and servers still supported SSLv3 until
quite recently.

When a TLS client connects to a TLS server, the two parties first
exchange some messages to decide which version of the TLS protocol to
use. In particular, the client will try to connect to the server using
the latest version of TLS it supports. The server will respond back with
either a confirmation (if the server supports the client's proposed TLS
version) or with garbage (if not). If garbage, the client will try to
connect to the server with an older version of TLS.

An important point is that \emph{none} of these negotiation messages are
authenticated---the client and server cannot start using authentication
(MACs, signatures, etc.) until they agree on which version of TLS to
use! So, an active in-network attacker can simply replace all of the
server responses in this protocol-negotiation phase with garbage until
the client downgrades all the way to SSLv3. Once the client and server
agree to use SSLv3, believing that this is their best available option,
the attacker can then monitor and decrypt their traffic using known
attacks on SSLv3.

The best defense against this attack is for both parties to disable
support for SSLv3 completely.

\subsection{TLS Structure}\label{tls-structure}

TLS consists of two main phases:

\begin{enumerate}
\def\labelenumi{\arabic{enumi}.}
\tightlist
\item
  \textbf{Handshake:} In this phase, the client and server use a
  key-exchange protocol to agree on a shared key to use to encrypt their
  application-layer traffic. This step uses public-key cryptography,
  since the client and server initially have no shared secret.
\item
  \textbf{Record protocol:} This phase is where the actual
  application-layer communication happens. This phases uses the secret
  key that the client and server agreed upon in the handshake phase for
  authentication and encryption.
\end{enumerate}

\subsection{TLS Handshake Properties}\label{tls-handshake-properties}

In our definitions of public-key encryption, we had only two (relatively
simple) properties: correctness and security. The TLS handshake,
however, has a much more complicated set of goals:

\begin{itemize}
\item
  \textbf{Correctness:} Both parties agree on the same session key at
  the conclusion of the handshake.
\item
  \textbf{Security:} adversary ``learns nothing'' about the secret key
  that the parties agree upon at the conclusion of the handshake.
\item
  \textbf{Peer authentication:} At the end of the handshake, each party
  believes that they are talking to the other party.
\item
  \textbf{Downgrade protection:} The parties agree on the same version
  of TLS that they would agree on absent an in-network attacker.
\item
  \textbf{Forward secrecy with respect to key compromise:} if an
  attacker steals the secrets stored on the client or the server, the
  attacker cannot decrypt past traffic.

  \marginnote{\begin{footnotesize}

  To provide forward secrecy, modern cryptographic protocols use
  long-term secrets only for signing---not for encryption. These
  protocols use \emph{ephemeral} (one-time-use) cryptographic keys for
  key exchange and encryption. Protocol participants delete these
  ephemeral keys on connection teardown and/or they rotate these keys
  often.

  This way, if an attacker compromises a party's secret key, the
  attacker can only sign messages on behalf of that party; the attacker
  cannot use the secret to decrypt past messages.

  \end{footnotesize}}
\item
  \textbf{Protection against key-compromise impersonation:} If an
  attacker steals a client's secret key, the attacker should not be able
  to impersonate other servers to the client.
\item
  \textbf{Protection of endpoint identities:} The public keys of the two
  parties should never flow over the wire in the clear. For example, if
  a client is connecting to a website that uses a content-distribution
  network, such as Akamai or Cloudflare, an attacker should not be able
  to tell which website the client is connecting to---only that it is
  hosted on Akamai or Cloudflare.
\end{itemize}

\subsection{TLS Handshake}\label{tls-handshake}

The TLS handshake is very carefully designed to achieve these
properties. A grossly simplified version looks something like the
following:

\begin{enumerate}
\def\labelenumi{\arabic{enumi}.}
\tightlist
\item
  At the start of the handshake, the TLS client holds the public
  \(\pk_{\text{CA}}\) of a certificate authority. The TLS server (for
  example, MIT) holds its secret signing key \(\sk_\text{MIT}\) and a
  public-key certificate \(\text{cert}_{MIT}\) binding its public key
  \(\pk_\text{MIT}\) to its domain \texttt{mit.edu}.
\item
  \textbf{Client Hello}: The client sends the following values to the
  server:

  \begin{itemize}
  \item
    a random nonce,
  \item
    list of supported ciphersuites, and

    \marginnote{\begin{footnotesize}

    A \emph{ciphersuite} contains all of the cryptographic parameters
    needed to perform key exchange, hashing, authentication, and
    encryption. For example, one possible ciphersuite for TLS 1.2 is
    \texttt{ECDHE-RSA-AES256-GCM-SHA384}. This indicates use of
    ephemeral elliptic-curve Diffie-Hellman key exchange (\texttt{DHE})
    with RSA signatures (\texttt{RSA}), 256-bit AES encryption in GCM
    mode (\texttt{AES256-GCM}), and SHA2-384 as the hash function
    (\texttt{SHA384}).

    \end{footnotesize}}
  \item
    an ephemeral Diffie-Hellman public key. (The client constructs this
    Diffie-Hellman public key using its preferred ciphersuite. If the
    server does not support the ciphersuite the client picked, the
    client will have to re-run this step using a different ciphersuite.)
  \end{itemize}
\item
  \textbf{Server Hello}: The server sends several values to the client,
  choosing a ciphersuite to use and completing the Diffie-Hellman key
  exchange. That is, the server sends:

  \begin{itemize}
  \tightlist
  \item
    a random nonce,
  \item
    a ciphersuite to use,
  \item
    and an ephemeral Diffie-Hellman public key
  \end{itemize}
\item
  Both partices compute a shared session key \(k\) using Diffie-Hellman
  key agreement on the ephemeral keys they exchanged.
\item
  Under encryption using the keys derived from the session key \(k\),
  the server sends the certificate for \texttt{mit.edu} as well as a
  signature over all messages sent so far, using its long-term secret
  key \(\sk_\text{MIT}\). The client then checks that:

  \begin{itemize}
  \tightlist
  \item
    the certificate has been signed by one of the client's trusted CAs,
    and
  \item
    the signature from the server matches their own record of the
    messages.
  \end{itemize}
\item
  Finally, the client and server run the TLS Record Protocol to exchange
  encrypted and authenticated application data.
\end{enumerate}

This simplified toy version of the TLS handshake does not provide many
of the features that the real TLS handshake provides. But it should give
you a flavor of what the real handshake looks like.

\section{Properties that TLS does not
provide}\label{properties-that-tls-does-not-provide}

\textbf{Authenticated End-of-File.} TLS does not provide any end-of-file
authentication, or ``clean closure.'' To explain what this means by
example:

A popular tool to install the toolchain for the trendy systems
programming language Rust is \texttt{rustup}. To use the tool and
install the Rust toolchain, the recommanded method is to run the command
``\texttt{curl\ https://sh.rustup.rs\ \textbar{}\ sh}.'' This downloads
a shell script from the internet over HTTPS and immediately runs it
using the shell \texttt{sh}.

Imagine that the contents of the downloaded script create a temporary
directory, copy things into it, install some things, and finally delete
the temporary directory with something like
\texttt{rm\ -r\ /tmp/install}. An in-network attacker, who knows the
structure of the \texttt{rustup} install script, could drop all of the
packets in the stream immediately after the characters
\texttt{rm\ -r\ /}. Eventually, the TLS connection will timeout, an
shell will run the command \texttt{rm\ -r\ /}, deleting the user's
entire file system.

To protect against this, script writers try to design their scripts such
that if the stream is cut off in the middle of a download, nothing
happens. For example, the install script might consist of a single
function definition that is called at the very end of the function.

\textbf{Plaintext Length Obfuscation.} As we have discussed, encryption
reveals exactly the length of the plaintext. If there is data that is
not encrypted that is then included inside the encrypted data as well,
this can cause a vulnerability---see the CRIME attack.

\chapter{Open Questions in
Encryption}\label{open-questions-in-encryption}

So far, we have established what may seem like a comprehensive set of
tools to transmit data over the network: we have schemes for verifying
the integrity of data and for hiding the contents of a transmission from
an adversary, both with and without a shared key.

Transport-layer security (TLS) effectively builds and ``encrypted pipe''
between a client and a server. Through the encrypted pipe that TLS
provides, we can run any of our favorite TCP-based protocols---HTTP,
SMTP, POP, IMAP, etc.---and can thus hide our data from an in-network
attacker. And yet, the security guarantees that TLS provides fall short
of the strongest possible security notions we could desire.

If we were to imagine the best possible security we could ask for
regarding network traffic, it might look (imprecisely) something like
the following:

\begin{tcolorbox}[enhanced jigsaw, arc=.35mm, bottomrule=.15mm, breakable, colback=white, colframe=quarto-callout-note-color-frame, left=2mm, leftrule=.75mm, opacityback=0, rightrule=.15mm, toprule=.15mm]

\emph{An attacker who controls many parties (clients and servers) as
well as the network should ``learn nothing'' about who the client is
talking to and what she is saying.}

\end{tcolorbox}

Unfortunately, the protocols we have for secure communication today fall
far short of this goal. In this section, we describe some of the
shortcomings of today's transport-security tools and some imperfect
solutions.

\section{Problem: Encryption does not hide the source and destination of
a
packet}\label{problem-encryption-does-not-hide-the-source-and-destination-of-a-packet}

In order to send IP packets over the internet, the Internet's routing
system relies on routers in the network knowing the source and
destination IP addresses of each packet: these are included,
unencrypted, in the packet header. (In some ways, the ``pipe'' analogy
fits here: anyone can see where the pipe starts and where it ends.)

\textbf{Solution Attempt: Tor.} The Tor system aims to allow a client to
connect to a server over the Internet while hiding---from certain types
of adversaries---which server the client is connecting to. For example,
a Tor client should be able to browse the web without anyone learning
which websites the client is visiting.

Tor's strategy is to bounce traffic around the Internet and hope that no
real-world attacker can gather enough information to figure out which
client is communicate with which server. Tor provides no precise
security guarantees, and there are scores of research papers
demonstrating various weaknesses in Tor's security plan. At the same
time, Tor is publicly available, is well supported, is widely used, and
seems to provide some meaningful privacy benefits in practice.

\marginnote{\begin{footnotesize}

It is difficult to evaluate the security of a tool like Tor, since
real-world attackers will not necessarily reveal that they can break the
tool's security guarantees. So using a tool like Tor requires taking a
leap of faith.

\end{footnotesize}}

Tor works by nesting several of these encrypted pipes: when opening a
connection, the Tor client will first select three \emph{relays} from
the Tor network \((A, B, C)\). The Tor client will then:

\begin{enumerate}
\def\labelenumi{\arabic{enumi}.}
\tightlist
\item
  Open an encrypted tunnel to the first relay \(A\) (the ``guard'').
\item
  Through that tunnel, open an encrypted tunnel to the second relay
  \(B\).
\item
  Through that tunnel-inside-a-tunnel, open an encrypted tunnel to the
  third relay \(C\) (the ``exit'').
\end{enumerate}

The client will then send its application-layer traffic through this
tunnel-inside-a-tunnel-inside-a-tunnel. So each byte of application data
will be encrypted first for relay \(C\), then for relay \(B\), then for
relay \(A\). When the client sends this ciphertext over the
\emph{circuit} from relay \(A\) to \(B\) to \(C\) to the real
destination, relay \(A\) will first strip off its layer of encryption
then forward the inner packet to relay \(B\). Relay \(B\) will do the
same, stripping off a layer of encryption and forwarding the packet to
relay \(C\). Finally, relay \(C\) will strip off the last layer of
encryption and be left with a normal IP packet that it can then send to
the destination server. As the response makes it back through the
network, each relay node will add a layer of encryption. The end result
of this is that no single relay can see the source and destination IP
addresses.

However, the security that Tor provides is imperfect. First, if an
attacker controls the guard node (relay \(A\)) and the exit node (relay
\(C\)), the attacker can correlate the timing of when a packet enters
the guard node and when a packet exits the exit node. Using this timing
an attacker can make a guess at the route traffic is taking through the
Tor network. Even without controlling relay nodes, if an attacker
controls certain key points in the Internet (e.g., Internet exchange
points or undersea fiber links) it may be able to perform this sort of
traffic analysis even without controlling relays.

\section{Problem: Attacker sees packet sizes and
timings}\label{problem-attacker-sees-packet-sizes-and-timings}

As we discussed, practical encryption schemes necessarily reveal the
length of the ciphertext. In the context of the Internet, this means
that an attacker can learn the size of each TCP packet that a client
sends, along with timing information. (Here the ``encrypted pipe''
analogy for TLS breaks down: an attacker can see how much traffic flows
through the encrypted pipe and when.)

Even without seeing the destination and source of packets sent to and
from a client's machine, an attacker can learn significant amounts about
the client's traffic. Some examples are:

\begin{itemize}
\tightlist
\item
  \emph{Watching a movie:} Video traffic has a distinct traffic pattern.
  By monitoring, for example, the length of time that a client spends
  watching a movie, a network attacker learn with fairly high accuracy
  \emph{which movie} the client is watching.
\item
  \emph{Using \texttt{ssh}:} Different commands will have different
  traffic patterns. An attacker may be able to infer what type of
  commands a client is running by inspecting traffic patterns.
\item
  \emph{Downloading a file:} The bitlength of a downloaded file can
  uniquely identify the file in many cases.
\item
  \emph{Browsing the web:} sizes leak individual pages.
\end{itemize}

\textbf{Example: New York Times.} The New York Times homepage
\texttt{nytimes.com/} downloads 1.56 MB of content, along with 76 total
assets (images, CSS, JavaScript). The webpage to submit a sensitive tip,
\texttt{nytimes.com/tips}, downloads 41.92 KB and only 15 assets. By
counting the number of HTTPS requests that a client makes over an
encrypted connection to \texttt{nytimes.com}, an attacker can easily
distinguish whether a client is visiting the homepage or the tips page,
even if the attacker cannot decrypt even a single bit of the HTTPS
traffic itself.

\textbf{Attempts at a solution.} There are several common ways that
people attempt to protect against this sort of traffic analysis. None of
these solutions works well.

\begin{enumerate}
\def\labelenumi{\arabic{enumi}.}
\item
  \textbf{Random Noise:} To try to hide the length of the packets it
  sends, a client can add a randomly chosen number of bytes of dummy
  data to the end of each packet. The hope is that by randomizing packet
  lengths, the client prevents the attacker from performing the traffic
  analysis.

  Unfortunately, a patient attacker can use \emph{averaging} to
  effectively eliminate the effect of the random noise. That is, if the
  attacker can trick the client into sending the same message a few
  times (as is often possible), the attacker can average the noised
  packet lengths to get a good estimate of the true length.
\item
  \textbf{Padding:} Another option is to just pad every packet (or
  webpage or encrypted message, etc.) to match the largest packet that
  the client will ever send. For example, whenever the client visits a
  page on \texttt{nytimes.com}, the client could download 50MB of page
  content and 100 fixed-size assets, even if the true page is tiny. This
  is somewhat secure, but incredibly costly and therefore not practical
  outside of very specific circumstances.
\end{enumerate}

\section{A Promising Direction: Metadata Privacy for
Messaging}\label{a-promising-direction-metadata-privacy-for-messaging}

Messaging apps like WhatsApp and iMessage are end-to-end encrypted, but
still may leak who you are talking to. This problem is more tractable
due to the circumstances of messaging:

\begin{itemize}
\tightlist
\item
  Messages are approximately fixed length.
\item
  Some latency is OK.
\item
  Total daily traffic per user is small.
\item
  Each user talks to few messaging partners.
\end{itemize}

Because of these constraints, it may be feasible to use techniques like
padding to greatly reduce the amount of data that messaging metadata
reveals and to do so in a way that provides strong formal guarantees
about security. But still, no widely used messaging app provides any
sort of metadata-privacy guarantees.

\section{Problem: Endpoint
Compromise}\label{problem-endpoint-compromise}

Say that we have a perfect scheme for transport security---one that
hides all data and metadata. Such a scheme is still not enough to
protect our data if an adversary can compromise the communication
\emph{endpoints}.

For example, in many applications, a client sends some sensitive data to
a server (e.g., its Google search queries). The server is free to lose
it in a breach, sell it, or turn it over to law enforcement agencies,
etc. Later on, we will discuss how to protect against server compromise
using software-engineering techniques. Here, we will give one example of
how cryptography can protect user data even against a compromised
server.

\subsection{Private Information
Retrieval}\label{private-information-retrieval}

In many applications, a client must read a record from a database stored
at a server. The client might like to perform such a database query
without the server learning which record it accessed.

\marginnote{\begin{footnotesize}

A concrete application of this is Google search---in order to give you
search results, Google necessarily learns what you are searching. With a
PIR scheme, it would be possible for Google to look up search results
without learning what you are searching for!

\end{footnotesize}}

In a \emph{private information retrieval} scheme:

\begin{itemize}
\tightlist
\item
  the server holds a public database of \(n\) bits:
  \(x_1, x_2, \ldots, x_n \in \bin\), and
\item
  the client holds a secret index \(i \in \{1, \dots, n\}\).
\end{itemize}

The client and server interact. At the end of the interaction we want
the following properties to hold:

\begin{itemize}
\tightlist
\item
  \emph{Correctness:} The client outputs \(x_i \in \bin\).
\item
  \emph{Security:} The server ``learns nothing'' about the client's
  secret index \(i\). In particular, we demand that the message that the
  client sends to the server is a CPA-secure encryption of its index
  \(i\).
\end{itemize}

\textbf{Naïve private information retrieval.} The simple scheme for
private information retrieval is to have the server send all \(n\) bits
of the database to the client. When the database is large, as it is for
Google search, this would be an infeasible amount of communication. A
surprising fact is that there are simple private-information-retrieval
protocols that involve much less than \(n\) bits of client-server
communication.

\subsection{A non-trivial private-information-retrieval
scheme.}\label{a-non-trivial-private-information-retrieval-scheme.}

To achieve this, we need a new tool called \emph{additively homomorphic
encryption}.

\textbf{Additively homomorphic encryption.} A secret-key additively
homomorphic encryption scheme is a CPA-secure secret-key encryption
scheme \((\Enc,
\Dec)\) over key space \(\calK\) and message space in \(\calM = \Z_p\)
with the added property that for all keys \(k \in \calK\) and all
messages \(m, \hat{m} \in \calM\),
\[ \Enc(k, m) \star \Enc(k, \hat{m}) = \Enc(k, m + \hat{m}), \] where
``\(\star\)'' is some fixed binary operation on ciphertexts.

In English: given two encrypted messages \(m\) and \(\hat m\), encrypted
under an additively homomorphic encryption scheme, anyone can compute
the encryption of \(m + \hat m\). Being able to add encrypted messages
also allows multiplying encrypted messages by public constants, since
\(m + m = 2m\) and \(2m + 2m = 4m\) and so on. Once we can add and
multiply by constants, and we can compute a matrix-vector product of an
encrypted vector and a public matrix.

It is possible to construct an additively homomorphic encryption scheme
from the DDH assumption, with only a slight tweak to ElGamal encryption.

\textbf{PIR Construction.} We can use additively homomorphic encryption
to construct a private-information-retrieval scheme. To do so, the
server represents its its database (the \(x_i\) values) as a \(\sqrt{n}
\times \sqrt{n}\) matrix \(D\). The client then tells the server which
column \(j\) it would like by supplying the encryption of a
\(\sqrt{n} \times 1\) vector \(\vec{m}\) with a 1 in the \(j\)th
location. The client sends this encryption (using a key only the client
knows) as \(\Enc(k, \vec{m})\). The server computes the matrix product
\(\Enc(k,
D\vec{m})\), which gives the \(j\)th column of the matrix, using
additively homomorphic encryption and returns the response to the
client, where the client can find the bit they are interested in in the
column.

This allows a client to retrieve a bit from a server's database without
the server learning anything about the desired bit, and to do so at the
communication cost of only \(2\sqrt{n}\) ciphertext. The server
computation cost is high---the server necessarily touches every bit of
the database---but at least it shows that making private queries is
feasible in theory.

\part{Platform Security}

\chapter{Architecting a secure
system}\label{architecting-a-secure-system}

So far, we have been focusing on security for network communication. We
have established many tools to achieve this, from message authentication
codes to public-key encryption.

Ultimately, however, applications need to make use of these tools. And
for our network security tools to provide meaningful security, the
applications themselves must be reasonably secure.

In discussing platform and application security, there are two classes
of problem that we want to defend against.

\textbf{1. Mistakes of various types.}

\begin{itemize}
\tightlist
\item
  Buggy systems: including hardware and software bugs
\item
  User mistakes: phishing, misconfiguration
\end{itemize}

\textbf{2. Malicious people or components.}

\begin{itemize}
\tightlist
\item
  Malicious components: malware, supply-chain attacks
\item
  Malicious users: what if the adversary gets the admin's password?
\item
  Attacker gets access to the system: insider attacks, the adversary
  guesses credentials,
\end{itemize}

In computer security, we tend to treat mistakes/bugs and malicious
software/components in the same way. We do that because (1) it's often
difficult to specify what it means for a component to be
``non-maliciously buggy'' and (2) an attacker can often leverage what
seems like a benign bug into full-fledged misbehavior.

Thus, a theme that will be present throughout this section is that we
will consider mistakes to be malicious: if we are prepared to handle
malicious components, we will similarly be prepared to handle our own
buggy code. If we are prepared to limit damage of a malicious user with
the admin password, we will also be limiting the damage that a
mistake-making admin can cause.

This multitude of threats makes designing secure applications quite
difficult. To make progress, we will seek to design systems that limit
damage when things go wrong.

When we design for security, we have essentially three goals:

\begin{enumerate}
\def\labelenumi{\arabic{enumi}.}
\tightlist
\item
  \textbf{Defend against known attacks.}
\item
  \textbf{Defend against unknown attacks.}
\item
  \textbf{Limiting damage.} In the cryptography part of this course,
  systems are either security or insecure. In systems security, things
  are much more gray. Attacks we care about often are outside of our
  threat model---even when this happens, we'd like to somehow contain
  the damage.
\end{enumerate}

\section{Isolation}\label{isolation}

One of the most effective strategies to limit damage is to split a
system into isolated components. If one of these components becomes
compromised, it should not be able to compromise the other components.
For example, if you run code in one virtual machine, it should not be
able to tamper with data in another virtual machine.

\marginnote{\begin{footnotesize}

When choosing what mechanism to use to isolate various components, we
think a lot about the \emph{performance overhead} of an isolation
mechanism. The challenge of building a good isolation mechanism is
ensuring strong isolation without slowing down the isolation components
too much (or taking up too much extra memory).

\end{footnotesize}}

These components will typically run on top of some \emph{host} that
enforces isolation. Importantly, this host must be correct! If there are
bugs in the host, malicious code in a component may be able to exploit a
bug to escape its isolation. The success of an isolation mechanism
depends on the correctness and configuration of the host.

\begin{longtable}[]{@{}rl@{}}
\caption{Some common types of
isolation}\label{tbl-isolationtypes}\tabularnewline
\toprule\noalign{}
\textbf{Examples} & \textbf{Host} \\
\midrule\noalign{}
\endfirsthead
\toprule\noalign{}
\textbf{Examples} & \textbf{Host} \\
\midrule\noalign{}
\endhead
\bottomrule\noalign{}
\endlastfoot
Docker Container & Operation System (e.g.~Linux) \\
Browser tabs & Browser \\
Language-Level (JavaScript, Wasm) & Language Runtime \\
Process & Linux kernel \\
Virtual machines (VMs) & VM Monitor \\
Physical (``air gap'') & Physics \\
\end{longtable}

In order for these isolated components to be useful, they will need to
be able to talk to each other in some form. For example, a client
component must be able to make requests to a database component, but we
would like to limit the power of the client to do damage. For this, we
would like to achieve \emph{controlled sharing}.

\subsection{Controlled Sharing}\label{controlled-sharing}

For an isolation mechanism to be useful, it additionally needs to have
some way to interact with other isolated components. For example, some
JavaScript code isolated in a browser tab still needs some means by
which to make requests over the network.

When a host decide whether to allow a request from a particular
component, it typically needs to do three things with each request:

\marginnote{\begin{footnotesize}

Since all three of these actions start with the letters ``Au,'' we
sometimes call this the \emph{gold standard} for controlled sharing.

\end{footnotesize}}

\begin{itemize}
\tightlist
\item
  Authenticate: Associate the request with some \emph{principal}. A
  principal could be a user name, an ``origin'' in the web context
  (e.g., \texttt{google.com}), a program, or some other entity in the
  system.
\item
  Authorize: Decide whether that principal is allowed to make the
  request.
\item
  Audit: Keep track of requests that each principal makes. Auditing is
  about limiting damage: often a host will mistakenly allow requests it
  shouldn't; audit logs make it easier to discover such mistakes and to
  clean up afterwards.
\end{itemize}

It is crucial that an isolation mechanism perform these three checks on
every single request---a single hold in the isolation boundary is often
enough to completely break any benefits isolation that would have
provided.

\section{Authentication}\label{authentication-1}

Since we already had an entire module on authentication using
signatures, MACs, passwords, and so on, we will not discuss
authentication further here.

\section{Authorization Policies}\label{authorization-policies}

In order to authorize requests, we need some sense of permissions---a
mapping from \emph{objects} to \emph{principals} that can access them.
We call these permissions the authorization \emph{policy}.

\textbf{Storing policies.} We can think of an authorization policy as a
gigantic matrix with one row per object and one column per principal.
For example, in a file system, we could have one row per file, and one
column per user in the system. The entry in column \(i\) and row \(j\)
lists the actions that user \(i\) can perform on file \(j\): read,
write, execute, etc. A common way of storing this gigantic matrix, for
example in AFS, is via an \emph{access control list} for each object.

\textbf{Setting policies.} There are many approaches to setting
authorization policies. As always in computer systems, there is no one
perfect solution:

\begin{itemize}
\item
  \textbf{Discretionary access control: ``Owner'' of each object sets
  the policy.} This approach is useful in file systems---each file has
  an owner and the owner can determine who has access to the file. A
  problem is that if an attacker hijacks the owner's account (or just
  one application that the owner runs), the attacker can tamper with the
  policy for all of the user's files. In addition, it may be difficult
  for non-expert users to set policies.
\item
  \textbf{Mandatory access control: Administrator sets policy.} This
  approach often is useful in a large organization, when administrators
  have opinions about which user should have access to which files or
  systems. A classic example of this is for systems that handle
  classified data in government systems. Normal users of the system
  cannot give unprivileged users access to a classified files.

  One limitation of this approach is that it is very coarse grained:
  administrators may not know exactly who should have access to what.
\end{itemize}

Systems in practice often use some combination of both of these
strategies.

\marginnote{\begin{footnotesize}

\emph{Role-based access control} tries to hit some midpoint between
discretionary and mandatory access control. In these systems, there is a
centrally defined set of ``roles.'' In a university, these could be
``Students,'' ``Faculty,'' and ``Staff.'' The security administrator
assigns users to roles. Then application developers determine which
roles have access to the application.

\end{footnotesize}}

Common issues are:

\begin{itemize}
\tightlist
\item
  It is difficult to keep policies up to date as the set of users
  evolves. Expiring permissions is one strategy.
\item
  Users will complain if they do not have enough permissions, but they
  will never complain if they have too many permissions. As a result,
  users often end up with more access than they need from a security
  standpoint.
\end{itemize}

\section{Auditing}\label{auditing}

We have relatively little to say about this. The most important thing to
remember about auditing is that a system should store the audit logs in
a container that is separate from the container holding application
logic. That is important because if the attacker compromises the
application, it should be difficult for the attacker to compromise the
logs as well.

\section{Delegation and Chained
Requests}\label{delegation-and-chained-requests}

Users often interact with systems \emph{indirectly}. For example, when
accessing Gmail, a user's browser first sends a request to the Gmail
server asking for new messages. The Gmail server then sends a request to
the database to fetch the message data.

For the first request, it is fairly clear that the principal should be
Alice: the request is coming from Alice's browser, and therefore should
have been initiated by Alice directly. Alice will send some credential
to the server, and the server can use this credential to verify that it
is really Alice on the other end. For the second request, however, it is
not as clear who the request should be from.

One option is to have the request come from Alice. This protects against
compromise of the Gmail server---the adversary cannot see all user data.
Systems like SSH and AFS follow a strategy like this. Another option is
for the principal of this second request to be the Gmail server itself.
This helps with isolation among services access the same database: if
the Google calendar code is buggy and gets compromised, the first plan
would allow an adversary to view Alice's gmail data even is the gmail
service was perfectly secure. However, it does not protect other users
from a buggy Gmail service.

\textbf{Compound Principal: ``B for A.''} To achieve something stronger,
we can create a new type of \emph{compound principal} that combines a
service or device with a user. For example, this server-to-database
request could carry a principal of ``Gmail Server for Alice''. This
provides protection against both gmail server compromise and against
compromise of other services.

However, it is not as clear how to actually implement this. One option
is to continue to have Alice send her credential to the server directly.
However, then the server can totally impersonate Alice and we gain
little protection. What we would really like is for Alice to give
permission to the Gmail server to fetch her emails, but not to do
anything else. This is called \emph{delegation}.

\textbf{Delegation with cryptography.} In interacting with the Gmail
server \(B\), we may like for Alice (\(A\)) to give \(B\) permission to
authenticate as ``\(B\) for \(A\)'' and to do so for only 60 seconds
into the future. To achieve this, \(A\) can sign a message that outlines
the permission it would like to give to \(B\). This signature becomes
the proof of authorization. As an example:
\[ \Sign(\sk-A, \text{“A delegates to B”}, \text{start}=\text{now}, \text{end}=\text{now}+60) \]

Google indeed uses a strategy like this. They have a global
DoS-resilient HTTP front-end that performs initial authentication. This
frontend is then responsible for generating these scoped delegation
signatures for each operation that the user would like to do and sending
them along to the individual services. These signatures are then used
for all following operations.

\textbf{Capabilities.} We may want more fine-grained access control. For
example, on Android, the Gmail app may like to delegate permission to a
PDF viewer to view an attachment. However, if all the attachments are
stored in some common database, we would like to avoid giving the PDF
viewer access to view everything in the database. To achieve this,
Android (and systems more generally) use a plan called cababilities.

A similar strategy can be seen, for example, in cloud file sharing: when
you share a file in google drive, it generates a long random link that
allows anyone with access to that link to view that file (and no
others). This link itself becomes a \emph{capability}---it allows anyone
that posesses it to perform some related action.

\chapter{Isolation}\label{isolation-1}

In many settings when building systems, it is useful to \emph{isolate}
different components of a computer system. For example, cloud providers
such as Amazon Web Services' EC2 run multiple virtual machines, your
phone runs several apps, and your browser runs many sites together. If
one of these virtual machines (or apps or websites) is malicious or
buggy, these platforms would like to protect the buggy component from
interfering with the execution of other code in the system. We often
refer to separate isolated components as running in separate
\emph{isolation domains}.

Ideally, we could have each isolation domain run on separate physical
computer. If this were the case, (modulo physical side-channel attacks)
one domain would clearly not be able to touch another's state.
Unfortunately, the cost of this physical ``air gap'' isolation is far
too high for most practical applications.

\marginnote{\begin{footnotesize}

Computer systems for high-stakes applications (e.g., classified
government data) do in fact use this type of ``air gap'' isolation.

\end{footnotesize}}

So, when creating isolation methods, we aim to make it appear as if each
domain is running on a separate computer but to do so all on the same
computer.

\section{Defining Isolation}\label{defining-isolation}

In order to define isolation, we will think about two key properties:
integrity and confidentiality.

\subsection{(Weak) Integrity}\label{weak-integrity}

One property that we would like an isolation scheme to provide is
\emph{integrity}: one domain cannot modifying the state of another
domain. To formalize this, let's consider two domains running on a
single host, an adversary domain \(A\) and a victim domain \(V\). We
would like to guarantee that \(A\) cannot change the execution of \(V\).
We can define something in terms of the \emph{state} of each domain,
\(S_A\) and \(S_V\). For any pair of starting states \((S_A, S_V)\),
after running \(A\), we would like the new pair of states to be
\((S_A', S_V)\): \[ (S_A, S_V) \xrightarrow{\text{run} A} (S_A', S_V) \]

\subsection{(Weak) Confidentiality}\label{weak-confidentiality}

We would also like \emph{confidentiality}: an adversarial domain should
not be able to read a data from any other domain. We can formalize this
by considering two worlds, each with a \emph{different} victim state.
After running \(A\) in each of these worlds, we would like the resulting
\(S_A'\) to be identical. That is, for all pairs of victim states
\(S_V^1, S_V^2\), running \((S_A, S_V^1)\) results in the same
adversarial state as running \((S_A, S_V^2)\):
\[ (S_A, S_V^1) \xrightarrow{\text{run} A} (S_A', -) \]
\[ (S_A, S_V^2) \xrightarrow{\text{run} A} (S_A', -) \] This definition
of confidentiality is often called ``non-leakage''.

\subsection{Non-interference: Strong confidentiality and
integrity}\label{non-interference-strong-confidentiality-and-integrity}

In our definitions of confidentiality and integrity so far, only the
adversary domain runs. In a real system, both the adversary and the
victim domain will run concurrently. Ideally, we would like to ensure
that our confidentiality and integrity properties hold under interleaved
execution of the adversary and victim.

To achieve this stronger notion of security, we can include an
interleaving of \(A\) and \(V\) in one world---we would like for the
resulting \(S_A'\) to be identical whether or not \(V\) runs.

\[ (S_A, S_V) \xrightarrow[A, V, V, A, A]{\text{run}\, A, V} (S_A', -) \]
\[ (S_A, S_V) \xrightarrow[A, A, \ldots, A]{\text{run only}\, A} (S_A', -) \]

This definition is often called \emph{non-interference}. We can
similarly strengthen our definition of integrity by requiring that
\(S_V'\) is identical after running \(V\) whether or not \(A\) is run.
\[ (S_A, S_V) \xrightarrow[A, V, V, A, A]{\text{run}\, A, V} (-, S'_V) \]
\[ (S_A, S_V) \xrightarrow[A, A, \ldots, A]{\text{run only}\, V} (-, S'_V) \]

\subsection{Non-interference is difficult to
achieve}\label{non-interference-is-difficult-to-achieve}

To achieve a non-interference style of isolation, an adversarial process
\(A\) must not be able to determine whether there is a victim \(V\)
process running alongside it concurrently. The challenge, though, is
that whenever the adversary and victim \emph{share limited
resources}---such as CPU, RAM, network bandwidth, hard disk space,
etc.---it is almost always possible for the adversary to determine
whether there is a victim process running concurrently.

\marginnote{\begin{footnotesize}

The information leakage across isolation boundaries as a results of
resource contention is one type of \emph{side channel} or \emph{covert
channel}. There is a vast literature on how to construct and exploit
various types of side channels that leak information from a victim to an
adversary.

\end{footnotesize}}

\textbf{Example: Memory Allocation.} A real system will have some bound
on the amount of memory available to it. After this memory is used, the
system will be unable to allocate any additional memory. Consider a
system with 16GB of memory and a victim process that allocates memory
based on the value of some secret:

\begin{verbatim}
int secret;
malloc(secret);
\end{verbatim}

An adversary could repeatedly try to allocate memory until the system
tells them they cannot. By keeping track of the amount of memory they
were able to allocate, the adversary can learn the secret: if the
adversary is able to allocate 15GB, the adversary will know that the
secret is \(2^{27}\), as the victim must have allocated 1GB.

\textbf{Example: Execution Time.} Consider another victim that runs some
computation that takes a variable amount of time to finish depending on
the value of a secret. By keeping track of how long the adversary takes
to finish, the adversary can learn how much execution time the victim
running on the same system takes to finish. Using this information, the
adversary may be able to learn information about the secret. This type
of information transfer are often called ``timing channels'', and can be
quite tricky to work with.

There are effectively three ways to deal with the fact that
non-interference is generally impossible to achieve with shared limited
resources:

\begin{itemize}
\item
  \textbf{Strictly partition resources to prevent contention.} Each
  isolation domain could run on a separate physical machine, or we can
  provision the resources on a machine are partitioned in such a way
  that there is never contention for resources between isolation
  domains.
\item
  \textbf{Prevent isolation domains from detecting resource contention.}
  Another (more practical) approach is to restrict the types of programs
  that can runs in such a way that prevents the programs in the system
  from detecting contention in shared resources. For example, if all
  programs in an isolated system are deterministic functions with no
  access to the outside world---no system calls, no networking,
  etc.---then programs may not be able to \emph{detect} resource
  contention when it exists. In most implementations of isolation (e.g.,
  virtual machines in a cloud environment), isolation domains absolutely
  need access to the outside world, so this approach is rarely useful.
\item
  \textbf{Give up on non-interference.} Most isolation mechanisms opt
  for this solution. Rather than trying to achieve strict
  non-interference, we aim for some ``good enough'' notion of isolation.
  Linux, for example, does not attempt to achieve strict
  non-interference between processes running on the same physical
  machine.
\end{itemize}

\section{Implementing Isolation}\label{implementing-isolation}

In principle, implementing isolation in a system involves three main
steps:

\begin{enumerate}
\def\labelenumi{\arabic{enumi}.}
\tightlist
\item
  identify the state for each domain,
\item
  identify operations that access state, and
\item
  ensure that state-modifying operations can only read/write the state
  within an isolation domain.
\end{enumerate}

The challenge is performance. An isolation mechanism will have to
perform checks to ensure that components in one isolate cannot influence
another. The game in isolation is to provide strong isolation at the
minimum possible cost.

We will start by considering simpler isolation mechanisms and then look
at more sophisticated ones.

\subsection{Emulation}\label{emulation}

One simple way to implement isolation is to have an interpreter that
executes isolated programs (e.g., x86 programs). The interpreter
inspects each opcode in the program one at a time, and then implements
the operations on the isolated state that the opcode indicates. While
processing each opcode, the interpreter enforces checks on the isolated
program to ensure that it can only modify its local state.

Some JavaScript engines (``runtimes'') in web browsers use emulation for
isolation. The Python interpreter is another example of emulation-based
isolation. The runtime for these languages is designed such that the
code can access only memory that belongs to the domain.

\textbf{Benefits.} Emulation can be simple to implement and provides
``good enough'' isolation for many applications.

\textbf{Downsides.} Emulation can be slow: to execute each logical
opcode in the emulated program, the emulator may have to run a large
number of physical instructions. Emulation can be inflexible: emulated
programs may not be able to take advantage of special-purpose hardware,
unless the emulator explicitly grants access to these devices to
emulated programs.

\subsection{Time Multiplexing}\label{time-multiplexing}

Another effective isolation strategy is to only allow the code from one
isolation domain to run on the hardware at once. When the isolation
mechanism switches from one isolation domain to another, it writes the
state of the hardware (e.g., registers, RAM) to storage, clears the
state of the hardware (e.g., zeros the contents of memory), and loads
the next program to run into the machine.

For example, gaming consoles implement this form of isolation: one game
runs at a time and has almost full control of the hardware during this
time. The state of the running game cannot tamper with the state of a
non-running game.

The security of this isolation scheme requires the isolation mechanism
to protect the stored state of each isolation domain, and to somehow
protect the code that switches between them.

\textbf{Benefits.} Time multiplexing is relatively simple to implement,
and it can give programs in each isolation domain almost full control of
the hardware (e.g., graphics hardware in game consoles)..

\textbf{Downsides.} There can be a high cost of switching execution
between between isolation domains. Since the isolation mechanism clears
the state of the hardware during context switches, storing and restoring
the state can be costly.

\subsection{Translation (Naming)}\label{translation-naming}

Translation is another common isolation mechanism in computer systems.
In a system using translation for isolation, an isolated program may not
access hardware resources (such as memory or files) directly. Isolated
programs can only access resources via pointers controlled by the
isolation mechanism. By construction, each process in an isolation
domain can only name resources inside of its isolated context.

The canonical example of naming/translation for isolation is
\emph{virtual memory}. In a system using virtual memory, the isolated
domains cannot read/write physical memory directly---they can only
read/write to virtual memory addresses..

To prevent one isolated program from accessing another's memory, the
isolation mechanism ensures that the valid virtual addresses in each
domain point to a separate portion of physical memory. Modern CPUs have
special support for implementing virtual memory to make the
virtual-address translation as fast as possible.

Other examples of naming/translation isolation in computer systems are:
file descriptors in Linux and virtual LANs in networking.

\subsection{Example: Isolation in Virtual
Machines}\label{example-isolation-in-virtual-machines}

When we run several virtual machines on one physical machine, we want
each virtual machine to run as if it had its own physical CPU, memory,
and devices, but we want to run all of these virtual machines on a
single physical machine. For performance, we would like to run
instructions from the VM directly on the host CPU---but we need to make
sure, for example, that the VM does not access memory belonging to
another VM. To achieve isolation and good performance, systems today use
several effective techniques.

Here is how a virtual-machine monitor handles the three questions that
an isolation mechanism must answer:

\begin{enumerate}
\def\labelenumi{\arabic{enumi}.}
\tightlist
\item
  \textbf{What is the isolation domain's state?} A few important pieces
  are:

  \begin{itemize}
  \tightlist
  \item
    the contents of memory,
  \item
    the values in the CPU registers,
  \item
    the data on disk.
  \end{itemize}
\item
  \textbf{What operations can a program perform on isolated state?}
  Virtual-machine monitors need to handle:

  \begin{itemize}
  \tightlist
  \item
    CPU instructions that modify register states ,
  \item
    CPU instructions that modify memory, and
  \item
    operations to read and write devices.
  \end{itemize}
\item
  \textbf{How does the virtual-machine monitor ensure isolation?}

  \begin{itemize}
  \tightlist
  \item
    To handle updates to the register state, the monitor uses \emph{time
    multiplexing}: one virtual machine runs on the CPU at a time and
    controls the CPU's registers.
  \item
    To handle accesses to memory, the monitor uses
    \emph{naming/translation} via virtual memory.
  \item
    To handle device operations, the monitor uses \emph{emulation}. When
    a virtual machine executes instructions that would cause device I/O,
    the CPU jumps into some dedicated code in the monitor that emulates
    these device operations. This is slow, but since device I/O is
    usually costly anyways, the overhead is isolation is tolerable.
  \end{itemize}
\end{enumerate}

\subsection{Software interposition}\label{software-interposition}

A final isolation technique that we will discuss is \emph{software
interposition}. Say that an isolated program wants to run the following
code:

\begin{verbatim}
var a = b[c];
var f = ....;
f();
\end{verbatim}

When using software isolation, a compiler can insert checks into this
isolated program to make sure that the memory access \texttt{b{[}c{]}}
does not access out-of-bounds memory:

\begin{verbatim}
if c >= b.size: error;
load b+c -> a             (***)
\end{verbatim}

After inserting these checks, the isolated code can run directly on
hardware, since the checks ensure that the isolated code cannot touch
any state outside of its isolated context.

A central challenge of using software interposition is handling function
calls. When calling the function \texttt{f} in the code snippet above,
the hardware might execute an instruction like:

\begin{verbatim}
jump *f
\end{verbatim}

which would execute the code at the location stored in memory location
\texttt{*f}.

If an isolated program can set the value in \texttt{*f} to the location
\texttt{(***)} in the snippet above, the isolated program can skip the
safety checks on the memory accesses.

\chapter{Software Trust}\label{software-trust}

The central question of this chapter is:

\begin{quote}
\emph{How do we know whether a system is running the software that we
expect it to be running?}
\end{quote}

This question comes up both when we are interacting with a machine in
person (e.g., typing a passcode into our phone) or across a network
(e.g., when sending sensitive data to a far-away server).

\marginnote{\begin{footnotesize}

A separate question, which we will discuss in future chapters is: How do
we ensure that the software itself is ``good'' or ``bug-free?''

\end{footnotesize}}

\textbf{Threats to software integrity.} There are a number of threats
that might cause a machine to be running unexpected software:

\begin{itemize}
\tightlist
\item
  \textbf{Malware.} An adversary may install bad software onto the
  laptop, such as a keylogger.
\item
  \textbf{User error.} A user may inadvertently install malware onto a
  machine.
\item
  \textbf{Software supply-chain problems.} An adversary may inject
  malware into a real app's libraries, by tricking or coercing a
  developer.
\item
  \textbf{Malicious updates.} An adversary may trick the software update
  process, converting a real piece of software into a malicious one.
\end{itemize}

\marginnote{\begin{footnotesize}

A classic way to trick users into installing malware is to show them a
warning (e.g., on a webpage) that says ``Your computer is infected.
Please download and install this anti-virus software.''

\end{footnotesize}}

\textbf{The software supply chain.} There are many steps that take place
between the development of a piece of software and its use:

\begin{enumerate}
\def\labelenumi{\arabic{enumi}.}
\tightlist
\item
  Developers write code. Their code may use many third-party libraries.
\item
  Compilers build and package the application.
\item
  The software vendor distributes binaries over the network.
\item
  Users download new software.
\item
  Users download software updates.
\item
  Users launch applications on their devices.
\item
  Running applications interact with remote servers running code.
\end{enumerate}

\section{Library Imports}\label{library-imports}

\subsection{Example: Python Imports}\label{example-python-imports}

In Python, importing a library requires downloading a library using the
\texttt{pip} package manager:

\begin{verbatim}
    pip install requests
\end{verbatim}

After that a user can import the package into their code like this:

\begin{verbatim}
    import requests
\end{verbatim}

Behind the scenes, the PyPi service maintains a database that maps
package names (e.g., \texttt{requests}) to a piece of code. When you
type \texttt{pip\ install\ requests}, the \texttt{pip} program fetches
the code from the PyPi repository and installs it on your machine.

\textbf{Benefits.} Benefits of this approach are:

\begin{itemize}
\tightlist
\item
  The centralized service makes it easy for users to identify packages.
\item
  It is relatively easily for users to discover and download software
  updates.
\item
  Developers do not need to run their own code-distribution service.
\end{itemize}

\textbf{Downsides.} Downsides of this approach are:

\begin{itemize}
\tightlist
\item
  The centralized update service is a single point of failure: if an
  attacker is able to change the code in PyPi, it can infect a large
  number of machines at once.
\item
  The end user has no idea who actually produced the library code. The
  user only is able to specify the package name.
\item
  In Python, the naming scheme is ambiguous: if there is a public
  package and a private package both with the name \texttt{requests},
  it's not clear when a user imports \texttt{requests}, which one the
  user wants to import.
\end{itemize}

\subsection{Example: Go Library
Imports}\label{example-go-library-imports}

The Go programming language takes a slightly different approach to
package management. In the Go programming language, a user imports a
library/package by specifying the URL of the package's Git repository.
For example, an import might look like this:

\begin{verbatim}
    import "github.com/grpc/grpc-go"
\end{verbatim}

When compiling the code, the user's PC will contact the server at the
given URL over HTTPS (verifying the server certificate via TLS) and
download the software bundle. On the other end, when a library developer
wants to update their library, they do so by interacting with the
hosting server via HTTPS and whatever authentication the server has set
up---credentials, maybe two-factor authentication, etc.

This has some good features: the server name is explicit so there is no
ambiguity about packages and the decentralized nature of specifying
individual URLs avoids the necessity for a central server that attracts
attacks. However, this requires trusting the server hosting the library
to secure the update process and distribute software honestly.

\subsection{More Explicit Trust: Code
Signing}\label{more-explicit-trust-code-signing}

In each of these approaches, if an attacker can cause the user to
download a bad package \emph{without} compromising the package
developer. In particular, if the attacker can compromise Github, it can
cause Github to distribute malware to end users.

To prevent this attack, a library developer could sign their software
using their private key and include the signature with their software
package. To verify that a package is authentic, the application
developer's PC can check that the signature is valid.

Of course, with any signature-based plan the mechanism for public key
distribution is crucially important. In the software distribution case,
the only reasonable plan is likely a Trust-on-First-Use based one which
accepts the first public key it sees but verifies that future software
updates use that same key. This protects against an adversary taking
control of, for example, the application's Github repository after the
end user installs the softare once. However, key management is hard, so
this is not widely used in practice.

\section{Building Binaries}\label{building-binaries}

In order to run software on our computer, it is necessary to convert the
source code (which is, at least in principle, manually auditable) into a
binary that is much more difficult to audit. Since compiling software is
computation-heavy, most application developers typically compile their
software and distribute the binary to their users. If an attacker
compromises the build server (or is able to backdoor the compiler), then
the attacker can cause users to execute bad code, even if the attacker
does not compromise the application developer itself.

\marginnote{\begin{footnotesize}

The XCodeGhost attack is an example of how an attacker can insert a
backdoor in a build system and exploit it to distribute malware.

\end{footnotesize}}

\textbf{Reproducible Builds.} One promising approach to the problem of
ensuring that a binary is the faithful compilation of a piece of
software is called ``reproducible builds.''

If a build process is reproducible, the function that turns a set of
source-code files into a binary is a deterministic function: if two
different people compile the same set of source-code files, they will
get exactly the same binary---the two will be bit-for-bit identical.
This allows anyone to \emph{audit} a build: to check that a build server
did its job correctly. In addition, having multiple independent parties
build the same piece of software (and sign the result) can give an end
user some assurance that the build server behaved correctly.

Implementing reproducible builds is not trivial. Traditional compilers
introduced many sources of non-determinism---not necessarily for any
particular reason, just for convenience. Creating reproducible builds
requires eliminating all of these sources of non-determinism, even
across multiple versions of the compiler.

As of today, the Go programming language now supports reproducible
builds.

\subsection{Juggling multiple versions of a
library}\label{juggling-multiple-versions-of-a-library}

Once a binary exists, the next step of the process is to distribute that
software to user devices. Typically, there are many different versions
of a piece of software around. When a user wants to install a piece of
software, they typically need to specify which version of the software
they want.

For example, in Python a user can specify a version of a package that
they would like to install when they run \texttt{pip\ install}. If an
attacker compromises the PyPi server, it can serve up any code it wants
to a user asking for a particular version of a library.

In contrast, in the Go programming language, when a users imports a
package, the \texttt{go\ get} software will store a hash of the
downloaded code in a file called \texttt{go.sum}. If an attacker later
compromises the server serving the package (e.g., Github), the
\texttt{go\ get} command will refuse to install the package unless its
code matches the stored hash value.

\marginnote{\begin{footnotesize}

This is an example of ``trust on first use'' in the context of code
installation.

\end{footnotesize}}

\section{Installing \& Updating
Software}\label{installing-updating-software}

Once a software developer finishes writing an application, it builds and
distributes it. When a user installs an application---e.g., by
downloading it from a website or fetching it using a package
manager---how does the user know that it got the authentic version of
the software? As usual in systems design, there are many possible
strategies.

\textbf{Application Developer Signs Package (Android Apps).} One
possible option is to have application developers sign the software that
they produce. When application developers distribute their software,
they attach a their signature to it. This way, it does not matter how a
user obtains the software---a user can download an application bundle
from any server and know that it came from the developer who owns the
corresponding secret key.

When a user first installs a piece of software they need to somehow
obtain the software developer's public key. Public-key distribution, as
always, is messy: trust on first use is a common strategy.

Once the user has the software developer's public key, the user can
easily verify that future updates to the software came from the same
developer. (To do this, the user can just check signatures on the
updates using the app-developer's public key.)

An important caveat is that signatures do not guarantee freshness: once
signed, a package is always valid.

\textbf{Repository Signs Packages.} For systems with a central
repository, another plan is for the repository to sign packages. This
again allows the user to fetch the signed packages from untrustworthy
sources---from a content-distribution network, for example.

In addition to signing the packages, the repository typically signs a
timestamped manifest of the latest package versions. This allows a user
to check that they are not only getting the right software but also that
they are getting the most up-to-date software.

Many Linux package managers, such as \texttt{apt}, \texttt{pacman}, and
\texttt{rpm}, use signatures to integrity-protect packages.

\textbf{Third-Party Validator Signs Packages.} Yet another option that
does not require a single central repository is to have a trusted
validator sign packages. This involves sending the source code and
package to a third party, who will then perform some inspection of the
package and, if it deems a package to be worthy, provide some signature
over that package that verifies that the validator thinks the package is
trustworthy.

A number of software platforms use this strategy for protecting
binaries. On Android, there is no requirement to install apps from the
Google Play Store, but Google provides a service that inspects packages
and attaches these signatures if the package passes. Similarly, Windows
uses a validation plan for its device drivers.

\textbf{Binary Transparency.} One different plan to help involves an
audit log that keeps track of all published binaries.

This helps prevent in particular targeted attacks---for example, if some
adversary has a specific target in mind and compromised the distribution
of the Linux kernel, they would likely be immediately noticed if they
introduced a backdoor into Linux for the whole world. However, if they
were able to introduce a backdoor and distribute that backdoored version
only to their target, the adversary would be much more likely to evade
detection. If clients check their received binary against the publicly
available one before installing it, this personalized attacks can be
avoided---if the attacker wants to change the binary for someone, they
will need to change it for everyone.

\section{Booting the System}\label{booting-the-system}

\subsection{Secure Boot}\label{secure-boot}

In order to actually run an application, we rely on large amounts of
software running on our computer, from the applications themselves to
the operating system that supports them. If the operating system itself
is compromised, for example, the modified OS could undermine all of the
defenses we just discussed. ``Secure Boot'' is one strategy for getting
some partial protection against these attacks.

Devices using secure boot have a small amount of read-only memory (ROM)
that contain a small piece of code that runs on boot. This boot-ROM code
has a signature-verification key baked into it. There is no way to
change this key---it is a fixed part of the hardware Booting then
involves several steps:

\begin{enumerate}
\def\labelenumi{\arabic{enumi}.}
\tightlist
\item
  On boot, the CPU will begin running the hardcoded Boot ROM code, which
  has a signature-verification key \(\vk_\text{ROM}\) hardcoded into it.
\item
  the Boot ROM will load the code for another layer called the
  \emph{boot loader}. The boot ROM will then verify that the boot loader
  code carries a correct signature that verifies under
  \(\vk_\text{ROM}\). If the signature is valid, the boot ROM code will
  begin executing the boot loader. The bootloader has another
  signature-verification key (\(\vk_\text{bootloader})\)) baked into it.
\item
  The boot loader will load the code for the operating system and verify
  it using \(\vk_\text{bootloader}\). If the signature is valid, the
  bootloader will redirect control to the operating system.
\end{enumerate}

This way, the system can verify that only boot loaders approved by the
hardware manufacturer can run on the machine. These boot loaders then
can verify that only trusted operating systems are executed.

Many systems use secure boot: iPhone, Android, chromebooks, game
consoles, and UEFI secure boot on PCs.

In some cases (e.g., UEFI secure boot) secure boot is a mechanism to
protect against malware that tampers with the operating-system code.
While the malware may be able to compromise the running machine, after
rebooting the machine, the user has some assurance that it is running an
uncompromised operating system.

In other cases (e.g., game consoles) secure boot is a mechanism to
prevent the device owner from installing a non-standard operating system
on the device. Game-console vendors often sell the console hardware at a
loss, and they make their money back by selling game software. They have
a strong incentive then to prevent users from buying game consoles and
using them for non-game purposes.

\marginnote{\begin{footnotesize}

A number of researchers have used PlayStation 3 consoles for brute-force
password-cracking and cryptanalysis (e.g., factoring). Game consoles
often have a large number of CPU/GPU cores, which make them appealing
hardware for applications that benefit from massive parallelism.

\end{footnotesize}}

\subsection{Measured Boot}\label{measured-boot}

Secure boot assumes you know which key is needed to sign the software.
But what if the hardware cannot determine which software is good and
which is bad? The idea is to measure what software is being booted. This
involves hashing the bootloader, then hashing the OS kernel, and finally
hashing the OS image. This idea is called measured boot.

While measured boot cannot prevent bad software from loading, different
software generates different secrets. The system contains a secret that
is durably stored in hardware and remains consistent across reboots.
When the system boots up, it derives a secret key based on this hardware
secret, using a derivation method that incorporates the hashes of the
bootloader, and OS kernel, and perhaps other software.

The OS then uses this secret key to decrypt its data and authenticate to
remote servers. Booting a different OS ---- perhaps due to malware
corruption ---- will result in a different key being generated, and this
is detected when the data cannot be decrypted.

\section{Secure attention key}\label{secure-attention-key}

When to approach a terminal and type your bank password into it, how do
you know that you are typing the password into the banking app or into
some other app (e.g., the flashlight app) on the machine?

A traditional approach to address this problem is a \emph{secure
attention key}: there is a special button or combination of buttons that
trap into the operating-system kernel code---interrupting whatever
application that may be running.

Windows workstations, for example, required users to type the keyboard
combination \texttt{CTRL-ALT-DEL} to bring up a login prompt. If a user
entered this combination while an app was running, the operating system
would interrupt the application and open the legitimate login screen.

Pressing the ``Home'' button on many smartphones has the same effect.

\chapter{Hardware Security}\label{hardware-security}

So far while discussing platform security, we have considered only our
software. However, software must run on some hardware, and the security
of this hardware is similarly vital---if an attacker can undermine the
security of our hardware, it does not matter how strong our software
security is. Luckily, hardware is typically much more difficult to
attack than software, but powerful attacks are still possible.

As a running example for this chapter, consider the example of a
certificate authority that signs certificates given that some policy is
met.

\marginnote{\begin{footnotesize}

An equivalent example is a cryptocurrency wallet---transactions in the
cryptocurrency are authorized by a signature over a transaction message.

\end{footnotesize}}

This server will accept requests and respond with a signature over that
request if some policy is met. For example, an MIT CA might enforce a
policy like ``I will sign only certificates for \texttt{*.mit.edu}
domains''. For this CA server, we can achieve some nice security
properties:

\begin{itemize}
\tightlist
\item
  We can prove the security of the signature scheme used under concrete
  assumptions such as the hardness of discrete log.
\item
  We can verify that the cryptography implementation faithfully
  implements the signature algorithm on some ideal hardware model.
\end{itemize}

In order to actually be used, however, we must first buy a computer,
load the code onto it, and run it (likely on a machine running other
software). This process is outside of the nice security properties with
achieved about our CA program. In cryptography and in software, we are
able to develop clean characterizations of the power of the attacker. In
cryptography, we assume that our adversaries are bounded by
probabilistic polynomial time. In software, we can model our attacker as
choosing arbitrary inputs to our software and accurately capture an
attacker's power. Once we consider the hardware, however, it becomes
much more difficult to meaningfully define the attacker's power.

When we consider the real hardware that the system is running on, there
are many attacks that we must consider.

\section{Hardware Bug}\label{hardware-bug}

Much of what we have discussed so far has been cryptography schemes that
rely on trusted parties knowing some randomness that the adversary does
not know. We have assumed that we have some source of ``true''
randomness to use, for example, as a seed for a PRF. When actually
implementing cryptographic algorithms in a computer, we need to actually
materialize this ``true'' randomness. However, it is not clear where
this randomness should come from: computers are designed specifically to
behave as they are instructed by the programs they run. Common solutions
are to measure statistics about the environment that should be hard for
an adversary to predict. For example, devices may use combinations of:

\begin{itemize}
\tightlist
\item
  Keypress timings
\item
  Packet timings
\item
  Clock
\item
  Temperature sensor
\end{itemize}

All of these require ``accumulating'' randomness from the environment. A
common hardware bug on embedded devices is to generate keys when the
randomness source has not accumulated enough random measurements from
the environment. For example, if a network card generates a
cryptographic key right after boot, this key may be predictable if an
attacker is able to accurately guess at the values of the randomness
source.

To help with this, many devices include specialized hardware that uses
some special circuit to generate randomness by measuring randomness
inherent to the universe.. Of course, developers must then use this
randomness---a common error is to use insufficient randomness, such as
the time, instead of this hardware randomness.

\marginnote{\begin{footnotesize}

On Intel CPUs, this takes the form of the \texttt{RDRAND} instruction

\end{footnotesize}}

\section{Attacks without Physical
Access}\label{attacks-without-physical-access}

Perhaps the most concerning attacks are those that do not require
physical access to a machine.

\subsection{Cache Timing Attacks}\label{cache-timing-attacks}

One major goal of operating systems is to provide isolation between
processes. Even if an attacker is able to run some software on the same
machine as our signing process, we would like to guarantee that an
attacker can not read, for example, the signing key used by our signing
process. However, the attacker and victim code both run on the same CPU,
and the victim may leave traces of secrets in the state of the CPU.

For example, consider that our signing process runs and, depending on
some secret value, either loads the value at memory address \(A\) or
does not. If the victim loads this address,the CPU will copy the value
into the cache to speed up future accesses to the value. An attacker
process that runs next can try to access this same memory address and
measure how long it takes to read the value. If the victim read that
value, the access will be fast since it comes from the cache, but if
not, the data will have to come from the much slower main memory. From
the difference in this timing, the attacker can learn about the victim's
access pattern, which may reveal data about the victim's secret.

\marginnote{\begin{footnotesize}

Because of attacks like this and others, it is important to write secure
code such that it does not branch on secret values.

\end{footnotesize}}

\marginnote{\begin{footnotesize}

This may seem like an unlikely attack since both the victim and attacker
must have access to the same memory location, which process-level
isolation should prevent. However, operating systems perform something
called deduplication that can be cleverly taken advantage of to achieve
this: if the victim process uses OpenSSL, the attacker process can also
use OpenSSL. The operating system will see that both processes are
linking the same library, and may map a section of virtual memory for
each process to the same physical memory.

\end{footnotesize}}

\subsection{Rowhammer}\label{rowhammer}

Data in a computers memory is stored in what is effectively a grid of
capacitors. These capacitors do not store values indefinitely, and so
their values must be refreshed every so often (commonly every 64
milliseconds). Reading a chunk of memory drains the corresponding
capacitors a bit, and they must then be rewritten. Memory is read one
row f this grid at a time. Reading a row drains the corresponding
capacitors, requiring them to be rewritten. This rewriting involves
voltage fluctuations, and since modern memory is so dense, these voltage
fluctuations can cause neighboring rows to discharge more quickly that
the refresh interval is equipped to handle. Surprisingly, reading a
single row repeatedly can cause bits in an adjacent row to flip.

An attacker could take advantage of this by ``hammering'' a memory
location, causing a bit to flip in memory that belongs to another
process or to the operating system. In some cases, this was enough to
allow the attacker to learn a secret or bypass isolation, etc.

\section{Physical Attacks}\label{physical-attacks}

If an attacker has physical access to a device, an entirely new class of
attacks becomes possible. They can measure the device, introduce faults
to the device, and more.

\subsection{Probing Attacks}\label{probing-attacks}

An attacker with physical access to a device can measure many things
about the device's behavior that a remote attacker could not. For
example:

\begin{itemize}
\tightlist
\item
  Place probes on the pins of a chip
\item
  Measure power consumption of a chip and watch for patterns
\item
  Measure optical emissions of a chip with an electron microscope
\item
  Measure RF emissions from a chip
\item
  Monitor the blinking light on an internet router
\end{itemize}

The information that an attacker can learn from attacks like this may be
limited, but even very slow information leakage can be enough to leak a
key in a relatively short amount of time. Protecting against these types
of attacks is difficult since the attacker can do such a broad range of
things. However, if we make certain assumptions about the attacker's
power, we can achieve principled solutions.

\textbf{Defense against Probing Attacks.} One assumption that may be
reasonable to make is that an attacker can probe at most \(t\) wires of
a circuit. By using techniques like secure multiparty communication, it
is possible to build a circuit that implements something like a
signature scheme, but that does so without leaking anything about the
secret given this assumption.

\subsection{Fault Attacks}\label{fault-attacks}

An attacker can also introduce faults that the system was not designed
to handle. For example, they can point a heat gun or a laser at the
chip, hoping to cause some bit flips. If they are successful, these bit
flips may leak a secret key or corrupt a kernel data structure, allowing
the attacker to take over the system.

\marginnote{\begin{footnotesize}

Designers of satellite systems have to think about similar attacks, but
in their case the attacker is the sun! Cosmic rays carry enough energy
to flip bits of CPU registers or memory.

\end{footnotesize}}

\subsection{Supply Chain Attacks}\label{supply-chain-attacks}

When we buy an device, we assume that the hardware inside is not working
against us. However, there is a long chain of steps that happens before
the device gets to us---it is built in the factory, packaged, mailed
across the world to a retailer, stored in a warehouse, packaged and
mailed again, and so on. An attacker that has control over any of these
steps could intercept the device on its way to you and modify it
somehow. For example, they could:

\begin{itemize}
\tightlist
\item
  Modify the randomness source to something predictable
\item
  Preload keys that the attacker knows
\item
  Add extra input/output interfaces
\item
  Add or enable management interfaces
\end{itemize}

There are no great solutions to defend against this---inspecting the
chips is impossible since they are so small, building a device yourself
is much too hard, and so on. One solution that can improve the
protection is to build a system out of \(n\) identical devices and use a
strategy like secure multiparty computation to protect against cases
where at most \(n-1\) of these devices is compromised by a supply chain
attacker.

\chapter{Case Study: iOS Security}\label{case-study-ios-security}

There are several things we might worry about when it comes to the
security of a smartphone, spanning software, platform, and hardware
security:

\begin{itemize}
\tightlist
\item
  Malicious apps that steal contacts, eavesdrop on calls, or steal you
  credit card data
\item
  Someone stealing a phone and extracting secret data
\item
  A non-Apple operating system that is loaded on a phone
  (i.e.~Jailbreaking)
\item
  Malicious chips installed in the factory, en route to the store, or
  during repairs
\item
  People selling fake phones as real iPhones
\end{itemize}

The iOS platform is a very well-designed integration of the topics that
have been covered in the platform security section and which addresses
many of these considerations effectively.

\section{App Security}\label{app-security}

In any platform, there is a balance between the ``openness'' of the
platform---who is able to run software on it---and the access that
software has to sensitive data. Of course, the most secure platform is
one that doesn't do anything at all: it is very closed and provides no
access to sensitive data. More practical systems have a different
balance:

\begin{itemize}
\tightlist
\item
  A typical laptop is very open---anyone can write software for it---and
  every application is able to access the entire filesystem.
\item
  The web is still very open, as any website can include JavaScript code
  that will execute on your machine, but access to sensitive data is
  tightly locked down via language sandboxing.
\end{itemize}

iOS improves its security story by locking this down: applications are
given some (checked) access to sensitive data, but are very
closed---applications must go through a review process and only approved
developers can write software for iOS.

\marginnote{\begin{footnotesize}

In addition to security policies, the review process checks that
applications follow policies in place for business reasons, such as the
restriction that digital purchases must go through apples In-App
Purchase mechanism.

\end{footnotesize}}

iOS apps run in a sandbox that does not provide access to a shared
filesystem and that allows communication between apps only though
limited APIs provided by the operating system. If an application wants
to interact with data that is controlled by the operating system, such
as VPN configurations or health data, the application developer must ask
for access to special APIs and have that access approved as part of the
review process.

\subsection{Things Can Still Go Wrong}\label{things-can-still-go-wrong}

Even with these checks and the sandboxing, malware can slip through the
gaps. As we discussed earlier, XCode Ghost was a compromised version of
the XCode compiler distributed via mirrors behind China's firewall that
inserted malware into apps developed by honest developers. This malware
was not found in App Store reviews, allowing apps that people expected
to be honest to run arbitrary malicious code. Even with the sandbox, an
app has quite a bit of access to do potentially sensitive actions:

\begin{itemize}
\tightlist
\item
  Learn Country
\item
  Learn Language
\item
  Learn UUID, before recent privacy changes
\item
  Read and modify clipboard contents, including copied passwords or
  credit card numbers
\item
  Open a URL that points to a phishing webpage
\end{itemize}

Despite this, isolation buys a lot of security. A malicious app that
makes it through review cannot access your text messages, browser
history, etc.

\section{iOS Secure Boot}\label{ios-secure-boot}

For security and business reasons, Apple would like to ensure that an
Apple-signed operating system is running on the phone. This prevents a
malicious actor from distributing a backdoored operating system and
convincing people to install it on their phone and against malware that
tries to persistently modify the operating system. It also prevents
users from installing a customized operating system on their phone,
bypassing the Apple restrictions on apps and other policies.

To achieve this, iOS uses a secure boot system as described in the last
chapter. Each phone ships with a Boot ROM that cannot be changed that is
burned with some public key for a secret key that Apple knows. This boot
ROM is responsible for verifying that this secret key signed the code
for a low-level bootloader and running that bootloader. This bootloader
will then verify and check the operating system kernel. This allows the
bootloader and the kernel to be updated as necessary, but places a root
of trust in the boot ROM that cannot be updated. If the bootloader is
updated by anyone besides Apple, however, the boot ROM's signature check
will fail and the boot ROM will refuse to run the bootloader.

Many device owners that wanted to customize their devices sought to
modify the operating system to add new features. This process generally
is referred to as ``jailbreaking''.

\subsection{Checkra1n Jailbreak}\label{checkra1n-jailbreak}

One set of exploits that constituted a jailbreak was named
\texttt{checkra1n}. This took advantage of complex code in the
unchangeable boot ROM---the devices supported running code directly via
USB, bypassing the low-level bootloader and the OS kernel, which meant
that the boot ROM contained code to act as a USB peripheral. USB code is
quite complex, and as with most complex code, it had bugs.
\texttt{checkra1n} took advantage of these USB bugs to trick the phone
into executing arbitrary code.

Because the Boot ROM can never be updated, it was impossible for Apple
to fix these bugs: this jailbreak will work forever on the devices that
had the bug.

However, these bugs did not allow the jailbreak to bypass the signature
check entirely, so this exploit needed to be run on every boot---if you
tried to reboot your phone without running this, it would no longer be
jailbroken.

\section{iOS protection for data at
rest}\label{ios-protection-for-data-at-rest}

If a phone is stolen, it would be nice if the theif could not learn any
sensitive information from the device. As with any time that we want to
hide data, encryption is the answer here---a simple solution is to
encrypt all data on the phone with 128-bit AES. However, this does not
tell the whole story: in order to decrypt the data, the key for this
encryption must be stored somewhere. We can't store the key in normal
flash memory, since then anyone could use it to decrypt the data. We
also can't use the only secret that the user knows, their 6-digit PIN,
since it is much too short to be an AES key.

Even an approach like using a PRF based on the key to extend it into key
for AES will not be secure---6 digits is so short that it can easily be
brute-forced. Instead, recent iPhones use a special chip called a
``secure enclave'' that holds the key and provides access to it only if
the user enters the correct PIN. Doing so allows the secure enclave to
enforce strict guess limits that prevent brute-forcing the PIN.

The secure enclave is essentially another processor that runs its own
operating system. It uses a similar secure boot system to prevent
tampering with the secure enclave's operating system, but it also uses
\emph{measured boot} to derive the secret encryption key from the
contents of the OS being run---if an attacker modifies the secure
enclave's operating system, the encryption key will change and the
attacker will not be able to decrypt the phone data.

Importantly, the secure enclave has access to two things that the main
application processor does not. First, on the first boot: the secure
enclave generates a long-term secret unique ID and burns it into
internal fuses. The enclave also has access to a secure NVRAM module
that has a limited amount of secure storage with support for real
deletion. This secure storage contains the root encryption key itself, a
hashed version of the user's PIN salted with the UID, and a guess
counter that keeps track of how many incorrect guesses have been made.

\marginnote{\begin{footnotesize}

In typical storage, deleting data does not really delete the
data---instead, it marks it as deleted and indicates that the operating
system should overwrite it in the future. However, if someone inspects
the storage directly, they are likely to be able to recover the data
that was deleted. Apple's ``effacable storage'' that is used for the
secure enclave supports deleting data such that it cannot be recovered.

\end{footnotesize}}

When put together, these elements enable the following process:

\begin{enumerate}
\def\labelenumi{\arabic{enumi}.}
\item
  User enters a PIN
\item
  iOS passes the PIN to the secure enclave
\item
  The enclave enforces some delay after each guess
\item
  The enclave passes \(H(\text{PIN}, \text{UID})\) to the secure
  storage.

  \marginnote{\begin{footnotesize}

  The communication between the secure enclave and the secure storage is
  also encrypted with a key burned into the enclave and into the secure
  storage. This prevents an attacker with a probe on the wire from
  learning the data.

  \end{footnotesize}}
\item
  Secure storage uses included logic to check whether the hashed PIN
  matches the stored hash.

  \begin{itemize}
  \tightlist
  \item
    If correct, return the root AES key and zero the guess counter
  \item
    If incorrect, increment the guess counter. If the guess counter is
    too high, erase the key from the effacable storage.
  \end{itemize}
\item
  If the PIN was correct, the enclave passes the key returned from the
  storage to the AES engine.
\item
  The application processor sends data to the AES engine to be encrypted
  or decrypted.

  \marginnote{\begin{footnotesize}

  Note that the application processor never sees the AES key---it is
  seen only by the secure enclave and the (hardware) AES engine.

  \end{footnotesize}}
\end{enumerate}

\subsection{Biometric Unlock}\label{biometric-unlock}

For convenience, however, iPhones do not require entering a PIN on every
unlock. A PIN unlock is always required on the first unlock after boot,
but afterwards they allow unlocking with a Biometric such as Face/Touch
ID. The phone includes dedicated biometrics hardware that is responsible
for reporting the ``hash'' of the measured face to the secure enclave.
The secure enclave then checks this against the stored one, and unlocks
the device if there is a match.

A possible attack, then, might look something like the following: steal
a phone that has been unlocked at least once and relocked (almost always
the case). Then, before it runs out of battery and shuts down, replace
the Face ID chip with a malicious one that always reports the correct
hash. Without protection against replacing the biometric hardware, an
attack like this would allow an attacker to access all the data on your
phone. To address this, iPhone include yet another shared key between
the secure enclave and the biometric hardware: if the Face ID module is
replaced, the keys will not match and the phone will refuse to unlock.

\part{Software Security}

\chapter{Software Security}\label{software-security-1}

Software vulnerabilities are at the root of some of the most serious
security failures in real-world software. A huge number of security
issues come from software-implementation bugs.

\marginnote{\begin{footnotesize}

In terms of practical consequences on deployed computer systems,
software issues are perhaps second only to phishing attacks.

\end{footnotesize}}

A rule of thumb to keep in mind is that, for reasonably
carefully-written code, there will be around one bug for every 1000
lines of code. Many of these bugs may seem minor, but surprisingly,
almost any kind of bug can be used in a security exploit and lead to
compromise---even bugs in code that may not seem security-critical. One
principle to remember is then:

\marginnote{\begin{footnotesize}

A common strategy to exploit seemingly benign software bugs is to string
together a number of benign bugs into an exploit that does something
very bad from a security perspective.

\end{footnotesize}}

\begin{quote}
``Any bug can be a security bug.''
\end{quote}

Protecting the security of computer system thus requires eliminating
software bugs.

Many bugs can lead to security exploits, but some of the most common
types of exploited bugs include memory corruption bugs (buffer overrun,
use after free, etc.), encoding and decoding errors, cryptographic
implementation bugs, race conditions and resource-consumption bugs.

\section{Memory Corruption}\label{memory-corruption}

Memory-corruption bugs show up in languages, such as C and C++, that do
not guarantee any sort of memory- or type-safety properites. The two
most common types of memory-corruption bugs are: (1) buffer overflows
and (2) use-after-free bugs.

\subsection{Buffer overflow}\label{buffer-overflow}

As an example of a memory-corruption bug, consider the following C code:

\begin{Shaded}
\begin{Highlighting}[]
\DataTypeTok{void}\NormalTok{ f}\OperatorTok{()} \OperatorTok{\{}
    \DataTypeTok{char}\NormalTok{ buf}\OperatorTok{[}\DecValTok{128}\OperatorTok{];}
\NormalTok{    gets}\OperatorTok{(}\NormalTok{buf}\OperatorTok{);} \CommentTok{// write bytes from stdin starting}
           \CommentTok{// at \&buf[0], followed by a \textquotesingle{}\textbackslash{}0\textquotesingle{}}
\OperatorTok{\}}
\end{Highlighting}
\end{Shaded}

This code allocates an array of 128 bytes and then uses \texttt{gets} to
read a string from standard input into the buffer. The \texttt{gets}
routine will read the input from until it reaches the end of the string
(indicated in C with a zero byte), then it will write the corresponding
string into the buffer. Arrays in C have no length information attached
to them, so the \texttt{gets} code will happily accept an input strength
of any length---possibly much larger than 127 bytes.

So, what happens if the input is longer than 128 bytes? The
\texttt{gets} function simply writes until it finds a NULL character
(\texttt{\textquotesingle{}\textbackslash{}0\textquotesingle{}}) in the
input. So if the input string is too long, \texttt{gets} will simply
write past the end of the memory allocated for \texttt{buf}. If an
attacker gives an input that is longer than 128 bytes, the attacker will
be able to write bytes of its choice into the memory after \texttt{buf}.

In C, the \texttt{buf} array will be allocated on the call stack, which
will be laid out in memory like this:

\begin{verbatim}
... rest of stack ...
---------------
return address of f's caller
---------------
buf[127]
...
buf[2]
buf[1]
buf[0]
---------------
\end{verbatim}

\marginnote{\begin{footnotesize}

You might think that the fact that this problem is caused in part by
having the stack in C grow down, so that running off the end of the
buffer can cause the attacker to overwrite the return address. It turns
out that similar attacks are possible even on architectures in which the
stack grows up. The fundamental problem is that the attacker can
scribble over data in the stack.

\end{footnotesize}}

Once the \texttt{f} function ends execution, the program will jump to
the return address sitting after at the end of there buffer. So if an
attacker can write data past the end of the \texttt{buf} array, the
attacker can overwrite the return address of \texttt{f} on the stack and
the attack can cause the program to jump to and begin executing code at
an arbitrary location in memory.

One simple mitigation is to modify the compiler to refuse to compile
programs that use routines such as \texttt{gets}. We discuss other
mitigations below.

To avoid this kind of \emph{buffer overlow} bug, we need to somehow
ensure that \texttt{gets} only writes within the bounds of the buffer.

\textbf{Mitigation: Runtime checks.} Modern compilers can try to check
in real time whether a memory access goes past the end of the buffer and
will crash the program if so. These defenses are imperfect but prevent
the most naive type of bug.

\textbf{Mitigation: Bounds Checking.} Another mitigation is to ensure
that the input is not too large before reading it in. To do this, we can
insert a check before writing.

Consider the following slightly more complex code, which receives \(n\)
records that are each 16 bytes long and writes them into the buffer:

\begin{Shaded}
\begin{Highlighting}[]
\DataTypeTok{void}\NormalTok{ f}\OperatorTok{()} \OperatorTok{\{}
    \DataTypeTok{char}\NormalTok{ buf}\OperatorTok{[}\DecValTok{256}\OperatorTok{];}
    \DataTypeTok{uint32\_t}\NormalTok{ n }\OperatorTok{=}\NormalTok{ get\_input}\OperatorTok{();}
    \ControlFlowTok{for} \OperatorTok{(}\DataTypeTok{uint32\_t}\NormalTok{ i}\OperatorTok{=}\DecValTok{0}\OperatorTok{;}\NormalTok{ i }\OperatorTok{\textless{}}\NormalTok{ n}\OperatorTok{;}\NormalTok{ i}\OperatorTok{++)} \OperatorTok{\{}
        \CommentTok{// read record i into}
        \CommentTok{// buf[i*16] .. buf[i*16+15]}
    \OperatorTok{\}}
\OperatorTok{\}}
\end{Highlighting}
\end{Shaded}

This code will write beyond the end of the buffer if \texttt{n*16} is
greater than \texttt{256}. We then may consider adding a check like the
following:

\begin{Shaded}
\begin{Highlighting}[]
\PreprocessorTok{\#define sz }\DecValTok{256}

\DataTypeTok{void}\NormalTok{ f}\OperatorTok{()} \OperatorTok{\{}
    \DataTypeTok{char}\NormalTok{ buf}\OperatorTok{[}\NormalTok{sz}\OperatorTok{];}
    \DataTypeTok{uint32\_t}\NormalTok{ n }\OperatorTok{=}\NormalTok{ get\_input}\OperatorTok{();}
    \ControlFlowTok{if} \OperatorTok{(}\NormalTok{n }\OperatorTok{*} \DecValTok{16} \OperatorTok{\textgreater{}}\NormalTok{ sz}\OperatorTok{)} \OperatorTok{\{}
        \CommentTok{// input too long!}
        \ControlFlowTok{return}
    \OperatorTok{\}}
    \ControlFlowTok{for} \OperatorTok{(}\DataTypeTok{uint32\_t}\NormalTok{ i}\OperatorTok{=}\DecValTok{0}\OperatorTok{;}\NormalTok{ i }\OperatorTok{\textless{}}\NormalTok{ n}\OperatorTok{;}\NormalTok{ i}\OperatorTok{++)} \OperatorTok{\{}
        \CommentTok{// read record i into buf[i*16] .. buf[i*16+15]}
    \OperatorTok{\}}
\OperatorTok{\}}
\end{Highlighting}
\end{Shaded}

However, consider an adversary that inputs data such that
\texttt{n\ =\ 2\^{}30}. If we were computing on paper, our check would
work just fine: \(2^{30} \cdot 16 = 2^{34}\), which is certainly greater
than \texttt{sz}. However, \texttt{uint32\_t} is a 32-bit value and
\(2^{34}\) will not fit into the 32 bits allocated for that
integer---the computation will \emph{overflow}. It turns out that if you
try to compute \texttt{n*16} in C when \texttt{n} is \(2^{30}\), the
answer is zero! Thus, our check will pass but our code will still write
beyond the end of the buffer.

To prevent this type of overflow, the program can explicitly check for
overflow---it's tricky to do, but important when accepting user-provided
input.

\subsection{Use after free}\label{use-after-free}

Another common type of bug is a use-after-free bug, in which a
programmer frees a chunk of heap memory and then reads or writes it
after freeing it.

In C this happens when a programmer uses \texttt{malloc} to allocate
some memory, then calls \texttt{free} to free it, and then accesses it.
An example piece of code with this bug is here:

\begin{Shaded}
\begin{Highlighting}[]
\DataTypeTok{void}\NormalTok{ f}\OperatorTok{()} \OperatorTok{\{}
  \DataTypeTok{char} \OperatorTok{*}\NormalTok{req }\OperatorTok{=}\NormalTok{ malloc}\OperatorTok{(}\DecValTok{1024}\OperatorTok{);}
  \DataTypeTok{int}\NormalTok{ err }\OperatorTok{=}\NormalTok{ read}\OperatorTok{(}\DecValTok{0}\OperatorTok{,}\NormalTok{ req}\OperatorTok{,} \DecValTok{1024}\OperatorTok{);}
  \ControlFlowTok{if} \OperatorTok{(}\NormalTok{err}\OperatorTok{)}\NormalTok{ free}\OperatorTok{(}\NormalTok{req}\OperatorTok{);}
                            \CommentTok{// ***}
  \ControlFlowTok{if} \OperatorTok{(}\NormalTok{err}\OperatorTok{)}\NormalTok{ printf}\OperatorTok{(}\StringTok{"Error }\SpecialCharTok{\%d}\StringTok{: }\SpecialCharTok{\%s\textbackslash{}n}\StringTok{"}\OperatorTok{,}\NormalTok{ err}\OperatorTok{,}\NormalTok{ req}\OperatorTok{);}
\OperatorTok{\}}
\end{Highlighting}
\end{Shaded}

If some other thread in the program calls \texttt{malloc} when the code
is a point \texttt{***}, the other thread of execution may use the array
\texttt{req} for something else. Then the \texttt{printf} line could
print out some other contents of memory---possibly exposing
cryptographic keys or other sensitive data.

These bugs are very difficult to track down since it is difficult for a
compiler to figure out which memory a piece of code should or should not
have access to. One way to defend against these bugs is to use a
programming language, such as Rust, that explicitly associates a
``lifetime'' with each piece of memory and can prevent code from
accessing free'd memory.

\section{Encoding Bugs}\label{encoding-bugs}

Another common source of security bugs comes from encoding and decoding
data from and to language data structures.

\subsection{SQL Injection}\label{sql-injection}

Most web sites that we interact with consist of some application
code---for example, a Flask app---that communicates with a database via
SQL queries. For example, a phone-number-to-name lookup site would
likely use SQL queries that look like

\begin{quote}
\texttt{\textquotesingle{}SELECT\ name\ FROM\ users\ WHERE\ phone\ =\ "6172536005"\textquotesingle{}}.
\end{quote}

When accepting a phone number, stored in variable \texttt{phone}, from a
user, the same query might look like:

\begin{quote}
\texttt{/*\ WRONG!!!\ */}\strut \\
\texttt{\textquotesingle{}SELECT\ name\ FROM\ users\ WHERE\ phone\ =\ "\textquotesingle{}\ +\ phone\ +\ \textquotesingle{}"\textquotesingle{}}.
\end{quote}

The problem with using string interpolation is that an adversarial user
can supply a phone number like

\begin{quote}
\texttt{123";\ DROP\ TABLE\ users;\ "}
\end{quote}

The SQL engine will then receive the query:

\begin{quote}
\texttt{\textquotesingle{}SELECT\ name\ FROM\ users\ WHERE\ phone\ =\ "123";\ DROP\ TABLE\ users;\textquotesingle{}},
\end{quote}

which will have the effect of deleting the \texttt{users} table.

The principled way to solve this problem is to have a strategy for
\emph{unambiguously encoding} data. In SQL, if you have a quote
character \texttt{"} in a data string, the programmer writes it as
\texttt{\textbackslash{}"}. This is called ``escaping'' a string. A SQL
library can automatically escape characters such as quotes, but escaping
is not as easy as replacing each quote with its escaped equivalent---you
need to worry about escaping the escape character
``\texttt{\textbackslash{}}'' and all sorts of other subtleties.

Modern libraries for interacting with databases perform escaping
automatically to avoid these ``SQL injection'' attacks.

\subsection{Cross-Site Scripting}\label{cross-site-scripting}

Another common behavior is to take input from the user via a form, save
it to the database, and later render that input to another user as HTML.
For example, a social media site will have each user enter their name
when the sign up, and may use some code like the following to render
another user's list of friends:

\begin{Shaded}
\begin{Highlighting}[]
\KeywordTok{def}\NormalTok{ render\_friends(friends: List[}\BuiltInTok{str}\NormalTok{]):}
    \BuiltInTok{print}\NormalTok{(}\StringTok{"\textless{}h3\textgreater{}Friends\textless{}/h3\textgreater{}"}\NormalTok{)}
    \BuiltInTok{print}\NormalTok{(}\StringTok{"\textless{}ul\textgreater{}"}\NormalTok{)}
    \ControlFlowTok{for}\NormalTok{ name }\KeywordTok{in}\NormalTok{ friends:}
        \BuiltInTok{print}\NormalTok{(}\StringTok{"\textless{}li\textgreater{}"} \OperatorTok{+}\NormalTok{ name }\OperatorTok{+} \StringTok{"\textless{}/li\textgreater{}"}\NormalTok{)}
    \BuiltInTok{print}\NormalTok{(}\StringTok{"\textless{}/ul\textgreater{}"}\NormalTok{)}
\end{Highlighting}
\end{Shaded}

If a friend's name is something expected, like ``Alice'' or ``Bob'',
this works fine. However, what if a friends sets their name to something
like
\texttt{\textless{}script\textgreater{}send\_to\_adversary(document.cookie)\textless{}/script\textgreater{}}?
Now, the rendered HTML will look something like the following:

\begin{Shaded}
\begin{Highlighting}[]
\DataTypeTok{\textless{}}\KeywordTok{h3}\DataTypeTok{\textgreater{}}\NormalTok{Friends}\DataTypeTok{\textless{}/}\KeywordTok{h3}\DataTypeTok{\textgreater{}}
\DataTypeTok{\textless{}}\KeywordTok{ul}\DataTypeTok{\textgreater{}}
    \DataTypeTok{\textless{}}\KeywordTok{li}\DataTypeTok{\textgreater{}}\NormalTok{Alice}\DataTypeTok{\textless{}/}\KeywordTok{li}\DataTypeTok{\textgreater{}}
    \DataTypeTok{\textless{}}\KeywordTok{li}\DataTypeTok{\textgreater{}}\NormalTok{Bob}\DataTypeTok{\textless{}/}\KeywordTok{li}\DataTypeTok{\textgreater{}}
    \DataTypeTok{\textless{}}\KeywordTok{li}\DataTypeTok{\textgreater{}\textless{}}\KeywordTok{script}\DataTypeTok{\textgreater{}}
        \FunctionTok{send\_to\_adversary}\NormalTok{(}\BuiltInTok{document}\OperatorTok{.}\AttributeTok{cookie}\NormalTok{)}
    \DataTypeTok{\textless{}/}\KeywordTok{script}\DataTypeTok{\textgreater{}\textless{}/}\KeywordTok{li}\DataTypeTok{\textgreater{}}
\DataTypeTok{\textless{}/}\KeywordTok{ul}\DataTypeTok{\textgreater{}}
\end{Highlighting}
\end{Shaded}

This \texttt{script} tag runs the contained Javascript code in the
viewer's browser. Anyone who views a friends list with this adversary in
it, then, will have their authentication cookie sent to the adversary,
potentially allowing the adversary to log in as that user.

The core of this issue is similar to the SQL injection attack: an
attacker is able to insert code that the victim's browser will run. To
prevent this, the solution is again escaping: we typically replace the
angle brackets (\texttt{\textless{}} and \texttt{\textgreater{}}) used
to denote HTML tags with the sequences \texttt{\&lt;} and
\texttt{\&gt;}. Now that \texttt{\&} has a special meaning, we also must
escape it as \texttt{\&amp;}. Modern web frameworks have ``templating''
systems that automate this escaping process.

\subsection{Decoding: Android Apps}\label{decoding-android-apps}

The application-installation process on Android also suffered from
decoding errors. Apps on Android (\texttt{.apk} files) are just renamed
ZIP files.

Apps shipped with a signature. When installing an app, Android would
first check that the contents of the ZIP file match the signature. If
the signature checked out, Android would then install the app. However,
the signature checking code and the installation code used different ZIP
decoders. An attacker was able to take advantage of a historical quirk
of the ZIP format that meant that ZIP holds two lists of files: the
signature checker used one list, while the installer used the other
list. By pointing to real files from one list but to malicious files in
the other list, an attacker was able to bypass this signature check and
cause a user to install malicious code.

\section{Concurrency Bugs}\label{concurrency-bugs}

When systems have code running in parallel, things become much more
difficult to reason about, and as a result, many bugs can occur.
Consider the following code running on a bank server:

\begin{Shaded}
\begin{Highlighting}[]
\KeywordTok{def}\NormalTok{ xfer(src, dst, amt):}
\NormalTok{    s }\OperatorTok{=}\NormalTok{ bal[src]}
\NormalTok{    d }\OperatorTok{=}\NormalTok{ bal[dst]}
    \ControlFlowTok{if}\NormalTok{ s }\OperatorTok{\textless{}}\NormalTok{ amt:}
        \ControlFlowTok{raise}\NormalTok{ InsufficientBalanceError}
\NormalTok{    balances[dst] }\OperatorTok{=}\NormalTok{ d }\OperatorTok{+}\NormalTok{ amt}
\NormalTok{    balances[src] }\OperatorTok{=}\NormalTok{ s }\OperatorTok{{-}}\NormalTok{ amt}
\end{Highlighting}
\end{Shaded}

By executing two of these \texttt{xfer} requests in parallel, an
attacker can cause unexpected behavior. For example, if an attacker
tries to send money from a single source to two destinations at once,
the check \texttt{s\ \textless{}\ amt} may pass in both executions, and
both destinations will then be updated to have money. However, the
source will only be deducted once, since \texttt{s} is stored before any
money is transferred.

The fix is to make sure that there is some kind of locking or
concurrency-control strategy in place. The important high-level bit is
to be thinking carefully about how concurrent execution can affect your
software, in cases when multiple threads of execution can access the
same data at the same time.

\subsection{File system races / Time-of-check-time-of-use (TOCTOU)
bugs}\label{file-system-races-time-of-check-time-of-use-toctou-bugs}

When deadline with files, symbolic links can cause all sorts of chaos. A
piece of code might want to check that it is deadline with a regular
file---rather than a symbolic link---before opening it. A bad way to
implement would be like this:

\begin{Shaded}
\begin{Highlighting}[]
\CommentTok{// WRONG}
\ControlFlowTok{if}\OperatorTok{(}\NormalTok{lstat}\OperatorTok{(}\NormalTok{path}\OperatorTok{,} \OperatorTok{\&}\NormalTok{st}\OperatorTok{)} \OperatorTok{\textless{}} \DecValTok{0}\OperatorTok{)}\NormalTok{ error}\OperatorTok{();}
\ControlFlowTok{if}\OperatorTok{(!}\NormalTok{S\_ISREG}\OperatorTok{(}\NormalTok{st}\OperatorTok{.}\NormalTok{st\_node}\OperatorTok{))}\NormalTok{ error}\OperatorTok{();}
\CommentTok{// ***}
\DataTypeTok{int}\NormalTok{ f }\OperatorTok{=}\NormalTok{ open}\OperatorTok{(}\NormalTok{path}\OperatorTok{,}\NormalTok{ O\_RDWR}\OperatorTok{);}
\OperatorTok{...}
\end{Highlighting}
\end{Shaded}

The problem is that if an attacker can replace the file when the code is
at line \texttt{***}, the attacker could swap out the regular file with
a symbolic link that points to somewhere else.

\marginnote{\begin{footnotesize}

The \texttt{openat()} API call gives a way to open a file while ensuring
that the file is of a particular type.

\end{footnotesize}}

The way to defend against this bug is to change the interface. Newer
versions of the POSIX file-system APIs enable checking for this type of
property in a way that defats race conditions.

\section{Resource Usage}\label{resource-usage}

Other bugs allow attackers to consume many resources on a system,
denying service to honest users.

For example, hash tables typically use a hash function designed to be
very fast when deciding which bin to place an input in. These hash
functions are typically not collision-resistant in the cryptographic
sense. Hash tables work well when inputs get distributed evenly across
all of the bins. However, if multiple inputs get mapped to the same hash
value, the performance of a hash table deviates from the constant-time
idea we have of a hash table---hash tables typically revert to using a
linked list of all the values for a given hash value.

If an attacker is able to predict which bin their input will end up in,
they can maliciously craft inputs that create a huge linked list,
resulting in very poor performance for the hash table.

The way to defend against this type of attack is to use a \emph{keyed}
hash function---essentially a pseudorandom function---where the
implementation does not reveal the secret key to the attacker (over the
network, etc.).

\section{Dealing with Software Bugs}\label{dealing-with-software-bugs}

As we have seen, there are many types of bugs so there are many ways to
defend against them. Here are some general rules:

\textbf{Clear Specification.} One way to avoid design-level bugs is to
have a clear and complete specification about what your program is
supposed to do. In particular, for things such as encoders and decoders,
a precise specification can make it easier to check that you have
handled all of the important corner cases.

\textbf{Design.} A simpler design is easier to understand, and thus bugs
are easier to find.

\textbf{Limit Bug Impact.} Some techniques that we will discuss, such as
privilege separation, allow us to protect security even when software
bugs arise.

\textbf{Find and Prevent Bugs at Development Time.} Techniques such as
fuzzing can find bugs before they hit production systems.

\textbf{Catch and Mitigate Bugs at Runtime.}

\textbf{Deploy Bug Fixes Quickly.} A big advance in browser security
came from mechanisms for pushing out software updates quickly.

\chapter{Privilege Separation}\label{privilege-separation}

The last chapter left off on a somewhat unsatisfying note: our software
is bound to have bugs and those bugs are likely to be exploitable.
Today, we will try to limit the impact that these bugs can have by
\emph{planning for compromise} and aiming to limit the damage that each
component can cause if an attacker compromises it.

The philosophy that we will employ is called the \emph{Principle of
Least Privilege}. The idea is to give each component in the system the
least privileges needed to do its job. By limiting the access that each
component of our software has to do sensitive actions, we can prevent
the compromise of our whole system when a single piece is attacked.

\marginnote{\begin{footnotesize}

In real systems, it is often costly to implement least privilege in the
strictest sense possible. You should think of the principle of least
privileges as a difficult-to-achieve design goal, rather than a rule
that every component of every system must satisfy.

\end{footnotesize}}

To summarize our strategy:

\begin{quote}
\textbf{Principle of Least Privilege:} Each component should have the
smallest set of privileges necessary to do its job.
\end{quote}

As a concrete example, a web server might have some networking code, it
might talk to a database server, and it might use some cryptographic
keys. One way to architect the system with least privilege would be to
split the server into multiple distinct processes---one that receives
network requests, one that interfaces with the database, and one that
uses the cryptographic keys. These components then would speak to each
other using narrow, restrictive APIs.

In some sense, a system that perfectly implements the principle of least
privilege gives one with the ``best security we could ask for,'' in the
face of component compromise. However, most real systems do not
completely implement the principle of least privilege for a number of
reasons:

\marginnote{\begin{footnotesize}

An orthogonal, but important, strategy is to reduce the number of
privileges that a component needs to do its job.

\end{footnotesize}}

\textbf{Challenge: Splitting Boundaries.} A large piece of software has
many components that often have many points of interaction. To
effectively implement least privilege, we need to find boundaries at
which we can separate our software. A few common ways to partition a
system or application are:

\marginnote{\begin{footnotesize}

If we don't split up a large piece of software into components, the
software just consists of one über-component. An attacker who breaks
into this one component can hijack the entire system.

\end{footnotesize}}

\begin{itemize}
\item
  \textbf{Isolating by user:} In an operating system, different users of
  the systems have different privileges. This way, if an attacker
  compromises one user's account it does not compromise the entire
  system.
\item
  \textbf{Splitting by feature:} In large business applications,
  different features (e.g., search, user management, mail) run as
  isolated components. This way, a bug in one component does not
  necessarily affect the entire system.
\item
  \textbf{Splitting off buggy code:} The Firefox web browser uses
  special sandboxes to compartmentalizes bug-prone video codec code
  (Section~\ref{sec-split-codec}). This way, if a malicious website is
  able to exploit the codec code, it is not so easy for the attacker to
  compromise the system.
\item
  \textbf{Separating exposed versus internal code:} Google's front-end
  HTTP servers are isolated from the servers that run core application
  logic. This way, an attacker on the network is one step removed from
  the servers with access to Google's internal resources (databases,
  etc.).
\item
  \textbf{Isolating sensitive data or keys:} Certificate authorities
  (CAs) use hardware security modules to isolate the cryptographic keys
  from all other code in the system. Laptops and servers sometimes have
  separate cryptographic co-processors for this purpose as well.
\end{itemize}

\textbf{Challenge: Interface/API.} If we chop our monolithic app into
different domains, we also need to clearly define an API that the
different domains can use to interact. This API must allow us to
preserve the functionality of the system while making it difficult for
an attacker that compromises one component from compromising an adjacent
one.

Some strategies for interfacing between different components of a system
might be

\begin{itemize}
\tightlist
\item
  remote procedure calls (RPCs) over a network,
\item
  message queues,
\item
  shared memory,
\item
  shared database, or
\item
  shared files or directories.
\end{itemize}

In designing interfaces between isolated components, we tend to worry
about:

\begin{itemize}
\tightlist
\item
  \textbf{Functionality:} Does the isolation plan allow the application
  to work as it should?
\item
  \textbf{Security:} Does the isolation strategy prevent the compromise
  of one component from affecting others?
\item
  \textbf{Performance:} How much performance overhead does the isolation
  mechanism add?
\item
  \textbf{Complexity:} How much code does the isolation mechanism
  require? If the compartmentalization strategy requires a large amount
  of code, this code might introduce new vulnerabilities.
\end{itemize}

\section{Example: Logging}\label{example-logging}

Most systems use logging to keep track of the actions that an
application has performed. That way, if an attacker compromises the
application, system administrators can look at the log to see what
happened and how to mitigate it. For a logging system to be useful, it
must be difficult for an attacker who compromises the application to
erase the log.

It is very natural to separate out the log into a separate component:
the app's functionality is almost entirely separate from the log, and
the app's interface to the log can be very simple. The app should be
allowed to append to the log and read from the log---the app should
\emph{not} be allowed to remove entries from the log. This way,
compromise of the application or log server will not be enough to erase
the log.

\begin{figure}

\centering{

\begin{verbatim}
Application  ----- Log entries -----> Log server
  server                                  |
     |                                    |   Append/Read API
     |                                    |   (No delete)
     |                                    v
     -------- No access -----> X         Log
\end{verbatim}

}

\caption{\label{fig-logging}One way to architect a logging system. An
attacker that compromises the application and/or log server may not be
able to erase the logs.}

\end{figure}%

\section{Example: Cryptography Keys}\label{example-cryptography-keys}

Many applications use cryptographic keys for authentication: certificate
authorities need to generate signed certificates for their clients,
cryptocurrency wallets need to sign transactions, and WebAuthn requires
an authenticator to sign a challenge from the relying party. In all of
these applications, we worry a lot about an attacker stealing the
application's secret signing key. Towards the goal of protecting
cryptographic secrets, we often isolate the code that uses cryptographic
keys into a separate software (or even hardware) component. In
particular, we might create a cryptography component whose only API is
\texttt{sign(msg)} and \texttt{get\_public\_key()}. In this way, even if
an attacker compromises the application, it cannot easily extract the
secret key.

This design does not completely protect us against the compromise of the
application. In particular, an attacker who compromises our app can
still call \texttt{sign()} as much as they like to sign any message they
like. At the same time, having the key isolated to this crypto component
allows us to add checks inside the crypto component to, for example,
reject messages of the wrong type. If our service is a certificate
authority, we could set up our crypto component to verify that each
message is an actual signature before signing it.

\marginnote{\begin{footnotesize}

Although in many cases the things we are worried about an attacker
asking for signatures of will pass this type check: for example, for a
CA, this would still allow an attacker to generate a certificate for
their malicious website.

\end{footnotesize}}

\begin{figure}

\centering{

\begin{verbatim}
Application  <----- Requests -----> Key manager
     |                                    ^
     |                                    |  Sign API
     |                                    |  (No extraction
     |                                    |     of key)
     |                                    v
     -------- No access -----> X     Signing key
\end{verbatim}

}

\caption{\label{fig-key-manager}A common architecture for isolating
cryptographic keys from an application.}

\end{figure}%

This division also allows our crypto code to do some meaningful logging:
imagine that the crypto module saves every signature it creates to a
log. Compromise of the main app would allow an attacker to generate
arbitrary signatures, but administrators could then see every signature
that was generated during the compromise. This would be very helpful in
recovery.

\section{Example: Media Codecs in Web Browsers}\label{sec-split-codec}

Media codecs (e.g., for JPEG decoding) are notoriously complex and
bug-prone: codec libraries have been at the source of many web-browser
exploits. If we were able to isolate these codecs, we could make it less
likely that a codec bug allows an attacker to access data other than the
media file being decoded.

\begin{figure}

\centering{

\begin{verbatim}
Browser tab  ----- Encoded data-----> JPG/Media codec
             <------- Bitmap --------
\end{verbatim}

}

\caption{\label{fig-codec}The Firefox browser isolates media codecs in a
separate sandbox.}

\end{figure}%

Isolating codecs may not always be as trivial as Figure~\ref{fig-codec}
makes it seem. In many cases, codecs require a sophisticated interface
to the browser tab: the codec library may progressively decode videos as
data arrives, for example. Even image decoders progressively decode
images, improving resolution as more data arrives. Supporting this
functionality without adding vulnerabilities or excessive latency
requires careful API design.

Firefox isolates these risky codecs using the language-level isolation
that the WebAssembly language provides.

\section{Example: Server for Network Time Protocol
(NTP)}\label{example-server-for-network-time-protocol-ntp}

Operating systems use the network-time protocol to fetch the current
time from time servers on the Internet, and to update the current time
to match. Setting the time requires root privilege on Unix-like systems,
but NTP also requires network accesses; the networking code can be
complicated and bug-prone. To prevent an attacker who finds a bug in
these network protocols from gaining root privilege on our machine, many
systems separate the two into a process that handles talking to the NTP
server on the network and a privileged process that accepts the time
from this other process and sets it. This way, an attacker who find a
bug in the network code can only set the time---they cannot perform
arbitrary actions as root. In addition, the time service may impose some
policy on time changes (e.g,. that time can never go backwards).

\begin{figure}

\centering{

\begin{verbatim}
Launcher --------> Network <---> Web
   |               Service
   |                  |
   |                  | "adjust clock"
   |                  v
   |-------------> Time Service [Policy enforcement]
\end{verbatim}

}

\caption{\label{fig-ntp}Modern operating systems implement a
network-time protocol (NTP) client as separate processes. This way, an
attacker who compromises the client over the network cannot arbitrarily
corrupt the system time.}

\end{figure}%

\section{Example: OpenSSH Server}\label{example-openssh-server}

A secure shell (SSH) server has access to many sensitive resources:
network port 22, a host secret key, the system's password file, and all
users' data. When it starts, the SSH server runs a ``monitor'' process
that listens for connections on port 22. When the monitor receives a
connection, the first thing it does is to spawn a new per-connection
worker process that communicates over the network. The worker client has
no access to the host key or password database---the worker process can
only ask the monitor for help in authenticating. The monitor-worker API
supports a few operations:

\begin{itemize}
\tightlist
\item
  a \emph{signing} operation, that instructs the monitor to sign a
  protocol transcript,
\item
  a \emph{password-authentication} operation, that instructs the monitor
  to check a password (this can be rate limited to prevent password
  guessing), and
\item
  a \emph{start-session} operation, that instructs the monitor to create
  a new process with a shell for the user.
\end{itemize}

\begin{figure}

\centering{

\begin{verbatim}
[ Host key  ]
[ Passwords ]
[ Port 22   ]
    |
    v
 Monitor --- Spawns per client --> Worker <--- Client1
    |                              Worker <--- Client2
    |                               ...
    --------> Session <----------> Worker <--- ClientN
              process
\end{verbatim}

}

\caption{\label{fig-openssh}The OpenSSH server is split into multiple
process to mitigate the compromise of the network-facing code.}

\end{figure}%

While this architecture adds quite a bit of complexity to the OpenSSH
server, it has paid off in terms of mitigating the impact of
vulnerabilities in client-facing code.

\section{Example: Web applications}\label{example-web-applications}

Companies will often implement Web applications (e.g., a photo-sharing
website) as a number of separate services, running on separate physical
machines. Client connections come into a front-end server that
terminates the TLS connection and proxies client data to a front-end
application server. The front-end server then routes requests to one or
more application services, each of which may have access to different
databases. For example, the login service may have access to the
password database, while the profile service may not.

\begin{figure}

\centering{

\begin{verbatim}
           [ TLS Key ]
               |
               v
Client ----> HTTPS  ---> Front-end
             server       service
                            ^
                            |---> Login   <--> Password DB
                            |     service
                            |
                            |---> Profile <--> User DB
                            |     service
                            |
                            |---> Photo   <--> Photo DB
                                  service
\end{verbatim}

}

\caption{\label{fig-web-services}Large Web services tend to isolate
different application features into different services, often on
different physical machines. The system gates access to these services
via minimal front-end client-facing servers.}

\end{figure}%

\section{Example: Web client}\label{example-web-client}

When you open a PDF attachment in Gmail , you might worry that the PDF
could exploit some bug in your browser that could steal your sensitive
Google data. To make this kind of attack more difficult, Gmail serves
attachments and other suspect files from a separate domain (``origin')):
\texttt{googleusercontent.com}. Code loaded from
\texttt{googleusercontent.com} cannot access cookies or data for
\texttt{google.com}. In this way, even if an attacker can somehow run
JavaScript in your browser, it cannot easily steal your Google cookies.

\section{Example: Web browser}\label{example-web-browser}

Web browsers today are extraordinarily complicated pieces of software.
The sensitive data that a browser is trying to protect are things, such
as user cookies, cached data, browser history, and other user data. The
browser may spawn new processes to handle rendering for each site from
each distinct domain/origin. In this way, if an attacker from one origin
can exploit a bug in the JavaScript engine, the attacker may still not
be able to compromise sensitive user data from other domains/origins.
GPU code, which is extremely complicated and bug-prone, may run in yet
another process. Today, compromising a browser entirely often requires
finding and exploiting a collection of bugs in multiple components.

\begin{figure}

\centering{

\begin{verbatim}
 codec                codec
   ^                    ^
   |                    |
   v                    v
mit.edu             nytimes.com
   |                    |
   |                    |
   ---- Browser core <------> UI <---> GPU
            |
            v
     [   Cookies   ]
     [ Cached data ]
\end{verbatim}

}

\caption{\label{fig-browser}Web browsers may isolate the execution of
each origin's code in a separate process. They further isolate
complicated and bug-prone codecs and GPU code in separate processes.}

\end{figure}%

\section{Example: Payment Systems}\label{example-payment-systems}

Processing credit-card transactions in web applications is risky: if a
vendor suffers a compromise, the credit-card network may fine them or
kick them off the network. To avoid ever having to handle credit-card
data, most websites use an external payment-processing service that
handles credit-card information. When the user makes a purchase, the
vendor redirects the user to the payment-processing service, who
collects the user's credit-card data. After payment, the
payment-processing service redirects the user back to the vendor's
website.

\begin{figure}

\centering{

\begin{verbatim}
   --------> Web app -----> Order DB
   |            ^
   |            |
Client    Payment data
   |            |
   |            |
   --------> Payment -----> Visa/MC
             service
                |
                v
        [ Credit card #s ]
\end{verbatim}

}

\caption{\label{fig-payment}Online vendors often use a separate payment
processor that handles the user's credit-card data.}

\end{figure}%

\chapter{Bug Finding}\label{bug-finding}

We are going to continue our treatment of security and bugs with a
discussion of how to find bugs. As we already have seen, bugs are a big
deal in terms of security problems. We have discussed how to architect a
system using privilege separation so that bugs do not matter so much.
But even with a very good privilege-separated design, we still want to
make sure that our system is as bug-free as possible.

The topic of this chapter is then: \emph{How do we find bugs?}

\textbf{Step 1: Define what is a bug.} Before we even begin talking
about how to find bugs, we need to answer the question: \emph{What is a
bug?} There are a number of application-independent ways to determine
when we have hit a bug:

\begin{itemize}
\tightlist
\item
  the program crashes,
\item
  the program makes an out-of-bounds memory access, or
\item
  the program jumps to an unknown or undefined point in the program.
\end{itemize}

More difficult types of bug to detect are ones that we can only detect
with knowledge of what the application is supposed to do:

\begin{itemize}
\tightlist
\item
  the program's output is incorrect, or
\item
  the program allows an attacker to access data it should not be able to
  access.
\end{itemize}

\textbf{Step 2: Find bugs.} To find a bug, we just need to identify a
possible execution of the program that leads to one of the buggy
outcomes we defined in Step 1. The reason this is difficult is:

\begin{itemize}
\tightlist
\item
  there are typically exponentially many possible inputs to the program
  and it may be difficult to find one the triggers the bug,
\item
  concurrent systems exhibit non-deterministic behavior, and
\item
  inputs from the environment (time, network state, etc.) can change the
  behavior of the program.
\end{itemize}

\section{Bug finding: A concrete
example}\label{bug-finding-a-concrete-example}

To frame our discussion of bugs, we will consider one specific example
of a C program that parses a binary packet that has the format depicted
in Figure~\ref{fig-packet}.

\begin{figure}

\centering{

\begin{verbatim}
------------------------------------------
| header | len | id | id | id | ... | id |
------------------------------------------
\end{verbatim}

}

\caption{\label{fig-packet}An example packet format.}

\end{figure}%

The packet starts with a header byte and a length byte and then has a
list of IDs.

\begin{figure}

\centering{

\begin{Shaded}
\begin{Highlighting}[numbers=left,,]
\CommentTok{// Parse header}
\DataTypeTok{char}\NormalTok{ in}\OperatorTok{[}\DecValTok{64}\OperatorTok{];}
\DataTypeTok{int}\NormalTok{ hdr }\OperatorTok{=}\NormalTok{ in}\OperatorTok{[}\DecValTok{0}\OperatorTok{];}
\ControlFlowTok{if} \OperatorTok{(}\NormalTok{hdr }\OperatorTok{!=}\NormalTok{ s}\OperatorTok{)}\NormalTok{ returnl}\OperatorTok{;}
\DataTypeTok{int}\NormalTok{ n }\OperatorTok{=}\NormalTok{ in}\OperatorTok{[}\DecValTok{1}\OperatorTok{];}
\ControlFlowTok{if} \OperatorTok{(}\NormalTok{n }\OperatorTok{\textgreater{}} \DecValTok{64}\OperatorTok{)} \ControlFlowTok{return}\OperatorTok{;}

\CommentTok{// Read ID fields}
\DataTypeTok{int}\NormalTok{ ctr}\OperatorTok{[}\DecValTok{32}\OperatorTok{];}
\DataTypeTok{char} \OperatorTok{*}\NormalTok{next }\OperatorTok{=} \OperatorTok{\&}\NormalTok{in}\OperatorTok{[}\DecValTok{2}\OperatorTok{];}
\ControlFlowTok{for} \OperatorTok{(}\DataTypeTok{int}\NormalTok{ i}\OperatorTok{=}\DecValTok{0}\OperatorTok{;}\NormalTok{ i }\OperatorTok{\textless{}}\NormalTok{ n}\OperatorTok{;}\NormalTok{ i}\OperatorTok{++)} \OperatorTok{\{}
\NormalTok{  ctr}\OperatorTok{[*}\NormalTok{next}\OperatorTok{]++;}
\NormalTok{  next}\OperatorTok{++;}
\OperatorTok{\}}
\end{Highlighting}
\end{Shaded}

}

\caption{\label{fig-buggy}An example of some buggy parsing code.}

\end{figure}%

\section{Manual testing}\label{manual-testing}

When we manually test code, we consider a \emph{specific execution} of
the program and predefine an \emph{expected result}. For the example of
Figure~\ref{fig-buggy}, we might tests to check that the following two
inputs give the following behavior:

\begin{Shaded}
\begin{Highlighting}[]
\NormalTok{in }\OperatorTok{=} \OperatorTok{\{}\DecValTok{6}\OperatorTok{\};}         \CommentTok{// Should have no effect.}
\NormalTok{in }\OperatorTok{=} \OperatorTok{\{}\DecValTok{5}\OperatorTok{,} \DecValTok{1}\OperatorTok{,} \DecValTok{2}\OperatorTok{\};}   \CommentTok{// Should cause ctr[2] == 1}
\end{Highlighting}
\end{Shaded}

\textbf{Benefits.} The main advantage of manual testing is that tests
can be targeted and can exercise application-specific logic. If you have
a precise correctness condition that your program should ensure, it is
often easiest to test it with manual testing.

An additional advantage is that manual tests are useful in
\emph{regression testing}: If you find a bug in your program today, you
can write a manual test that tests that the buggy condition does not
occur. As you update your program later on, this regression test can
determine whether the same bug occurs again.

\textbf{Downsides.} The downsides of manual tests are that they are
expensive to write, the test cases can themselves be buggy (especially
when the program is complicated), it requires a lot of program-specific
understanding, and it is difficult to write enough tests to cover a
large fraction of the program's behavior.

\section{Fuzzing}\label{fuzzing}

Fuzzing is the process of running a program on a very large number of
\emph{randomly generated inputs}. As soon as a random input causes the
program to crash, we know that we have detected a bug.

When using a fuzzer to test a piece of code, we will typically ask the
compiler to instrument the code with extra instructions to test whether
the code behaved improperly. In C, for example, we will instruct the
compiler to check for out-of-bounds memory accesses.

\marginnote{\begin{footnotesize}

With the GCC compiler, you can compile your code with the
\texttt{-fsanitize=address} flag to insert extra checks for
memory-access errors.

\end{footnotesize}}

To test for application bugs with a fuzzer, we can instrument our code
with assertions that will crash the program if the program ever violates
certain programmer-specified invariants.

Fuzzers may have to test a large number of inputs before finding one
that triggers a bug. In the code of Figure~\ref{fig-buggy}, if a
well-formed packet has \texttt{n=64}, the program will write off the end
of the \texttt{in} array. (The check in Line 6 should test whether
\texttt{n\ \textgreater{}=\ 64} instead of
\texttt{n\ \textgreater{}\ 64}.) For a fuzzer to hit this bug, it will
need to choose a random input value of the form:

\begin{Shaded}
\begin{Highlighting}[]
\NormalTok{in }\OperatorTok{=} \OperatorTok{\{}\DecValTok{5}\OperatorTok{,} \DecValTok{64}\OperatorTok{,} \OperatorTok{...} \DecValTok{64}\NormalTok{ arbitrary values }\OperatorTok{...\};}
\end{Highlighting}
\end{Shaded}

The probability that a uniform randomly bitstring of the appropriate
length hits this bug is \(2^{-8} \cdot 2^{-8} = 2^{-16}\), which is
quite small.

\textbf{Coverage-guided fuzzing.} Real fuzzers do not just feed uniform
random bitstrings to programs that they are trying to fuzz. Instead,
real fuzzers try to pick inputs in a way that maximizes the fuzzer's
\emph{code coverage}: the fraction of the lines of the program's code
that the fuzzed programs have executed.

\marginnote{\begin{footnotesize}

Most random bitstrings will probably not cause the program to exercise a
large fraction of the program's code, since the program will reject them
early most of the time.

\end{footnotesize}}

To implement this strategy, the fuzzer maintains a corpus of bitstrings.
Each time the fuzzer runs the program, it picks an input from the corpus
and randomly mutates it in some way (e.g., by changing or adding a
byte). If running the new string on the program causes the program to
execute some new lines of code (i.e., the coverage increases), the
fuzzer adds the newly mutated string to the corpus.

While this strategy does not have a robust theory to support it,
coverage-guided fuzzing works shockingly well in practice. Many major
software projects use fuzzing extensively to find bugs.

A coverage-guided fuzzer run on Figure~\ref{fig-buggy} might find the
following input that causes the program to crash with an out-of-bounds
write:

\begin{Shaded}
\begin{Highlighting}[]
\NormalTok{in }\OperatorTok{=} \OperatorTok{\{}\DecValTok{5}\OperatorTok{,} \DecValTok{1}\OperatorTok{,} \DecValTok{100}\OperatorTok{\};}
\end{Highlighting}
\end{Shaded}

\textbf{Benefits.} A major benefit of fuzzers is that they are almost
completely automated---they require very little input from the
programmer. Since programmer time is more expensive than machine time,
finding bugs using fuzzers is often much cheaper than finding bugs via
manual test cases. Since fuzzers execute the program on billions of
inputs, they will often find tricky bugs that a human might never find.

\textbf{Drawbacks.} A drawback of fuzzers is that they cannot find
application-specific bugs: they are essentially limited to only finding
violated assertions in a program. So while fuzzers are a useful tool for
finding bugs, they are typically only useful in conjunction with manual
tests.

\textbf{Generalizations of fuzzing.} The first fuzzers used random
bitstrings as their initial pool of inputs. More recent fuzzers have
application-specific logic for handling HTML, JSON, or other file
formats---these fuzzers are better at catching higher-level logic
errors. Some languages, such as Go, have support for fuzzing in their
test infrastructure.

\section{Symbolic execution}\label{symbolic-execution}

A weakness of fuzzing is that a fuzzer may not be able to trigger bugs
that hide behind \texttt{if} conditions that are very very rarely true.
Symbolic execution is a testing strategy that can find bugs in these
difficult-to-fuzz programs.

\marginnote{\begin{footnotesize}

For example, if a parser first checks a CRC32 checksum on a packet, a
fuzzer will almost never find an input that causes the program to run
past the checksum check. Another example of difficult-to-fuzz code might
be some HTML parsing code that checks that every \texttt{\textless{}}
symbol is followed by an \texttt{\textgreater{}} symbol.

\end{footnotesize}}

The idea of symbolic execution is that we will run the program. But
instead of running the program on actual concrete input values, we will
run the program on \emph{symbolic variables} that represent arbitrary
values. For example, we might want to use symbolic execution to run the
following simple snippet of code:

\begin{Shaded}
\begin{Highlighting}[]
\NormalTok{c }\OperatorTok{=}\NormalTok{ a }\OperatorTok{+}\NormalTok{ b}\OperatorTok{;}
\NormalTok{e }\OperatorTok{=}\NormalTok{ d }\OperatorTok{+}\NormalTok{ c}\OperatorTok{;}
\NormalTok{f }\OperatorTok{=}\NormalTok{ a }\OperatorTok{*}\NormalTok{ d}\OperatorTok{;}
\end{Highlighting}
\end{Shaded}

When executing this code, the state of memory might look like this,
where we replace the value of the variable \texttt{d} with a variable
\(X\):

\begin{verbatim}
--------------------------------------------------
...    |  5  |  7  |  12  |  X  |     |     |  ...
--------------------------------------------------
          a     b      c     d     e     f
\end{verbatim}

The major headache when using symbolic execution is \emph{control flow}.
For example, we might have code that looks like this:

\begin{Shaded}
\begin{Highlighting}[]
\ControlFlowTok{if} \OperatorTok{(}\NormalTok{e }\OperatorTok{==}\NormalTok{ f}\OperatorTok{)} \OperatorTok{\{}
\NormalTok{  BUG}\OperatorTok{();}
\OperatorTok{\}} \ControlFlowTok{else} \OperatorTok{\{}
  \CommentTok{// Something else}
\OperatorTok{\}}
\end{Highlighting}
\end{Shaded}

If running with the symbolic variable \(\mathtt{d} = X\), then we will
have the condition: \texttt{(d+c\ ==\ a*d)}, which simplifies to:
\texttt{(X+12\ ==\ 5*X)}.

To handle this sort of case, we can use a program called a \emph{SAT
Solver} to search for inputs that cause the condition to be true. For
example, a SAT solver run on the branch condition
\texttt{(X+12\ ==\ 5*X)} will likely find the value \texttt{X==3} that
causes the condition to be true. Then the symbolic-execution engine can
continue executing the program down the \texttt{true} branch of the
program with the constraint \texttt{X==3} on the symbolic variable
\(X\). In parallel, the engine can search for values of \texttt{X} that
cause the program to go down the \texttt{false} branch of the program. A
SAT solver might find \texttt{X==908234} as one input that causes the
program to traverse the \texttt{false} branch.

\marginnote{\begin{footnotesize}

An important thing to know is that there is \emph{no guarantee} that a
SAT solver will actually be able to find an input that causes the
program to execute one branch or another. The reason is that we know of
no efficient algorithm for finding an assignment of the variables for an
arbitrary condition that causes the condition to be true. (This problem
is NP complete.) At the same time, for finding branch conditions in
``reasonable'' programs, SAT solves work surprisingly well.

\end{footnotesize}}

Running a symbolic-execution engine on the code of
Figure~\ref{fig-buggy} might produce the following output:

\begin{verbatim}
ERROR: buggy.c:10 memory error
ERROR: buggy.c:12 memory error
\end{verbatim}

The symbolic-execution engine may consider the following tree of program
executions, branching on each condition:

\begin{verbatim}
// Input data
in = {i0, i1, i2, i3, i4, ...};

[i0 == 5]
  /
  | FALSE -> return
  | TRUE  -> [i1 > 64]
  \           /
              | TRUE  -> return
              | FALSE -> [0 < i1]
              \           /
                          | FALSE -> return
                          | TRUE  -> [i2 < 0 || i2 >= 32]
                          \             /
                                        | TRUE  -> BUG!
                                        | FALSE -> [1 < i1]
                                        \           /
                                                    | ...continue ...
                                                    | ...execution...
                                                    \
\end{verbatim}

\textbf{Benefits.} Symbolic execution can find tricky bugs involving
complicated branch conditions that fuzzers and manual tests may not
find. In addition, symbolic execution may be able to find bugs that do
not crash the program. In addition, a symbolic-execution engine can in
principle consider \emph{all possible inputs} to a program and can give
a guarantee that the program has no bugs of a certain type (e.g.,
out-of-bounds read).

\textbf{Drawbacks.} Symbolic-execution engines can be very slow to run
and often work poorly on very large pieces of code. As the length of an
execution grows, the symbolic-execution engine accumulates more and more
symbolic variables (representing values in memory) and the number of
constraints on each symbolic variables grows as well. When there are
many variables and many constraints, the SAT solver may not be able to
determine---in a reasonable amount of time---whether there is or is not
a satisfying assignment to the variables. For these reasons, symbolic
execution can work well for small snippets of code; in large programs
(such as a web browser), symbolic execution may not be able to progress
very deep into the program.

\marginnote{\begin{footnotesize}

To address these drawbacks in symbolic executions, one approach is to
write a \emph{scheduler} that guides the search that the symbolic
execution makes through the program state. A second approach is to
define \emph{loop invariants} or \emph{function invariants} that aim
simplify the job that the SAT-solver must perform by giving it more
information about the program's expected behavior.

\end{footnotesize}}

\chapter{Runtime Defenses}\label{runtime-defenses}

We have considered a number of ways to improve software security. First,
we explored \emph{privilege separation} as a way to architect a system
so that bugs do not lead to catastrophic security failures. Next, we
looked at \emph{bug finding}---techniques to find and eliminate bugs in
source code before we use the code in production. Now, we will discuss
\emph{runtime defenses}: how to detect buggy behavior as it occurs in a
piece of production software so that we can halt the program when a bug
occurs.

There are a number of reasons why it is difficult to build runtime
defenses. We must:

\begin{itemize}
\tightlist
\item
  determine the classes of bugs that we want to find,
\item
  identify which components of the system that we want to monitor,
\item
  figure out how to avoid \emph{false positives} (erroneously halting a
  program when there is no bug), and
\item
  try to implement these runtime defenses with minimal overhead, since
  we will apply these defenses to production code.
\end{itemize}

\section{Defenses against buffer
overflows}\label{defenses-against-buffer-overflows}

We have discussed the pervasive buffer overflow attack, which takes
advantage of a missing bounds check to overwrite memory beyond the
bounds of an array, often modifying the current function's return
address to cause the attacked system to run attacker-specified code
which is placed in the buffer itself.

An example of C code that is vulnerable to a buffer-overflow attack is
this:

\begin{Shaded}
\begin{Highlighting}[numbers=left,,]
\DataTypeTok{void}\NormalTok{ f}\OperatorTok{()} \OperatorTok{\{}
  \DataTypeTok{char}\NormalTok{ buf}\OperatorTok{[}\DecValTok{128}\OperatorTok{];}
\NormalTok{  gets}\OperatorTok{(}\NormalTok{buf}\OperatorTok{);}
\OperatorTok{\}}
\end{Highlighting}
\end{Shaded}

The call-stack layout for this program will look like this:

\begin{verbatim}
        |                 |
        |-----------------|
        | Return address  |
        |-----------------|
        | buf[127]        |
        | ....            |
        | buf[1]          |
        | buf[0]          |
        |-----------------|
        |                 |
        |                 |
        |     ....        |
        |-----------------|
\end{verbatim}

The \texttt{gets} function in C will read from \texttt{stdin} until it
finds a NULL (\texttt{\textbackslash{}0}) character in the input string.
If the string has length \(>127\), the \texttt{gets} function could copy
adversarially generated input byte onto the stack, overwriting the
return address of the function \texttt{f}.

There are three steps involved in executing a buffer-overflow attack:

\begin{enumerate}
\def\labelenumi{\arabic{enumi}.}
\tightlist
\item
  write past the end of a buffer,
\item
  cause the victim process to jump to an adversary-controlled address,
  and
\item
  cause the victim process to run adversarial code.
\end{enumerate}

Runtime defenses against buffer overflows can try to disrupt each of
these three steps.

\subsection{Non-Executable Stack}\label{sec-overflow-nx}

A first defense against buffer-overflow attacks is to prevent the CPU
from executing code on the C call stack. In traditional buffer-overflow
attacks, an attacker will somehow place some adversarial code on the
stack (e.g., in the buffer \texttt{buf} in the example above) and then
cause the victim process to jump to that code on the stack.

However, C usually puts normal program code in a separate region of
memory---not on the stack. Modern CPUs allow us to add permissions to
different memory regions: the OS can mark each region as read (R), write
(W), and/or execute (X). To make buffer-overflow attacks more difficult,
we can prevent the CPU from running code from the stack at all by
marking the stack as RW only.

\marginnote{\begin{footnotesize}

This is sometimes called an NX defense---for ``no execute.''

\end{footnotesize}}

Marking the stack as non-executable seems to eliminate a major piece of
an effective buffer overflow attack: the attacker can no longer supply
code of their choice and point to it with the return address. However,
modern attackers have worked around this with something called
\emph{return-oriented programming}: with the ability to supply their own
code removed, attackers must find code that already exists to do what
they like. This may be full existing functions, but more likely
attackers will set the return address to point into the middle of some
function and execute just a fragment that does something useful. It
turns out that with more work, it is possible to perform many attacks
using only code that already exists in a victim process.

\subsection{Stack Canary}\label{stack-canary}

To try to remove the adversary's ability to overwrite the return address
in the presence of a buffer overflow, another defense is to insert a
\emph{stack canary} in every stack frame between the function variables
and the return address. A stack canary is typically a secret random
value, chosen when the program starts running.

\marginnote{\begin{footnotesize}

A buffer-overflow attack does not directly allow the attacker to read
arbitrary memory, so the attacker has no direct way to read the canary
before overflowing the buffer.

\end{footnotesize}}

\marginnote{\begin{footnotesize}

A clever non-random canary includes a collection of the
string-terminating characters, such as
\texttt{\textbackslash{}0\textbackslash{}n\textbackslash{}r}. Many C
functions that read strings from input, such as \texttt{gets} will stop
reading once they reach one of these values, so it could be difficult
for an attacker to feed in an overflowing string that includes these
characters.

\end{footnotesize}}

At the start of each function, the compiler inserts instruction that
write this canary to some value. At the end of the function before
returning, the compiler also adds some code that checks that the canary
value has not changed. If the canary has changed, the attacker must have
overflowed a buffer and the program should exit to avoid running unknown
code.

\begin{verbatim}
        |                 |
        |-----------------|
        | Return address  |
        |-----------------|
        | Stack canary    |
        |-----------------|
        | buf[127]        |
        | ....            |
        | buf[1]          |
        | buf[0]          |
        |-----------------|
        |                 |
        |                 |
        |     ....        |
        |-----------------|
\end{verbatim}

This is effective because in a buffer-overflow scenario, the attacker
needs to write memory sequentially until the address they care about
writing is reached: if the canary is between the function variables and
the return address, the attacker must overwrite the canary to modify the
return address. However, this is not a perfect defense: if the attacker
writes the same value to the canary as was already there, it will go
undetected. Therefore, the canary value must be hard for the attacker to
guess.

This defense is still not perfect---for example, it does not prevent an
attacker from overwriting function pointers. However, it does make a
successful attack significantly harder.

There are a few ways to subvert canaries:

\begin{itemize}
\tightlist
\item
  An attacker can corrupt \emph{data} on the stack, even if it does not
  corrupt the return address. This could be very bad for the program's
  behavior.
\item
  An attacker might find a way to read the canary from memory (e.g., if
  there is some other bug in the program) and then execute the
  traditional buffer-overflow attack to overwrite the return pointer.
\item
  In a forking web server, the child process may have the same canary
  value as the parent process. An attacker can potentially exploit this
  to learn the canary (See (\citeproc{ref-bittau2014hacking}{Bittau et
  al. 2014})).
\end{itemize}

\subsection{Address Space Layout Randomization
(ASLR)}\label{address-space-layout-randomization-aslr}

Another approach to defend against buffer overflow-style attacks is to
make it more difficult for the adversary to guess a useful address to
jump to. To do this, many modern systems randomize the locations of
code, stack, and heap memory regions when a process starts. With this
defense in place, an attacker needs to learn the location of the code
memory region in order to mount a return-oriented programming attack
(Section~\ref{sec-overflow-nx}).

\marginnote{\begin{footnotesize}

Implementing ASLR requires compiler support. Most modern compilers do.

\end{footnotesize}}

A weakness of ASLR schemes is that they typically shift the location of
an entire region of memory: the entire heap, stack, and code sections
move around in memory, but the layout within each section is typically
fixed at compile time. Thus if the adversary can learn the location in
memory of the code for a single function, it can mount an effective
return-oriented programming attack.

\subsection{Bounds Checking with Fat
Pointers}\label{bounds-checking-with-fat-pointers}

All of the defenses so far have attempted only to minimize the damage of
a buffer overflow after an attacker has exploited it. However, we could
prevent a more comprehensive suite of attacks if we could make sure that
our code never reads or writes a pointer that is outside the bounds of a
given buffer. Memory-safe languages such as Go, Rust, and Python have
this bounds checking built in. For older languages, such as C, we can
try to retrofit the C compiler to achieve this bounds checking. As we
now discuss, implementing bounds checking in C is challenging.

One way to implement bounds checking in C is a technique called
\emph{fat pointers}. When using fat pointers, the compiler changes the
representation of a pointer to include not only an address, but also the
\emph{base} and the \emph{limit} of the buffer the pointer points to.
This base and limit are initialized on an allocation, and the compiler
inserts bounds checks on each pointer dereference to guarantee that the
dereferenced value is within the array bounds that the base and limit
specify. Pointer arithmetic preserves the base and limit but modifies
the pointer itself as before, allowing the pointer to possibly go out of
bounds.

For example, if a programmer allocates an array of 128 bytes using
\texttt{void*\ ptr\ =\ malloc(128)}, the compiler will associate values
\texttt{base\ =\ ptr} and \texttt{limit\ =\ ptr+128} with the pointer
\texttt{ptr}. If we then assign \texttt{ptr3\ =\ ptr\ +\ 3}, then the
new pointer \texttt{ptr3} will have the same base and limit as
\texttt{ptr}:

\begin{verbatim}
              |                 |
   limit -->  |-----------------|
              | buf[127]        |
              | ....            |
              | buf[4]          |
   ptr3  -->  | buf[3]          |
              | buf[2]          |
              | buf[1]          |
ptr,base -->  | buf[0]          |
              |-----------------|
              |                 |
              |                 |
              |     ....        |
              |-----------------|
\end{verbatim}

When the program deferences \texttt{ptr3}, the compiler will insert
checks to ensure that
\(\mathtt{base} \leq \mathtt{ptr3} < \mathtt{limit}\) and crash the
program otherwise.

One limitation of fat pointers is that they only check overflowing an
\emph{allocated region of memory}---not overflows \emph{within} a region
of memory. For example, if we use \texttt{malloc} to allocate a C
\texttt{struct}, overflowing the \texttt{buf} member of the struct could
allow the attacker to overwrite the values of other elements of the
struct:

\begin{Shaded}
\begin{Highlighting}[]
          \KeywordTok{struct} \OperatorTok{\{}
            \DataTypeTok{int}\NormalTok{ x}\OperatorTok{;}
            \DataTypeTok{char}\NormalTok{ buf}\OperatorTok{[}\DecValTok{16}\OperatorTok{];}
            \DataTypeTok{void} \OperatorTok{(*}\NormalTok{f}\OperatorTok{)();}
          \OperatorTok{\}}
\end{Highlighting}
\end{Shaded}

If there are function pointers in the \texttt{struct}, the result of
this within-region overflow could be as bad as overflowing the return
address.

Unlike a nonexecutable stack, stack canaries, and ASLR, fat pointers are
not widely used. This is largely because the modified ``fat'' pointers
can break the functionality of existing C code. In particular, a fat
pointer on a 64-bit architecture will typically take more than 64 bits
to represent. If the programmer cast a pointer to an \texttt{int} and
back again, the behavior of the program could change when using fat
pointers versus when using unmodified 64-bit pointers. In addition, a C
program may implicitly require that a pointer takes exactly 64 bits to
represent (e.g., as a field within a \texttt{struct}); a compiler using
fat pointers will break such programs.

\subsection{Control-Flow Integrity
(CFI)}\label{control-flow-integrity-cfi}

The goal of techniques for \emph{control-flow integrity} is to limit the
set of addresses that a program will jump to when returning. In this
way, a victim process may be able to detect when it is about to jump to
a return address that an adversary modified with a buffer overflow.
Implementing CFI requires compiler support---many modern compilers
support some form of CFI.

To implement CFI, we need to add several checks on different kinds of
jumps.

\begin{itemize}
\item
  For direct jumps (e.g,. ``\texttt{call\ gets}'', we know at compile
  time that these jumps are valid since there is nothing that the
  attacker can control. There is no need to implement any CFI checks on
  these jumps.
\item
  For indirect jumps that use some variable in the jump target (e.g.,
  function pointers, function returns), the compiler inserts checks to
  insure that the return address is a valid jump target. To do this, the
  compiler inserts into the program a data structure that maintains a
  set of all valid indirect-jump targets.

  \marginnote{\begin{footnotesize}

  Often a compiler will implement this data structure using some sort of
  simple Bloom filter.

  \end{footnotesize}}

  Checking whether a jump point is valid takes a bit of extra
  computation---and thus imposes some runtime cost---it may tolerable in
  practice.
\end{itemize}

\subsection{Unsolved program: Use-after-free
bugs}\label{unsolved-program-use-after-free-bugs}

The defenses in this section have made buffer overflows extremely
challenging to exploit in modern code on modern systems. These defenses
do very little to prevent against \emph{use-after-free} bugs, in which a
program \texttt{free}s allocated memory and then inadvertently reads or
writes the \texttt{free}d pointer.

\section{Taint Tracking to Defend against Input-Sanitization
Bugs}\label{taint-tracking-to-defend-against-input-sanitization-bugs}

We have discussed SQL injection and cross-site scripting, which allow an
attacker to run code by adding special characters, such as
``\texttt{"}'' or ``\texttt{\textgreater{}},'' into their input.

Unlike buffer-overflow bugs, input-sanitization bugs arise not because a
bug at one particular point in the program---they are more due to a
systematic failure to consider certain types of inputs. As a result, the
defenses are more systematic as well.

\subsection{Taint tracking in
libraries}\label{taint-tracking-in-libraries}

A common approach to check for sanitization failures at runtime is
called \emph{taint tracking}. In taint tracking, some infrastructure in
the program (e.g., a SQL library) marks any data coming from user input
as \emph{tainted}. At functions that perform escaping, the
infrastructure removes the taint label. The infrastructure marks
sensitive functions, such as the HTML renderer, are marked as
\emph{sinks}. Any time the program runs sink code, the taint-tracking
infrastructure checks that the data headed into the sink is not tainted.
If the data is tainted, some part of the input must not have been
sanitized since it came from the user, and the infrastructure will halt
the program to avoid an exploit. See Figure~\ref{fig-taint} for an
example.

\begin{figure}

\centering{

\begin{Shaded}
\begin{Highlighting}[]
\NormalTok{name }\OperatorTok{=}\NormalTok{ read\_from\_user()     }\CommentTok{\# name is tainted}
\NormalTok{first,last }\OperatorTok{=}\NormalTok{ name.split()   }\CommentTok{\# first,last are tained}
\CommentTok{\# query is tainted}
\NormalTok{query }\OperatorTok{=} \StringTok{"SELECT * FROM USERS WHERE last\_name = \textquotesingle{}"} \OperatorTok{+}\NormalTok{ last }\StringTok{"\textquotesingle{}"}
\NormalTok{query\_database(query)       }\CommentTok{\# will cause an exception}

\NormalTok{qesc }\OperatorTok{=}\NormalTok{ escape(query)        }\CommentTok{\# qesc is not tainted}
\NormalTok{query\_database(qesc)        }\CommentTok{\# will succeed}
\end{Highlighting}
\end{Shaded}

}

\caption{\label{fig-taint}A hypothetical Python example of how
taint-tracking might prevent SQL injection}

\end{figure}%

Many browsers implement taint tracking to prevent cross-site scripting
using \emph{Trusted Types}: for JavaScript calls that update the
displayed HTML, such as \texttt{innerHTML\ =\ foo}, browsers may
restrict the type of \texttt{foo} to ensure that the code explicity
converts its type to something like \texttt{TrustedHTML}. This does not
guarantee that the sanitization was done correctly, but does ensure that
the programmer acknowledged the risk in their code.

\subsection{Taint tracking in language
runtimes}\label{taint-tracking-in-language-runtimes}

Some programming languages associate a \emph{taint bit} with every
string in the program. The application developer can set or clear the
taint bit manually. In addition, the language runtime will automatically
set the taint bit on strings that certain functions (e.g., those that
read from user input) return. An application developer then can
implement some taint-checking policy to systematically prevent against
escaping bugs. For example, if there is one function that makes a SQL
query to the database, the application could check the taint bit is
cleared on any query string that a developer passes to this function.

\subsection{Taint tracking in operating
systems}\label{taint-tracking-in-operating-systems}

Operating systems also use taint tracking, often to flag suspicious
files---typically those that the user downloaded from the Internet.
MacOS, for example, will set a taint bit on any executable that the user
downloaded from the Internet. The OS will maintain this taint bit when
the user copies or moves a tainted file. If the user ever tries to
execute a tainted file, the OS will raise a warning to the user before
executing it.

\part{Privacy}

\chapter{Privacy and zero-knowledge
proofs}\label{privacy-and-zero-knowledge-proofs}

So far, we have discussed several cryptography primitives, from hash
functions to encryption schemes, and explored many applications of those
primitives to systems security. These primitives have provided security,
but in a very all-or-nothing sense: in order to provide broadly
applicable security, our definitions required that a certain party (with
the key) could either completely decrypt the message, learning the
message contents, or cannot do anything with the message at all.

\section{Zero-knowledge proofs of knowledge}\label{sec-zkdefn}

Zero-knowledge proofs allow one party (called a prover) to convince
another party (called a verifier) that the prover ``knows'' some secret,
without revealing anything else about the secret. More formally, say
that both parties hold a function \(f\), and a value \(y = f(x)\). The
prover may want to convince the verifier that it ``knows'' an \(x\) such
that \(y=f(x)\). We call such a prover-verifier interaction a
\emph{proof system}. If the prover reveals no information about \(x\),
apart from the fact that \(y=f(x)\), to the verifier, we say that the
proof system is a \emph{zero-knowledge proof of knowledge}.

The proof takes the form of an interaction between the prover and the
verifier, in which the two parties exchange a sequence of messages. At
the end of the interaction, the verifier either accepts the proof or
rejects it. This is an \emph{interactive proof}: the prover and verifier
may exchange many messages. In contrast, in a classical
(non-interactive) proof, the prover writes down a proof and sends it to
the verifier in a single message.

\marginnote{\begin{footnotesize}

Probably all of the proofs you have seen in your life up until this
point, including the ones in your math textbooks, are classical proofs.

\end{footnotesize}}

In our example, both parties hold a function \(f\) and a value \(y\).
The prover wants to convince the verifier that it knows a value \(x\)
such that \(y=f(x)\). To be a zero-knowledge proof of knowledge, the
proof system must have three properties. We state this informally:

\marginnote{\begin{footnotesize}

The literature usually refers to the secret value \(x\) as the
\emph{witness}.

\end{footnotesize}}

\begin{itemize}
\item
  \emph{Completeness.} If the prover ``knows'' \(x\), it will always
  convince an honest verifier to accept.
\item
  \emph{Soundness.} If the verifier accepts the proof, then prover must
  really ``know'' \(x\). The soundness property is about protecting an
  honest verifier from a cheating prover: if the prover does not know
  \(x\), then no matter how it cheats it will not be able to convince
  the verifier that it does.
\item
  \emph{Zero-knowledge.} The verifier ``learns nothing'' about \(x\),
  apart from the fact that \(y=f(x)\), as a result of interacting with
  the prover. The zero-knowledge property is about protecting the prover
  from a malicious verifier: if the prover is honest, then no matter how
  the verifier misbehaves, it should learn nothing about the prover's
  secret witness \(x\).
\end{itemize}

\todobox{\textbf{TODO (HCG):} Insert formal definitions here?}

One important question here is: How do we define the notion of
knowledge? That is, what does it mean for the prover to really ``know''
\(x\)? A naïve definition might say that \(x\) is stored somewhere in
the prover's memory, but that might be too strong: a prover might
``know'' \(x\) without storing it literally in its memory, perhaps
storing it in base64-encoded form. A more meaningful definition of
knowledge, which is the one we will use, thinks of knowledge in terms of
extractability. That is, we say that a prover knows \(x\) if there is an
efficient algorithm that extracts \(x\) by interacting with the prover
and observing the prover's internal state.

We also must pin down what it means for the verifier to learn nothing
through an interaction with the prover. A good way to define this, it
turns out, is to think of ``learning nothing'' in terms of the verifier
being able to simulate its interaction with the prover without actually
talking to the real prover. Specifically, if the verifier can produce a
transcript of its hypothetical interaction with the prover---that is
cryptographically indistinguishable from a real transcript---without
actually communicating with the prover, we say the verifier has learned
nothing about \(x\).

\marginnote{\begin{footnotesize}

There is a rich theory of zero-knowledge proof systems that we will not
be able to cover here. One surprising fact is that there exist
zero-knowledge proofs of knowledge for \emph{any} choice of the public
function \(f\), provided that it is an efficiently computable function.

\end{footnotesize}}

One immediate application of zero-knowledge proofs of knowledge is to
authentication. The prover's witness \(x\) becomes the user's secret
key, the verifier's \(y\) is the user's public key, and the function
\(f\) is a public parameter of the authentication scheme. The
authentication protocols we've seen so far, such as challenge-response
protocols using MACs or signatures, do not quite satisfy the notions of
soundness and zero-knowledge we use here. For example, in a MAC-based
challenge-response protocol, the server learns the MAC of a
server-chosen challenge value under the client's (prover's) secret key.
Provided that the MAC is secure, the verifier could not produce a
simulated transcript that is identical to a real one without actually
interacting with the real prover.

\section{Discrete-log problem and Schnorr
signatures}\label{discrete-log-problem-and-schnorr-signatures}

As an example of a zero-knowledge proof of knowledge, we will show how a
prover can convince a verifier that it knows a discrete log, without
revealing it.

\section{Reminder: The discrete-log
problem}\label{reminder-the-discrete-log-problem}

We recall the discrete-log problem. The problem is parameterized by a
group \(\G\) of prime order \(q\), generated by a value \(g \in \G\).
Then, given a group element \(y = g^x \in \G\), for \(x \getsr \Z_q\),
the discrete-log problem is to find the value of \(x\in \Zq\). We have
already seen the discrete-log problem in the context of Diffie-Hellman
key exchange and elliptic-curve signatures.

\marginnote{\begin{footnotesize}

In cryptographic settings, \(q\) is a big prime.

\end{footnotesize}}

\section{Schnorr's protocol for proving knowledge of discrete
logs}\label{schnorrs-protocol-for-proving-knowledge-of-discrete-logs}

In Schnorr's protocol, the prover and verifier both hold the generator
\(g \in \G\) and the group order \(q\). In addition the prover holds a
secret value \(x \in \Z_q\), the verifier holds a public value
\(y \in \G\). The prover must convince the verifier that it knows the
discrete log of \(y\) base \(g\): that is, that it knows \(x\) such that
\(y=g^x \in \G\). The prover wants to prove this to the verifier this
without revealing anything about else about its secret value \(x\).

The protocol goes as follows:

\begin{itemize}
\item
  The prover picks a random \(r\in \Z_q\), and sends \(R=g^r\in \G\) to
  the verifier.
\item
  The verifier picks a challenge bit \(c\) at random, either 0 or 1, and
  sends the challenge to the prover.
\item
  The prover sends back a response \(z \in \Z_q\), chosen as follows:

  \marginnote{\begin{footnotesize}

  In the true Schnorr protocol, the verifier samples \(c \getsr \Z_q\),
  as we discuss in Section~\ref{sec-hvzk}

  \end{footnotesize}}

  \begin{itemize}
  \tightlist
  \item
    If \(c=0\), the prover sets \(z \gets r \in \Z_q\).
  \item
    If \(c=1\), the prover sets \(z \gets r+x\in \Z_q\).
  \end{itemize}

  Mathematically, \(z=r+cx\in \Z_q\).
\item
  The verifier checks that \(g^z=R y^{c}\in \G\). If equal, the verifier
  accepts, otherwise the verifier rejects.
\end{itemize}

Why is interaction necessary here? It turns out that this protocol
critically depends on the prover not knowing the challenge bit \(c\) in
advance. If the prover knew \(c\) in advance, the prover could pick
\(z \getsr \Z_q\) at random, and just compute \(R\) as
\((g^z)(y^c)^{-1}\in \G\).

\textbf{Reducing soundness error.} In any given interaction of the
protocol, the prover could falsely convince the verifier that it knows
\(x\), using the attack we just described. The prover only will succeed
in this attack with probability \(\frac 12\), which is exactly the
probability that the prover guesses the verifier's challenge \(c\) in
advance. By repeating this protocol \(\lambda\) times and accepting only
if all \(\lambda\) iterations accept, the verifier can reduce the
probability of accepting a cheating prover (called the ``soundness
error'') to \(2^{-\lambda}\).

A more efficient way to reduce the soundness error is to have the
verifier sample the challenge value \(c \getsr \Z_q\) as a random
\(\Z_q\) element instead of a random bit. With this transformation, the
soundness error becomes roughly \(1/q\), which is \(\approx 2^{-256}\)
for the cryptographic groups we use in practice. An important caveat is
that if we increase the challenge space in this way, we cannot prove
that the resulting protocol is zero knowledge. This modified protocol
only satisfies a weaker flavor of zero knowledge called
``honest-verifier zero knowledge.'' We won't discuss those details here.

\section{Analysis of Schnorr's
protocol}\label{analysis-of-schnorrs-protocol}

We will now argue that Schnorr's protocol is really a zero-knowledge
proof of knowledge of discrete log. To do this, we must show that the
protocol satisfies completeness, soundness, and zero knowledge, as we
defined them in Section~\ref{sec-zkdefn}.

The completeness follows by construction. The more challenging steps are
proving soundness and zero knowledge.

\subsection{Soundness using an
extractor}\label{soundness-using-an-extractor}

We first need to convince ourselves that the protocol is \emph{sound}:
that is, if the verifier accepts, the prover really knows \(x\). We can
prove soundness by showing that, if there is a (possibly adversarial)
prover \(P^*\) that can convince our honest verifier to accept, there
exists an efficient extractor that can use the cheating prover \(P^*\)
to extract \(x\) with some high probability (say, \(\gg 1/2\)).

For simplicity, assume here that the prover \(P^*\) manages to convince
the verifier to accept with probability \(1\)---i.e., for all choices of
the verifier's challenge bit \(c\) and all choices of the prover's
randomness. Then, given a prover algorithm \(P^*\), the extractor for
Schnorr's protocol works as follows:

\begin{itemize}
\item
  The extractor runs the prover \(P^*\) to get some initial protocol
  transcript with challenge \(c=0\). The transcript consists of
  \((R, c=0, z)\).
\item
  The extractor rewinds the state of the prover \(P^*\) to just after
  the point when it sent the first protocol message \(R\) to the
  verifier.
\item
  The extractor runs \(P^*\) again, but this time feeds its the
  challenge \(c=1\). This produces a second transcript \((R, c=1, z')\).
\item
  The extractor outputs \(z'-z\) as the discrete log \(x\).
\end{itemize}

By our assumption---namely, that \(P^*\) can convince our honest
verifier to accept with probability one---both runs of the prover
\(P^*\) by the extractor must convince the verifier to accept.
(Otherwise, the prover \(P^*\) would not convince the verifier in at
least one of the two runs of the protocol.) That means that both
\(g^z=R \in \G\) and that \(g^{z'}=RX \in \G\). Thus,
\(g^{z'-z}=y \in \G\), which is exactly the definition of what it means
for something (\(z'-z\) in particular) to be the discrete log of
\(y \in \G\).

Here we have glossed over the details of what happens if the prover
convinces the verifier with some non-negligible probability that is
nonetheless less than 1. A similar but more involved argument handles
that case.

\subsection{Zero-knowledge using a
simulator}\label{zero-knowledge-using-a-simulator}

Our second task is to show that Schnorr's protocol satisfies zero
knowledge. That is, that even a malicious verifier \(V^*\) that deviates
from the protocol ``learns nothing'' about the prover's witness \(x\) as
a result of interacting with the prover in the Schnorr protocol. To do
this, given a malicious verifier algorithm \(V^*\), we construct a
\emph{simulator} algorithm. The simulator's job is to produce a
transcript of the prover-verifier interaction that is indistinguishable
from a true prover-verifier interaction. Crucially, in the simulated
interaction, the simulator \emph{does not know} the witness \(x\), while
in the real prover-verifier interaction, the prover \emph{does know} the
witness \(x\).

Given the verifier \(V^*\), we construct this simulator as follows:

\begin{itemize}
\item
  Make a guess \(c' \getsr \{0,1\}\) of the verifier's challenge bit.
\item
  Choose the third protocol message \(z \getsr \Z_q\) at random.
\item
  Set the first message to \(R \gets g^z y^{-c'} \in \G\).
\item
  Run the verifier \(V^*\) on first message \(R\). The verifier outputs
  a challenge \(c\).
\item
  If \(c=c'\), output \((R, c, z)\) as the simulated transcript. This
  happens with probability \(1/2\).
\item
  Otherwise, retry. This happens with probability \(1/2\) as well.
\end{itemize}

Both the extractor and simulator constructions rely crucially on being
able to rewind the prover or verifier and try again. In a real
prover-verifier interaction, the verifier cannot rewind the prover and
the prover cannot rewind the verifier. This is why the existence of an
extractor does \emph{not} mean that in a real protocol interaction the
verifier can extract the witness \(x\) from the prover.

\section{Fiat-Shamir heuristic and Schnorr signatures}\label{sec-hvzk}

One challenge with Schnorr's protocol as we have presented it, and
zero-knowledge proof systems more generally, is that they are
interactive. They require the prover and verifier to send messages back
and forth. A clever trick, called the Fiat-Shamir heuristic, allows us
to turn interactive zero-knowledge proofs into non-interactive proofs,
provided that

\begin{itemize}
\item
  we use a cryptographic hash function that we model as a random oracle,
  and
\item
  the verifier only sends the prover independent random values.

  \marginnote{\begin{footnotesize}

  In Schnorr's protocol, the verifier's challenge is just a random bit,
  so it satisfies this second property.

  \end{footnotesize}}
\end{itemize}

The idea is that, whenever the protocol expects the verifier to send a
challenge to the prover, we can instead derive the challenge just as a
hash of the protocol transcript so far. The prover can then execute the
zero-knowledge protocol against a synthetic verifier (the hash
function), and produce a transcript. The prover can now send this
transcript to the verifier, and if the verifier can confirm that indeed
all of the challenges were correctly computed using a cryptographic hash
function, such as SHA-256, it's sound to accept this transcript as a
proof.

We can use this Fiat-Shamir heuristic to construct a (regular,
non-interactive) signature scheme from an (interactive) zero-knowledge
proof of knowledge for the discrete-log problem. The elliptic-curve
signatures used in practice today effectively take the Schnorr protocol
for proof of knowledge of discrete log using challenge space \(\Z_q\),
and apply the Fiat-Shamir heuristic to it. The resulting signature
scheme is called the \emph{Schnorr signature scheme}.

\chapter{Differential Privacy}\label{differential-privacy}

In many of the systems we have seen so far---encryption, iOS security,
etc---we try to prevent an adversary from learning \emph{any
information} about our sensitive data. Today we will be talking about
how to deal with situations in which we \emph{want} to leak some
information about sensitive user data, but we want to do it in a
``privacy-preserving'' way.

For example:

\begin{itemize}
\tightlist
\item
  the U.S. Census collects sensitive demographic information about U.S.
  citizens and then publishes aggregate statistics about it,
\item
  Internet Service Providers may want to publish aggregate statistics
  about network traffic,
\item
  Google may want to publish anonymized data about user search queries,
  or
\item
  public-health officials may want to publish metrics about a
  population's health at large without revealing any individual user's
  health status.
\end{itemize}

In each of these settings, we many \(n\) users, where each user \(i\)
has some data \(x_i\). There is a central party that has collected
\((x_1, \dots, x_n)\) and wants to publish an aggregate statistic
\(f(x_1, \dots, x_n)\) in a way that ``protects the privacy'' of each
user's data.

\marginnote{\begin{footnotesize}

The techniques we describe here also apply when there is only a single
user and the user wants to publish some function about their private
data.

\end{footnotesize}}

\begin{verbatim}
  user x_1 -----
               |
  user x_2 -----
               |
   ...        ...
               |                         Publish
  user x_n ---------> Database ---> f(x_1, ..., x_n)
\end{verbatim}

\begin{itemize}
\tightlist
\item
  \textbf{Q1.} What functions of private data can we publish without
  compromising user privacy? How do we even define ``privacy'' in this
  context?
\item
  \textbf{Q2.} Once we have a definition of privacy, how do we achieve
  it in practice?
\end{itemize}

Differential privacy gives us \emph{one possible} answer to the question
Q1 here---it is a \emph{definition} of privacy that has an number of
appealing properties. For reasons we will discuss, it is not a perfect
definition of privacy, but it is essentially the best one we have.

Once we buy the privacy definition that differential privacy gives us,
there is a rich literature that explains mechanisms for achieving
differential privacy (i.e., the answer to question Q2).

\section{A bad idea: Attempt to anonymize the
data}\label{a-bad-idea-attempt-to-anonymize-the-data}

A common \textbf{bad} strategy for attempting to release data with
privacy protections is to try to \emph{anonymize} the data. For example,
if a hospital wants to release medical-records data, they could publish
the records with the names redacted. One surprise---that has bitten many
companies and governments---is that it is shockingly easy to re-identify
data in datasets that are supposedly anonymized.

\textbf{Example: Re-identification by Linking (Sweeney 1997).} For
example, in Massachusetts, the Group Insurance Commission released a
dataset that they believed to be anonymized. This dataset contained
health data, including patient ethnicity, ZIP code, birth date, sex,
date of visit, diagnosis, procedure, and medication given.

In 1997, Latanya Sweeney demonstrated that this anonymization strategy
completely failed. In particular, she purchased (for \$20!) the public
voter-registration data from the government of Cambridge, MA. This
second dataset included voter name, address, ZIP code, birth date, and
gender. Crucially, this other database included each voter's zip code,
birth date, and gender! Using the voter-registration dataset---which
linked names to zip codes and birth dates---and the medical data
set---which linked zip codes and birth dates to medical
conditions---Sweeney was able to determine which people had which
medical conditions. That is, Sweeney was able to completely de-anonymize
the GIC dataset and reveal private medical information for everyone up
to the governor of Massachusetts.

\textbf{Netflix Competition.} In 2006, Netflix aimed to improve their
movie recommendation system. To do this, they planned to have
researchers compete to come up with the best movie-recommendation
algorithms. For the purposes of running the competition, Netflix
published a supposedly anonymized dataset that included a randomized
user id, movie id, rating, and date. Netflix released a portion of the
dataset, and the research group that could produce a model that best
predicted ratings for the unreleased set of movies would win a million
dollars.

One group of researchers was able to link individual records in the
Netflix database to records from the public IMDb database. This allowed
them to de-anonymize the data from Netflix even though the database
contained no identifying information whatsoever!

\textbf{The bottom line.} Anonymizing data sets simply does not work.

\section{\texorpdfstring{A flawed idea:
\(k\)-Anonymity}{A flawed idea: k-Anonymity}}\label{a-flawed-idea-k-anonymity}

The \(k\)-anonymity approach to privacy is to think of a dataset release
as being sufficiently private if there are at least \(k\) users' data in
the dataset that are identical. For example, if there are \(k=10\)
people in a medical-records database with the same birth date, zip code,
and medical condition, we might be okay with publishing the dataset with
those three fields only when there are at least \(k=10\) users whose
medical records are identical on those three fields.

The flawed intuition here is that even if there is still some leakage,
maybe it is not so bad. For example, even if you can link the
voter-registration database to the medical-records database, you might
hope that it will be difficult to figure out exactly which user has
exactly which medical condition.

There are a few issues with the \(k\)-anonymity approach to privacy:

\begin{itemize}
\tightlist
\item
  It is not robust to \emph{side information}. For example, if an
  attacker knows a few people with a given medical condition, it can use
  a process of elimination to de-anonymize the dataset still.
\item
  The anonymized dataset may still leak \emph{some private information}.
  For example, say that every student in a class of 100 got the same
  grade. The teacher might be happy publishing the grades of all
  students (without names) since this release satisfies
  \(k=100\)-anonymity. The problem is that the anonymized release
  reveals the \emph{grade of every student in the class}---definitely
  not exactly what we would expect from a privacy-preserving release.
\end{itemize}

\section{A flawed idea: Publish only aggregate
statistics}\label{a-flawed-idea-publish-only-aggregate-statistics}

Another approach we might consider is to publish only summary statistics
with the aim that these summary statistics compress the data so much
that it is impossible to learn anything meaningful about an individual
from them. However, we have to be careful: even a few statistics can
reveal sensitive information.

For example, consider a company that released the average salary of its
employees regularly. If the company releases this average before and
after the resignation of one individual, anyone who knows that that
individual resigned can learn his salary.

So, releasing only statistics does not cleanly protect privacy either.

\section{A new definition: Differential privacy}\label{sec-dp-defn}

Differential privacy is a very mathematically clean definition of
privacy that addresses many of the concerns of the flawed ideas above.
It is not perfect---for reasons we will see---but many regard it as the
best formal definition of privacy we have.

Say that \((x_1, \dots, x_n)\) is a dataset of \(n\) users' data.
Differential privacy, roughly, says that a it is safe to release an
algorithm \(A(x_1, \dots, x_n)\)---called a \emph{mechanism} in the
language of differential privacy---applied to the dataset if the value
of the algorithm \(A\)'s output is roughly the same if we release any
user \(i\)'s data \(x_i\) with some other data value \(x^*_i\):
\[ A(x_1, \dots, x_n) \approx A(x_1, \dots, x_{i-1}, x^*_i, x_{i+1}, \dots, x_n). \]
In other words, dataset-release mechanism provides differential privacy
if its output looks roughly the same, whether or not a particular user's
data is included in the original dataset.

In picture form, the situation looks like this:

\begin{verbatim}
                      Dataset -----> A(Dataset)
                         |              ^
      Drop Alice's data  |              |  Should be "roughly equal"
                         v              v
                      Dataset' ----> A(Dataset')
\end{verbatim}

To formalize this, we will require that a pair of datasets that differ
in only a single row---that is, one contains a row and the other does
not---are \emph{close}. The formal definition is as follows:

\marginnote{\begin{footnotesize}

Notice that a mechanism \(A\) that outputs nothing trivially satisfies
differential privacy. For a mechanism to be useful, it should add as
little noise as possible to maintain a given level of differential
privacy.

\end{footnotesize}}

\begin{definition}[\(\epsilon\)-differential
privacy]\protect\hypertarget{def-dp}{}\label{def-dp}

An mechanism \(A\) provides \(\epsilon\)-differentially private if, for
all neighboring datasets \(x\) and \(x'\) that differ in only one row,
for all subsets \(S\) of outputs of \(A\):
\[ \text{Pr}[A(x) \in S] \leq e^\epsilon \cdot \text{Pr}[A(x') \in S]. \]

\marginnote{\begin{footnotesize}

We have that \(e^\epsilon \approx 1+\epsilon\) when \(\epsilon\) is very
small.

\end{footnotesize}}

This is approximately saying that whether or not a given row is included
in the dataset, the probability of the output being detectably different
is less than \(\epsilon\).

\end{definition}

\marginnote{\begin{footnotesize}

In the differential-privacy literature, this parameter \(\epsilon\) is
often called the ``privacy budget.''

\end{footnotesize}}

Notice that the definition of differential privacy is parameterized by a
real number \(\epsilon\) that specifies how close the two outputs of the
mechanism \(A\) must be on neighboring datasets. The smaller
\(\epsilon\) is, the stronger the privacy guarantee is.

\marginnote{\begin{footnotesize}

One of the delicate points in practice is figuring out what actual value
of \(\epsilon\) a system should aim to achieve.

\end{footnotesize}}

In order for a non-trivial mechanism to provide differential privacy,
the algorithm must be randomized. Take our salary example from before:
if we want to reveal an average salary of all employees, achieving
differential privacy requires adding \emph{noise} in order to make sure
that the output of the release algorithm looks roughly the same whether
or not a particular employee's salary is included in the computation.

\textbf{Differential privacy is closed under post-processing.} A really
convenient property of the differential-privacy view of privacy is that
the definition is \emph{closed under post processing}. That is, if a
mechanism \(A\) is \(\epsilon\)-differentially private, then for
\emph{every} function \(B\), the composed mechanism \(B(A(\cdot))\) is
also \(\epsilon\)-differentially private.

\textbf{Differential privacy is composes nicely.} In many applications,
we want to publish many functions of the same dataset. An very powerful
property of the definition of differential privacy is that it can nicely
handle this situation. In particular, say that a mechanism \(A_1\) is a
mechanism that satisfies \(\epsilon_1\)-differential privacy. Say that a
mechanism \(A_2\) is a mechanism that satisfies
\(\epsilon_2\)-differential privacy. Then the mechanism
\(B(x) = (A_1(x) \| A_2(x))\) that publishes the output of both
mechanisms is \((\epsilon_1 + \epsilon_2)\)-differentially private.

\marginnote{\begin{footnotesize}

There are more sophisticated composition theorems that give tighter
bounds on the differential privacy parameter of the composed mechanism
when the number of mechanisms you are composing is large.

\end{footnotesize}}

The value of \(\epsilon\) increases with each statistic we
publish---which nicely captures our intuition that the more statistics
you release about a dataset, the more information you are leaking about
it.

\section{Achieving differential privacy: Adding
noise}\label{achieving-differential-privacy-adding-noise}

A common approach to achieving differential is exactly this: adding
random noise to the data in order to add uncertainty about the real
data. The first instance of this was by Warner in 1965, who proposed
using random noise to allow individuals to answer sensitive survey
questions (i.e., revealing some private data) while also giving them
some privacy protection. Warner's approach is called \emph{randomized
response}.

To give an example of randomized response: consider a professor that
wants to learn the fraction of students who cheated on a test. Obviously
no student wants to admit that they cheated on an exam. Using Warner's
approach, the professor could ask students to answer honestly with
probably 2/3, but to lie with probability 1/3. This allows a student who
did cheat to claim that they were lying, as per the directions. However,
with many students, the noise should average out, and the professor can
still learn approximately how many students cheated on the test.

We can apply this same approach if we want to make a certain algorithm
differentially private. Instead of publishing a dataset (or a function
of a dataset) directly, we can publish a \emph{noisy version of it}.
Much of the technical complexity in differential privacy goes into
trying to design exactly how to choose the noise in a way that maintains
as much of the ``signal'' of the underlying dataset as possible.

\section{The Laplace mechanism}\label{the-laplace-mechanism}

The \emph{Laplace mechanism} gives a very general way to take an
aggregate statistic \(f\) and modify it to have
\(\epsilon\)-differential privacy for any choice of \(\epsilon\).

First, we need to define the \emph{global sensitivity} of the function
\(f\). This aims to capture the maximal change in \(f\) that results
from changing one of the function's \(n\) inputs. Intuitively, the
higher the sensitivity of a function is, the more each user's data can
affect the functions output and therefore the more noise we will have to
add to achieve differential privacy.

\begin{definition}[Global
Sensitivity]\protect\hypertarget{def-global-sensitivity}{}\label{def-global-sensitivity}

~

\marginnote{\begin{footnotesize}

The definition of sensitivity and the Laplace mechanism generalize
nicely to functions that output multiple real numbers.

\end{footnotesize}}

The global sensitivity of a function \(f \colon \calD^n \to \R\), for
some domain \(\calD\), is:
\[ GS_f = \max_{\text{neighbors}~x, x' \in \calD^n} \| f(x) - f(x') \|. \]
Here, ``neighbors'' inputs in \(\calD^n\) that differ in only a single
coordinate.

\end{definition}

We will also need to define the zero-mean real-valued Laplace
distribution, which is parameterized by a value \(b \in \R\). The
probability distribution function is:
\[ h(y) = \frac{1}{2b} e^{-\frac{\abs{y}}{b}} \]

Now, the Laplace mechanism for achieving a differential privacy release
of the function \(f(x_1, \dots, x_n)\) is:

\begin{enumerate}
\def\labelenumi{\arabic{enumi}.}
\tightlist
\item
  Compute the true statistic \(y \gets f(x_1, \dots, x_n)\).
\item
  Let \(b \gets \frac{GS_f}{\epsilon}\) be the global sensitivity of
  \(f\).
\item
  Sample a random noise value \(\nu\) from the Laplace distribution with
  mean \(0\) and parameter \(b\).
\item
  Output the noised statistic \(y + b\).
\end{enumerate}

\textbf{Handling non-numerical data.} The Laplace mechanism gives a very
clean way to provide differential privacy for releasing real-valued
aggregate statistics about a dataset---where the aggregate statistic is
just a number. But in many applications, we need to publish
\emph{strings}, such as names or zip codes or birth dates.

One way to handle strings is to convert them into numerical statistics
and them apply differential privacy. For example, we could publish (1)
the number of users with name ``Alice,'' (2) the number of users with
name ``Bob,'' and so on. Then we are back to working in a world of
numerical statistics and we can against use the Laplace mechanism.

\section{Challenges with differential
privacy}\label{challenges-with-differential-privacy}

As we discussed in Section~\ref{sec-dp-defn}, the more statistics you
publish about a database, the larger the differential-privacy parameter
\(\epsilon\) grows. A consequence is that if you want to publish many
statistics about a dataset, you will have to add a huge amount of noise
to the statistics you publish to achieve \(\epsilon\)-differential
privacy any reasonable value of \(\epsilon\). This is a major headache
in practice.

A second issue is that the designer of the system has to choose what
value of \(\epsilon\) to use. What value is good enough? The theory of
differential privacy cannot tell you---it's really a subjective
question.

Yet a third issue is that in many applications, a company has to publish
an aggregate statistic every day or week. Since the \(\epsilon\) values
``add up,'' after a short amount of time the effective \(\epsilon\)
value can end up being so large as to render the differential privacy
guarantee meaningless.

\part{Conclusions}

\chapter{Conclusions}\label{conclusions-1}

To conclude, we will explore five case studies that exemplify what we
have covered so far.

\section{Authentication: OPM Hack}\label{authentication-opm-hack}

In order to get security clearance to view classified documents, it is
necessary to fill out a form called an SF86 that covers all kinds of
personal details from relationships and mental health to drug use and
finances---some of the most sensitive data there is. Many, many people
have filled these out---around 2.8 million people have some level of
current security clearance. The goal of this invasive background check
was to understand people's exposure to blackmailing.

These records are stored in a database at the Office of Personnel
Management. In 2015, the OPM announced that 20 million of these records
have been exposed. This was a big problem for the US Government. Not
only did it expose exactly how to blackmail every person with security
clearance, but also records of CIA employees were not stored in this OPM
database---this meant that if you knew someone had a security clearance
but they were not in the database, you had an idea that they might be a
CIA agent.

\subsection{How it Happened}\label{how-it-happened}

\textbf{Learn Contractor Credentials.} First, the attacker somehow got a
contractor's credentials. This could have been through phishing or some
other means. The system did not use two-factor authentication, and only
around \(1\%\) of OPM users used smart cards. Importantly, the
contractor did not need to have many privileges on the system---perhaps
only enough permissions to log in.

\textbf{Compromise Root Account.} Then the attacker likely compromised
the root account on a local machine. Given that these systems were old,
this was likely easy---old version of Windows were not particularly
careful about protecting access to the root account. One way to do this
was by scheduling a job for the future. Windows allowed unprivileged
users to schedule jobs that would run as the system user, so a command
like \texttt{at\ 16:05\ /interactive\ "cmd.exe"} would open a command
prompt as the system user at 4:05 PM.

\textbf{Learn Administrator Credentials.} Even with access to the root
account on a local machine, the attacker needed to learn credentials to
be able to log into one of the machines with access to the database.
Windows did not store the password of the currently logged in user, but
it did store the \emph{hash} of the password and uses that to
authenticate to the server. This means that the client does not need the
cleartext password to log in!

As we have covered many times, passwords are terrible for
authentication. Signatures, as used with standards like FIDO2/WebAuthn,
provide much stronger security and should be used whenever possible.

\section{Transport Security: POODLE}\label{transport-security-poodle}

TLS is used all over the web. Many servers, from data centers to
embedded hardware devices, like doorbells, implement TLS. Some of these
servers take a long time to upgrade, so even if a client supports the
latest version of TLS (e.g., TLS 1.3), it might be unable to use TLS 1.3
to communicate with some servers. To handle this situation, TLS
implements version negotiation; the client and server should agree to
use the most recent version of TLS supported by both sides. However, the
way this negotiation was implemented was subtly insecure.

When an old server receives a message from a new client speaking a new
version of TLS that the server does not support, the server might reply
back with garbage data. This could be due to a bug in the server
implementation, which was never found because the developers didn't
think in advance to test against various ways that future clients
sending messages from a future version of TLS. Despite the fact that
this is a server bug, users want to still be able to communicate with
this device; after all, communicating with the device worked fine before
the client upgraded to support a new version of TLS. As a result,
clients that receive garbage data in response to their TLS connection
attempt will try to downgrade the version of TLS that they use to
connect in their next attempt. However, the ``garbage'' that the client
is reacting to was not authenticated. Therefore, an adversary could
inject garbage to force a client to downgrade to an older TLS/SSL
version.

Old versions of TLS/SSL have well-known weaknesses in their
cryptographic protocol. In particular, SSL 3.0 computed the MAC of the
message (for authentication) before encrypting the message
(MAC-then-encrypt), rather than the encrypt-then-MAC approach that we
discussed in the authenticated encryption lecture. This allowed the
attacker to send corrupted ciphertexts and see how the client or server
respond to them. In combination with MAC-then-encrypt, SSL 3.0 also used
a particular encryption construction, called cipher-block-chaining (CBC)
mode, whereby changing bits in one part of the ciphertext caused
corresponding changes to another part of the plaintext. Finally, the
plaintext payload included padding, constructed with a well-known
scheme, which was checked after decryption to make sure it was not
corrupted; if the padding was corrupted, the connection was terminated.
However, this gives a signal to the adversary as to whether their
corrupted ciphertext happens to match the expected padding or not. We
saw the CBC padding oracle attack earlier in the lecture on
authenticated encryption.

The POODLE attack was a combination of these two weaknesses (downgrade
attack and CBC padding oracle), together with running adversary-supplied
code in Javascript in the victim user's web browser, as a way of sending
partially-adversary-chosen requests over the TLS/SSL connection. The
result is that the adversary can recover other data sent to the server
over the same connection, such as the victim's cookie.

\section{Platform Security: Sony PS3
Hack}\label{platform-security-sony-ps3-hack}

The Sony PS3 originally could boot Linux and Windows. Since the hardware
was subsidized by the games that they sold, PS3s were cheaper than
comparable PCs, and for that reason PS3s were a popular option for a
cheap PC. In a software update, Sony disabled the ability to run a
custom operating system. Like the iPhone and other systems we discussed,
PS3s then used secure boot to ensure that they only boot Sony-signed
operating systems.

Sony used EC-DSA for their signatures, which resulted in signatures
along the lines of:

\[ \sigma = (g^r, r + H(\pk || g^r || m)\cdot \sk) \pmod{q} \]
\[ \sigma' = (g^r, r + H(\pk || g^r || m')\cdot \sk) \pmod{q} \]

In EC-DSA, \(r\) is supposed to be a long random number, serving as a
nonce. However, Sony re-used these nonces in pairs of signatures. This
meant that their signatures revealed their secret key! This allowed
others to sign their own operating systems.

Importantly, Sony had a plan for updating the PS3 firmware that allowed
them to ship a fix for this attack. However, attackers quickly found
flaws in every other update they shipped---it is very hard to secure a
device that the attacker has unconstrained access to.

\section{Software Security: WannaCry
Ransomware}\label{software-security-wannacry-ransomware}

Ransomware is a type of malware that encrypt important-looking files on
the infected system and demands a payment in Bitcoin to decrypt the
files. This is inconvenient and upsetting for personal computers, but
for enterprise computer systems this can cause huge monetary losses. For
hospital systems, this can even lead to loss of life. This WannaCry
randomware infected hundreds of thousands of computers and caused
billions of dollars of damange, but did not make much money due to bad
payment systems and slow decryption.

The bugs used to enable this were part of an exploit developed and kept
secret by the NSA called EternalBlue. It took advantage of several C
bugs, including an invalid cast, a parser bug, and an allocation bug, to
eventually achieve remote code execution over the network on a Windows
system.

The NSA intended to keep these exploits to themselves, and thus did not
inform Microsoft (or anyone else) about the bugs in their software.
However, an NSA contractor took terabytes of NSA data home with him,
including this exploit. On his home computer, he ran Kaspersky
antivirus. Importantly, Kaspersky sends suspicious files home for
analysis. Wall Street Journal reported that this was likely how the
exploit leaked and became a part of malware.

In order to spread to many computers, WannaCry looked something like the
following:

\begin{enumerate}
\def\labelenumi{\arabic{enumi}.}
\tightlist
\item
  Connect to a website at a random-looking address and exit if it
  succeeds. Security teams would often analyze software that they
  thought may include malware by running it in a VM, allowing network
  requests to succeed, and watching what happens. This random-looking
  domain did not really exist, so this may have been a way to try to
  detect when the software is running in a VM and make it behave
  ``normally``.
\item
  Install Tor and connect to command-and-contral infrastructure.
\item
  Encrypt all files with a fixed set of extensions with RSA and AES.
\item
  Demand a ransome to be paid to one of four static Bitcoin addresses.
\item
  Spread itself by trying to perform the exploit on all IPs in the local
  network.
\end{enumerate}

When designing a system, it is prudent to have the system as simple as
possible, since less software leads to fewer bugs. And of course, any
bug is a security bug---each of the bugs used in the EternalBlue exploit
did not look like a security bug, but the combination of them allowed
for a powerful exploit.

\section{Privacy: US Census}\label{privacy-us-census}

The US Census, performed every decade, collects data used to allocate
seats in the House of Representatives and by many researchers. The
Census Bureau is mandated to make this information public, but is
forbidden by law from publishing any data that allows individuals to be
identified.

In the 2020 census, the bureau used differential privacy to protect
released data from de-identification. However, they used
\(\epsilon = 19.61\) to avoid adding so much noise that the data lost
its utility. This meant that if the probability of some event happening
to an individual without the release of these data was \(p\), the
release of these data with \(\epsilon = 19.61\) was guaranteed to make
the probility at most \(\approx e^{19.61} p\), or around a million times
the original probability.

\bookmarksetup{startatroot}

\chapter*{References}\label{references}
\addcontentsline{toc}{chapter}{References}

\markboth{References}{References}

\protect\phantomsection\label{refs}
\begin{CSLReferences}{1}{1}
\bibitem[\citeproctext]{ref-BR93}
Bellare, Mihir, and Phillip Rogaway. 1993. {``Random Oracles Are
Practical: A Paradigm for Designing Efficient Protocols.''} \emph{ACM
Conference on Computer and Communications Security}.

\bibitem[\citeproctext]{ref-bittau2014hacking}
Bittau, Andrea, Adam Belay, Ali Mashtizadeh, David Mazières, and Dan
Boneh. 2014. {``Hacking Blind.''} \emph{IEEE Symposium on Security and
Privacy}, 227--42.

\bibitem[\citeproctext]{ref-DH76}
Diffie, Whitfield, and Martin E Hellman. 1976. {``New Directions in
Cryptography.''} \emph{Transactions on Information Theory} 22 (6).

\bibitem[\citeproctext]{ref-DH77}
Diffie, Whitfield, and Martin E. Hellman. 1977. {``Exhaustive
Cryptanalysis of the {NBS} Data Encryption Standard.''} \emph{Computer}
6 (10): 74--84.

\bibitem[\citeproctext]{ref-flickr}
Duong, Thai, and Juliano Rizzo. 2009. \emph{{Flickr's} {API} Signature
Forgery Vulnerability}.
\href{https://vnhacker.blogspot.com/2009/09/flickrs-api-signature-forgery.html}{Https://vnhacker.blogspot.com/2009/09/flickrs-api-signature-forgery.html}.

\bibitem[\citeproctext]{ref-granville2008smooth}
Granville, Andrew. 2008. {``Smooth Numbers: Computational Number Theory
and Beyond.''} \emph{Algorithmic Number Theory: Lattices, Number Fields,
Curves and Cryptography} 44: 267--323.

\bibitem[\citeproctext]{ref-HILL99}
Hastad, J., R. Impagliazzo, L. A. Levin, and M. Luby. 1999. {``A
Pseudorandom Generator from Any One-Way Function.''} \emph{SIAM Journal
on Computing} 28 (4): 1364--96.

\bibitem[\citeproctext]{ref-hill:spectre-meltdown}
Hill, Mark D., Jon Masters, Parthasarathy Ranganathan, Paul Turner, and
John L. Hennessy. 2019. {``On the {Spectre} and {Meltdown} Processor
Security Vulnerabilities.''} \emph{IEEE Micro} 39 (2): 9--19.

\bibitem[\citeproctext]{ref-honan}
Honan, Mat. 2012. {``How {Apple} and {Amazon} Security Flaws Led to My
Epic Hacking.''} In \emph{Wired}.
\href{https://www.wired.com/2012/08/apple-amazon-mat-honan-hacking/}{Https://www.wired.com/2012/08/apple-amazon-mat-honan-hacking/}.

\bibitem[\citeproctext]{ref-koscher2010experimental}
{Koscher, Karl, Alexei Czeskis, Franziska Roesner, et al.} 2010.
{``Experimental Security Analysis of a Modern Automobile.''} \emph{IEEE
Symposium on Security and Privacy}.

\bibitem[\citeproctext]{ref-L79}
Lamport, Leslie. 1979. \emph{Constructing Digital Signatures from a One
Way Function}.

\bibitem[\citeproctext]{ref-lenstra2000integer}
Lenstra, Arjen K. 2000. {``Integer Factoring.''} \emph{Designs Codes,
and Cryptography} 19 (2/3).

\bibitem[\citeproctext]{ref-lenstra:universal}
Lenstra, Arjen K, Thorsten Kleinjung, and Emmanuel Thomé. 2013.
{``Universal Security.''} \emph{Number Theory and Cryptography}.

\bibitem[\citeproctext]{ref-motoyama:recaptchas}
Motoyama, Marti, Kirill Levchenko, Chris Kanich, Damon McCoy, Geoffrey
M. Voelker, and Stefan Savage. 2010. {``Re: {CAPTCHAs}---Understanding
{CAPTCHA}-Solving Services in an Economic Context.''} \emph{Proceedings
of the 19th {USENIX} Security Symposium} (Washington, DC), August.

\bibitem[\citeproctext]{ref-muller2019practical}
Müller, Jens, Fabian Ising, Vladislav Mladenov, Christian Mainka,
Sebastian Schinzel, and Jörg Schwenk. 2019. {``Practical Decryption
Exfiltration: Breaking Pdf Encryption.''} \emph{ACM CCS}.

\bibitem[\citeproctext]{ref-nemec2017return}
Nemec, Matus, Marek Sys, Petr Svenda, Dusan Klinec, and Vashek Matyas.
2017. {``The Return of {Coppersmith's} Attack: Practical Factorization
of Widely Used {RSA} Moduli.''} \emph{CCS}.

\bibitem[\citeproctext]{ref-RSA}
Rivest, Ronald L., Adi Shamir, and Leonard Adleman. 1978. {``A Method
for Obtaining Digital Signatures and Public-Key Cryptosystems.''}
\emph{Communications of the ACM} 21 (2): 120--26.

\bibitem[\citeproctext]{ref-thompson1984reflections}
Thompson, Ken. 1984. {``Reflections on Trusting Trust.''}
\emph{Communications of the ACM} 27 (8).

\bibitem[\citeproctext]{ref-icloud}
Trew, J. 2014. \emph{{``{Find My iPhone}''} Exploit May Be to Blame for
Celebrity Photo Hacks (Update)}.
\href{https://www.engadget.com/2014-09-01-find-my-iphone-exploit.html}{Https://www.engadget.com/2014-09-01-find-my-iphone-exploit.html}.

\bibitem[\citeproctext]{ref-CW81}
Wegman, Mark N, and J Lawrence Carter. 1981. {``New Hash Functions and
Their Use in Authentication and Set Equality.''} \emph{Journal of
Computer and System Sciences} 22 (3).

\bibitem[\citeproctext]{ref-yilek2009private}
Yilek, Scott, Eric Rescorla, Hovav Shacham, Brandon Enright, and Stefan
Savage. 2009. {``When Private Keys Are Public: Results from the 2008
{Debian OpenSSL} Vulnerability.''} \emph{SIGCOMM}.

\bibitem[\citeproctext]{ref-palin}
Zetter, Kim. 2008. {``{Palin} e-Mail Hacker Says It Was Easy.''} In
\emph{Wired}.
\href{https://www.wired.com/2008/09/palin-e-mail-ha/}{Https://www.wired.com/2008/09/palin-e-mail-ha/}.

\end{CSLReferences}

\cleardoublepage
\phantomsection
\addcontentsline{toc}{part}{Appendices}
\appendix

\chapter{Factoring integers}\label{sec-fact}

The problem of integer factorization was central to 20th-century
cryptography. Breaking the one-wayness of the RSA trapdoor one-way
function (Chapter~\ref{sec-rsa}), for example, is no harder than
factoring integers. In this chapter, we will see a couple of
surprisingly powerful algorithms for factoring integers.

We will only consider factoring numbers of the form \(N=pq\), for
distinct odd primes \(p\) and \(q\). (The general case is not too much
more challenging.) Throughout, let \(n = \lceil \log_2 N \rceil\) be the
bitlength of the number to factor.

\section{Background}\label{sec-fact-bg}

\textbf{Trial division.} We can factor \(N\) by trying to divide \(N\)
by each of the primes of size \(\leq \sqrt N\) and checking whether the
result is an integer. If so, we have found a factor of \(N\). Since at
least one of the two factors of \(N\) is in \(\{1, \dots, \sqrt N\}\),
this algorithm (``trial division'') runs in time roughly
\(\sqrt{N} = 2^{n/2}\).

Trial division is an \emph{exponential time} algorithm, since it runs in
time \(2^{\Omega(n)}\), where the bitlength \(n\) is the size of the
number to be factored. The best known factoring algorithms run in
\emph{sub-exponential time} \(2^{O(n^c)}\), for some constant \(c < 1\).

\textbf{Euclid's algorithm.} An important subroutine in almost all
factoring algorithms is Euclid's polynomial-time algorithm for computing
the greatest common divisor of two integers \(x\) and \(y\).

\marginnote{\begin{footnotesize}

In this discussion, we draw on Arjen Lenstra's very nice survey on
factoring (\citeproc{ref-lenstra2000integer}{Lenstra 2000}).

\end{footnotesize}}

The principle of Euclid's algorithm is that
\[ \gcd(x, y) = \gcd(x, y \bmod x) \qquad \text{and} \qquad \gcd(x, 0) = x.\]
So, for example, if we want to compute \(\gcd(46, 12)\), we can compute
it as: \[ \gcd(46,12) = \gcd(12, 10) = \gcd(10, 2) = 2. \]

\textbf{Difference of squares.} The second key idea is that, if we can
find two numbers \(x,y \in \Z\) whose squares are congruent modulo
\(N\), we can use these numbers to factor \(N\): \[\begin{aligned}
x^2 &= y^2 &\pmod N \\
x^2 - y^2 &= 0 &\pmod N\\
(x+y)(x-y) &= 0 &\pmod N\\
\end{aligned}\]

If \(x = \pm y\), then this relation is not helpful to us.

\marginnote{\begin{footnotesize}

It is not necessarily obvious that the useful pairs \((x,y)\) will ever
exist. The key idea is that, modulo \(N=pq\), ever number in \(\Z^*_N\)
either has four square roots or has none. If an element in \(\Z^*_N\)
has four square roots then the roots are of the form \(r,-r, s, -s\). In
this case, a pair \((\pm r, \pm s)\) yields the sort of relation that we
need to factor.

\end{footnotesize}}

But if \(x \neq \pm y\), then we know that \(x+y \neq 0 \bmod N\) and
\(x-y \neq 0 \bmod N\). So we have: \[ (x+y)(x-y) = kN \quad \in \Z,\]
for some positive integer \(k \in \Z\). Then \(x+y\) must be a multiple
of one of the factors of \(N\) (but not both), and \(\gcd(x+y, N)\)
reveals a factor of \(N\).

The goal of many factoring algorithms---including the one we will see
today---is finding these integers \(x\) and \(y\) whose squares are
congruent modulo \(N\).

\section{Dixon's algorithm}\label{dixons-algorithm}

Dixon's algorithm is one of the simplest sub-exponential-time factoring
algorithms. It gives a fast method for finding two numbers whose squares
are congruent modulo \(N\). Once we have these squares, we can use them
to factor as in Section~\ref{sec-fact-bg}

\subsection{The idea}\label{the-idea}

The principle of Dixon's algorithm is that we will pick many random
numbers \(r \in \Z^*_N\) and square them modulo the integer \(N\) we
would like to factor.

Say that we are somehow able to find numbers \(r, r'\) such that
\[\begin{aligned}
  r^2 &= 2 \cdot 3^2 \cdot 5 &\pmod N\\
  r'^2 &= 2 \cdot 5 &\pmod N,\\
\end{aligned}\] then we know that: \[\begin{aligned}
(rr')^2 &= 2^2 \cdot 3^2 \cdot 5^2 &\pmod N\\
(rr')^2 &= (2 \cdot 3 \cdot 5)^2 &\pmod N
\end{aligned}\] and now we have two numbers whose squares are congruent
modulo \(N\):
\[ x = r r' \qquad \text{ and }\qquad y = 2 \cdot 3 \cdot 5.\] If we are
lucky, this is the useful type of congruence that we can use to factor
\(N\) (i.e., \(r r' \neq \pm 2 \cdot 2 \cdot 5 \bmod N\)).

The principle of Dixon's algorithm is to generate many such \(r\)s and
then use linear algebra to find a subset of them whose product modulo
\(N\) is a perfect square.

\subsection{The algorithm}\label{the-algorithm}

\textbf{Input:} An integer \(N = pq\) for odd primes \(p\) and \(q\). A
parameter \(B \in \N\), which we refer to as ``the size of the factor
base.''

\textbf{Output:} The factors \((p,q)\) of \(N\).

\begin{enumerate}
\def\labelenumi{\arabic{enumi}.}
\item
  \textbf{Collect linear relations.} Maintain a list \(L\) of pairs of

  \begin{enumerate}
  \def\labelenumii{(\alph{enumii})}
  \tightlist
  \item
    an element in \(\Z^*_N\) and (b) vectors over \(\Z_2^B\). Repeat
    until \(L\) contains \(B+1\) pairs:
  \end{enumerate}

  \begin{itemize}
  \item
    Sample \(r \getsr \Z^*_N\).
  \item
    Compute \(s \gets (r^2 \bmod N)\).
  \item
    Attempt to write \(s\) as a product of the first \(B\) primes:
    \[ s = 2^{e_2} 3^{e_3} 5^{e_5} \dots \]

    \marginnote{\begin{footnotesize}

    If a number completely splits into prime factors \(\leq B\), we say
    that the number is ``\(B\)-smooth.''

    \end{footnotesize}}
  \item
    If successful, add the pair \((r, (e_2, e_3, e_5, \dots))\) to the
    list \(L\).
  \end{itemize}
\item
  \textbf{Solve linear system.} Let
  \(L = \{ (r_1, \vec v_1), (r_2, \vec v_2), \dots \}\). Find a non-zero
  combination of the vectors in \(L\) that sums to zero modulo \(2\).
  That is, find \(S \subseteq [B+1]\) such that
  \[ \sum_{i \in S} \vec v_i = (0, 0, 0, \cdots, 0) \qquad \in \Z_2^B. \]

  \marginnote{\begin{footnotesize}

  We can find the set \(S\) using Gaussian elimination in roughly
  \(B^3\) time. Since the set of vectors will be extremely sparse, there
  are faster methods that implementers use in practice.

  \end{footnotesize}}

  Letting \[ (e_2, e_3, e_5, \dots) \gets \sum_{i \in S} \vec v_i,\] we
  then have a difference of squares:
  \[ \left(\sum_{i \in S} r_i \right)^2 = (2^{(\frac{e_2}{2})} \cdot 3^{(\frac{e_3}{3})} \cdot 5^{(\frac{e_5}{2})} \cdots )^2 \qquad \pmod N.\]
\item
  \textbf{Use Euclid's algorithm to try to factor \(N\).} We can take:
  \[\begin{aligned}
    x &= \left(\sum_{i \in S} r_i \right)\\
    y&= 2^{(\frac{e_2}{2})} \cdot 3^{(\frac{e_3}{3})} \cdot 5^{(\frac{e_5}{2})} \cdots )
  \end{aligned}\] and compute \(\gcd(x+y, N)\). With probability roughly
  \(1/2\), over the random choice of the \(r\)s, this will yield a
  factor of \(N\).
\end{enumerate}

\subsection{The analysis.}\label{the-analysis.}

The costs of the three steps of the algorithms are:

\begin{enumerate}
\def\labelenumi{\arabic{enumi}.}
\item
  Each iteration of the loop requires us to try to factor a number into
  primes \(\leq B\). We can factor in this way by trial division using
  time roughly \(B\).

  \marginnote{\begin{footnotesize}

  I'm ignoring any \(\log B\) factors, which do matter very much in
  practice.

  \end{footnotesize}}

  The question then is how many trials it will take for us to find a
  single smooth number. For a smoothness bound \(B\), let's say for now
  that it takes \(T(B)\) trials---we will look into the precise value of
  \(T(B)\) in a moment..
\item
  Solving the linear system using Gaussian elimination takes roughly
  \(B^3\) time.
\item
  Run Euclid's algorithm---the time required here is negligible compared
  to the time of the first two steps. This step runs in time
  \(\poly(\log N) = \poly(n)\).
\end{enumerate}

Putting everything together, we have that factoring an \(n\)-bit number
with a factor base of size \(B\) takes time:
\begin{equation}\protect\phantomsection\label{eq-time}{
B \cdot T(B) + B^3 + \poly(n).
}\end{equation}

\textbf{Smoothness probabilities.} The key question that we need to
answer to complete the analysis is

\begin{quote}
``If we pick an integer uniformly at random from \(\{1, \dots, N\}\),
what is the probability that the integer will be \(B\)-smooth?''
\end{quote}

The convention is to denote the number of \(B\)-smooth numbers in
\(\{1, \dots, N\}\) as \(\Psi(N,B)\). When \(B\) is ``not too small,''
we have:

\marginnote{\begin{footnotesize}

For many more details on these estimates, take a look at Granville's
very nice survey on smooth numbers.
(\citeproc{ref-granville2008smooth}{Granville 2008})

\end{footnotesize}}

\[ \Psi(N, B) \approx N \cdot u^{-u + o(1)}\qquad \text{for } u = \frac{\log N}{\log B}.\]

The probability of a random number modulo \(N\) being \(B\)-smooth is
then \(\Psi(N,B)/N\) and the expected number of tries it will take for
us to find a smooth number is: \[ T(B) = 1/\Psi(N,B).\]

\marginnote{\begin{footnotesize}

We are being slightly imprecise---we actually need to know the number of
squares (quadratic residues) modulo \(N\) that are smooth. But
heuristically, we can assume that quadratic residues behave like random
integers modulo \(N\) for the purposes of smoothness.

\end{footnotesize}}

Now we can plug this estimate into the expression
(Equation~\ref{eq-time}) for the running time of Dixon's algorithm and
we can solve for the value of \(B\) that minimizes the running time. In
particular, to minimize the running time we want: \[\begin{aligned}
  B \approx T(B) \approx N / \Psi(N,B) = u^u,
\end{aligned}\] for \(u = (\log N)/(\log B)\).

\[\begin{aligned}
  B &= u^u\\
  \log B &= u \log u\\
  \log B &= \frac{\log N}{\log B} \log \frac{\log N}{\log B}\\
  \log^2 B &\approx \log N \log \log N\\
\log B &\approx \sqrt{\log N \log \log N}\\
B &\approx \exp(\sqrt{\log N \log \log N}).\\
\end{aligned}\]

If we plug this value of \(B\) into Dixon's algorithm, we get a running
time of
\[ \exp(O(\sqrt{ \log N \log \log N}) = 2^{O(\sqrt{ n \log n })}.\]


\backmatter


\end{document}
